\kantbook{Tweede Boek: Analytiek van het Sublieme}

\section{Overgang van het vermogen om het schone te oordelen naar dat om het verhevene te oordelen.}

Het schone valt samen met het verhevene doordat beide op zichzelf behagen. En verder doordat beide noch een zintuiglijk oordeel noch een logisch bepalend oordeel veronderstellen, maar een reflecterend oordeel: bijgevolg hangt de bevrediging niet af van een sensatie, zoals bij het aangename, noch van een bepaald begrip, zoals de bevrediging bij het goede; maar zij is niettemin nog steeds verwant aan begrippen, hoewel onbepaald welke, vandaar dat de bevrediging verbonden is aan de loutere voorstelling of aan het vermogen daartoe, waardoor het voorstellingsvermogen of de verbeelding, in het geval van een gegeven aanschouwing, beschouwd wordt als in overeenstemming met het begripsvermogen van het verstand of de rede, als het laatste bevorderend.
Derhalve zijn beide soorten oordelen ook enkelvoudig, en toch oordelen die pretenderen universeel geldig te zijn ten aanzien van elk subject, hoewel zij slechts aanspraak maken op het genotgevoel en niet op enige kennis van het object.

Maar opvallende verschillen tussen beide treffen ook het oog. Het schone in de natuur betreft de vorm van het object, die bestaat in beperking; het verhevene daartegen is te vinden in een vormloos object voor zover grenzeloosheid daarin voorgesteld wordt, of bij gelegenheid daarvan, en toch ook als een totaliteit gedacht wordt: zodat het schone schijnt te worden opgenomen als de voorstelling van een onbepaald begrip van het verstand, maar het verhevene als dat van een dergelijk begrip van de rede.
Derhalve is de bevrediging in het eerste geval verbonden met de voorstelling van kwaliteit, maar in dit geval met die van kwantiteit. Ook is het laatste genot van geheel andere aard dan het eerste, doordat het eerste (het schone) direct een gevoel van bevordering van het leven met zich meebrengt, en derhalve verenigbaar is met bekoringen en een spelende verbeelding, terwijl het laatste (het gevoel van het verhevene) een genot is dat slechts indirect ontstaat, namelijk opgewekt door het gevoel van een ogenblikkelijke remming van de levenskrachten en de onmiddellijk daarop volgende en des te krachtiger uitstorting daarvan; derhalve als een aandoening schijnt het geen spel, maar iets ernstigs in de werkzaamheid van de verbeelding te zijn.
Derhalve is het ook onverenigbaar met bekoringen, en, aangezien de geest niet slechts door het object wordt aangetrokken, maar ook altijd wederkerig door het wordt afgestoten, bevat de bevrediging in het verhevene niet zozeer een positief genot als wel bewondering of eerbied, d.w.z. het verdient negatief genot genoemd te worden.

Het belangrijkste en wezenlijke verschil tussen het verhevene en het schone is echter dit: dat wij, indien wij, zoals passend is, hier eerst slechts het verhevene in objecten van de natuur beschouwen (dat in de kunst immers altijd beperkt is tot de voorwaarden van overeenstemming met de natuur), de natuurlijke schoonheid (de op zichzelf staande soort) een doelmatigheid in haar vorm met zich meebrengt, waardoor het object als het ware vooraf bepaald schijnt voor ons oordeelsvermogen, en aldus een object van bevrediging op zichzelf uitmaakt, terwijl datgene wat, zonder enige rationalisering, louter in de aanschouwing, het gevoel van het verhevene in ons opwekt, weliswaar in zijn vorm voor ons oordeelsvermogen doelloos kan schijnen, ongeschikt voor ons voorstellingsvermogen, en als het ware geweld aandoet aan onze verbeelding, maar desondanks des te verhevener wordt geoordeeld.
Maar hieruit ziet men onmiddellijk dat wij ons in het algemeen onjuist uitdrukken indien wij een object van de natuur verheven noemen, hoewel wij er zeer vele terecht schoon kunnen noemen; want hoe kunnen wij met een uitdrukking van instemming aanduiden datgene wat op zichzelf als doelloos wordt aanschouwd?
Wij kunnen niet meer zeggen dan dat het object dient voor de voorstelling van een verhevenheid die in de geest kan worden gevonden; want wat eigenlijk verheven is, kan in geen zinnelijke vorm bevatten worden, maar betreft slechts ideeën van de rede, die, hoewel geen aan hen adequate voorstelling mogelijk is, juist door deze onvoldoende voorstelling worden opgeroepen en in herinnering gebracht, welke toch een zinnelijke voorstelling toelaat.
Zo kan de wijde oceaan, door stormen vertoornd, niet verheven genoemd worden. Zijn aanblik is afschuwelijk; en men moet de geest reeds met allerlei ideeën hebben gevuld indien door middel van zo’n aanschouwing hij in de stemming moet worden gebracht voor een gevoel dat op zichzelf verheven is, doordat de geest wordt aangespoord de zinnelijkheid te verlaten en zich bezig te houden met ideeën die een hogere doelmatigheid bevatten.
De op zichzelf staande schoonheid van de natuur openbaart ons een techniek van de natuur, die het mogelijk maakt haar voor te stellen als een systeem in overeenstemming met wetten waarvan het principe wij nergens in ons gehele verstandsvermogen aantreffen, namelijk dat van een doelmatigheid met betrekking tot het gebruik van het oordeelsvermogen ten aanzien van verschijnselen, zodat dit beoordeeld moet worden als niet slechts toebehorend aan de natuur in haar doelloze mechaniek, maar veeleer ook aan de analogie met de kunst.
Zo breidt het niet onze kennis van natuurlijke objecten uit, maar ons begrip van de natuur, namelijk als louter mechaniek, tot het begrip van de natuur als kunst: wat diepgaande onderzoeken uitnodigt naar de mogelijkheid van zo’n vorm.
Maar in datgene wat wij gewoon zijn verheven in de natuur te noemen, is er zo weinig dat leidt tot bijzondere objectieve principes en vormen van de natuur die daaraan beantwoorden, dat het meestal veeleer in haar chaos of in haar wildste en meest ongeregelde wanorde en verwoesting is, indien het slechts een glimp van grootsheid en macht toelaat, dat het de ideeën van het verhevene opwekt.

Hieruit zien wij dat het begrip van het verhevene in de natuur verre van zo belangrijk en rijk aan gevolgen is als dat van haar schoonheid, en dat het in het algemeen niets doelmatigs in de natuur zelf aangeeft, maar slechts in het mogelijke gebruik van haar aanschouwingen om in ons een doelmatigheid waarneembaar te maken die geheel onafhankelijk is van de natuur.
Voor het schone in de natuur moeten wij een grond buiten onszelf zoeken, maar voor het verhevene slechts een in onszelf en in de wijze van denken die verhevenheid in de voorstelling van het eerste inbrengt, een zeer noodzakelijke inleidende opmerking, die de ideeën van het verhevene geheel scheidt van die van een doelmatigheid van de natuur, en van de theorie van het verhevene een louter aanhangsel maakt van het esthetisch oordelen over de doelmatigheid van de natuur, aangezien hierdoor geen bijzondere vorm in het laatste wordt voorgesteld, maar slechts een doelmatig gebruik dat de verbeelding van haar voorstelling maakt, ontwikkeld wordt.

De op zichzelf staande schoonheid van de natuur openbaart ons een techniek van de natuur, die het mogelijk maakt haar voor te stellen als een systeem in overeenstemming met wetten waarvan het principe wij nergens in ons gehele verstandsvermogen aantreffen, namelijk dat van een doelmatigheid met betrekking tot het gebruik van het oordeelsvermogen ten aanzien van verschijnselen, zodat dit beoordeeld moet wordena als niet slechts toebehorend aan de natuur in haar doelloze mechaniek, maar veeleer ook aan de analogie metb de kunst.
Zo breidt het niet onze kennis van natuurlijke objecten uit, maar ons begrip van de natuur, namelijk als louter mechaniek, tot het begrip van de natuur als kunst: wat diepgaande onderzoeken uitnodigt naar de mogelijkheid van zo’n vorm.
Maar in datgene wat wij gewoon zijn verheven in de natuur te noemen, is er zo weinig dat leidt tot bijzondere objectieve principes en vormen van de natuur die daaraan beantwoorden, dat het meestal veeleer in haar chaos of in haar wildste en meest ongeregelde wanorde en verwoesting is, indien het slechts een glimp van grootsheid en macht toelaat, dat het de ideeën van het verhevene opwekt.
Hieruit zien wij dat het begrip van het verhevene in de natuur verre van zo belangrijk en rijk aan gevolgen is als dat van haar schoonheid, en dat het in het algemeen niets doelmatigs in de natuur zelf aangeeft, maar slechts in het mogelijke gebruik van haar aanschouwingen om in ons een doelmatigheid waarneembaar te maken die geheel onafhankelijk is van de natuur.
Voor het schone in de natuur moeten wij een grond buiten onszelf zoeken, maar voor het verhevene slechts een in onszelf en in de wijze van denken die verhevenheid in de voorstelling van het eerste inbrengt, een zeer noodzakelijke inleidende opmerking, die de ideeën van het verhevene geheel scheidt van die van een doelmatigheid van de natuur, en van de theorie van het verhevene een louter aanhangsel maakt van het esthetisch oordelen over de doelmatigheid van de natuur, aangezien hierdoor geen bijzondere vorm in het laatste wordt voorgesteld, maar slechts een doelmatig gebruik dat de verbeelding van haar voorstelling maakt, ontwikkeld wordt.

\begin{mdframed}[leftmargin=0pt,rightmargin=0pt,backgroundcolor=blue!6,linecolor=white,innertopmargin=6pt,innerbottommargin=6pt]

Het schone en het verhevene hebben gemeen dat ze beide op zichzelf bevallen en niet afhankelijk zijn van zintuiglijke waarneming of logische oordelen, maar van een reflecterend oordeel. Beide roepen een bevrediging op die niet gebaseerd is op sensatie (zoals bij het aangename) of op een vastomlijnd begrip (zoals bij het goede), maar die wel gerelateerd is aan begrippen, ook al zijn die niet precies gedefinieerd. Daardoor is de bevrediging verbonden met de loutere voorstelling of het vermogen daartoe, waarbij de verbeelding in overeenstemming lijkt met het begripsvermogen van het verstand of de rede.
Beide soorten oordelen zijn uniek, maar claimen universele geldigheid voor ieder subject, ook al gaat het slechts om een gevoel van genot en niet om kennis van het object.

Het schone in de natuur heeft betrekking op de vorm van het object, die beperkt en afgebakend is. Het verhevene daarentegen vind je in vormloze objecten, waarin grenzeloosheid wordt voorgesteld, maar die toch als een totaliteit worden gedacht. Het schone lijkt een onbepaald begrip van het verstand te presenteren, terwijl het verhevene een vergelijkbaar begrip van de rede weergeeft. Bij schoonheid gaat de bevrediging over kwaliteit, bij het verhevene over kwantiteit.
Het genot van schoonheid is direct en levensbevorderend, verenigbaar met bekoring en een speelse verbeelding. Het verhevene daartegen roept een genot op dat indirect ontstaat: eerst is er een gevoel van remming van de levenskrachten, gevolgd door een krachtige uitstorting daarvan. Dit maakt het verhevene serieus en onverenigbaar met bekoring; het wekt eerbied of bewondering op, en kan zelfs als een ‘negatief genot’ worden beschouwd.

Het grootste verschil is dat natuurlijke schoonheid een doelmatigheid in haar vorm toont, alsof het object van nature past bij ons oordeelsvermogen. Het verhevene daartegen lijkt juist doelloos, ongeschikt voor onze voorstellingskracht, en doet geweld aan aan onze verbeelding, maar wordt juist daardoor als des te verhevener ervaren.

We drukken ons eigenlijk verkeerd uit als we een natuurobject ‘verheven’ noemen. Wat verheven is, zit niet in het object zelf, maar in de ideeën van de rede die het oproept. Een woeste oceaan is op zichzelf niet verheven; zijn aanblik is afschrikwekkend. Pas als onze geest gevuld is met ideeën, kan zo’n aanschouwing een verheven gevoel oproepen, doordat we ons losmaken van zintuiglijke waarneming en ons richten op hogere, doelmatige ideeën.

Natuurlijke schoonheid toont ons een techniek in de natuur, alsof ze volgens wetten functioneert die ons verstand niet kent: een doelmatigheid die past bij ons oordeelsvermogen. Dit breidt ons begrip van de natuur uit, van louter mechaniek naar een soort kunst. Het verhevene daartegen vind je niet in specifieke natuurvormen, maar juist in chaos, wildheid en verwoesting, als die maar een glimp van grootsheid en kracht tonen. Het verhevene zegt niets over doelmatigheid in de natuur zelf, maar over hoe wij onze aanschouwingen kunnen gebruiken om een doelmatigheid in onszelf te ervaren die onafhankelijk is van de natuur.

Voor schoonheid zoeken we een grond buiten onszelf, maar voor het verhevene ligt die in onszelf en in onze manier van denken. Dit onderscheidt het verhevene van de doelmatigheid van de natuur en maakt het tot een aanhangsel van het esthetisch oordeel over natuurlijke doelmatigheid.

\bigskip
Waarom ervaren we natuur als doelmatig zonder dat we haar een doel toeschrijven? Kant wijst erop dat onze ervaring van natuurlijke schoonheid niet voortkomt uit een bewuste toekenning van een doel of bedoeling, maar uit een spontane overeenstemming tussen de vorm van natuurverschijnselen en onze eigen kenvermogens. Het is alsof de natuur, zonder dat we haar een expliciete functie of intentie toedichten, toch past bij de manier waarop wij haar waarnemen en beoordelen. Deze ervaring is geen kennis over de natuur zelf, maar een subjectieve wijze van oordelen: we ervaren de natuur als “geschikt” voor ons vermogen om haar te aanschouwen en te beoordelen, zonder dat we daarbij een concreet doel of een externe bedoeling hoeven aan te nemen. Het is een gevoel van harmonie dat ontstaat in de reflectie, niet in de analyse.
Kant benadrukt dat deze ervaring van doelmatigheid niet berust op een logische afleiding of een zintuiglijke reactie, maar op een reflecterend oordeel. De vorm van natuurverschijnselen, hun beperking, hun evenwicht, hun structuur, spreekt ons aan alsof ze voor ons zijn “ontworpen”, zonder dat we daarbij een ontwerper of een doel hoeven te veronderstellen. Dit is geen objectieve eigenschap van de natuur, maar een subjectieve ervaring: de natuur lijkt doelmatig omdat ze aansluit bij onze wijze van waarnemen en denken. Het is een ervaring die ons uitnodigt om de natuur niet louter als een mechanisch geheel te zien, maar als iets dat in harmonie is met onze eigen geestelijke vermogens. Hierin ligt de kern van Kants idee van “zuivere” schoonheid: de natuur bevalt ons niet om wat ze doet, maar om wie we zijn als we haar aanschouwen.
In de Kritik der Urteilskraft vormt deze gedachte een cruciale schakel. Kant bereidt hiermee de weg voor naar het idee dat natuurlijke schoonheid een symbolische betekenis kan dragen, een brug tussen onze zintuiglijke ervaring en morele ideeën. Door de natuur als doelmatig te ervaren, zonder haar een doel toe te schrijven, openen we de mogelijkheid om haar te zien als een teken van een diepere orde, een orde die niet mechanisch is, maar die aansluit bij onze vermogens tot oordeel en verbeelding. Dit is de eerste stap richting Kants latere koppeling van esthetische ervaring aan moraliteit: als de natuur ons aanspreekt alsof ze voor ons is “bestemd”, dan kunnen we haar ook gaan zien als een weerspiegeling van de morele wetten die wij in onszelf dragen. De ervaring van schoonheid wordt zo niet alleen een kwestie van smaak, maar een opening naar een bredere, meer fundamentele harmonie tussen de mens en de wereld.
De natuur is voor Kant een sterker voorbeeld van “zuivere” schoonheid dan kunst, omdat ze ons confronteert met een schoonheid die niet het product is van menselijke intentie. Kunst is altijd gebonden aan een bedoeling, die van de kunstenaar, en roept daardoor onvermijdelijk vragen op over doel en functie. De natuur daartegenover lijkt doelmatig zonder dat we een doel hoeven te veronderstellen; ze spreekt ons aan alsof ze voor ons is, zonder dat we haar kunnen reduceren tot een menselijk ontwerp. Dit zegt iets diepgaands over onze houding tegenover een wereld die ons niet “bedoeld” lijkt: ook als we geen doel of bedoeling kunnen aanwijzen, ervaren we de natuur als betekenisvol. Ze nodigt ons uit om haar niet alleen te begrijpen, maar ook te ervaren als iets dat in harmonie is met wie we zijn. In die ervaring ligt misschien wel de meest fundamentele openbaring van Kants esthetica: dat schoonheid ons leert om de wereld niet alleen te zien als een mechanisme, maar als een weerspiegeling van onze eigen vermogens tot betekenisgeving.

\end{mdframed}

\section{Over de verdeling van een onderzoek naar het gevoel van het verhevene.}

Wat de verdeling van de momenten van het esthetische oordeel over voorwerpen in relatie tot het gevoel van het verhevene betreft, zal de analyse volgens hetzelfde principe kunnen verlopen als bij de analyse van smaakoordelen. Want als oordeel van de esthetische reflecterende oordeelskracht moet de bevrediging in het verhevene, net als die in het schone, worden voorgesteld als algemeen geldig in haar \textbf{hoedanigheid}, als belangloos in haar \textbf{kwaliteit}, als subjectieve doelmatigheid in haar \textbf{betrekking}, en deze laatste, wat haar \textbf{modaliteit} aangaat, als noodzakelijk. De methode hier zal dus niet afwijken van die in de voorgaande paragraaf, hoewel er rekening moet worden gehouden met het feit dat we daar, waar het esthetische oordeel de vorm van het voorwerp betrof, begonnen met het onderzoek naar de kwaliteit, maar hier, gelet op de vormloosheid die het verhevene kan kenmerken, met de hoedanigheid als het eerste moment van het esthetische oordeel over het verhevene zullen beginnen; de reden hiervoor is echter te vinden in de voorgaande §.

Maar één verdeling is noodzakelijk bij de analyse van het verhevene die bij het schone niet vereist was, namelijk die in het \textbf{wiskundig} verhevene en het \textbf{dynamisch} verhevene. 

Want aangezien het gevoel van het verhevene als kenmerkend teken een \textbf{beweging} van de geest met zich meebrengt die verbonden is met het oordelen over het voorwerp, terwijl de smaak voor het schone de geest in \textbf{kalme} beschouwing veronderstelt en behoudt, moet deze beweging toch als subjectief doelmatig worden beoordeeld (omdat het verhevene bevalt); deze beweging staat dus via de verbeelding of in relatie tot de \textbf{kennisvermogens} of tot de \textbf{begeertevermogens}, maar in beide relaties wordt de doelmatigheid van de gegeven voorstelling slechts beoordeeld met betrekking tot dit \textbf{vermogen} (zonder doel of belang): in het eerste geval wordt deze aan het voorwerp toegeschreven als een \textbf{wiskundige}, in het tweede als een \textbf{dynamische} gesteldheid van de verbeelding, en zo wordt het voorwerp op de bedoelde tweevoudige wijze als verheven voorgesteld.

\begin{mdframed}[leftmargin=0pt,rightmargin=0pt,backgroundcolor=blue!6,linecolor=white,innertopmargin=6pt,innerbottommargin=6pt]

Kant legt in deze paragraaf uit hoe het onderzoek naar het \textbf{verhevene} (sublieme) is opgebouwd. Net als het oordeel over het \textbf{schone} (mooie) moet het verhevene universeel geldig, belangloos, subjectief doelmatig en noodzakelijk worden voorgesteld. Een cruciaal verschil is dat de analyse van het verhevene begint met \textbf{hoeveelheid} (kwantiteit), vanwege de vormloosheid die het verhevene vaak kenmerkt. Het verhevene wordt onderverdeeld in twee soorten:
\begin{enumerate}
    \item \textbf{Wiskundig verhevene}: gerelateerd aan kennisvermogens, zoals grootsheid of oneindigheid.
    \item \textbf{Dynamisch verhevene}: gerelateerd aan begeertevermogens, zoals kracht of overweldiging.
\end{enumerate}

In tegenstelling tot het schone, dat kalme beschouwing uitlokt, roept het verhevene een \textbf{beweging van de geest} op. Deze beweging wordt als doelmatig ervaren, zonder dat er sprake is van een concreet doel of belang.

\bigskip
Deze paragraaf vormt een sleutelmoment in Kants betoog over esthetische oordelen. Waar het schone de geest in harmonie brengt, zet het verhevene de geest in beweging door confrontatie met grootsheid of overweldigende kracht. Kant benadrukt hier dat het verhevene niet alleen een esthetische ervaring is, maar ook een \textbf{reflectie op de menselijke geest} en haar vermogens, zoals verbeelding en rede. Deze paragraaf bereidt de verdere analyse van het wiskundige en dynamische verhevene voor en toont hoe esthetica verbonden is met zelfreflectie.

De kern van deze paragraaf is dat het verhevene een \textbf{actieve, dynamische ervaring} is. Het confronteert ons met wat ons overstijgt, of het nu de oneindigheid van het heelal (wiskundig) of de kracht van de natuur (dynamisch) is. Deze ervaring is niet passief, maar vereist een actieve betrokkenheid van de geest, die zich bewust wordt van haar eigen grenzen én haar vermogen om betekenis te geven aan het onbegrijpelijke. Het verhevene is daarmee niet alleen een esthetische categorie, maar ook een \textbf{filosofische uitdaging}: hoe vinden we betekenis in wat ons te boven gaat?

Wat deze paragraaf bijzonder maakt, is Kants inzicht dat esthetische ervaringen niet alleen over waarneming gaan, maar ook over \textbf{zelfreflectie}. Het verhevene confronteert ons met onze kleinheid, maar tegelijkertijd met ons vermogen om die te overstijgen. Dit roept fundamentele vragen op over de rol van de mens in een wereld die vaak overweldigend lijkt.

Kants analyse van het verhevene is nog steeds relevant, vooral in een tijd van complexe uitdagingen zoals klimaatverandering, technologische vooruitgang of informatie-overload. Het verhevene herinnert ons eraan dat ervaringen die ons overweldigen niet per se negatief zijn, maar juist kunnen leiden tot een dieper begrip van onszelf en onze plaats in de wereld. Kants idee dat het verhevene een \textbf{beweging van de geest} teweegbrengt, moedigt ons aan om actief betekenis te zoeken in wat ons te boven gaat.

Bovendien laat deze paragraaf zien dat esthetische ervaringen niet alleen over schoonheid gaan, maar ook over \textbf{kracht, grootsheid en menselijkheid}. Het is een uitnodiging om na te denken over hoe we omgaan met wat ons overstijgt --- of dat nu de natuur is, de wetenschap, of de complexiteit van het moderne leven.

\end{mdframed}

\subdiv{Het wiskundig sublieme}

\section{Nominale definitie van het sublieme}

We noemen \textbf{subliem} datgene wat \textbf{absoluut groot} is. Maar ‘groot’ en ‘grootte’ zijn heel verschillende begrippen (\textit{magnitudo} en \textit{quantitas}). Evenzo is het \textbf{simpelweg zeggen} dat iets ‘groot’ is iets heel anders dan zeggen dat het \textbf{absoluut groot} is (\textit{absolute, non comparative magnum}). Dat laatste is wat \textbf{zo groot is dat het elke vergelijking te boven gaat}.

Wat betekent het eigenlijk als we zeggen dat iets ‘groot’, ‘klein’ of ‘gemiddeld’ is? Het is geen zuiver begrip van het verstand, en ook geen zintuiglijke waarneming, en evenmin een begrip van de rede, omdat het geen enkel kennisprincipe in zich draagt. Het moet dus een begrip van de oordeelskracht zijn, of daaruit voortkomen, en geworteld zijn in een subjectieve doelmatigheid van de voorstelling ten opzichte van de oordeelskracht.
Dat iets een grootte (\textit{quantum}) is, kunnen we vaststellen aan de hand van het ding zelf, zonder vergelijking met iets anders, als tenminste een veelheid van homogene elementen samen een eenheid vormen. Maar \textbf{hoe groot} iets is, vereist altijd iets anders, wat eveneens een grootte is, als maatstaf. Omdat bij het beoordelen van grootte niet alleen de veelheid (het aantal) maar ook de grootte van de eenheid (van de maatstaf) een rol speelt, en de grootte van die eenheid op zijn beurt weer een andere maatstaf nodig heeft om mee vergeleken te worden, zien we dat elke bepaling van de grootte van verschijnselen absoluut ongeschikt is om een absoluut begrip van grootte te geven, ze kan hoogstens een vergelijkend begrip opleveren.

Als ik nu simpelweg zeg dat iets ‘groot’ is, lijkt het alsof ik helemaal geen vergelijking in gedachten heb, in elk geval niet met een objectieve maatstaf, want er wordt niet bepaald hoe groot het object is. Toch eist het oordeel, zelfs als de vergelijkingsstandaard louter subjectief is, universele instemming: de uitspraken ‘Deze man is mooi’ en ‘Hij is groot’ beperken zich niet tot het oordelend subject, maar eisen, net als theoretische oordelen, de instemming van iedereen op.

Maar als ik in een oordeel iets simpelweg ‘groot’ noem, zeg ik niet alleen dat het object een grootte heeft, maar ik ken het die grootte ook toe in hogere mate dan aan vele andere van dezelfde soort, zonder dat die superioriteit precies wordt bepaald. Dit oordeel berust dus op een standaard die we veronderstellen dat die voor iedereen hetzelfde is, maar die niet bruikbaar is voor een logische (wiskundig bepaalde) beoordeling van grootte, maar alleen voor een esthetische, omdat het een louter subjectieve standaard is die het reflecterende oordeel over grootte grondvest.
Deze standaard kan overigens empirisch zijn, zoals bij de gemiddelde grootte van de mensen die we kennen, van dieren van een bepaalde soort, van bomen, huizen, bergen, enzovoort. Of het kan een \textit{a priori} gegeven standaard zijn, die door de beperkingen van het oordelend subject wordt teruggebracht tot subjectieve voorwaarden van een voorstelling \textit{in concreto}, zoals in de praktische sfeer de grootte van een bepaalde deugd, of van openbare vrijheid en rechtvaardigheid in een land; of in de theoretische sfeer de grootte van de nauwkeurigheid of onnauwkeurigheid van een waarneming of meting, en dergelijke.

Opmerkelijk is hier dat zelfs als we geen enkel belang hebben in het object, dat wil zeggen, als het ons onverschillig laat of het bestaat, de loutere grootte ervan, zelfs als we die als vormloos beschouwen, een bevrediging kan opleveren die universeel meedeelbaar is. Dit betekent dat ze een bewustzijn kan bevatten van een subjectieve doelmatigheid in het gebruik van onze kennisfaculteiten, niet een bevrediging in het object, zoals bij het schone (want het kan vormloos zijn), waar de reflecterende oordeelskracht zich doelmatig voelt ten opzichte van kennis in het algemeen, maar eerder in de vergroting van de verbeeldingskracht op zichzelf.

Als we (onder de hierboven genoemde beperking) van een object zeggen dat het absoluut groot is, is dat geen wiskundig-bepalend oordeel, maar een louter reflecterend oordeel over zijn voorstelling, dat subjectief doelmatig is voor een bepaald gebruik van onze kenniskrachten bij het schatten van grootte. In dat geval verbinden we altijd een zekere eerbied met de voorstelling, net zoals we minachting verbinden met wat we ‘absoluut klein’ noemen. Bovendien geldt het oordelen van dingen als groot of klein voor alles, zelfs voor al hun eigenschappen; zo noemen we ook schoonheid ‘groot’ of ‘klein’. De reden hiervoor moet gezocht worden in het feit dat alles wat we in de aanschouwing kunnen presenteren volgens het voorschrift van de oordeelskracht (en dus esthetisch voorstellen) louter verschijning is, en dus ook een quantum (grootte).

Maar als we iets niet alleen ‘groot’ noemen, maar eenvoudigweg, absoluut groot, groot in alle opzichten (boven elke vergelijking uit), dat wil zeggen subliem, dan zien we meteen dat we geen passende maatstaf daarbuiten zoeken, maar alleen binnenin het object zelf. Het is een grootte die alleen aan zichzelf gelijk is. Dat het sublieme dus niet in de dingen van de natuur gezocht moet worden, maar alleen in onze ideeën, volgt hieruit, maar waar in die ideeën het ligt, moet worden bewaard voor de deductie.

De bovenstaande uitleg kan ook als volgt worden geformuleerd: \textbf{Subliem is datgene waarbij alles anders klein is in vergelijking}. Hieruit blijkt meteen dat er in de natuur niets kan worden gegeven, hoe groot we het ook mogen vinden, dat niet, vanuit een ander perspectief, tot het oneindig kleine kan worden gereduceerd, en omgekeerd, er is niets zo klein dat niet, vergeleken met nog kleinere maatstaven, voor onze verbeelding tot de grootte van een wereld kan worden vergroot. De telescoop heeft ons rijk materiaal gegeven voor de eerste observatie, de microscoop voor de laatste. Niets wat een object van de zintuigen kan zijn, is dus, vanuit dit oogpunt, subliem te noemen.
Maar juist omdat er in onze verbeelding een streven ligt om tot het oneindige vooruit te gaan, terwijl in onze rede een claim ligt op absolute totaliteit als een echte idee, wekt de onvoldoende capaciteit van ons vermogen om de grootte van de dingen in de zintuiglijke wereld te schatten juist het gevoel van een bovenzintuiglijke faculteit in ons. En het is het gebruik dat de oordeelskracht van nature maakt ten behoeve van die faculteit (het gevoel), niet het object van de zintuigen, dat absoluut groot is, terwijl elke andere toepassing in vergelijking daarme klein is. Het is dus de gemoedsgesteldheid die voortkomt uit een bepaalde voorstelling die de reflecterende oordeelskracht in beslag neemt, niet het object zelf, die subliem moet worden genoemd.

We kunnen aan de voorgaande formulering van de uitleg van het sublieme dus ook het volgende toevoegen: \textbf{Subliem is datgene waarover zelfs maar nadenken een vermogen van de geest aantoont dat elke maatstaf van de zintuigen te boven gaat.}

\begin{mdframed}[leftmargin=0pt,rightmargin=0pt,backgroundcolor=blue!6,linecolor=white,innertopmargin=6pt,innerbottommargin=6pt]

Kant definieert het sublieme als iets dat absoluut groot is, zo groot dat het elke vergelijking te boven gaat. Het gaat niet om een objectieve grootte, maar om de ervaring dat iets onze verbeeldingskracht overweldigt en dat we geen passende maatstaf kunnen vinden om het te meten. Terwijl we bij gewone grootte altijd een referentiepunt nodig hebben, zoals de gemiddelde lengte van een mens of de hoogte van een boom, ontbreekt zo’n referentie bij het sublieme. Het is geen kwestie van meetbare grootheid, maar van een gevoel dat iets onze voorstellingskracht overstijgt.

Het sublieme ligt niet in de dingen zelf, zoals een berg of een storm, maar in onze reactie erop. Niets in de natuur is op zichzelf subliem, want alles kan vanuit een ander perspectief klein lijken; een berg is immens voor ons, maar nietig vergeleken met de aarde. Toch roept het sublieme een universeel gevoel van eerbied op, niet omdat het object indrukwekkend is, maar omdat het ons bewust maakt van de beperkingen van onze zintuigen en tegelijkertijd van het vermogen van onze geest om daaroverheen te denken. Het is geen eigenschap van het waargenomene, maar van de wijze waarop ons oordeel ermee omgaat.

Kant benadrukt dat het sublieme ons confronteert met het oneindige en ons laat voelen dat onze rede verder reikt dan wat we kunnen waarnemen. Het is geen esthetische kwaliteit zoals schoonheid, maar een ervaring die ons subjectieve vermogen om het onmetelijke te vatten aantoont, en juist daardoor een diep gevoel van verhevenheid oproept.

\bigskip
De paragraaf markeert een cruciaal onderscheid tussen het schone en het sublime, en toont hoe het sublieme een brug slaat tussen onze zintuiglijke ervaring en ons morele bewustzijn. Het sublieme is geen kwestie van smaak of harmonie, zoals schoonheid, maar van een conflict: onze verbeeldingskracht wordt overweldigd door iets dat zo groot, krachtig of oneindig is dat we het niet kunnen bevatten (zoals een storm, een afgrond, of de sterrenhemel). Dat overweldigende gevoel, een mix van angst, ontzag en zelfs afschuw, wordt echter niet als louter negatief ervaren. Integendeel, het roept uiteindelijk een gevoel van trots of verhevenheid op, omdat we beseffen dat onze rede (ons vermogen om het oneindige te denken) boven de beperkingen van onze zintuigen uitstijgt. Het sublieme is dus geen aangename ervaring, maar een uitdagende: het confronteert ons met onze kwetsbaarheid als zintuiglijke wezens, terwijl het tegelijkertijd onze vrijheid als redelijke wezens bevestigt. Dit past perfect in Kants grotere project, waarin hij laat zien hoe de mens niet alleen een natuurlijk, maar ook een moreel wezen is. Het sublieme is als het ware een symbool van onze morele bestemming: net zoals we het oneindige niet kunnen zien maar wel kunnen denken, zo kunnen we morele idealen (zoals absolute rechtvaardigheid of deugd) niet volledig verwezenlijken, maar wel als richtinggevend erkennen.

De aantrekking tot angstwekkende dingen, zoals horrorfilms, extreme hoogtes, of zelfs existentiële vragen over de eindigheid van het leven, kan indien als een moderne echo van het sublieme worden gezien. We zoeken deze ervaringen op omdat ze ons uit onze comfortzone halen en ons even doen voelen hoe klein en kwetsbaar we zijn, maar ook hoe uniek ons vermogen is om betekenis te geven aan wat ons overstijgt. Een horrorfilm confronteert ons met het oncontroleerbare, een afgrond met het niets, en de ruimte met onze nietigheid, maar juist daarin ligt de kick: we weten dat we veilig zijn (of dat we de angst kunnen stoppen), en dat geeft ons een gevoel van overwinning over wat ons bedreigt. Het is alsof we, net als bij Kant, genieten van het besef dat we de angst kunnen overstijgen met onze verbeelding of rede. Het is geen loutere vlucht in sensatie, maar een oefening in transcendentie, een manier om onze eigen grenzen te verkennen en te bevestigen dat we meer zijn dan louter zintuiglijke wezens.

In het algehele betoog van de Kritik der Urteilskraft fungeert het sublieme als een scharnier tussen de theorie van kennis (wat we weten) en de moraal (wat we behoren te doen). Schoonheid toont hoe vorm en verstand samen kunnen vallen in een harmonische ervaring; het sublieme daentegen onthult de kloof tussen onze zintuigen en onze rede, en nodigt ons uit die kloof te overbruggen. Het is geen toeval dat Kant later in het boek de verbinding legt tussen het sublieme en morele ideeën: beide gaan over het streven naar het oneindige, het absolute. Het sublieme herinnert ons eraan dat we niet alleen deel uitmaken van de natuur, maar ook van een morele wereldorde die we niet kunnen zien, maar wel kunnen denken. Zo wordt een dreigende storm niet alleen een spectaculair natuurverschijnsel, maar ook een metafoor voor de uitdagingen die ons morele leven kenmerken, en voor het vermogen om die aan te gaan.

\end{mdframed}

\section{Over de schatting van de grootte van natuurlijke dingen die vereist is voor de idee van het sublieme.}

De schatting van grootte door middel van numerieke begrippen (of hun symbolen in de algebra) is wiskundig, maar die in louter aanschouwing (met het oog gemeten) is esthetisch. Nu kunnen we weliswaar bepaalde begrippen verkrijgen van \textbf{hoe groot} iets is alleen door middel van getallen (of in elk geval door benaderingen door middel van numerieke reeksen die tot in het oneindige voortschrijden), waarvan de eenheid de maat is; en in die mate is alle logische schatting van grootte wiskundig. Maar aangezien de grootte van de maat nog steeds als bekend moet worden aangenomen, dan, als deze op zijn beurt alleen door middel van getallen moet worden geschat waarvan de eenheid weer een andere maat zou moeten zijn, en dus wiskundig, kunnen we nooit een oorspronkelijke of fundamentele basismaat hebben, en dus ook nooit een bepaald begrip van een gegeven grootte. Zo bestaat de schatting van de grootte van de basismaat er simpelweg in dat men deze onmiddellijk in een aanschouwing kan vatten en door middel van de verbeeldingskracht kan gebruiken voor de voorstelling van numerieke begrippen, dat wil zeggen, uiteindelijk is alle schatting van de grootte van natuurlijke objecten esthetisch (dat wil zeggen, subjectief en niet objectief bepaald).

Voor de wiskundige schatting van grootte is er zeker geen grootste (want de kracht van getallen gaat tot in het oneindige); maar voor de esthetische schatting van grootte is er zeker wel een grootste; en hierover zeg ik dat, als het als een absolute maat wordt beoordeeld, waarboven geen grotere subjectief (voor het oordelende subject) mogelijk is, het de idee van het sublieme met zich meebrengt en die emotie voortbrengt die geen wiskundige schatting van grootheden door middel van getallen kan voortbrengen (behalve voor zover die esthetische basismaat levendig in de verbeelding bewaard blijft), aangezien die laatste altijd alleen maar relatieve grootte door vergelijking met andere van dezelfde soort voorstelt, maar de eerste grootte absoluut voorstelt, voor zover de geest het in één aanschouwing kan vatten.

Om een quantum in de verbeelding aanschouwelijk op te nemen, om het als maat of eenheid te kunnen gebruiken voor de schatting van grootte door middel van getallen, omvat twee handelingen van deze faculteit: \textbf{apprehensie} (\textit{apprehensio}) en \textbf{esthetische begripvorming} (\textit{comprehensio aesthetica}). Er is geen moeite met apprehensie, want die kan tot in het oneindige doorgaan; maar begripvorming wordt steeds moeilijker naarmate de apprehensie vordert, en bereikt spoedig zijn maximum, namelijk de esthetisch grootste basismaat voor de schatting van grootte. Want wanneer de apprehensie zo ver is gevorderd dat de gedeeltelijke voorstellingen van de zintuiglijke aanschouwing die eerst waren opgevat al beginnen te vervagen in de verbeelding terwijl deze doorgaat met de apprehensie van verdere, dan verliest zij aan de ene kant evenveel als zij aan de andere wint, en er is in het begrip een grootste punt waarboven zij niet kan gaan.

Dit maakt het mogelijk om een punt te verklaren dat Savary opmerkt in zijn verslag over Egypte: dat men om het volle emotionele effect van de grootte van de piramides te krijgen, er niet te dichtbij mag komen en ook niet te veraf mag staan. Want in het laatste geval worden de delen die worden opgevat (de op elkaar gestapelde stenen) alleen maar vaag voorgesteld, en hun voorstelling heeft geen effect op het esthetische oordeel van het subject. In het eerste geval echter vereist het oog enige tijd om zijn apprehensie van het grondvlak tot de top te voltooien, maar tijdens deze tijd vervaagt het eerste gedeelte altijd deels voordat de verbeelding het laatste heeft opgenomen, en het begrip is nooit compleet. --
Precies hetzelfde kan ook dienen om de verwarring of een soort verlegenheid te verklaren die, naar men zegt, de toeschouwer bevangt bij het betreden van de Sint-Pieter in Rome. Want hier is er een gevoel van ontoereikendheid van zijn verbeelding om de ideeën van een geheel voor te stellen, waarin de verbeelding haar maximum bereikt en, in de poging om zich uit te breiden, terugzakt in zichzelf, maar daardoor in een emotioneel ontroerende bevrediging wordt gebracht.

Ik zal nog niets toevoegen over de grond voor deze bevrediging, die verbonden is met een voorstelling waarvan men het minst zou verwachten, namelijk een voorstelling die ons de ontoereikendheid, en dus ook de subjectieve ondoelmatigheid van de voorstelling voor de oordeelskracht bij de schatting van grootte doet opmerken; ik merk alleen op dat, als het esthetische oordeel \textbf{zuiver} moet zijn (\textbf{niet vermengd met iets teleologisch} als oordelen van de rede) en als er een voorbeeld van moet worden gegeven dat volledig geschikt is voor de kritiek van de \textbf{esthetische} oordeelskracht, dan moet het sublieme niet worden getoond in producten van kunst (bijvoorbeeld gebouwen, zuilen, enz.), waar een menselijk doel de vorm en de grootte bepaalt, noch in natuurlijke dingen \textbf{waarvan het begrip al een bepaald doel met zich meebrengt} (bijvoorbeeld dieren met een bekende natuurlijke bestemming), maar eerder in de ruwe natuur (en zelfs daarin alleen voor zover deze op zichzelf noch bekoring noch emotie van echte gevaar met zich meebrengt), louter voor zover deze grootte bevat.
Want in deze soort voorstelling bevat de natuur niets dat \textbf{monsterlijk} (of prachtig of verschrikkelijk) zou zijn; de grootte die wordt opgevat kan zo groot worden als men wil, zolang deze in één geheel door de verbeelding kan worden begrepen. Een object is monsterlijk als het door zijn grootte het doel vernietigt dat zijn begrip vormt.
De loutere voorstelling van een begrip, echter, dat bijna te groot is voor elke voorstelling (dat grenst aan het relatief monsterlijke), wordt \textbf{kolossaal} genoemd, omdat het doel van de voorstelling van een begrip moeilijker wordt gemaakt als de aanschouwing van het object bijna te groot is voor ons vermogen tot apprehensie.,- Een zuiver oordeel over het sublieme, echter, mag geen doel van het object als zijn bepalende grond hebben als het esthetisch moet zijn en niet vermengd met enig oordeel van het verstand of de rede.

\bigskip
Omdat alles wat louter de reflecterende oordeelskracht zonder belangstelling moet behagen, in zijn voorstelling subjectieve en als zodanig universeel geldige doelmatigheid moet bevatten, hoewel hier geen doelmatigheid van de \textbf{vorm} van het object (zoals in het geval van het schone) de grond is voor het oordelen, rijst de vraag: wat is deze subjectieve doelmatigheid? en hoe wordt deze voorgeschreven als een norm die een grond biedt voor universeel geldige bevrediging in de loutere schatting van grootte, en wel waar deze bijna tot het punt van de ontoereikendheid van ons vermogen tot verbeelding in de voorstelling van het begrip van een grootte is gedreven?

De verbeelding, op zichzelf, zonder dat iets haar hindert, schrijdt voort tot in het oneindige in de samenstelling die vereist is voor de voorstelling van grootte; het verstand, echter, leidt deze door middel van numerieke begrippen, waarvoor de eerste het schema moet leveren; en in deze procedure, die behoort tot de logische schatting van grootte, is er zeker iets objectief doelmatigs in overeenstemming met het begrip van een doel (zoals alle meting is), maar niets dat doelmatig en bevallig is voor de esthetische oordeelskracht. Er is ook in deze intentionele doelmatigheid niets dat zou vereisen dat de grootte van de maat en dus het \textbf{begrip} van de vele in één aanschouwing tot de grenzen van het vermogen tot verbeelding wordt gedreven en zo ver als deze in voorstellingen zou kunnen reiken. Want in de schatting van grootheden door het verstand (in de rekenkunde) komt men even ver, of men nu de samenstelling van de eenheden tot het getal 10 (in het decimale stelsel) of slechts tot 4 (in het tetradische stelsel) voortzet; de verdere voortzetting van de grootte in samenstelling, of, als het quantum in de aanschouwing gegeven is, in apprehensie, verloopt slechts progressief (niet omvattend) volgens een aangenomen principe van voortgang. In deze wiskundige schatting van grootte wordt het verstand even goed gediend en bevredigd, of de verbeelding nu voor haar eenheid een grootte kiest die in één oogopslag kan worden begrepen, bijvoorbeeld een voet of een roede, of dat zij een Duitse mijl of zelfs een aarddiameter kiest, waarvan de apprehensie weliswaar, maar niet de samenstelling in een aanschouwing van de verbeelding mogelijk is (niet door \textit{comprehensio aesthetica}, hoewel zeker door \textit{comprehensio logica} in een numeriek begrip). In beide gevallen verloopt de logische schatting van grootte ongehinderd tot in het oneindige.

Maar nu hoort de geest in zichzelf de stem van de rede, die totaliteit eist voor alle gegeven grootheden, zelfs voor die welke nooit geheel kunnen worden opgevat, hoewel zij (in de zintuiglijke voorstelling) als geheel gegeven worden beoordeeld, dus begrip in \textbf{één} aanschouwing, en het eist een \textbf{voorstelling} voor alle leden van een progressief toenemende numerieke reeks, en onttrekt zich niet aan deze eis, zelfs niet voor het oneindige (ruimte en verleden tijd), maar maakt het juist onvermijdelijk voor ons om het (in het oordeel van de gewone rede) als \textbf{geheel gegeven} te denken (in zijn totaliteit).

Het oneindige, echter, is absoluut (niet slechts relatief) groot. Vergeleken hiermee is alles anders (van dezelfde soort grootte) klein. Maar wat het belangrijkste is, is dat zelfs het kunnen denken ervan als \textbf{een geheel} wijst op een vermogen van de geest dat elke maatstaf van de zintuigen te boven gaat. Want dit zou een begrip vereisen dat als maat een eenheid zou opleveren die een bepaalde relatie tot het oneindige heeft, uitdrukbaar in getallen, wat onmogelijk is. Maar \textbf{zelfs het kunnen denken} van het gegeven oneindige zonder tegenspraak vereist een vermogen in de menselijke geest dat op zichzelf bovenzintuiglijk is. Want het is alleen door middel hiervan en zijn idee van een noumenon, dat op zichzelf geen aanschouwing toelaat, hoewel het wordt verondersteld als het substraat van de aanschouwing van de wereld als louter verschijning, dat het oneindige van de zintuiglijke wereld \textbf{volledig} wordt begrepen in de zuivere intellectuele schatting van grootte \textbf{onder} een begrip, hoewel het nooit volledig kan worden gedacht in de wiskundige schatting van grootte \textbf{door middel van numerieke begrippen}. Zelfs een vermogen om het oneindige van bovenzintuiglijke aanschouwing als gegeven te kunnen denken (in zijn intelligibele substraat) gaat elke maatstaf van de zintuiglijkheid te boven, en is groot boven elke vergelijking, zelfs met het vermogen van wiskundige schatting, niet natuurlijk vanuit een theoretisch oogpunt, ten behoeve van het kennisfaculteit, maar toch als een vergroting van de geest die zich bekwaam voelt om de grenzen van de zintuiglijkheid vanuit een ander (praktisch) oogpunt te overschrijden.

De natuur is dus subliem in die van haar verschijnselen waarvan de aanschouwing de idee van haar oneindigheid met zich meebrengt. Nu kan dit laatste niet gebeuren behalve door de ontoereikendheid van zelfs de grootste inspanning van onze verbeelding bij de schatting van de grootte van een object. Nu echter is de verbeelding weliswaar geschikt voor de wiskundige schatting van elk object, dat wil zeggen, om een adequate maat ervoor te geven, omdat de numerieke begrippen van het verstand, door middel van progressie, elke maat adequaat kunnen maken voor elke gegeven grootte. Zo moet het de \textbf{esthetische} schatting van grootte zijn waarin de inspanning tot begrip wordt gevoeld die het vermogen van de verbeelding te boven gaat om de progressieve apprehensie in één geheel van aanschouwing te begrijpen, en waarin tegelijkertijd de ontoereikendheid van deze faculteit, die onbegrensd is in haar voortgang, wordt waargenomen om een basismaat te grijpen die geschikt is voor de schatting van grootte met de minste inspanning van het verstand en om deze te gebruiken voor de schatting van grootte. Nu is de eigenlijke onveranderlijke basismaat van de natuur haar absolute geheel, dat, in het geval van de natuur als verschijning, oneindigheid omvat. Maar aangezien deze basismaat een zelf-tegensprekend begrip is (wegens de onmogelijkheid van de absolute totaliteit van een eindeloze progressie), moet die grootte van een natuurlijk object waarop de verbeelding vruchteloos al haar vermogen tot begrip uitgeeft, het begrip van de natuur leiden tot een bovenzintuiglijk substraat (dat zowel deze als tegelijkertijd ons vermogen om te denken grondvest), dat groot is boven elke maatstaf van de zintuigen en dus niet zozeer het object als wel eerder de gezindheid van de geest bij het beoordelen ervan \textbf{subliem} doet oordelen.

Zoals de esthetische oordeelskracht bij het oordelen over het schone de verbeelding in haar vrije spel relateert aan het \textbf{verstand}, om overeen te stemmen met zijn \textbf{begrippen} in het algemeen (zonder bepaling daarvan), zo wordt in het oordelen van een ding als subliem dezelfde faculteit gerelateerd aan de \textbf{rede}, om subjectief overeen te stemmen met haar \textbf{ideeën} (hoewel onbepaald), dat wil zeggen, om een gezindheid van de geest voort te brengen die in overeenstemming is met deze ideeën en verenigbaar met wat de invloed van bepaalde (praktische) ideeën op het gevoel zou voortbrengen. Het is ook duidelijk hieruit dat de ware sublimiteit alleen moet worden gezocht in de geest van degene die oordeelt, niet in het object in de natuur, waarvan het oordelen deze gezindheid daarin teweegbrengt. En wie zou wilden noemen subliem de vormloze bergmassa’s die in wilde wanorde boven elkaar uittorenen met hun ijs piramides, of de donkere en woeste zee, enzovoort? Maar de geest voelt zich verheven in zijn eigen oordelen als het, bij de beschouwing van dergelijke dingen, zonder acht te slaan op hun vorm, zich overgeeft aan de verbeelding en aan een rede die, hoewel zij geheel zonder enig bepaald doel met haar verbonden is, haar slechts uitbreidt, en niettemin het gehele vermogen van de verbeelding ontoereikend vindt voor haar ideeën.

Voorbeelden van het wiskundig sublieme in de natuur in louter aanschouwing worden ons geboden door alle die gevallen waarin ons niet zozeer een groter numeriek begrip wordt gegeven als wel een grote eenheid als maat (om de numerieke reeks in te korten) voor de verbeelding. Een boom die we schatten naar de lengte van een mens kan dienen als standaard voor een berg, en als deze laatste, laten we zeggen, een mijl hoog is, kan deze dienen als eenheid voor het getal dat de diameter van de aarde uitdrukt, om deze aanschouwelijk te maken; de diameter van de aarde zou kunnen dienen als eenheid voor het planetenstelsel zoals wij dat kennen, dit voor de Melkweg, en de onmeetbare menigte van dergelijke Melkwegstelsels, nebulae genoemd, die waarschijnlijk onderling weer zo’n systeem vormen, laat hier geen grenzen verwachten. Nu ligt in het esthetische oordelen van zo’n onmeetbaar geheel de sublimiteit niet zozeer in de grootte van het getal als wel in het feit dat we, naarmate we vorderen, steeds grotere eenheden tegenkomen; de systematische indeling van de structuur van de wereld draagt hiertoe bij, door ons alles wat groot is in de natuur op zijn beurt als klein voor te stellen, maar eigenlijk onze verbeelding in al haar grenzeloosheid, en daarmee de natuur, als verblekend tot nietigheid naast de ideeën van de rede als zij een voorstelling zou moeten geven die daaraan adequaat is.

\begin{mdframed}[leftmargin=0pt,rightmargin=0pt,backgroundcolor=blue!6,linecolor=white,innertopmargin=6pt,innerbottommargin=6pt]

Kant legt uit hoe we de grootte van natuurlijke dingen beoordelen om tot het gevoel van het sublieme te komen. Hij maakt onderscheid tussen twee manieren van grootte inschatten: wiskundig (met getallen en meetbare eenheden) en esthetisch (door directe aanschouwing, met het oog). Wiskundige schattingen zijn oneindig en objectief, maar esthetische schattingen stuiten op grenzen, namelijk op het vermogen van onze verbeeldingskracht om iets in één keer te overzien.

Kant beschrijft hoe onze verbeelding werkt bij het waarnemen van grote objecten, zoals bergen of gebouwen. Eerst vat ze delen van het object (apprehensie), maar naarmate het object groter wordt, lukt het steeds minder om al die delen samen te vatten in één geheel (esthetische begripvorming, comprehensio aesthetica). Op een gegeven moment raakt onze verbeelding overbelast: terwijl we nieuwe delen waarnemen, vervagen de eerdere delen alweer in ons bewustzijn. Dit is het moment waarop we de grootte niet meer kunnen bevatten, en dat roept het gevoel van het sublieme op. Een voorbeeld is de ervaring van de piramides: staan we te ver weg, dan zien we ze niet goed; staan we te dichtbij, dan kunnen we ze niet in één keer overzien. Ook bij het betreden van de Sint-Pieter in Rome ervaren mensen een soort verwarring, omdat de grootte van het gebouw hun verbeelding te boven gaat.

Het sublieme zit dus niet in het object zelf, maar in onze reactie daarop: het besef dat onze verbeelding tekortschiet, terwijl onze rede wel het oneindige kan denken. Dit geldt vooral voor ruwe natuur (zoals woeste bergen of een stormachtige zee), niet voor kunstmatige dingen (zoals gebouwen) of levende wezens (zoals dieren), omdat die een duidelijk doel of vorm hebben. Het sublieme ontstaat wanneer we geconfronteerd worden met iets zo enorm dat we het niet kunnen bevatten, zoals de oneindige ruimte van het heelal. Onze verbeelding probeert steeds grotere eenheden te vinden (een boom als maat voor een berg, een berg voor de aarde, de aarde voor het zonnestelsel, enzovoort), maar uiteindelijk beseffen we dat elke maatstaf te klein is. Dat gevoel van ontoereikendheid leidt tot het sublieme: niet omdat het object zelf indrukwekkend is, maar omdat het ons bewust maakt van de grenzen van onze waarneming, en tegelijkertijd van het vermogen van onze geest om daaroverheen te denken.

Kort samengevat: het sublieme is de ervaring dat iets zo groot is dat we het niet kunnen overzien, en dat ons bewust maakt van de oneindige mogelijkheden van onze rede, terwijl onze zintuigen tekortschieten. Het gaat niet om het object, maar om onze reactie daarop.

\bigskip
Kant onthult hoe het wiskundig sublieme werkt als een spiegel van onze cognitieve grenzen en mogelijkheden. Het gaat niet om de objectieve grootte van een object, maar om het moment waarop onze verbeeldingskracht faalt in haar poging om iets in één aanschouwing te vatten. Neem het voorbeeld van de Grand Canyon: als je aan de rand staat, probeert je verbeelding de afmetingen te begrijpen, de diepte, de breedte, de lagen van gesteente die zich over miljoenen jaren hebben gevormd. Je oog glijdt langs de wanden, maar terwijl je de ene kant probeert te bevatten, vervaagt de andere kant alweer in je bewustzijn. Je kunt niet tegelijkertijd de volledige diepte, de horizon, en de details van de rotsformaties overzien. Je verbeelding raakt overbelast: ze kan de canyon niet als één geheel "grijpen", zoals ze dat wel kan bij een kleiner landschap of een voorwerp dat binnen je blikveld past (zoals een kristal of een boom).

Wat dit uitzicht over de Grand Canyon wiskundig subliem maakt, is niet de absolute grootte (want vanuit een satelliet is het slechts een kras in het aardoppervlak), maar de ervaring van het falen van je verbeelding om het te bevatten. Je weet \textit{rationeel} dat de canyon meetbaar is, kilometers breed, honderden meters diep, maar je \textit{kunt} die metingen niet \textit{voelen}. Je verbeelding zoekt wanhopig naar een "maatstaf": een boom als referentie voor een rotswand, een rotswand voor de hele canyon, maar elke maatstaf die je kiest, blijkt ontoereikend. Dat gevoel van ontoereikendheid, het besef dat je de canyon niet in één keer kunt "denken" zoals je een tafel of een huis kunt overzien, is precies waar het sublieme ontstaat.

Tegelijkertijd, en dat is het cruciale punt bij Kant, overschrijdt je rede deze beperking. Terwijl je verbeelding worstelt, weet je \textit{wel} dat de canyon eindig is, dat hij gemeten \textit{kan} worden, dat hij deel uitmaakt van een groter geologisch systeem. Je kunt je zelfs het \textit{oneindige} voorstellen, niet als een beeld, maar als een \textit{idee}. Dat dubbele besef, dat je zintuigen tekortschieten, maar je rede niet, is wat het sublieme zo krachtig maakt. De Grand Canyon wordt niet subliem door zijn schoonheid (hoewel die er zeker is), maar door het conflict tussen wat je ziet en wat je \textit{weet}: je ziet een afgrond die je niet kunt bevatten, maar je begrijpt \textit{wel} dat hij bestaat, dat hij deel uitmaakt van een systeem dat je rede wel kan omvatten. Het is alsof de canyon je dwingt om je eigen kennisfaculteiten te ervaren: je voelt je klein als waarnemer, maar tegelijkertijd verheven als redelijk wezen dat het oneindige kan \textit{denken}.

Dit voorbeeld illustreert waarom het sublieme volgens Kant niet in het object zelf ligt, maar in onze interactie ermee: de Grand Canyon is niet \textit{op zich} subliem, maar wordt dat pas in de ontmoeting met een menselijke geest die probeert hem te bevatten. Het is een ervaring van transcendentie, niet in mystieke zin, maar in de zin dat je je bewust wordt van de grenzen van je zintuigen \textit{en} van het vermogen van je rede om daarvoorbij te gaan. Dat is waarom zo’n ervaring zowel ontzag als trots oproept: ontzag voor wat je niet kunt vatten, trots omdat je het toch kunt \textit{begrijpen}.

\end{mdframed}

\section{Over de aard van de bevrediging bij het oordeel over het sublieme.}

Het gevoel van ontoereikendheid van ons vermogen om een idee, \textbf{dat voor ons als wet geldt}, te bereiken, is \textbf{eerbied}. Nu is de idee van het begrijpen van elke verschijning die ons gegeven kan worden in de aanschouwing van een geheel een idee dat ons door een wet van de rede wordt opgelegd. Deze wet erkent geen andere bepaalde, voor iedereen geldige en onveranderlijke maatstaf dan het absolute geheel. Maar onze verbeeldingskracht, zelfs in haar grootste inspanning om een gegeven voorwerp in een geheel van aanschouwing (en dus voor de voorstelling van de idee van de rede) te bevatten, toont haar grenzen en ontoereikendheid, maar tegelijkertijd haar roeping om die idee als wet adequaat te verwezenlijken. Zo is het gevoel van het sublieme in de natuur eerbied voor onze eigen roeping, die we tonen aan een voorwerp in de natuur door een bepaalde subreptie (vervanging van eerbied voor het voorwerp in plaats van voor de idee van de mensheid in ons subject), die als het ware de superioriteit van de redelijke roeping van onze kennisfaculteit boven de grootste faculteit der zintuiglijkheid aanschouwelijk maakt.

Het gevoel van het sublieme is dus een gevoel van misoegen over de ontoereikendheid van de verbeelding bij de esthetische schatting van grootte voor de schatting door middel van de rede, en een genoegen dat tegelijkertijd wordt opgewekt door de overeenstemming van dit oordeel over de ontoereikendheid van de grootste zintuiglijke faculteit in vergelijking met de ideeën van de rede, voor zover het streven naar die ideeën toch een wet voor ons is. Dat wil zeggen, het is een wet (van de rede) voor ons en deel van onze roeping om alles wat groot is in de natuur, als een voorwerp der zintuigen, voor ons als klein te schatten in vergelijking met de ideeën van de rede; en alles wat het gevoel van deze bovenzintuiglijke roeping in ons opwekt, is in overeenstemming met die wet. Nu is de grootste inspanning van de verbeelding bij de voorstelling van de eenheid voor de schatting van grootte een relatie tot iets \textbf{absoluut groots}, en bijgevolg ook een relatie tot de wet van de rede om dit alleen als de opperste maatstaf van grootte aan te nemen. Zo komt de innerlijke waarneming van de ontoereikendheid van elke zintuiglijke maatstaf voor de schatting van grootte door de rede overeen met de wetten van de rede, en is een misnoegen dat het gevoel van onze bovenzintuiglijke roeping in ons opwekt, overeenkomstig waarmee het doelmatig en dus een genoegen is om elke maatstaf der zintuiglijkheid ontoereikend te vinden voor de ideeën van het verstand.

De geest voelt zich \textbf{bewogen} bij de voorstelling van het sublieme in de natuur, terwijl hij bij het esthetische oordeel over het schone in de natuur in \textbf{kalme} beschouwing verkeert. Deze beweging (met name in haar begin) kan vergeleken worden met een trilling, dat wil zeggen, een snel wisselende afstoting van en aantrekking tot één en hetzelfde voorwerp. Wat te veel is voor de verbeelding (waarnaar zij gedreven wordt bij de apprehensie van de aanschouwing) is als het ware een afgrond, waarin zij vreest zichzelf te verliezen, maar voor de idee van de rede van het bovenzintuiglijke is een dergelijke inspanning van de verbeelding niet te veel, maar wetmatig, en is dus precies zo aantrekkelijk als het afstotend was voor louter zintuiglijkheid. Ook in dit geval blijft het oordeel zelf echter slechts esthetisch, omdat het, zonder een bepaald begrip van het voorwerp als zijn grond, slechts het subjectieve spel der geesteskrachten (verbeelding en rede) voorstelt als harmonieus, zelfs in hun contrast. Want zoals verbeelding en \textbf{verstand} in het oordeel over het schone door hun eensgezindheid een subjectieve doelmatigheid der geesteskrachten voortbrengen, zo brengen verbeelding en \textbf{rede} in hun conflict een subjectieve doelmatigheid voort: namelijk een gevoel dat wij zuivere, zelfgenoegzame rede bezitten, of een faculteit om grootte te schatten, waarvan de superioriteit niet door iets anders intuitable gemaakt kan worden dan door de ontoereikendheid van die faculteit die zelf onbegrensd is in de voorstelling van grootheden (van zintuiglijke voorwerpen).

Het meten van een ruimte (als apprehensie) is tegelijkertijd het beschrijven ervan, dus een objectieve beweging in de verbeelding en een voortgang; daartegenover staat het begrijpen van veelheid in de eenheid niet van het denken, maar van de aanschouwing, dus het begrijpen in één moment van datgene wat opeenvolgend wordt opgevat, een teruggang, die op haar beurt de tijdvoorwaarde in de voortgang van de verbeelding opheft en \textbf{gelijktijdigheid} aanschouwelijk maakt. Het is dus (aangezien tijdelijke opeenvolging een voorwaarde is van het innerlijk zintuig en van een aanschouwing) een subjectieve beweging van de verbeelding, waardoor zij geweld aandoet aan het innerlijk zintuig, wat des te opvallender moet zijn naarmate het quantum groter is dat de verbeelding in één aanschouwing begrijpt. Zo is de inspanning om in één aanschouwing een maat voor grootheden op te nemen, die een waarneembare tijd voor haar apprehensie vereist, een soort apprehensie die, subjectief beschouwd, tegendelmatig is, maar die objectief, voor de schatting van grootte, noodzakelijk en dus doelmatig is; op deze wijze wordt echter precies hetzelfde geweld dat de verbeelding het subject aandoet beoordeeld als doelmatig \textbf{voor de gehele roeping} van de geest.

De \textbf{kwaliteit} van het gevoel van het sublieme is dat het een gevoel van misoegen is met betrekking tot de esthetische faculteit van het oordelen over een voorwerp, dat toch tegelijkertijd als doelmatig wordt voorgesteld, hetgeen mogelijk is omdat de eigen onmacht van het subject het bewustzijn onthult van een onbegrensde bekwaamheid van hetzelfde subject, en de geest die laatste slechts esthetisch kan beoordelen door middel van de eerste.

In de logische schatting van grootte werd de onmogelijkheid om ooit absolute totaliteit te bereiken door de voortgang van de meting der dingen van de zintuiglijke wereld in tijd en ruimte als objectief erkend, dat wil zeggen als een onmogelijkheid om het oneindige zelfs maar als gegeven te \textbf{denken}, en niet slechts als subjectief, dat wil zeggen als een onbekwaamheid om het te \textbf{bevatten}; want daar komt het geheel niet aan op de graad van begrip in één aanschouwing als maatstaf, maar alles komt neer op een numeriek begrip. Maar in een esthetische schatting van grootte moet het numerieke begrip wegvallen of gewijzigd worden, en het begrip van de verbeelding met betrekking tot de eenheid van de maatstaf (zodat het begrip van een wet van de opeenvolgende voortbrenging van groottebegrippen wordt vermeden) is alleen voor haar doelmatig. Nu, als een grootte bijna de uiterste grens van ons vermogen tot begrip in één aanschouwing bereikt, en toch de verbeelding door middel van numerieke begrippen (waarvoor wij ons bewust zijn dat ons vermogen onbegrensd is) wordt opgeroepen tot esthetisch begrip in een grotere eenheid, dan voelen wij ons in onze geest als esthetisch binnen grenzen opgesloten; maar met betrekking tot de noodzakelijke uitbreiding van de verbeelding tot het punt van adequaatheid aan datgene wat onbegrensd is in ons vermogen van de rede, namelijk de idee van het absolute geheel, wordt het misnoegen en dus de tegendelmatigheid van de faculteit der verbeelding toch voorgesteld als doelmatig voor de ideeën van de rede en hun ontwaken. Het is precies op deze wijze echter dat het esthetische oordeel zelf doelmatig wordt voor de rede, als bron van ideeën, dat wil zeggen voor een intellectueel begrip waarvoor alle esthetisch begrip klein is; en het voorwerp wordt als subliem opgenomen met een genoegen dat alleen mogelijk is door middel van een misnoegen.

\begin{mdframed}[leftmargin=0pt,rightmargin=0pt,backgroundcolor=blue!6,linecolor=white,innertopmargin=6pt,innerbottommargin=6pt]

Kant onderzoekt hoe het sublieme ons raakt op een manier die fundamenteel verschilt van hoe we schoonheid ervaren. Het sublieme is geen prettig of harmonisch gevoel, maar een complexe emotie die voortkomt uit een diepe spanning tussen wat we kunnen waarnemen en wat we kunnen denken. Centraal staat het idee dat het gevoel van ontoereikendheid, het besef dat onze verbeeldingskracht iets niet kan bevatten, eigenlijk een vorm van eerbied is, niet voor het object zelf, maar voor onze eigen redelijke vermogens. Wanneer we geconfronteerd worden met iets oneindig groots, zoals een eindeloze oceaan of de sterrenhemel, probeert onze verbeelding dit in één keer te vatten, maar faalt daarin. Onze zintuigen kunnen de omvang niet overzien, maar ons verstand weet wel dat dit oneindige bestaat en dat we het kunnen denken, zelfs als we het niet kunnen zien. Dit dubbele besef, dat we het niet kunnen bevatten, maar wel kunnen begrijpen, vormt de kern van het sublieme. Het is geen eigenschap van de natuur, maar een weerspiegeling van onze eigen geest, die ons zowel onze beperkingen als onze mogelijkheden toont.

Het sublieme roept een tegenstrijdig gevoel op: aan de ene kant ervaren we onbehagen, omdat onze verbeeldingskracht tekortschiet. We kunnen de grootte van iets als de Grand Canyon of het heelal niet in één keer overzien, wat een gevoel van desoriëntatie en zelfs angst kan oproepen, alsof we dreigen te verdwijnen in een afgrond die we niet kunnen bevatten. Aan de andere kant voelen we een diep genoegen, omdat we weten dat deze falende verbeelding juist wetmatig is. Onze rede eist van ons dat we het oneindige proberen te denken, ook al kunnen we het niet zien. Het mislukken van onze verbeelding is dus niet alleen normaal, maar zelfs doelmatig: het dwingt ons om ons bewust te worden van onze redelijke vermogens, die wel in staat zijn het oneindige te omvatten. Dit is het tegenovergestelde van de ervaring van schoonheid, waar we in kalme contemplatie verkeren. Bij het sublieme voelen we ons bewogen, alsof we heen en weer worden geslingerd tussen afstoting ("ik kan dit niet bevatten!") en aantrekking ("maar ik weet dat het bestaat, en ik kan erover nadenken!"). Kant vergelijkt dit met een trilling: een snel wisselend gevoel van afschuw en fascinatie.

De vraag is dan: waarom voelen we ons toch bevredigd door deze ervaring, ondanks het onbehagen? Volgens Kant komt dit doordat het sublieme ons laat zien dat we meer zijn dan onze zintuigen. Onze verbeelding kan de Grand Canyon niet in één keer overzien, maar onze rede kan wel begrijpen dat hij bestaat, dat hij deel uitmaakt van een groter geologisch systeem, en zelfs dat het heelal oneindig is. Het sublieme confronteert ons met onze eigen vrijheid: we zijn niet beperkt tot wat we kunnen zien, maar kunnen denken wat we niet kunnen ervaren. Dit gevoel is doelmatig, niet voor het object zelf, maar voor onszelf. Het herinnert ons eraan dat we niet alleen zintuiglijke wezens zijn, maar ook redelijke wezens die het oneindige kunnen denken. Het is alsof het sublieme ons uitdaagt om onze eigen grenzen te overschrijden.

Kant benadrukt dat het sublieme een brug vormt tussen onze zintuigen en onze rede. Bij schoonheid werken verbeelding en verstand samen in harmonie: we zien iets moois en begrijpen het direct. Bij het sublieme zijn verbeelding en rede juist in conflict met elkaar. De verbeelding faalt, omdat ze het oneindige niet kan vatten, maar de rede wint, omdat ze het oneindige wel kan denken. Dit conflict is productief: het dwingt ons om ons bewust te worden van onze bovenzintuiglijke vermogens. Het sublieme is dus geen louter esthetische ervaring, zoals schoonheid, maar een morele ervaring. Het herinnert ons eraan dat we niet alleen deel uitmaken van de natuur, maar ook van een morele wereldorde die we niet kunnen zien, maar wel kunnen denken. Het sublieme is daarom geen genoegen zoals schoonheid, maar een diepere bevrediging: het bewijs dat we meer zijn dan onze zintuigen, en dat we deel hebben aan een werkelijkheid die we niet kunnen waarnemen, maar wel kunnen begrijpen.

\bigskip

Het sublieme en het schone voelen fundamenteel anders omdat ze verschillende psychische dynamieken in beweging zetten. Bij het schone, zoals een mooi uitgevoerd ritueel met zijn harmonie, symmetrie en doelmatige vorm, ervaren we kalme bevrediging. Onze verbeelding en ons verstand werken samen als in een vloeiende dans: wat we waarnemen, past naadloos in onze verwachtingen en begrippen. Het ritueel voelt juist aan, alsof het precies klopt met wat we als ordening en betekenis herkennen. Er is geen spanning, alleen een gevoel van evenwicht en instemming, alsof de wereld even perfect aansluit bij ons innerlijke kader.

Bij het sublieme, zoals de ervaring van de Grand Canyon, is er juist disbalans en een fysiek gevoel van overweldiging. Waar schoonheid ons in rust brengt, schokt het sublieme ons. Je staat aan de rand van de afgrond en voelt je letterlijk duizelig: je verbeelding kan de diepte en uitgestrektheid niet in één beeld vatten, je lichaam reageert met een lichte angst ("ik zou kunnen vallen"), maar je rede weet tegelijkertijd dat je veilig bent en dat de canyon een meetbaar, begrijpelijk geheel is. Dat dubbele gevoel, van falen (ik kan dit niet bevatten) en overwinnen (maar ik begrijp het wel), is wat het sublieme zijn emotionele lading geeft: een mengeling van ontzag, angst en trots. Het is geen genoegen, maar een uitdaging die ons bewust maakt van onze eigen beperkingen en mogelijkheden.

Deze paragraaf is een scharniermoment in Kants betoog, omdat hij hier de morele dimensie van het sublieme blootlegt. Kant heeft hiervoor (§§23–26) uitgelegd wat het sublieme is (de ervaring van het oneindige) en hoe het werkt (via het falen van de verbeelding). Nu laat hij zien waarom dat falen juist waardevol is: omdat het ons confronteert met onze redelijke roeping. Het sublieme is geen esthetische categorie zoals schoonheid, maar een brug naar de moraal. Het herinnert ons eraan dat we niet alleen zintuiglijke wezens zijn, maar ook vrije, redelijke wezens die het oneindige kunnen denken; en dat is precies wat ons in staat stelt om morele ideeën (zoals vrijheid, plicht, of het absolute) te omarmen. Vooruitwijzend bereidt dit de weg voor Kants latere discussie over het moreel sublieme (in het tweede deel van de Kritik der Urteilskraft), waar hij zal laten zien hoe het esthetisch sublieme een voorproefje is van de eerbied voor de morele wet.

Het kernpunt van §27 is dat het sublieme ons twee dingen tegelijk leert: onze eigen kleinheid (als zintuiglijke wezens) en onze eigen grootheid (als redelijke wezens). Dat dubbele besef, dat we de wereld niet kunnen bevatten, maar wel kunnen begrijpen, is wat het sublieme zijn morele kracht geeft. Het is geen ervaring van harmonie, maar van conflict (tussen verbeelding en rede), en juist dat conflict is productief: het dwingt ons om ons bewust te worden van onze vrijheid ten opzichte van de natuur. Schoonheid bevestigt de orde die we al kennen; het sublieme daagt ons uit om voorbij die orde te denken.

\end{mdframed}

\subdiv{Het dynamisch sublieme in de natuur}

\section{Over de natuur als een macht.}

\textbf{Macht} is een vermogen dat superieur is aan grote obstakels. Hetzelfde wordt \textbf{heerschappij} genoemd als het ook superieur is aan de weerstand van iets dat zelf over macht beschikt. De natuur, beschouwd in esthetisch oordeel als een macht die geen heerschappij over ons heeft, is \textbf{dynamisch verheven}.

Als de natuur door ons dynamisch als verheven moet worden beoordeeld, moet zij worden voorgesteld als angst opwekkend (hoewel, omgekeerd, niet elk voorwerp dat angst opwekt in ons esthetisch oordeel verheven wordt gevonden). Want in esthetisch oordelen (zonder begrip) kan de superioriteit over obstakels alleen worden beoordeeld naar de grootte van de weerstand. Wat wij echter trachten te weerstaan is een kwaad, en als wij ons vermogen daartegen niet opgewassen vinden, is het een voorwerp van angst. Dus kan de natuur voor het esthetisch oordeelsvermogen alleen als een macht, en dus als dynamisch verheven, worden beschouwd voor zover zij als een voorwerp van angst wordt gezien.

Wij kunnen echter een voorwerp als \textbf{angstaanjagend} beschouwen zonder er bang \textbf{voor} te zijn, als wij het namelijk zo beoordelen dat wij alleen \textbf{denken} aan het geval waarin wij het zouden willen weerstaan en denken dat in dat geval alle weerstand volkomen vruchteloos zou zijn. Zo vreest de deugdzame mens God zonder bang voor Hem te zijn, omdat hij niet denkt aan het geval van God en zijn geboden te willen weerstaan als iets dat \textbf{hem} zorgen baart. Maar omdat hij zo’n geval niet als onmogelijk op zich beschouwt, erkent hij God als angstaanjagend.

Iemand die bang is, kan net zomin oordelen over het verhevene in de natuur als iemand die in de greep is van neiging en begeerte over het schone. De eerste vlucht voor het gezicht van een voorwerp dat hem ontzetting inboezemt, en het is onmogelijk bevrediging te vinden in een ernstig bedoelde terreur. Daarom is de aangenaamheid in het ophouden van iets verontrustends \textbf{vreugde}. Maar deze vreugde over bevrijding van een gevaar gaat gepaard met de voorwaarde dat men nooit meer aan dat gevaar wordt blootgesteld; men kan zelfs terughoudend zijn om terug te denken aan die sensatie, laat staan de gelegenheid ervoor op te zoeken.

Steile, dreigende, als het ware bovenhangende rotsen, onweerswolken die zich tot in de hemel verheffen en bliksemflitsen en donderslagen met zich meebrengen, vulkanen met hun allesvernietigende geweld, orkanen met de verwoesting die zij achterlaten, de grenzenloze oceaan in woede, een hoge waterval op een machtige rivier, enzovoort, maken ons vermogen om weerstand te bieden tot een nietig kleinigheidje in vergelijking met hun macht. Maar hun aanblik wordt des te aantrekkelijker naarmate hij angstaanjagender is, zolang wij ons in veiligheid bevinden, en wij noemen deze voorwerpen graag verheven, omdat zij de kracht van onze ziel boven haar gewone niveau verheffen en ons in staat stellen om in onszelf een vermogen tot weerstand van een geheel andere aard te ontdekken, dat ons de moed geeft om ons te meten met de schijnbare almacht van de natuur.

Want net zoals wij onze eigen beperking vonden in de onmetelijkheid van de natuur en de ontoereikendheid van ons vermogen om een maatstaf aan te nemen die evenredig is aan de esthetische schatting van de grootte van haar \textbf{domein}, maar toch tegelijkertijd in ons eigen redenvermogen een andere, niet-zintuiglijke maatstaf vonden, die die oneindigheid zelf als eenheid onder zich heeft, waartegen alles in de natuur klein is, en zo in onze eigen geest een superioriteit over de natuur zelf, zelfs in haar onmetelijkheid, vonden: zo maakt de onweerstaanbaarheid van haar macht ons, beschouwd als natuurlijke wezens, weliswaar ons lichamelijke machteloosheid bewust, maar tegelijkertijd openbaart zij een vermogen om onszelf als onafhankelijk daarvan te beoordelen en een superioriteit over de natuur, waarop een geheel andere vorm van zelfbehoud is gebaseerd dan die welke door de natuur buiten ons kan worden bedreigd en in gevaar gebracht, waardoor de menselijkheid in onze persoon onverminderd blijft, ook al moet het menselijke wezen zich aan die heerschappij onderwerpen. Op deze manier wordt de natuur in ons esthetisch oordeel als verheven beoordeeld, niet voor zover zij angst opwekt, maar omdat zij onze kracht (die geen deel uitmaakt van de natuur) oproept om die dingen waarover wij ons zorgen maken (goederen, gezondheid en leven) als onbeduidend te beschouwen, en dus haar macht (waarvoor wij, wat deze dingen betreft, weliswaar onderworpen zijn) niet als een heerschappij over onszelf en onze autoriteit te beschouwen waaraan wij ons zouden moeten onderwerpen als het aankomt op onze hoogste principes en hun bevestiging of verlating. Zo wordt de natuur hier alleen verheven genoemd omdat zij de verbeelding tot het punt brengt waarop de geest zichzelf de verhevenheid van zijn eigen roeping, zelfs boven de natuur, kan voorstellen.

Dit zelfrespect wordt niet verminderd door het feit dat wij ons veilig moeten voelen om deze verheffende bevrediging te kunnen ervaren, waarbij (zo zou het kunnen lijken) omdat het gevaar niet ernstig is, de verhevenheid van ons geestelijk vermogen ook niet serieus genomen hoeft te worden. Want de bevrediging hier betreft alleen de \textbf{roeping} van ons vermogen zoals die zich aan ons openbaart in zo’n geval, net zoals de aanleg daartoe in onze natuur ligt; terwijl de ontwikkeling en uitoefening daarvan aan ons wordt overgelaten en onze verantwoordelijkheid blijft. En hierin ligt de waarheid, hoezeer de persoon, als hij zijn reflectie zo ver doorvoert, zich ook bewust kan zijn van zijn huidige werkelijke machteloosheid.

Zeker, dit principe lijkt ver gezocht en subtiel, en dus overdreven voor een esthetisch oordeel; maar de observatie van mensen toont het tegendeel, namelijk dat het het principe kan zijn voor de meest gewone oordelen, ook al is men zich daar niet altijd van bewust. Want wat is het dat zelfs bij de wilde het grootste bewondering wekt? Iemand die niet bang is, die geen angst kent, die dus niet terugdeinst voor gevaar maar met volle overtuiging en krachtig aan het werk gaat. En zelfs in de meest beschaafde omstandigheden blijft deze uitzonderlijk hoge achting voor de krijger bestaan, alleen wordt nu ook geëist dat hij tegelijkertijd alle deugden van de vrede toont: zachtmoedigheid, mededogen en zelfs de juiste zorg voor zijn eigen persoon, precies omdat op deze manier de onoverwinnelijkheid van zijn geest door gevaar kan worden erkend. Hoe veel discussie er ook mag zijn over de vraag of het de staatsman of de generaal is die meer respect verdient in vergelijking met de ander, het esthetisch oordeel beslist in het voordeel van de laatste. Zelfs oorlog, als hij met orde en respect voor de rechten van burgers wordt gevoerd, heeft iets verhevens, en maakt tegelijkertijd de mentaliteit van het volk dat hem op deze manier voert des te verhevener, naarmate hij aan meer gevaren is blootgesteld en daartegen zijn moed heeft kunnen bewijzen; terwijl een lange vrede ertoe leidt dat de geest van louter handel drijven de overhand krijgt, samen met laf selfishheid, lafheid en zwakte, en gewoonlijk de mentaliteit van het volk verlaagt.

Deze analyse van het begrip van het verhevene, voor zover het aan macht wordt toegeschreven, lijkt in tegenspraak met het feit dat wij God gewoonlijk voorstellen als zich in toorn vertonend, maar tegelijkertijd in zijn verhevenheid in donder, storm, aardbeving, enzovoort, waarbij het zich voorstellen dat onze geest enige superioriteit zou hebben over de effecten en, zoals het lijkt, zelfs over de bedoelingen van zo’n macht, zowel dwaasheid als godslastering zou lijken. Hier lijkt het niet te gaan om een gevoel van de verhevenheid van onze eigen natuur, maar eerder om onderwerping, neerslachtigheid en een gevoel van volledige machteloosheid, dat de gepaste gemoedsgesteldheid van de geest is bij de verschijning van zo’n voorwerp, en dat ook gewoonlijk wordt geassocieerd met de idee ervan in het geval van natuurverschijnselen van deze aard. In religie in het algemeen lijken onderwerping, aanbidding met gebogen hoofd, en berouwvolle en angstige gebaren en stem, het enige gepaste gedrag in de aanwezigheid van de Godheid te zijn, en zo zijn ze door de meeste mensen aangenomen en nog steeds in acht genomen. Maar deze gemoedsgesteldheid is verre van intrinsiek en noodzakelijkerwijs verbonden met de idee van de \textbf{verhevenheid} van een religie en haar voorwerp. Iemand die werkelijk bang is omdat hij daarvoor reden in zichzelf vindt, omdat hij zich bewust is met zijn verachtelijke gezindheid een macht te hebben beledigd wier wil onweerstaanbaar en tegelijkertijd rechtvaardig is, bevindt zich zeker niet in de juiste gemoedsgesteldheid om zich te verwonderen over de grootheid van God, waartoe een stemming van kalme beschouwing en een geheel vrij oordeel vereist is. Alleen wanneer hij zich bewust is van zijn oprechte, God welgevallige gezindheid, dienen die effecten van macht om in hem de idee van de verhevenheid van dit wezen op te wekken, voor zover hij in zichzelf een verhevenheid van gezindheid herkent die past bij Gods wil, en zo boven de angst voor dergelijke effecten van de natuur wordt verheven, die hij niet beschouwt als uitbarstingen van Gods toorn. Zelfs nederigheid, als het meedogenloze oordelen over eigen tekortkomingen, die anders, gegeven het bewustzijn van goede gezindheden, gemakkelijk bedekt zouden kunnen worden met de mantel van de broosheid van de menselijke natuur, is een verheven gemoedstoestand, namelijk die van zich vrijwillig onderwerpen aan de pijn van zelfverwijt om geleidelijk de oorzaken daarvan weg te nemen. Alleen op deze manier onderscheidt religie zich intern van bijgeloof, dat geen basis biedt in de geest voor eerbied voor het verhevene, maar alleen voor angst en bezorgdheid voor het wezen van superieure macht, aan wiens wil de verschrikte persoon zich onderworpen ziet zonder hem in hoge achting te houden; waaruit natuurlijk niets anders kan voortkomen dan pogingen om gunsten te winnen en zich in te gratie te brengen, in plaats van een religie van goed gedrag in het leven.

Dus is verhevenheid niet in iets in de natuur zelf bevatten, maar alleen in onze geest, voor zover wij ons bewust kunnen worden van superieur te zijn aan de natuur in ons en dus ook aan de natuur buiten ons (voor zover deze ons beïnvloedt). Alles wat dit gevoel in ons opwekt, inclusief de \textbf{macht} van de natuur die onze eigen krachten oproept, wordt dus (hoewel onjuist) verheven genoemd; en alleen onder de vooronderstelling van deze idee in ons en in relatie daartoe zijn wij in staat om te komen tot de idee van de verhevenheid van dat wezen dat innerlijke eerbied in ons opwekt, niet alleen door zijn macht, die hij in de natuur toont, maar nog meer door het vermogen dat in ons is gelegd om de natuur zonder angst te beoordelen en onze roeping als verheven in vergelijking daarmee te beschouwen.

\begin{mdframed}[leftmargin=0pt,rightmargin=0pt,backgroundcolor=blue!6,linecolor=white,innertopmargin=6pt,innerbottommargin=6pt]

Macht is het vermogen om grote obstakels te overwinnen. Als die macht ook weerstand kan bieden aan iets dat zelf krachtig is, noemen we dat heerschappij. Wanneer we de natuur esthetisch bekijken als een kracht die geen controle over ons heeft, ervaren we haar als dynamisch verheven.

Om de natuur als verheven te beoordelen, moeten we haar zien als iets dat angst oproept — hoewel niet alles wat angst oproept ook als verheven wordt ervaren. Bij een esthetisch oordeel (zonder begrip) kunnen we de superioriteit over obstakels alleen afmeten aan de grootte van de weerstand. Wat we proberen te weerstaan, is iets slechts. Als we ons niet sterk genoeg voelen, wordt het een bron van angst. Daarom kan de natuur alleen als verheven worden gezien als ze ons angst inboezemt.

We kunnen iets als angstaanjagend beschouwen zonder er echt bang voor te zijn. Bijvoorbeeld: als we ons voorstellen dat we weerstand zouden willen bieden, maar weten dat elke poging zinloos zou zijn. Zo vreest een deugdzaam mens God zonder bang voor Hem te zijn, omdat hij niet van plan is om Hem te weerstaan. Toch erkent hij Gods kracht als angstaanjagend, omdat hij weet dat weerstand onmogelijk is.

Iemand die echt bang is, kan niet oordelen over het verhevene in de natuur, net zoals iemand die wordt overheerst door begeerte niet kan oordelen over schoonheid. Wie bang is, vlucht weg van wat hem angst aanjaagt. Echte angst geeft geen ruimte voor bevrediging. De opluchting die komt als het gevaar voorbij is, is puur vreugde. Maar die vreugde gaat gepaard met de wens om nooit meer in zo’n situatie terecht te komen — men wil er niet eens aan terugdenken.

Steile, dreigende kliffen, donkere onweerswolken met bliksem en donder, vulkanen met hun vernietigende kracht, orkanen die alles verwoesten, een woeste oceaan, een hoge waterval op een machtige rivier: al deze natuurverschijnselen maken ons duidelijk hoe klein ons vermogen is om weerstand te bieden. Toch worden ze des te aantrekkelijker naarmate ze angstaanjagender zijn — zolang we zelf in veiligheid zijn. We noemen ze graag verheven, omdat ze onze geest boven het alledaagse tillen en ons een innerlijke kracht laten zien die we niet kenden. Die kracht geeft ons de moed om ons te meten met de schijnbare almacht van de natuur.

Net zoals we onze beperkingen ervaren door de oneindigheid van de natuur, vinden we in ons verstand een andere maatstaf: een niet-zintuiglijk perspectief dat die oneindigheid relatief maakt. Alles in de natuur lijkt klein in vergelijking met die maatstaf. Zo ontdekken we een superioriteit in onze geest, zelfs ten opzichte van de onmetelijke natuur. De onweerstaanbaarheid van de natuurlijke kracht toont ons onze fysieke machteloosheid, maar tegelijkertijd zien we dat we ons kunnen bevrijden van die afhankelijkheid. Onze menselijkheid blijft onaangetast, ook al moeten we ons onderwerpen aan de natuur. Op deze manier beoordelen we de natuur als verheven, niet omdat ze angst oproept, maar omdat ze ons uitdaagt om onze eigen kracht te herkennen — een kracht die niet afhankelijk is van de natuur. Ze daagt ons uit om dingen die ons normaal gesproken bezighouden (zoals bezittingen, gezondheid en leven) als onbelangrijk te zien. Zo ervaren we dat de natuurlijke macht geen absolute controle over ons heeft, vooral niet als het gaat om onze hoogste principes en waarden.

Daarom noemen we de natuur verheven: ze prikkelt onze verbeelding om ons de verhevenheid van onze eigen bestemming te laten zien, zelfs boven de natuur uit.

Dit zelfrespect wordt niet minder waardevol omdat we ons veilig moeten voelen om deze verheffende ervaring te kunnen hebben. Het gaat niet om de ernst van het gevaar, maar om het besef van onze eigen mogelijkheden — mogelijkheden die in onze natuur liggen, maar die we zelf moeten ontwikkelen en toepassen. Dat is onze verantwoordelijkheid, ook al zijn we ons bewust van onze huidige beperkingen.

Dit principe lijkt misschien ingewikkeld en te subtiel voor een esthetisch oordeel, maar de praktijk leert anders. Het is juist de basis voor veel alledaagse oordelen, ook al zijn we ons daar niet altijd van bewust. Wat wekt bijvoorbeeld de grootste bewondering op, zelfs bij de meest ‘wilde’ mensen? Iemand die geen angst kent, die niet terugdeinst voor gevaar, maar met overtuiging en moed handelt. Ook in beschaafde samenlevingen blijft de krijger hoog in achting staan — mits hij ook vredevolle deugden toont, zoals mededogen en zorg voor zichzelf. Dat laat zien dat zijn geest niet overwonnen kan worden door gevaar. Of het nu gaat om een staatsman of een generaal: esthetisch gezien verdient de laatste meer respect. Zelfs oorlog, als hij met orde en respect voor burgers wordt gevoerd, heeft iets verhevens. Hij maakt de mentaliteit van een volk des te verhevener, naarmate ze meer gevaren hebben doorstaan en hun moed hebben getoond. Een lange vrede daartegenover kan leiden tot een mentaliteit van handel drijven, egoïsme, lafheid en zwakte, wat de geest van een volk vaak verlaagt.

Deze uitleg van het verhevene als kracht lijkt in tegenspraak met hoe we God vaak voorstellen: als een toornige, almachtige kracht in donder, storm en aardbevingen. Het idee dat onze geest superieur zou zijn aan Gods effecten — of zelfs aan Zijn bedoelingen — lijkt dwaas en respectloos. Hier gaat het niet om een gevoel van verhevenheid, maar om onderwerping, neerslachtigheid en machteloosheid. Dat is ook de gebruikelijke houding in religie: onderdanigheid, aanbidding met gebogen hoofd, berouw en angst. Maar deze houding hoort niet noodzakelijk bij het idee van de verhevenheid van religie. Iemand die echt bang is omdat hij zich schuldig voelt en weet dat hij een rechtvaardige, almachtige kracht heeft beledigd, kan zich niet verwonderen over Gods grootheid. Daarvoor is kalmte en een vrij oordeel nodig. Alleen wie zich bewust is van een oprechte, God welgevallige houding, kan in Gods kracht de verhevenheid zien. Hij herkent dan in zichzelf een verhevenheid die past bij Gods wil en voelt zich boven de angst voor natuurlijke effecten verheven, zonder ze als uitingen van Gods toorn te zien. Zelfs nederigheid — het streng oordelen over eigen fouten — is een verheven toestand: vrijwillig de pijn van zelfkritiek ondergaan om geleidelijk beter te worden. Zo onderscheidt echte religie zich van bijgeloof, dat alleen angst en onderdanigheid kweekt, zonder respect voor het verhevene. Bijgeloof leidt alleen tot pogingen om gunsten af te smelen, in plaats van een leven volgens goede principes.

Verhevenheid zit niet in de natuur zelf, maar in onze geest. We ervaren het als we ons bewust worden van onze superioriteit ten opzichte van de natuur — zowel in ons als buiten ons. Alles wat dit gevoel oproept, inclusief de kracht van de natuur die onze eigen krachten activeert, noemen we (misschien ten onrechte) verheven. Alleen met dit besef kunnen we het idee van verhevenheid begrijpen: niet als iets dat ons angst aanjaagt, maar als iets dat ons uitdaagt om onze eigen bestemming als verheven te zien, boven de natuur uit.

\bigskip
Kant beschrijft hier hoe de natuur als "dynamisch verheven" wordt ervaren wanneer ze ons confronteert met een overweldigende kracht (bijv. stormen, vulkanen, oceanen). Deze ervaring roept angst op, maar alleen als we ons veilig voelen, kunnen we deze angst omzetten in een gevoel van verhevenheid.
De natuur is een macht die obstakels overwint, maar geen heerschappij over ons heeft. Het verhevene ontstaat wanneer we ons bewust worden van onze eigen innerlijke kracht om deze natuurlijke overweldiging te weerstaan, niet fysiek, maar moreel en geestelijk.
Het verhevene ervaren we alleen als we ons veilig voelen. Pas dan kunnen we de angst ombuigen tot bewondering en het besef van onze eigen geestelijke superioriteit.
Kant linkt het verhevene ook aan religie en moraliteit. Echte verhevenheid is niet gebaseerd op angst voor God of de natuur, maar op het besef van onze eigen morele kracht en bestemming, die boven de natuur uitstijgt.
Hij onderscheidt echte religieuze verhevenheid (gebaseerd op morele autonomie) van bijgeloof (gebaseerd op angst en onderwerping).

Tijdens een rondje golfen in de Zwitserland trok plotseling een zwaar onweer op. De lucht werd donker, de wind stak op, en je hoorde de donder dichterbij komen. Ik zocht snel beschutting in een klein, houten schuilhutje — niet meer dan een afdak met vier muren en een bliksemafleider op het dak. Buiten, op slechts enkele meters afstand, sloeg de bliksem met een explosief, oorverdovend geweld in de grond. De initiële knal was zo luid en scherp dat hij door merg en been ging, gevolgd door een sissend geluid waar de hitte van de bliksem de natte aarde deed verdampen. De aarde trilde onder mijn voeten, en de lucht vulde zich met de scherpe, metalige geur van ozon. Het was een overweldigend schouwspel: de pure, ongeremde kracht van de natuur, die elk moment alles had kunnen vernietigen — en toch, door de bliksemafleider en het besef van mijn eigen nietigheid, voelde ik me vreemd genoeg veilig genoeg om het te aanschouwen.

Toch was er geen angst. Niet omdat ik moedig was, maar omdat het besef daalde dat een dodelijke inslag zo snel en willekeurig zou zijn, dat ik het niet eens zou merken. De bliksemafleider bood een technisch, rationeel houvast — een symbool van menselijke ingeniositeit temidden van de chaos. Maar het echte gevoel van veiligheid kwam voort uit iets diepers: het inzicht dat ik, als individu, in dit immense natuurgeweld eigenlijk nietig was. Dat besef maakte de ervaring niet kleiner, maar juist groter. Het was alsof de natuur me uitdaagde om mijn eigen plaats te begrijpen; niet als slachtoffer, maar als getuige van iets dat mij oversteeg. In die stilte, tussen de donderklappen door, voelde ik een vreemde rust, een mengeling van ontzag en een soort trots: de natuur kon mij fysiek vernietigen, maar niet mijn bewustzijn van deze ervaring zelf. Dat was het verhevene.

Kant benadrukt dat veiligheid essentieel is om het verhevene te kunnen ervaren, omdat angst anders overheerst. Als je werkelijk in gevaar bent, ben je gericht op overleven; er is geen ruimte voor esthetische beschouwing of morele reflectie. Het verhevene ontstaat juist wanneer je veilig bent, maar je je tegelijkertijd bewust bent van de overweldigende kracht van de natuur. Zonder veiligheid blijft alleen angst over, en angst blokkeert het vermogen om de ervaring als verheven te interpreteren. In mijn ervaring was de bliksemafleider een symbool van die veiligheid: hij stelde me in staat om de storm niet als een bedreiging, maar als een spectakel te ervaren.

Kant stelt dat angst ons bewust maakt van onze fysieke kwetsbaarheid, maar morele moed ons in staat stelt om boven die angst uit te stijgen. Het verhevene ontstaat wanneer we angst overwinnen door ons te richten op onze innerlijke kracht: ons vermogen om betekenis te geven aan de ervaring, zelfs als we fysiek machteloos zijn. In de Alpen was er geen angst, maar wel een diep respect voor de kracht van de natuur. Dat respect was geen onderwerping, maar een erkenning van iets groters; en tegelijkertijd een bevestiging van mijn eigen vermogen om dat te ervaren zonder eraan ten onder te gaan.

Kant maakt een scherp onderscheid tussen echte religieuze verhevenheid en bijgeloof:
\begin{compactitem}
\item Echte religieuze verhevenheid is gebaseerd op het besef van onze eigen morele autonomie en waardigheid, zelfs ten opzichte van een almachtige God. Het is een gevoel van respect en bewondering, dat voortkomt uit het idee dat we, ondanks onze fysieke beperkingen, een morele kracht bezitten die boven de natuur uitstijgt.
\item Bijgeloof daarentegen is gebaseerd op angst en onderwerping. Het reduceert religie tot een poging om gunsten af te smeken of straf te vermijden, zonder echte morele reflectie. Bijgeloof ziet God als een willekeurige, dreigende macht, terwijl echte verhevenheid God ziet als een bron van inspiratie voor onze eigen morele groei.
\end{compactitem}
In mijn ervaring was het gebrek aan angst tijdens het onweer vergelijkbaar met Kants idee van echte verhevenheid: ik voelde me niet onderworpen aan de storm, maar juist verbonden met iets groters: een ervaring die mijn eigen bewustzijn verrijkte.

"Der Wanderer über das Nebelmeer" (Caspar David Friedrich, 1818) belichaamt het verhevene perfect: een eenzame wandelaar staat op een rots, uitkijkend over een zee van wolken en bergen. De afgrond voor hem is zowel letterlijk als metaforisch: een symbool van de oneindigheid en de kracht van de natuur. Toch staat de wandelaar stevig, niet bang, maar in contemplatie. Het schilderij roept een gevoel van ontzag op, maar ook van innerlijke kracht: de mens is klein, maar niet machteloos. De veiligheid van de wandelaar (hij staat op vaste grond) stelt hem in staat om de grootsheid van de natuur te ervaren zonder eraan ten onder te gaan.

\bigskip
Kant laat zien dat esthetische ervaringen (zoals het verhevene) niet alleen over schoonheid of angst gaan, maar ook over zelfreflectie. Ze confronteren ons met onze eigen grenzen en met ons vermogen om die te overstijgen. De wiskundige verhevenheid doet dat via de oneindigheid, de dynamische verhevenheid via de kracht. Beide ervaringen leiden tot hetzelfde inzicht: we zijn meer dan alleen natuur. We zijn wezens die betekenis kunnen geven aan wat ons te boven gaat — en dat is de kern van onze morele autonomie.

\end{mdframed}

\section{Over de modaliteit van het oordeel over het verhevene in de natuur.}

Er zijn ontelbare dingen in de mooie natuur waarover wij onmiddellijk instemming met ons eigen oordeel van ieder ander verlangen en ook, zonder bijzonder vatbaar te zijn voor dwaling, kunnen verwachten; maar wij kunnen niet van onszelf beloven dat ons oordeel over het verhevene in de natuur zo gemakkelijk aanvaarding bij anderen zal vinden. Want een veel hogere cultuur, niet alleen van het esthetisch oordeelsvermogen, maar ook van de kennisvermogens waarop dat gebaseerd is, lijkt vereist te zijn om een oordeel over deze uitmuntendheid van de voorwerpen van de natuur te kunnen vellen.

De gezindheid van de geest voor het gevoel van het verhevene vereist haar ontvankelijkheid voor ideeën; want het is juist in de ontoereikendheid van de natuur ten opzichte van deze ideeën, dus alleen onder de vooronderstelling daarvan en van de inspanning van de verbeelding om de natuur als een schema daarvoor te behandelen, dat bestaat wat voor de gevoeligheid afstotend is, maar wat tegelijkertijd aantrekkelijk voor haar is, omdat het een heerschappij is die de rede uitoefent over de gevoeligheid, alleen om deze op een wijze uit te breiden die geschikt is voor haar eigen domein (het praktische) en haar te laten uitkijken over het oneindige, dat voor de gevoeligheid een afgrond is. Inderdaad, zonder de ontwikkeling van morele ideeën zal wat wij, voorbereid door cultuur, verheven noemen, voor de onontwikkelde persoon slechts afstotend lijken. Hij zal in de bewijzen van de heerschappij van de natuur, gegeven door haar vernietigende kracht en de enorme omvang van haar macht, waartegen zijn eigen tot niets verbleekt, alleen de nood, het gevaar en de ellende zien die de persoon zou omringen die daarheen verbannen zou zijn. Zo had de goede en verder verstandige Savoyardse boer (zoals de heer de Saussure vertelt) geen aarzeling om alle liefhebbers van de ijzige bergen voor gek te verklaren. En wie weet of dat geheel onterecht zou zijn geweest als die waarnemer de gevaren waaraan hij zich daar blootstelde, zoals de meeste reizigers gewoonlijk doen, slechts als een hobby of om ze ooit met pathos te kunnen beschrijven, had ondernomen? Maar zijn bedoeling was de veredeling van de mensheid, en deze uitstekende man ervoer het verheffende gevoel dat hij aan de lezers van zijn reizen meegaf als onderdeel van de koop.

Maar juist omdat het oordeel over het verhevene in de natuur cultuur vereist (meer dan dat over het schone), wordt het daarom niet pas door cultuur voortgebracht en zo te zeggen slechts als een kwestie van conventie in de samenleving geïntroduceerd; maar het heeft zijn grond in de menselijke natuur, en wel in datgene wat van iedereen kan worden vereist en van hem gevraagd samen met een gezond verstand, namelijk in de aanleg voor het gevoel voor (praktische) ideeën, d.w.z. voor datgene wat moreel is.

Dit is de grond voor de noodzaak van de instemming van het oordeel van andere mensen met het onze over het verhevene, die wij tegelijkertijd in ons eigen oordeel insluiten. Want net zoals wij iemand die onverschillig is in het oordelen over een voorwerp in de natuur dat wij mooi vinden, een gebrek aan \textbf{smaak} verwijten, zo zeggen wij van iemand die onbewogen blijft door datgene wat wij verheven oordelen, dat hij geen \textbf{gevoel} heeft. We eisen beide echter van elk menselijk wezen, en veronderstellen het ook bij iedereen die enige cultuur heeft, alleen met dit verschil, dat wij de eerste onmiddellijk van iedereen verlangen omdat daarin het oordeelsvermogen de verbeelding slechts relateert aan het verstand, als het vermogen van begrippen, maar omdat de laatste de verbeelding relateert aan de rede, als het vermogen van ideeën, eisen wij het alleen onder een subjectieve vooronderstelling (die wij echter menen gerechtigd te zijn van iedereen te verlangen), namelijk die van het morele gevoel in de mens, en zo schrijven wij ook noodzaak toe aan dit esthetisch oordeel.

In deze modaliteit van esthetische oordelen, namelijk hun veronderstelde noodzakelijkheid, ligt een hoofdmoment voor de kritiek van het oordeelsvermogen. Want het maakt ons bewust van een \textit{a priori} principe in hen, en verheft hen uit de empirische psychologie, waarin zij anders begraven zouden blijven tussen de gevoelens van genot en pijn (slechts met het betekenisloze epitheton van een \textbf{meer verfijnd} gevoel), om hen en daardoor het oordeelsvermogen in de klasse van diegene te plaatsen die als grond een \textit{a priori} principe hebben, en als zodanig hen over te brengen naar de transcendentale filosofie.

\begin{mdframed}[leftmargin=0pt,rightmargin=0pt,backgroundcolor=blue!6,linecolor=white,innertopmargin=6pt,innerbottommargin=6pt]

Kant wijst erop dat het oordeel over het verhevene in de natuur minder vanzelfsprekend wordt gedeeld dan het oordeel over schoonheid. Terwijl we bij een mooie zonsondergang vaak zonder aarzeling aannemen dat anderen het met ons eens zullen zijn, geldt dat niet voor het verhevene. Niet iedereen ervaart een storm, een afgrond of een oneindige woestijn op dezelfde manier als verheven. Dit komt doordat het verhevene meer vraagt dan alleen een esthetische reactie; het vereist een diepere cultuur en ontwikkeling, niet alleen in onze zintuigen, maar ook in ons vermogen om na te denken en te reflecteren. Het schone spreekt direct tot onze verbeelding, maar het verhevene daagt ons uit om verder te gaan: het vraagt om een vermogen om abstracte, morele of rationele ideeën te begrijpen en ermee om te gaan.

Het gevoel voor het verhevene is nauw verbonden met onze ontvankelijkheid voor ideeën, met name morele ideeën. Wanneer we worden geconfronteerd met iets in de natuur dat onze zintuigen en verbeelding overweldigt, zoals een storm of een afgrond, roept dat aanvankelijk een gevoel van afstoting op — angst, onbehagen. Maar tegelijkertijd trekt het ons ook aan, omdat het onze rede uitdaagt om betekenis te geven aan wat we ervaren. De natuur toont haar overweldigende kracht, die ons fysiek zou kunnen vernietigen, en dat confronteert ons met onze eigen kwetsbaarheid. Toch ervaren we ook dat onze rede deze kracht kan overstijgen. We beseffen dat we, ondanks onze fysieke beperkingen, een innerlijke kracht hebben die boven de natuur uitgaat. Dit is het verheffende gevoel: we voelen ons klein, maar ook trots op ons vermogen om deze ervaring te dragen en er betekenis aan te geven. Kant haalt hier het voorbeeld aan van een Savoyardse boer die bergbeklimmers "gek" noemt, omdat ze zich zonder duidelijke reden blootstellen aan gevaar. Maar als die klimmers niet alleen voor de sensatie gaan, maar om de mensheid te veredelen — bijvoorbeeld door wetenschappelijke kennis te vergaren of morele lessen te trekken — dan ervaren ze het verhevene wel. Het gaat dus om de intentie en het vermogen om de ervaring te koppelen aan hogere ideeën, zoals moed, wetenschap of morele groei.

Het verhevene is volgens Kant geen kwestie van conventie of opvoeding, maar heeft zijn wortels in de menselijke natuur zelf. Cultuur helpt ons weliswaar om het verhevene te herkennen en te waarderen, maar de aanleg ervoor ligt diep in ons morele gevoel — ons vermogen om na te denken over wat goed, rechtvaardig of betekenisvol is. Zonder enige cultuur of ontwikkeling zullen mensen het verhevene vaak alleen maar als bedreigend ervaren, zoals een storm die enkel gevaarlijk lijkt of een afgrond die enkel angst aanjaagt. Maar met cultuur leren we deze ervaringen te koppelen aan ideeën, zoals de grootsheid van de natuur, de kracht van de menselijke geest of morele moed. Toch is het verhevene geen uitvinding van de cultuur; het is een natuurlijke aanleg die bij iedereen aanwezig is, maar die alleen tot uiting komt als we ons morele en rationele vermogens ontwikkelen.

Kant stelt dat we van anderen verwachten dat ze instemmen met ons oordeel over het verhevene, net zoals we van iemand die een mooi landschap niet kan waarderen zeggen dat hij "geen smaak heeft." We verwachten van elk menselijk wezen dat het zowel schoonheid als verhevenheid kan ervaren, maar er is een belangrijk verschil. Schoonheid vereist alleen dat we onze verbeelding relateren aan het verstand, wat een relatief eenvoudig proces is. Het verhevene daartegen vereist dat we onze verbeelding relateren aan de rede — ons vermogen om ideeën zoals vrijheid, oneindigheid of moraliteit te denken. Dit is complexer en vereist een morele ontwikkeling. We veronderstellen daarom dat iedereen die enige cultuur heeft, ook dit vermogen bezit, maar alleen onder de voorwaarde dat ze een moreel gevoel hebben ontwikkeld. Kant schrijft dus noodzakelijkheid toe aan het oordeel over het verhevene, maar deze noodzakelijkheid is niet absoluut, zoals bij wiskundige waarheden. Het is voorwaardelijk en geldt alleen voor hen die openstaan voor morele ideeën.

Tot slot benadrukt Kant dat de wijze waarop we oordelen over het verhevene een belangrijk moment is in zijn kritiek van het oordeelsvermogen. Het verhevene maakt ons bewust van een \textit{a priori} principe — een waarheid die niet afhankelijk is van ervaring, maar diep in ons denken geworteld is. Hiermee tilt hij het verhevene uit de sfeer van de empirische psychologie, waar gevoelens van genot en pijn centraal staan, en plaatst het in de transcendentale filosofie. Het verhevene is dus niet slechts een verfijnd gevoel, maar een fundamenteel vermogen van de menselijke geest om zich te verhouden tot wat ons overstijgt. Het toont aan dat we niet alleen natuurlijke wezens zijn, maar ook morele en rationele wezens die betekenis kunnen geven aan de wereld.

Kants betoog in §29 laat zien dat het verhevene niet voor iedereen even toegankelijk is. Het vereist cultuur, ontwikkeling en een openheid voor morele ideeën. Maar tegelijkertijd is het geen elitair concept, want de aanleg ervoor ligt in de menselijke natuur zelf. Het verhevene confronteert ons met onze eigen grenzen, maar toont ook ons vermogen om die grenzen te overstijgen. Dat is precies wat Kant bedoelt met morele autonomie: een oproep om niet alleen te voelen, maar ook te denken en te handelen als vrije, morele wezens.

\end{mdframed}

\section*{Algemene opmerking bij de uiteenzetting van esthetische reflecterende oordelen}

Wat betreft het gevoel van genot kan een voorwerp worden gerekend tot het \textbf{aangename}, het \textbf{mooie}, het \textbf{verhevene}, of het absoluut \textbf{goede} (\textit{iucundum}, \textit{pulchrum}, \textit{sublime}, \textit{honestum}).

Het \textbf{aangename}, als prikkel voor begeren, is overal van dezelfde aard, ongeacht waar het vandaan komt en hoe verschillend de voorstelling (van zintuig en gevoel, objectief beschouwd) ook mag zijn. Bij het oordelen over zijn invloed op de geest gaat het dus alleen om het aantal charmes (gelijktijdig en opeenvolgend), en als het ware alleen om de massa van aangename sensatie; en dus kan dit niet begrepen worden anders dan door \textbf{kwantiteit}. Het draagt ook niet bij aan cultuur, maar is louter een kwestie van genot.

Het \textbf{mooie}, daarentegen, vereist de voorstelling van een bepaalde \textbf{kwaliteit} van het voorwerp, die zich ook begrijpelijk maakt en tot begrippen gebracht kan worden (hoewel dit in het esthetische oordeel niet gebeurt); en het draagt wel bij aan cultuur, doordat het ons leert aandacht te schenken aan doelmatigheid in het genotsgevoel.

Het \textbf{verhevene} bestaat slechts in de verhouding waarin het zintuiglijke in de voorstelling van de natuur wordt beoordeeld als geschikt voor een mogelijk bovenzintuiglijk gebruik ervan.

Het \textbf{absoluut goede}, subjectief beoordeeld naar het gevoel dat het opwekt (het voorwerp van het morele gevoel), als de bepaalbaarheid van de vermogens van het subject door middel van de voorstelling van een \textbf{absoluut noodzakelijke} wet, onderscheidt zich vooral door de \textbf{modaliteit} van een noodzakelijkheid die berust op \textit{a priori} begrippen, die niet alleen een \textbf{aanspraak}, maar ook een \textbf{gebod} bevat dat iedereen ermee instemt, en die niet behoort tot de esthetische, maar tot de zuiver intellectuele oordeelskracht; het wordt ook niet toegeschreven in een louter reflecterend, maar in een bepalend oordeel, niet aan de \textbf{natuur}, maar aan de \textbf{vrijheid}.

Maar de \textbf{bepaalbaarheid van het subject} door middel van deze idee, en wel van een subject dat in zichzelf \textbf{obstakels} in de gevoeligheid kan waarnemen, maar tegelijkertijd superioriteit over deze obstakels door ze te overwinnen als een \textbf{wijziging van zijn toestand}, d.w.z. het morele gevoel, is toch verwant aan de esthetische oordeelskracht en haar \textbf{formele voorwaarden}, voor zover het kan dienen om de wetmatigheid van de handeling uit plicht tegelijkertijd voorstelbaar te maken als esthetisch, d.w.z. als verheven of ook als mooi, zonder de zuiverheid ervan prijs te geven; wat niet het geval zou zijn als men het in natuurlijke combinatie met het gevoel van het aangename zou plaatsen.

Als men de uitkomst trekt uit de uiteenzetting tot nu toe van de twee soorten esthetische oordelen, zou het resultaat de volgende beknopte uitleggen opleveren:

\textbf{Mooi} is wat bevalt in het louter oordelen (dus niet door middel van zintuiglijke sensatie, noch volgens een begrip van het verstand).
Hieruit volgt vanzelf dat het zonder enige belangstelling moet behagen.

\textbf{Verheven} is wat onmiddellijk bevalt door zijn weerstand tegen de belangstelling van de zintuigen.

Beide, als uitleggen van esthetisch universeel geldige oordelen, zijn gerelateerd aan subjectieve gronden, namelijk aan de ene kant aan die van de gevoeligheid, voor zover deze doelmatig is ten behoeve van het beschouwende verstand, en aan de andere kant, \textbf{in tegenstelling} daartoe, als doelmatig voor de doelen van de praktische rede; en toch zijn beide, verenigd in hetzelfde subject, doelmatig in relatie tot het morele gevoel. Het mooie bereidt ons voor om iets te liefhebben, zelfs de natuur, zonder belangstelling; het verhevene, om het te waarderen, zelfs tegen onze (zintuiglijke) belangstelling in.

Men kan het verhevene als volgt beschrijven: het is een voorwerp (van de natuur) \textbf{waarvan de voorstelling de geest bepaalt om te denken aan de onbereikbaarheid van de natuur als een presentatie van ideeën}.

Letterlijk genomen, en logisch beschouwd, kunnen ideeën niet worden gepresenteerd. Maar als we onze empirische voorstellingsvermogen (wiskundig of dynamisch) uitbreiden voor de intuïtie van de natuur, dan komt de rede onvermijdelijk binnen als een vermogen van de onafhankelijkheid van de absolute totaliteit, en brengt de inspanning van de geest voort, hoewel tevergeefs, om de voorstelling van de zintuigen adequaat te maken aan die idee. Deze inspanning, en het gevoel van de onbereikbaarheid van de idee door middel van de verbeelding, is zelf een presentatie van de subjectieve doelmatigheid van onze geest in het gebruik van de verbeelding voor haar bovenzintuiglijke roeping, en dwingt ons om de natuur zelf in haar totaliteit, als de presentatie van iets bovenzintuiglijks, subjectief te \textbf{denken}, zonder deze presentatie \textbf{objectief} te kunnen voortbrengen.

Want we beseffen snel dat de natuur volledig tekortschiet ten opzichte van het onvoorwaardelijke in ruimte en tijd, en dus ten opzichte van absolute grootte, hoewel dit door de gewoonste rede wordt geëist. En juist hierdoor worden we eraan herinnerd dat we alleen met een natuur als verschijning te maken hebben, en dat deze zelf beschouwd moet worden als louter de presentatie van een natuur op zich (die de rede in de idee heeft). Deze idee van het bovenzintuiglijke, die we natuurlijk niet verder kunnen bepalen, zodat we de natuur niet als een presentatie ervan kunnen \textbf{kennen}, maar alleen kunnen \textbf{denken}, wordt in ons gewekt door middel van een voorwerp waarvan het esthetische oordeel de verbeelding tot het uiterste rekkt, of dat nu van vergroting (wiskundig) of van zijn macht over de geest (dynamisch) is, doordat het gegrond is in het gevoel van een roeping van de geest die geheel voorbijgaat aan het domein van de eerste (het morele gevoel), waarvoor de voorstelling van het voorwerp als subjectief doelmatig wordt beoordeeld.

In feite kan een gevoel voor het verhevene in de natuur niet eens worden bedacht zonder het te verbinden met een gezindheid van de geest die lijkt op de morele gezindheid; en hoewel het mooie in de natuur eveneens een bepaalde \textbf{vrijgevigheid} in de wijze van denken veronderstelt, d.w.z. onafhankelijkheid van de bevrediging van louter zintuiglijk genot, wordt de vrijheid toch meer voorgesteld als in \textbf{spel} dan als onderworpen aan een wetmatig \textbf{bedrijf}, wat de eigenlijke eigenschap is van de menselijke moraliteit, waar de rede heerschappij moet uitoefenen over de gevoeligheid; het is juist dat in het esthetische oordeel over het verhevene deze heerschappij wordt voorgesteld als uitgeoefend door de verbeelding zelf, als een instrument van de rede.

De bevrediging in het verhevene in de natuur is dus ook alleen \textbf{negatief} (terwijl die in het mooie \textbf{positief} is), namelijk een gevoel van de ontneming van de vrijheid van de verbeelding door zichzelf, voor zover deze doelmatig wordt bepaald volgens een wet anders dan die van empirisch gebruik. Hierdoor verkrijgt het een uitbreiding en kracht die groter is dan die het opoffert, maar waarvan de grond voor het verborgen blijft, terwijl het de opoffering of ontneming en tegelijkertijd de oorzaak waaraan het onderworpen is, \textbf{voelt}. De \textbf{verbazing} die grenst aan ontzetting, de afschuw en de ontzagwekkende huivering, die de toeschouwer bevatten bij het aanschouwen van bergketens die tot in de hemel reiken, diepe ravijnen en de woeste stroomversnellingen daarin, diep beschaduwde woestijnen die melancholische reflectie opwekken, enz., is, gezien de veiligheid waarin hij zich weet, geen werkelijke angst, maar slechts een poging om onszelf erin te betrekken door middel van de verbeelding, om de kracht van die faculteit zelf te voelen, om de beweging van de geest die daardoor wordt opgewekt te combineren met zijn kalmte, en zo boven de natuur in onszelf uit te stijgen, en dus ook boven die buiten ons, voor zover deze invloed kan hebben op ons gevoel van welzijn. Want de verbeelding, volgens de wet van associatie, maakt onze toestand van tevredenheid fysiek afhankelijk; maar dezelfde verbeelding, volgens de principes van het schematisme van de oordeelskracht (dus voor zover het onderworpen is aan de vrijheid), is een instrument van de rede en haar ideeën, maar als zodanig een kracht om onze onafhankelijkheid te bevestigen tegenover de invloeden van de natuur, om de waarde van wat groot is volgens deze te verminderen, en zo het absoluut grote alleen te plaatsen in zijn (het subject) eigen roeping. Deze reflectie van de esthetische oordeelskracht, die zich verheft tot adequaatheid aan de rede (toch zonder een bepaald begrip van de laatste), stelt het voorwerp, juist door middel van de objectieve ontoereikendheid van de verbeelding in haar grootste uitbreiding tot de rede (als een vermogen van ideeën), voor als subjectief doelmatig.

Hier moet men vooral op letten wat al eerder is opgemerkt, dat in de transcendentale esthetiek van de oordeelskracht het strikt om zuiver esthetische oordelen gaat, dus de voorbeelden niet mogen worden ontleend aan die mooie of verhevene voorwerpen van de natuur die het begrip van een doel veronderstellen; want in dat geval zou het óf teleologisch óf louter gebaseerd zijn op sensaties van een voorwerp (bevrediging of pijn), en dus in het eerste geval geen esthetische doelmatigheid zijn en in het tweede geval geen louter formele doelmatigheid. Zo moet, als iemand de aanblik van de sterrenhemel \textbf{verheven} noemt, hij dit oordeel niet baseren op begrippen van werelden bewoond door redelijke wezens, waarbij de heldere punten waarmee we de ruimte boven ons gevuld zien als hun zonnen, waaromheen ze in hun doelmatig aangewezen banen bewegen, maar moet het nemen zoals we het zien, slechts als een breed, alomvattend gewelf; en het moet louter onder deze voorstelling zijn dat we de verhevenheid plaatsen die een zuiver esthetisch oordeel aan dit voorwerp toekent. Op dezelfde manier moeten we de aanblik van de oceaan niet nemen zoals we hem \textbf{denken}, verrijkt met allerlei kennis (die echter niet in de onmiddellijke intuïtie zijn opgenomen), bijvoorbeeld als een uitgestrekt rijk van waterwezens, als de grote voorraad water voor de verdamping die de lucht met wolken verzadigt ten bate van het land, of als een element dat delen van de wereld van elkaar scheidt maar tegelijkertijd de grootste gemeenschap tussen hen mogelijk maakt, want dit zou slechts teleologische oordelen opleveren; in plaats daarvan moet men de oceaan beschouwen zoals de dichters dat doen, overeenkomstig wat zijn verschijning toont, bijvoorbeeld wanneer hij in periodes van kalmte wordt beschouwd als een heldere waterige spiegel slechts begrensd door de hemel, maar ook wanneer hij turbulent is, een afgrond die alles dreigt te verzwelgen, en toch nog steeds verheven gevonden kan worden. Hetzelfde geldt voor het verhevene en het mooie in de menselijke gedaante, waar we niet naar begrippen van de doelen \textbf{waarvoor} al zijn leden bestaan moeten kijken voor bepalende gronden van ons oordeel, en niet mogen toelaten dat overeenstemming daarmee ons esthetische oordeel \textbf{beïnvloedt} (wat in dat geval niet langer zuiver zou zijn), hoewel het feit dat ze daar niet mee in strijd zijn natuurlijk een noodzakelijke voorwaarde is, zelfs van esthetische bevrediging.
Esthetische doelmatigheid is de wetmatigheid van de oordeelskracht in haar \textbf{vrijheid}. De bevrediging in het voorwerp hangt af van de verhouding waarin we de verbeelding zouden plaatsen: namelijk dat deze de geest door zichzelf in vrije activiteit vermaakt. Als daarentegen iets anders het oordeel bepaalt, of het nu een sensatie van de zintuigen is of een begrip van het verstand, dan is het zeker wetmatig, maar niet het oordeel van een \textbf{vrije} oordeelskracht.

Als men dus spreekt van een intellectuele schoonheid of verhevenheid, dan zijn deze uitdrukkingen \textbf{ten eerste} niet geheel correct, omdat het soorten van esthetische voorstelling zijn die bij ons helemaal niet zouden voorkomen als we louter zuivere intelligenties waren (of zelfs als we ons in onze gedachten tot zulke zouden transformeren); \textbf{ten tweede}, hoewel beide, als voorwerpen van een intellectuele (morele) bevrediging, zeker verenigbaar zijn met het esthetische voor zover ze niet op enige belangstelling \textbf{berusten}, zijn ze toch moeilijk te verenigen met het esthetische omdat ze een belangstelling moeten \textbf{opwekken} die, als de voorstelling moet overeenstemmen met de bevrediging in esthetisch oordelen, nooit zou voorkomen behalve door middel van een zintuiglijke belangstelling, die ermee verbonden is in de voorstelling, waardoor echter schade wordt toegebracht aan de intellectuele doelmatigheid en het onzuiver wordt.

Het voorwerp van een zuivere en onvoorwaardelijke intellectuele bevrediging is de morele wet in al zijn kracht, die hij in ons uitoefent over elke prikkel van de geest \textbf{die eraan voorafgaat}; en aangezien deze kracht zich esthetisch alleen kenbaar maakt door opofferingen (wat een ontneming is, hoewel ten behoeve van innerlijke vrijheid, maar ook een onpeilbare diepte van deze bovenzintuiglijke faculteit samen met haar gevolgen die verder reiken dan wat gezien kan worden), is de bevrediging aan de esthetische kant (in relatie tot de gevoeligheid) negatief, d.w.z. tegengesteld aan deze belangstelling, maar beschouwd vanuit de intellectuele kant is het positief, en verbonden met een belangstelling. Hieruit volgt dat het intellectuele, intrinsiek doelmatige (morele) goede, esthetisch beoordeeld, niet zozeer als mooi maar eerder als verheven moet worden voorgesteld, zodat het meer het gevoel van respect (dat charmes veracht) opwekt dan dat van liefde en innige genegenheid, aangezien de menselijke natuur niet van nature overeenstemt met dat goede, maar alleen door de heerschappij die de rede uitoefent over de gevoeligheid. Omgekeerd wordt zelfs datgene wat we verheven noemen in de natuur buiten ons of zelfs in onszelf (bijv. bepaalde affecten) alleen voorgesteld als een kracht van de geest om boven \textbf{bepaalde} obstakels van de gevoeligheid uit te stijgen door middel van morele principes, en daardoor interessant te worden.

Ik zou graag even stilstaan bij het laatste punt. Het idee van het goede met affect wordt \textbf{enthousiasme} genoemd. Deze gemoedstoestand lijkt verheven, zozeer zelfs dat algemeen wordt aangenomen dat zonder dit niets groots kan worden verricht. Nu is echter elk affect\footnote{\textbf{Affecten} zijn specifiek verschillend van \textbf{passies}. De eerste hebben betrekking uitsluitend op het gevoel; de laatste behoren tot het begeertevermogen en zijn neigingen die elke bepaalbaarheid van het keuzevermogen door middel van principes moeilijk of onmogelijk maken. Affecten zijn stormachtig en onvoorbedacht, passies zijn aanhoudend en overwogen; zo is verontwaardiging, als woede, een affect, maar als haat (wraakzucht) is het een passie. Deze laatste kan nooit, onder welke omstandigheden dan ook, verheven worden genoemd, omdat in het geval van een affect de vrijheid van de geest weliswaar wordt \textbf{belemmerd}, maar in het geval van een passie wordt opgeheven.} blind, ofwel in de keuze van zijn doel, ofwel, zelfs als dit door de rede wordt gegeven, in de uitvoering ervan; want het is die beweging van de geest die hem onbekwaam maakt om zich vrij te bezinnen over principes, om zich daaraan te bepalen. Zo kan het op geen enkele wijze een bevrediging van de rede verdienen. Toch is enthousiasme esthetisch verheven, omdat het een uitrekking van de vermogens door ideeën is, die de geest een momentum geven dat veel krachtiger en aanhoudender werkt dan de impuls die door zintuiglijke voorstellingen wordt gegeven. Maar (wat vreemd lijkt) zelfs \textbf{affectloosheid} (\textit{apatheia}, \textit{phlegma in significatu bono}) in een geest die nadrukkelijk zijn onveranderlijke principes nastreeft, is verheven, en wel op een veel superieure wijze, omdat het ook de bevrediging van de zuivere rede aan zijn zijde heeft. Alleen zo'n mentaliteit wordt \textbf{edel} genoemd, een uitdrukking die vervolgens ook wordt toegepast op dingen, bijv. gebouwen, kleding, een literaire stijl, een lichaamshouding, enz., als deze niet zozeer \textbf{verbazing} wekken (een affect in de voorstelling van nieuwheid die de verwachting overtreft), maar \textbf{eerbied} (een verbazing die niet ophoudt wanneer de nieuwheid verdwijnt), wat gebeurt wanneer ideeën in hun presentatie onbedoeld en zonder kunstgrepen overeenstemmen met esthetische bevrediging.

Elk affect van de \textbf{moedige soort} (dat is, dat het bewustzijn van onze kracht oproept om elke weerstand te overwinnen (\textit{animi strenui})) is \textbf{esthetisch verheven}, bijv. woede, zelfs wanhoop (dat is, de \textbf{verbolgen}, niet de \textbf{neerslachtige} soort). Affect van de \textbf{toegevende} soort (dat de inspanning tot weerstand zelf tot een voorwerp van ongenoegen maakt (\textit{animum languidum})) heeft niets \textbf{edels}, hoewel het tot de schoonheid van de zintuiglijke soort gerekend kan worden. Daarom zijn de \textbf{emoties} die de kracht van een affect kunnen bereiken ook zeer divers. We hebben \textbf{moedige} en \textbf{tedere} emoties. De laatste, als ze het niveau van een affect bereiken, zijn goed voor niets; de neiging daartoe wordt \textbf{overgevoeligheid} genoemd. Een medelijdende pijn die zich niet laat troosten, of waarmee we ons, wanneer het om verzonnen kwaden gaat, bewust inlaten, tot het punt dat we erdoor worden meegesleept alsof het echt is, bewijst en vormt een teergevoelig maar tegelijkertijd zwak karakter, dat een mooie kant toont, en dat zeker fantastisch genoemd kan worden, maar niet eens enthousiast. Romans, sentimentale toneelstukken, oppervlakkige morele voorschriften, die spelen met (valselijk) zogenaamde edele gezindheden, maar in feite het hart verzwakken en het ontoegankelijk maken voor het strenge voorschrift van de plicht en onbekwaam voor alle respect voor de waardigheid van de mensheid in onze eigen persoon en het recht van de mens (wat iets geheel anders is dan hun geluk), en in het algemeen onbekwaam voor alle vaste principes; zelfs een religieuze preek die een kruiperige, laag-bij-de-grondse vleierij en zelfverloochening predikt, die alle vertrouwen in onze eigen weerstandsvermogen tegen het kwaad opgeeft, in plaats van de energieke beslissing om de krachten die nog in ons zijn, ondanks al onze zwakheid, te zoeken om neigingen te overwinnen; de valse nederigheid die de enige manier ziet om de opperwezen te behagen in zelfverachting, in gejammer, geveinsde berouw en een louter passieve houding van de geest, al deze dingen hebben niets te maken met wat als de schoonheid, laat staan de verhevenheid, van een mentaliteit kan worden beschouwd.

Maar zelfs stormachtige bewegingen van de geest, of ze nu geassocieerd worden met religieuze ideeën onder de naam van stichting, of, als louter behorend tot cultuur, met ideeën die een sociaal belang bevatten, hoezeer ze ook de verbeelding uitrekken, kunnen op geen enkele wijze de eer opeisen van een \textbf{verheven} voorstelling te zijn, als ze niet een gezindheid van de geest achterlaten die, zelfs als het alleen indirect is, invloed heeft op het bewustzijn van zijn kracht en vastberadenheid met betrekking tot datgene wat intellectuele doelmatigheid met zich meebrengt (het bovenzintuiglijke). Anders behoren al deze emoties slechts tot de \textbf{beweging} die we graag hebben omwille van de gezondheid. De aangename uitputting die volgt op zo'n opwinding door het spel van affecten is een genot van het welzijn dat voortkomt uit het evenwicht van de verschillende levenskrachten dat zo in ons wordt geproduceerd, wat uiteindelijk neerkomt op hetzelfde als wat de genieters van het Oosten zo troostend vinden wanneer ze hun lichaam als het ware laten kneden, en al hun spieren en gewrichten zachtjes laten drukken en buigen; alleen is in het eerste geval het bewegende principe voor het grootste deel in ons, terwijl het in het laatste geheel buiten ons is. Nu gelooft menig mens dat hij gesticht wordt door een preek waarin echter niets (geen systeem van goede maximen) is opgebouwd, of veredeld door een tragedie wanneer hij slechts blij is met een gelukkige ontsnapping aan verveling. Zo moet het verhevene altijd een relatie hebben tot de \textbf{wijze van denken}, d.w.z. tot maximen om het intellectuele en de ideeën van de rede superieur te maken aan de gevoeligheid.

Er hoeft geen angst te bestaan dat het gevoel van het verhevene iets zal verliezen door zo'n abstracte voorstelling, die geheel negatief is ten opzichte van het zintuiglijke; want de verbeelding, hoewel het zeker niets buiten het zintuiglijke vindt waaraan het zich kan hechten, voelt zich juist door deze eliminatie van de grenzen van de gevoeligheid onbegrensd; en die scheiding is dus een voorstelling van het oneindige, die om die reden nooit iets anders kan zijn dan een louter negatieve voorstelling, die desondanks de ziel uitbreidt. Misschien is er geen verhevenere passage in het Joodse Wetboek dan het gebod: Gij zult u geen gesneden beeld maken, noch enige gelijkenis van hetgeen in de hemel boven is, noch van hetgeen op de aarde beneden is, enz. Dit gebod alleen kan de enthousiaste gevoelens verklaren die het Joodse volk in zijn beschaafde periode voor zijn religie voelde wanneer het zichzelf met andere volkeren vergeleek, of de trots die de islam inspireerde. Precies hetzelfde geldt voor de voorstelling van de morele wet en de aanleg tot moraliteit in ons. Het is volkomen verkeerd om te vrezen dat, als deze van alles wat de zintuigen kunnen aanbevelen zou worden beroofd, het dan niets anders dan koude, levenloze instemming en geen bewegende kracht of emotie zou meebrengen. Het is juist het omgekeerde: want waar de zintuigen niets meer voor zich zien, maar de onmiskenbare en onuitblusbare idee van moraliteit blijft, zou het nodig zijn om de vaart van een onbegrensde verbeelding te matigen om niet tot enthousiasme te laten komen, in plaats van, uit angst voor de machteloosheid van deze ideeën, hulp voor hen te zoeken in beelden en kinderlijke hulpmiddelen. Daarom hebben zelfs regeringen graag toegestaan dat de religie rijkelijk werd uitgerust met zulke aanvullingen en zo de onderdaan van de moeite, maar tegelijkertijd ook van het vermogen om de krachten van zijn ziel uit te breiden voorbij de grenzen die willekeurig voor hem zijn gesteld, en waardoor hij, louter passief, gemakkelijker kan worden behandeld.

Deze zuivere, verheffende, louter negatieve voorstelling van de moraliteit brengt daarentegen geen risico met zich mee van \textbf{visionaire extase}, die een \textbf{waan is om iets voorbij alle grenzen van de gevoeligheid te kunnen zien}, d.w.z. om te dromen volgens principes (om met de rede te razen), precies omdat de voorstelling in dit geval louter negatief is. Want de \textbf{ondoorgrondelijkheid van de idee van vrijheid} sluit elke positieve voorstelling geheel uit; maar de morele wet is op zichzelf in ons voldoende en oorspronkelijk bepalend, zodat hij ons niet eens toestaat om elders naar een bepalende grond te zoeken. Als enthousiasme vergeleken kan worden met de \textbf{waan van de zintuigen}, dan moet visionaire extase vergeleken worden met de \textbf{waan van de geest}, die het minst van alles verenigbaar is met het verhevene, omdat het broedend en absurd is. In enthousiasme, als een affect, is de verbeelding ongetemd; in visionaire extase, als een diepgewortelde, onderdrukkende passie, is het ongeregeerd. De eerste is een voorbijgaand ongeluk, dat af en toe het gezondste verstand aantast; de laatste is een ziekte die het vernietigt.

\textbf{Eenvoud} (kunstloze doelmatigheid) is als het ware de stijl van de natuur in het verhevene, en dus ook van de moraliteit, die een tweede (bovenzintuiglijke) natuur is, waarvan we alleen de wetten kennen, zonder door intuïtie het bovenzintuiglijke vermogen in onszelf te kunnen bereiken dat de grond van deze wetgeving bevat.

Het moet verder worden opgemerkt dat, hoewel de bevrediging in het mooie, evenals die in het verhevene, zich niet alleen duidelijk onderscheidt van andere esthetische oordelen door middel van universele \textbf{mededeelbaarheid}, maar ook door middel van deze eigenschap een belangstelling verkrijgt in relatie tot de samenleving (waar het gecommuniceerd kan worden), toch de \textbf{afzondering van alle samenleving} ook als iets verhevens wordt beschouwd als het berust op ideeën die voorbij alle zintuiglijke belangstelling gaan. Zelfvoorzienend zijn, dus de samenleving niet nodig hebben, maar zonder onmaatschappelijk te zijn, d.w.z. haar te ontvluchten, is iets dat dicht bij het verhevene komt, net als elke superioriteit boven behoeften. In tegenstelling daarmee is het ontvluchten van mensen uit \textbf{mensenhaat}, omdat men vijandig tegenover hen staat, of uit \textbf{antropofobie} (angst voor mensen), omdat men hen vreest als vijanden, deels hatelijk en deels verachtelijk. Toch is er een soort mensenhaat (zeer onjuist zo genoemd), waarvoor de aanleg vaak wordt gevonden in de geest van veel weldenkende mensen naarmate ze ouder worden, die zeker filantropisch genoeg is wat hun \textbf{welwillendheid} betreft, maar door lange, droevige ervaring ver verwijderd is van enig \textbf{genot} in mensen; bewijs hiervan is te vinden in de neiging om zich terug te trekken uit de samenleving, de fantastische wens naar een geïsoleerd landgoed, of zelfs (bij jongeren) de droom van geluk in het kunnen doorbrengen van hun leven op een eiland onbekend bij de rest van de wereld met een klein gezin, wat de romanciers of dichters die Robinsonades schrijven zo goed weten te exploiteren. Leugenachtigheid, ondankbaarheid, onrechtvaardigheid, de kinderlijkheid in doelen die we zelf belangrijk en groot vinden, in de nasleep waarvan mensen elkaar elk denkbaar kwaad aandoen, zijn zo in strijd met de idee van wat ze zouden kunnen zijn als ze wilden, en zijn zo tegengesteld aan de levendige wens om een beter beeld van hen te krijgen, dat, om ze niet te haten, omdat men ze niet kan liefhebben, het ontberen van alle sociale vreugden slechts een kleine opoffering lijkt. Deze droefheid, niet over het kwaad dat het lot anderen aandoet (wat door medelijden wordt veroorzaakt), maar over datgene wat ze zichzelf aandoen (wat berust op antipathie in fundamentele principes), is, omdat het op ideeën berust, verheven, terwijl de eerste hoogstens als mooi kan worden beschouwd. Saussure, zo geïnspireerd als hij grondig is, zegt in de beschrijving van zijn reizen in de Alpen over Bonhomme, een van de bergen van Savoye: ‘‘Daar heerst een zekere \textbf{eentonige droefheid}.’’ Maar hij kende ook een \textbf{interessante} droefheid, die wordt ingegeven door de aanblik van een woestenij waarnaar mensen zouden verhuizen om niets meer van de wereld te horen of te ervaren, maar die toch niet zo volkomen onherbergzaam mag zijn dat het mensen slechts een uiterst belastend toevluchtsoord zou bieden. Ik maak deze opmerking alleen met de bedoeling om in herinnering te brengen dat zelfs verdriet (niet neerslachtige droefheid) tot de \textbf{krachtige} affecten gerekend kan worden als het gegrond is in morele ideeën, maar als het gegrond is in medelijden, en als zodanig ook liefdevol is, behoort het slechts tot de \textbf{verzachtende} affecten, alleen om de aandacht te vestigen op de gezindheid van de geest die alleen in het eerste geval \textbf{verheven} is.

\bigskip
De transcendentale uiteenzetting van esthetische oordelen die nu is voltooid, kan worden vergeleken met de fysiologische uiteenzetting, zoals die is uitgewerkt door een \textbf{Burke} en vele scherpe mannen onder ons, om te zien waar een louter empirische uiteenzetting van het verhevene en het mooie toe zou leiden. \textbf{Burke}, die verdient genoemd te worden als de voornaamste auteur in deze benadering, brengt in deze zin naar voren (p. 223 van zijn werk) ‘‘dat het gevoel van het verhevene berust op de drang tot zelfbehoud en op angst, d.w.z. een pijn, die, omdat het niet zo ver gaat als de daadwerkelijke vernietiging van lichaamsdelen, bewegingen voortbrengt die, omdat ze de fijnere of grovere vaten van gevaarlijke en lastige verstoppingen zuiveren, in staat zijn om aangename sensaties op te wekken, niet zeker genot, maar een soort aangename afschuw, een zekere rust die vermengd is met terreur.’’ Het mooie, dat hij baseert op liefde (die hij echter gescheiden zou moeten weten van begeerte), leidt hij af (pp. 251–52) ‘‘tot ontspanning, losmaking en verslapping van de vezels van het lichaam, dus tot een verzachting, een oplossing, een verzwakking, een wegzinken, een sterven van genot.’’ En nu bevestigt hij deze soort uitleg door gevallen waarin de verbeelding in staat is het gevoel van het mooie evenals het verhevene niet alleen in verband met het verstand, maar zelfs in verband met zintuiglijke sensaties op te wekken., Als psychologische opmerkingen zijn deze analyses van de verschijnselen van onze geest uiterst fijn, en bieden rijke materialen voor de geliefde onderzoeken van de empirische antropologie. Bovendien kan niet worden ontkend dat alle voorstellingen in ons, of ze nu objectief louter zintuiglijk zijn of geheel intellectueel, toch subjectief verbonden kunnen zijn met genot of pijn, hoe onopvallend ook (omdat ze allemaal het levensgevoel aantasten, en geen van hen, voor zover het een wijziging van het subject is, onverschillig kan zijn), of zelfs dat, zoals Epicurus beweerde, \textbf{genot} en \textbf{pijn} altijd uiteindelijk lichamelijk zijn, of ze nu voortkomen uit de verbeelding of zelfs uit voorstellingen van het verstand: omdat leven zonder het gevoel van het lichamelijke orgaan slechts bewustzijn van het eigen bestaan is, maar geen gevoel van wel- of onwelzijn, d.w.z. de bevordering of remming van de levenskrachten; omdat de geest op zichzelf geheel leven is (het principe van het leven zelf), en belemmeringen of bevorderingen buiten hem moeten worden gezocht, hoewel in de mens zelf, dus in combinatie met zijn lichaam.

Als men echter de bevrediging in het voorwerp geheel situeert in het feit dat het genot verschaft door middel van charme en emotie, dan mag men niet verwachten dat \textbf{anderen} zullen instemmen met de esthetische oordelen die wij vellen; want daarover is iedereen gerechtigd om slechts zijn eigen persoonlijke zin te raadplegen. In dat geval houdt echter alle kritiek van smaak geheel op; want men zou dan het voorbeeld dat anderen geven door middel van de toevallige overeenstemming onder hun oordelen tot een \textbf{gebod} voor instemming van ons moeten maken, in strijd waarmee principe we waarschijnlijk zouden worstelen en een beroep zouden doen op het natuurlijke recht om het oordeel dat berust op het onmiddellijke gevoel van ons eigen welzijn aan onze eigen zin te onderwerpen, en niet aan die van anderen.

Als het oordeel van smaak dus niet als \textbf{egoïstisch} mag worden beschouwd, maar noodzakelijkerwijs, volgens zijn innerlijke aard, d.w.z. vanzelf, niet omwille van de voorbeelden die anderen geven van hun smaak, als \textbf{pluralistisch}, als men het evalueert als een oordeel dat tegelijkertijd kan eisen dat iedereen ermee instemt, dan moet het berusten op een soort \textit{a priori} principe (of het nu objectief of subjectief is), waar men nooit toe kan komen door te speuren naar empirische wetten van de veranderingen van de geest: want deze laten ons alleen zien hoe dingen worden beoordeeld, maar nooit voorschrijven hoe ze beoordeeld moeten worden, met name op een \textbf{onvoorwaardelijke} wijze; hoewel het iets van deze aard is dat de oordelen van smaak veronderstellen wanneer ze de bevrediging willen laten kennen als \textbf{onmiddellijk} verbonden met een voorstelling.

Dus kan de empirische uiteenzetting van esthetische oordelen altijd een begin maken met het leveren van materiaal voor een hoger onderzoek, maar een transcendentale discussie van deze faculteit is nog steeds mogelijk en essentieel voor de kritiek van de smaak. Want tenzij deze \textit{a priori} principes heeft, zou het onmogelijk zijn om de oordelen van anderen te leiden en aanspraak te maken op het recht om ze goed of af te keuren.

Wat behoort tot het overige van de analytiek van de esthetische oordeelskracht bevat allereerst de:

\begin{mdframed}[leftmargin=0pt,rightmargin=0pt,backgroundcolor=blue!6,linecolor=white,innertopmargin=6pt,innerbottommargin=6pt]

In de "Algemene opmerking" bij de uiteenzetting van esthetische reflecterende oordelen in Kants \textit{Kritik der Urteilskraft} schetst Kant een systematisch overzicht van hoe voorwerpen zich verhouden tot het gevoel van genot en onderscheidt hij vier categorieën: het aangename, het mooie, het verhevene en het absoluut goede. Het aangename is louter zintuiglijk en draagt niet bij aan cultuur, maar is slechts een kwestie van genot, terwijl het mooie wel bijdraagt aan cultuur doordat het ons leert aandacht te schenken aan doelmatigheid in het genotsgevoel. Het verhevene daentegen bestaat in de verhouding waarin het zintuiglijke in de voorstelling van de natuur wordt beoordeeld als geschikt voor een mogelijk bovenzintuiglijk gebruik. Het absoluut goede, subjectief beoordeeld, onderscheidt zich door een noodzakelijkheid die berust op \textit{a priori} begrippen en behoort tot de zuivere intellectuele oordeelskracht, gerelateerd aan vrijheid in plaats van natuur.

Kant benadrukt dat het verhevene niet kan worden begrepen zonder een verband met een gezindheid die lijkt op de morele gezindheid. Het verhevene roept een gevoel op van weerstand tegen de belangstelling van de zintuigen en wordt gekenmerkt door een negatieve bevrediging, namelijk een gevoel van ontneming van de vrijheid van de verbeelding, die zichzelf echter tegelijkertijd als doelmatig ervaart. Dit gevoel wordt opgewekt door voorstellingen die de verbeelding tot haar uiterste rekken, zoals de oneindigheid van de natuur of de overweldigende kracht ervan. Het verhevene confronteert de mens met zijn eigen grenzen, maar toont ook het vermogen om deze te overstijgen door middel van rede en morele reflectie.

Het mooie bereidt de mens erop voor om iets zonder belangstelling lief te hebben, terwijl het verhevene ertoe aanzet om het te waarderen, zelfs tegen de zintuiglijke belangstelling in. Beide, het mooie en het verhevene, zijn esthetisch universeel geldig en gerelateerd aan subjectieve gronden: het mooie aan de gevoeligheid die doelmatig is voor het beschouwende verstand, het verhevene aan de gevoeligheid die doelmatig is voor de doelen van de praktische rede. Samen zijn ze doelmatig in relatie tot het morele gevoel.

Kant gaat vervolgens in op het onderscheid tussen affecten en passies. Affecten zijn stormachtig en onvoorbedacht, terwijl passies aanhoudend en overwogen zijn. Passies kunnen nooit verheven worden genoemd, omdat ze de vrijheid van de geest opheffen, terwijl affecten deze slechts belemmeren. Het verhevene wordt gekoppeld aan een gevoel van respect en een bewustzijn van morele kracht, terwijl het mooie meer gerelateerd is aan liefde en genegenheid.

Ten slotte waarschuwt Kant tegen een louter empirische benadering van het verhevene en het mooie, zoals voorgesteld door denkers als Burke, die het verhevene baseert op angst en zelfbehoud, en het mooie op lichamelijke ontspanning. Kant stelt dat een dergelijke benadering onvoldoende is, omdat esthetische oordelen \textit{a priori} principes vereisen om universeel geldig te zijn. Hij benadrukt dat een transcendentale discussie van de esthetische oordeelskracht essentieel is voor een kritiek van de smaak, omdat zonder \textit{a priori} principes esthetische oordelen niet kunnen claimen universeel te zijn.

\bigskip
Kant benadrukt dat het verhevene niet slechts een esthetische ervaring is, maar een morele oproep. Het confronteert ons met onze eigen grenzen en toont tegelijkertijd ons vermogen om die grenzen te overstijgen. Dit is geen louter gevoel, maar een reflectie op onze morele autonomie — het vermogen om ons te verhouden tot wat ons overstijgt, zoals de oneindigheid van de natuur of de eis van de morele wet. Het verhevene is dus geen kwestie van conventie of cultuur, maar van een \textit{a priori} principe: een fundamenteel vermogen van de menselijke geest.

Kant maakt een cruciaal onderscheid tussen affecten (zoals woede of verontwaardiging) en passies (zoals haat of wraakzucht). Affecten zijn stormachtig en tijdelijk, maar passies zijn aanhoudend en ondermijnen de vrijheid van de geest. Het verhevene kan alleen worden ervaren in een staat van vrije reflectie, niet in een toestand van passie. Dit betekent dat het verhevene ons uitnodigt om ons te verheffen boven onze zintuiglijke neigingen en ons te richten op morele principes.

Het verhevene is niet voor iedereen even toegankelijk. Het vereist cultuur, ontwikkeling en morele reflectie. Zonder deze kan de mens het verhevene slechts als bedreigend ervaren (bijv. een storm als enkel gevaarlijk). Met cultuur leren we deze ervaringen te koppelen aan hogere ideeën, zoals de grootsheid van de natuur of de kracht van de menselijke geest. Het verhevene is dus een oproep tot zelfontwikkeling.

In tegenstelling tot het mooie, dat positieve bevrediging geeft, is de bevrediging in het verhevene negatief: het is een gevoel van ontneming van de vrijheid van de verbeelding, maar tegelijkertijd een bewustzijn van onze morele kracht. Dit negatieve karakter is essentieel, omdat het ons confronteert met onze eigen beperkingen en ons uitdaagt om deze te overstijgen.

Neem de rituele arbeid in een vrijmetselaarsloge. Vrijmetselarij is een traditie die sterk gericht is op morele ontwikkeling, symboliek en zelfreflectie. Tijdens een logebijeenkomst worden leden geconfronteerd met symbolen van oneindigheid (zoals de "alziendheid" van het Oog van de Voorzienigheid), de broosheid van het menselijk leven (bijv. het gebruik van de doodskop in de Kamer van Overdenking), en de noodzaak om de innerlijke passies te overwinnen (bijv. het "temmen van de ruwe steen" als symbool voor zelfverbetering).

De ervaring van de rituele ruimte, met zijn indrukwekkende symboliek en plechtige sfeer, kan een gevoel van ontzag oproepen — een typisch kenmerk van het verhevene. Dit ontzag is geen angst, maar een bewustzijn van iets dat ons overstijgt: de morele eis om een beter mens te worden.
Een vrijmetselaar die tijdens een ritueel wordt geconfronteerd met zijn eigen tekortkomingen (bijv. door een morele les of een symbolische "rechterstoel"), kan een affect ervaren, zoals schaamte of verontwaardiging. Dit is stormachtig, maar tijdelijk. Als hij deze ervaring echter gebruikt om zich te richten op morele groei, overstijgt hij de passie en bereikt hij een staat van vrije reflectie — precies wat Kant bedoelt met het verhevene.
Het ritueel kan ook een gevoel van "ontneming" oproepen: het loslaten van egoïsme, vooroordelen of wereldse bekommernissen. Dit is negatief in de zin van Kant, maar het opent de weg naar een diepere, morele bevrediging.

\end{mdframed}

\kantbook{Deductie van zuiver esthetische oordelen.}

\section{De deductie van esthetische oordelen over de voorwerpen van de natuur mag niet gericht zijn op datgene wat wij onder hen verheven noemen, maar alleen op het mooie.}

De aanspraak van een esthetisch oordeel op universele geldigheid voor elk subject, als een oordeel dat op een of ander \textit{a priori} principe moet berusten, behoeft een deductie (d.w.z. een legitimatie van zijn veronderstelling), die aan de uiteenzetting ervan moet worden toegevoegd, indien het tenminste een bevrediging of onbevrediging in de \textbf{vorm van het voorwerp} betreft. De smaakoordelen over het mooie in de natuur zijn van deze aard. Want in dit geval heeft de doelmatigheid haar grond in het voorwerp en zijn vorm, zelfs als het niet de relatie van het voorwerp tot anderen volgens begrippen aangeeft (voor kennisoordelen), maar eerder algemeen slechts de opvatting van deze vorm betreft, voor zover deze zich in de geest toont als geschikt voor zowel het \textbf{vermogen} tot begrippen als tot de voorstelling daarvan (wat één en hetzelfde is als dat van de opvatting). Daarom kan men ook vele vragen stellen over het mooie in de natuur, met betrekking tot de oorzaak van deze doelmatigheid van haar vormen: bijv. hoe moet men verklaren dat de natuur schoonheid zo overdadig overal heeft verspreid, zelfs op de bodem van de oceaan, waar het menselijk oog (voor wie het immers doelmatig is) slechts zelden doordringt? en zo voort.

Alleen het verhevene in de natuur, als wij een zuiver esthetisch oordeel erover vellen, dat niet vermengd is met begrippen van perfectie als objectieve doelmatigheid, in welk geval het een teleologisch oordeel zou zijn, kan als geheel vormloos of gedaanteloos worden beschouwd, maar desondanks als voorwerp van een zuivere bevrediging, en kan subjectieve doelmatigheid aantonen in de gegeven voorstelling; en de vraag rijst nu of in het geval van deze soort esthetisch oordeel, naast de uiteenzetting van wat erin wordt gedacht, ook een deductie van zijn aanspraak op een zekere (subjectieve) \textit{a priori} principe gevraagd zou kunnen worden.

Als antwoord hierop dient dat het verhevene in de natuur eigenlijk onjuist zo wordt genoemd en eigenlijk alleen aan de wijze van denken, of liever aan de grond ervan in de menselijke natuur, moet worden toegeschreven. De opvatting van een verder vormloos en niet-doelmatig voorwerp biedt slechts de gelegenheid om ons hiervan bewust te worden, wat op deze manier subjectief doelmatig wordt \textbf{gebruikt}, maar niet \textbf{als zodanig} en omwille van zijn vorm (als het ware \textit{species finalis accepta, non data}) wordt beoordeeld. Daarom was onze uiteenzetting van de oordelen over het verhevene in de natuur tegelijkertijd hun deductie. Want toen wij de reflectie van de oordeelskracht in deze oordelen analyseerden, vonden wij daarin een doelmatige relatie van de kennisvermogens, die het vermogen tot doelen (de wil) \textit{a priori} moet gronden, en dus zelf \textit{a priori} doelmatig is, wat dan onmiddellijk de deductie, d.w.z. de rechtvaardiging van de aanspraak van een dergelijk oordeel op universeel noodzakelijke geldigheid, bevat.

Wij zullen dus alleen de deductie van smaakoordelen, d.w.z. van de oordelen over de schoonheid van dingen in de natuur, moeten zoeken, en daardoor de taak voor de gehele esthetische oordeelskracht in haar geheel voltooien.

\begin{mdframed}[leftmargin=0pt,rightmargin=0pt,backgroundcolor=blue!6,linecolor=white,innertopmargin=6pt,innerbottommargin=6pt]

Kant legt uit dat we alleen een filosofische rechtvaardiging (deductie) hoeven te zoeken voor esthetische oordelen over schoonheid in de natuur, niet voor het verhevene. Hier is waarom:
Wanneer we iets "mooi" noemen, claimen we dat iedereen het met ons eens moet zijn. Dit oordeel is niet willekeurig, maar berust op een \textit{a priori} principe, een universele, objectieve basis in de vorm of structuur van het voorwerp zelf. Bij schoonheid in de natuur (zoals een bloem of een landschap) kunnen we vragen stellen over waarom die vorm als doelmatig of harmonieus wordt ervaren. Bijvoorbeeld: waarom is de natuur zo overdadig mooi, zelfs op plekken waar mensen haar nauwelijks zien, zoals de bodem van de oceaan? Dit roept de vraag op naar een objectieve grond voor ons oordeel, een deductie is dus nodig.

Het verhevene daarentegen (zoals een overweldigende storm of een oneindige woestijn) is vormloos en roept geen universele instemming op zoals schoonheid. Het verhevene is geen eigenschap van het voorwerp zelf, maar van onze manier van denken erover. Wanneer we iets verheven vinden, gaat het niet om de vorm, maar om hoe wij als mens reageren op wat ons overstijgt: we ervaren onze eigen morele kracht om boven onze zintuiglijke beperkingen uit te stijgen. Omdat het verhevene geen objectieve vorm heeft, maar een subjectieve ervaring is, hoeft Kant hier geen aparte deductie voor te geven. De analyse van het verhevene is al zijn eigen rechtvaardiging: het toont hoe onze geest zichzelf doelmatig ervaart in de confrontatie met het oneindige of overweldigende.

Kortom: voor schoonheid moeten we een universele basis zoeken, omdat we claimen dat iedereen het met ons eens moet zijn. Voor het verhevene is dat niet nodig, omdat het gaat om een persoonlijke, morele reactie. Kant richt zich daarom alleen op de deductie van smaakoordelen over schoonheid om zijn esthetische theorie af te ronden.

\bigskip
In §30 van de Kritik der Urteilskraft legt Immanuel Kant uit waarom alleen schoonheid — en niet het verhevene — een filosofische rechtvaardiging, ofwel een deductie, behoeft. Dit onderscheid vormt de kern van zijn esthetische theorie. Wanneer we iets mooi noemen, claimen we niet alleen een persoonlijke voorkeur, maar verwachten we dat iedereen het met ons eens is. Deze aanspraak op universele instemming vereist een \textit{a priori} principe: een objectieve basis in de vorm of structuur van het voorwerp zelf. Daarom moet Kant aantonen waarom deze ervaring van harmonie niet louter subjectief is, maar een noodzakelijke, universele geldigheid heeft. Dit is precies wat de deductie beoogt: het bewijs dat schoonheid geen kwestie is van individuele smaak, maar van een gedeelde, rationele ervaring.

Het verhevene daarentegen is fundamenteel anders. Het is niet gebonden aan de eigenschappen van het voorwerp — zoals een storm, een berg of de oneindige sterrenhemel — maar aan onze morele reactie erop. Het verhevene confronteert ons met onze eigen beperkingen, met wat ons overstijgt, maar toont tegelijkertijd ons vermogen om boven onze zintuiglijke ervaring uit te stijgen. Omdat het geen objectieve vorm heeft, maar een persoonlijke, morele ervaring is, hoeft Kant hier geen aparte deductie voor te leveren. De analyse van het verhevene is al zijn eigen rechtvaardiging: ze toont hoe onze geest zichzelf als doelmatig ervaart wanneer ze geconfronteerd wordt met het oneindige of overweldigende. Het verhevene is dus geen esthetische categorie die om universele instemming vraagt, maar een subjectieve, morele beweging die ons uitdaagt om na te denken over onze plaats in de wereld.

De confrontatie met klimaatverandering illustreert dit onderscheid treffend. De omvang van het probleem — smeltende ijskappen, extreme weersomstandigheden, het uitsterven van soorten — is overweldigend en vormloos, precies zoals het verhevene. Het roept gevoelens van angst en onmacht op, maar confronteert ons ook met onze collectieve morele verantwoordelijkheid. Dit is het verhevene in actie: een ervaring die ons uitnodigt om boven onze individuele bekommernissen uit te stijgen en gezamenlijk te handelen. Tegelijkertijd ervaren we de natuur zelf, die we willen beschermen, vaak als mooi. Deze schoonheid roept een universeel gevoel van waardering op, dat ons motiveert om haar te behouden. Hier wordt zichtbaar hoe schoonheid en het verhevene samenwerken: schoonheid trekt ons aan, spreekt tot ons verlangen naar harmonie, terwijl het verhevene ons uitdaagt en in beweging zet.

Ditzelfde samenspel zien we terug in rituelen, zoals die van de vrijmetselarij. Schoonheid schept daarin een harmonieuze ruimte — door symboliek, architectuur of muziek — die contemplatie en openheid mogelijk maakt. Maar het is het verhevene dat de deelnemer confronteert met diepere vragen, zoals sterfelijkheid, morele plicht of de zoektocht naar betekenis. Deze confrontatie wekt ontzag en een gevoel van morele urgentie op. Zonder schoonheid zou het ritueel te confronterend zijn, te overweldigend; zonder het verhevene zou het louter een esthetische ervaring blijven, zonder diepgang. Pas in hun samenwerking vormen ze een krachtig middel voor zelfreflectie en morele groei.

Uiteindelijk laat §30 zien dat Kant schoonheid en het verhevene niet als gescheiden, maar als complementaire krachten beschouwt. Schoonheid vereist een deductie omdat ze universele instemming opeist; het verhevene niet, omdat het een subjectieve, morele reactie is. Toch werken ze in zowel moderne ervaringen — zoals de uitdagingen van klimaatverandering — als in traditionele rituelen nauw samen. Schoonheid spreekt tot ons verlangen naar verbinding en genot, terwijl het verhevene ons uitdaagt, confronteert en oproept tot actie en reflectie. Samen vormen ze de twee zijden van een volledige esthetische en morele ervaring.

\end{mdframed}

\section{Over de methode van de deductie van smaakoordelen.}

De verplichting om een deductie te leveren, d.w.z. de garantie van de legitimiteit, van een bepaalde soort oordeel doet zich alleen voor als het oordeel een aanspraak maakt op noodzakelijkheid, wat ook het geval is wanneer het subjectieve universaliteit eist, d.w.z. de instemming van iedereen, ondanks het feit dat het geen kennisoordeel is, maar slechts een oordeel over genot of ongenoegen in een gegeven voorwerp, d.w.z. een veronderstelling van een subjectieve doelmatigheid die voor iedereen geldend is, die niet gegrondvest mag zijn in enig begrip van het ding, omdat het een smaakoordeel is.

Aangezien wij in het laatste geval geen kennisoordeel voor ons hebben, noch een theoretisch, gegrondvest in het begrip van een \textbf{natuur} in het algemeen door het verstand, noch een (zuiver) praktisch, gegrondvest in het idee van \textbf{vrijheid} zoals \textit{a priori} door de rede gegeven, en wij dus de geldigheid van noch een oordeel dat voorstelt wat een ding is, noch een oordeel dat ik, om het voort te brengen, iets moet verrichten, \textit{a priori} hoeven te rechtvaardigen, is het alleen de \textbf{universele geldigheid} van een \textbf{enkelvoudig} oordeel, dat de subjectieve doelmatigheid van een empirische voorstelling van de vorm van een voorwerp uitdrukt, die voor de oordeelskracht in het algemeen aangetoond moet worden, om te verklaren hoe het mogelijk is dat iets louter in het oordelen (zonder een sensatie van de zintuigen of een begrip) kan behagen en dat, net zoals het oordelen over een voorwerp ter wille van kennis in het algemeen universele regels heeft, ook de bevrediging van de ene als regel voor alle anderen kan worden aangekondigd.

Nu, als deze universele geldigheid niet gegrondvest mag zijn op het verzamelen van stemmen en het rondvragen bij andere mensen naar de soort sensaties die zij hebben, maar als het ware moet berusten op een autonomie van het subject dat oordeelt over het gevoel van genot in de gegeven voorstelling, d.w.z. op zijn eigen smaak, maar toch ook niet afgeleid mag worden uit begrippen, dan heeft een dergelijk oordeel — zoals het smaakoordeel inderdaad heeft — een tweevoudige en zelfs logische eigenaardigheid: namelijk, \textbf{ten eerste}, universele geldigheid \textit{a priori}, maar niet een logische universaliteit in overeenstemming met begrippen, maar de universaliteit van een enkelvoudig oordeel; \textbf{ten tweede}, een noodzakelijkheid (die altijd op \textit{a priori} gronden moet berusten), die echter niet afhangt van enige \textit{a priori} gronden van bewijs, door middel van de voorstelling waarvan de instemming die het smaakoordeel van iedereen vereist, afgedwongen zou kunnen worden.

De oplossing van deze logische eigenaardigheden, waarin een smaakoordeel verschilt van alle kennisoordelen, als wij hier aanvankelijk abstractie maken van al zijn inhoud, namelijk het gevoel van genot, en slechts de esthetische vorm vergelijken met de vorm van objectieve oordelen, zoals de logica die voorschrijft, zal op zichzelf voldoende zijn voor de deductie van deze ongebruikelijke vermogen. Wij zullen daarom allereerst een voorstelling geven van deze kenmerkende eigenschappen van de smaak, verduidelijkt aan de hand van voorbeelden.

\begin{mdframed}[leftmargin=0pt,rightmargin=0pt,backgroundcolor=blue!6,linecolor=white,innertopmargin=6pt,innerbottommargin=6pt]

Een deductie (een rechtvaardiging van de geldigheid) is alleen nodig is als een oordeel claimt dat het noodzakelijk en universeel geldt, zelfs als het geen kennisoordeel is, maar een oordeel over genot of ongenoegen, zoals bij smaak. Een smaakoordeel is bijzonder omdat het geen kennis over een voorwerp geeft, maar wel eist dat iedereen het ermee eens is. Het is geen theoretisch oordeel (over hoe de wereld is) of een praktisch oordeel (over wat we moeten doen), maar een subjectief oordeel over hoe iets ons bevalt.

Het probleem is: hoe kan zo’n oordeel universeel geldig zijn, zonder gebaseerd te zijn op feiten, zintuiglijke sensaties of begrippen? Kant stelt dat de universele geldigheid van smaak niet kan worden bepaald door bij anderen te peilen wat zij voelen. In plaats daarvan moet het berusten op de autonomie van degene die oordeelt: zijn eigen smaak. Toch heeft het smaakoordeel twee opvallende kenmerken: het claimt universele geldigheid (alsof het voor iedereen zou gelden), maar niet op basis van logische regels of bewijzen. Het eist ook noodzakelijkheid (alsof iedereen het ermee eens moet zijn), zonder dat dit afgedwongen kan worden met argumenten.

Om dit te verklaren, wil Kant eerst de logische structuur van smaakoordelen onderzoeken, los van het gevoel van genot zelf. Hij zal laten zien hoe deze oordelen verschillen van gewone kennisoordelen, en dat zal voldoende zijn om hun bijzondere aard te begrijpen. Hiervoor zal hij de kenmerken van smaak uitleggen en met voorbeelden verduidelijken.

\bigskip
Smaak is niet afhankelijk van externe autoriteit (zoals stemmen of regels), maar op een interne, vrije beoordeling berust die toch universeel geldig claimt te zijn. Het smaakoordeel is zelfwetgevend: het stelt een regel op zonder deze te kunnen afdwingen. Dit is vergelijkbaar met morele oordelen, die ook een universele aanspraak maken zonder empirisch bewijs. Kant wil laten zien dat smaak, net als moraliteit, een vermogen is dat ons verbindt; niet door dwang, maar door een gedeelde structuur van de menselijke geest.

Kant ziet smaak als een voorbereiding op moraliteit: het leert ons om zonder begrippen of regels te oordelen, maar wel met een gevoel voor wat universeel geldig is. In de vrijmetselarij is dit precies wat rituele beoefening doet: het traint het vermogen om zonder dogma’s betekenis te geven aan symbolen en handelingen, terwijl men toch een gedeelde waarheid ervaart.

Het is geen toeval dat zowel Kant als de vrijmetselarij benadrukken dat vrijheid (autonomie) en plicht (universele geldigheid) samengaan. Een vrijmetselaar oefent in het ritueel zijn vermogen om vrij te oordelen, maar altijd met het besef dat zijn ervaring aansluit bij een groter geheel.

\end{mdframed}

\section{Eerste eigenaardigheid van het smaakoordeel.}

Het smaakoordeel bepaalt zijn voorwerp met betrekking tot bevrediging (als schoonheid) met een aanspraak op de instemming van \textbf{iedereen}, alsof het objectief was.

Wanneer iemand zegt: ‘‘Deze bloem is mooi’’, dan is dat hetzelfde als te beweren dat iedereen er bevrediging in zou moeten vinden. Wat de aangenaamheid van zijn geur betreft, heeft de bloem echter helemaal geen aanspraak op algemene instemming. De ene persoon raakt er namelijk in vervoering van, terwijl een ander er duizelig van wordt. Wat moeten we hieruit concluderen, anders dan dat schoonheid als een eigenschap van de bloem zelf beschouwd moet worden, die niet afhangt van de verschillen tussen mensen en hun zintuigen, maar waaraan deze juist moeten beantwoorden als ze erover willen oordelen? Toch is dit niet het geval. Want het smaakoordeel bestaat juist hierin: dat het een ding alleen mooi noemt op grond van die eigenschap waardoor het overeenkomt met onze manier van waarnemen.

Bovendien wordt van elk oordeel dat de smaak van het subject moet bewijzen, geëist dat het subject zelfstandig oordeelt, zonder eerst bij anderen te hoeven rondvragen naar hun oordeel over hetzelfde voorwerp, en zonder zich te baseren op hun bevrediging of onbevrediging. Het oordeel moet dus \textit{a priori} worden uitgesproken, alsof het voorwerp werkelijk universeel bevalt. Men zou denken dat een \textit{a priori} oordeel een begrip van het voorwerp moet bevatten, waarvoor het het principe levert; het smaakoordeel is echter helemaal niet op begrippen gebaseerd en is al helemaal geen kennis, maar louter een esthetisch oordeel.

Daarom laat een jonge dichter zich niet afbrengen van zijn overtuiging dat zijn gedicht mooi is, zelfs niet door het oordeel van het publiek of dat van zijn vrienden. Als hij naar hen luistert, is dat niet omdat hij zijn oordeel verandert, maar omdat hij, zelfs als hij denkt dat het hele publiek slechte smaak heeft, toch (tegen zijn eigen oordeel in) redenen vindt om zich aan de algemene misvatting aan te passen, uit verlangen naar erkenning. Pas later, wanneer zijn oordeelsvermogen door oefening is gescherpt, wijzigt hij zijn eerdere oordeel uit vrije wil, net zoals hij dat zou doen met oordelen die geheel op rede berusten. Smaak eist autonomie. Het oordeel van anderen tot grond van het eigen oordeel maken, zou heteronomie zijn.

Dat de werken van de oudheid terecht worden geprezen als voorbeelden en hun auteurs als klassieken worden beschouwd — als een soort adel onder schrijvers, die door hun voorbeeld wetten voorschrijven aan het volk — lijkt erop te wijzen dat smaak a posteriori bronnen heeft, en dat spreekt de autonomie van smaak in elk subject tegen. Maar men zou net zo goed kunnen zeggen dat de oude wiskundigen, die tot nu toe als bijna onmisbare voorbeelden van de grootste grondigheid en elegantie van de synthetische methode worden gezien, ook een imiterende rede bij ons aantonen, en onze onvermogen om uit eigen bronnen strenge bewijzen te leveren met de grootste aanschouwelijke evidentie, door middel van de constructie van begrippen. Er is helemaal geen gebruik van onze vermogens, hoe vrij ook, en zelfs niet van de rede (die al haar oordelen uit de gemeenschappelijke \textit{a priori} bron moet putten), dat niet in dwaling zou vervallen als anderen niet met hun eigen pogingen waren voorgegaan — niet om hun opvolgers tot loutere navolgers te maken, maar om hen door hun methode op het juiste spoor te zetten bij het zoeken naar principes in zichzelf, en zo hun eigen, vaak betere, weg te laten volgen.
Zelfs in de religie, waar ieder zeker de regel voor zijn gedrag uit zichzelf moet afleiden — omdat hij zelf verantwoordelijk blijft en de schuld voor zijn overtredingen niet op anderen kan afschuiven, of het nu leraren of voorgangers zijn — bereiken algemene voorschriften, die men of van priesters of filosofen heeft geleerd of uit zichzelf heeft afgeleid, nooit zoveel als een voorbeeld van deugd of heiligheid, dat in de geschiedenis is vastgelegd. Zo’n voorbeeld maakt de autonomie van de deugd, voortkomend uit het eigen oorspronkelijke idee van moraliteit (\textit{a priori}), niet overbodig en verandert deze niet in een mechanisme van imitatie. \textbf{Opvolging}, in relatie tot een voorganger, en niet imitatie, is de juiste uitdrukking voor elke invloed die de werken van een voorbeeldig auteur op anderen kunnen hebben. Dit betekent niets anders dan dat men uit dezelfde bronnen put waaruit de voorganger putte, en van hem alleen leert hoe men te werk moet gaan. Maar van alle vermogens en talenten is smaak juist datgene wat, omdat zijn oordeel niet door begrippen en voorschriften bepaald kan worden, het meest behoefte heeft aan voorbeelden van wat in de loop der cultuur het langst goedkeuring heeft genoten, wil het niet snel terugvallen in barbarisme en verworden tot de ruwheid van zijn eerste pogingen.

\begin{mdframed}[leftmargin=0pt,rightmargin=0pt,backgroundcolor=blue!6,linecolor=white,innertopmargin=6pt,innerbottommargin=6pt]

De eerste opvallende eigenschap van het smaakoordeel: het doet een aanspraak op universele instemming, alsof schoonheid een objectieve eigenschap is van het voorwerp zelf, zoals wanneer iemand zegt: "Deze bloem is mooi" en daarmee impliciet stelt dat iedereen die bloem mooi zou moeten vinden. Dit in tegenstelling tot iets subjectiefs als geur, waar de ene persoon van geniet en de andere misselijk van wordt. Toch is schoonheid geen objectieve eigenschap, zoals kleur of grootte. Het smaakoordeel gaat niet over hoe het voorwerp is, maar over hoe het door ons wordt waargenomen en hoe het overeenkomt met onze manier van ervaren. Het is dus geen kennisoverdragend oordeel (zoals "deze bloem is rood"), maar een esthetisch oordeel dat claimt dat iedereen het ermee eens moet zijn.

Een cruciaal punt is dat een smaakoordeel autonoom moet zijn: degene die oordeelt, doet dat uit zichzelf, zonder afhankelijk te zijn van de meningen van anderen. Een jonge dichter die overtuigd is van de schoonheid van zijn gedicht, zal zich niet zomaar laten verleiden door kritiek van het publiek of vrienden. Als hij naar anderen luistert, is dat niet omdat hij zijn eigen oordeel bijstelt, maar misschien uit praktische overwegingen, bijvoorbeeld om erkenning te krijgen, zelfs als hij denkt dat het publiek geen smaak heeft. Pas later, wanneer zijn oordeelsvermogen zich heeft ontwikkeld, kan hij zijn eerdere mening herzien, net zoals hij dat zou doen bij een oordeel dat op rede berust. Het overnemen van andermans oordeel zou heteronomie zijn: het eigen oordeel laten bepalen door extern gezag, in plaats van door eigen inzicht.

Kant vergelijkt dit met andere gebieden, zoals wiskunde en religie. Ook daar leren we van voorbeelden — zoals de klassieke wiskundigen of heilige figuren uit de geschiedenis — maar het gaat niet om blindelings navolgen. Een voorbeeld dient als leidraad, niet als voorschrift. Het helpt ons om onze eigen principes te ontdekken en onze eigen weg te vinden. In de wiskunde leert men van voorgangers hoe men strenge bewijzen kan opbouwen, maar men imitieert ze niet; men leert de methode om zelfstandig tot inzichten te komen. Hetzelfde geldt voor moraliteit: hoewel we ons kunnen laten inspireren door voorbeelden van deugd, blijft ieder zelf verantwoordelijk voor zijn eigen handelen. De waarde van een voorbeeld ligt niet in het klakkeloos overnemen, maar in het leren hoe men zelf tot morele inzichten kan komen.

Voor smaak geldt dit nog sterker dan voor andere vermogens. Omdat smaak niet kan worden vastgelegd in regels of begrippen, is het afhankelijk van voorbeelden uit de cultuur — zoals de klassieke werken uit de oudheid — om niet terug te vallen in onbeholpenheid. Deze voorbeelden fungeren niet als wetten, maar als richtingwijzers die helpen om het eigen oordeelsvermogen te scherpen. Zonder dergelijke voorbeelden zou smaak snel vervallen tot willekeur of barbarisme. Het gaat erom dat men leert om uit dezelfde bronnen van schoonheid en betekenis te putten als de grootmeesters, zonder hun stijl een-op-een te kopiëren. Op die manier blijft smaak een vrij, maar geïnformeerd vermogen.

\bigskip
Kants §32 onthult een fascinerende spanning in het smaakoordeel: hoewel het subjectief is, claimt het universele geldigheid, alsof schoonheid een objectieve eigenschap is die iedereen moet herkennen. Deze paradox wijst op een diepgeworteld vertrouwen in een gedeelde structuur van waarneming, een idee dat later in Kants deductie centraal zal staan. Het smaakoordeel is geen willekeurige voorkeur, maar een vrije ervaring die toch aanspraak maakt op intersubjectiviteit — alsof onze persoonlijke beleving van schoonheid aansluit bij een gemeenschappelijke menselijke grond.

De nadruk op autonomie is cruciaal: smaak vereist dat we zelfstandig oordelen, zonder ons te laten leiden door de meningen van anderen. Dit is geen arbitraire koppigheid, maar een wezenlijk kenmerk, vergelijkbaar met morele autonomie. Net zoals we in de ethiek leren om zelfstandig morele keuzes te maken, leert smaak ons om betekenis te geven zonder afhankelijk te zijn van externe autoriteit. De jonge dichter die vasthoudt aan zijn oordeel, zelfs tegen de stroom in, illustreert dit: smaak is geen kwestie van meerderheidsstemming, maar van een innerlijk kompas dat vertrouwt op een gedeelde, maar niet afdwingbare waarheid.

Voorbeelden — zoals klassieke kunstwerken of literaire meesterwerken — spelen hierin een dubbele rol. Ze zijn geen voorschriften, maar leidraden die ons helpen ons eigen oordeelsvermogen te ontwikkelen. Net zoals in wiskunde of moraliteit gaat het niet om imitatie, maar om het ontdekken van principes die in ons zelf liggen. Voor smaak is dit zelfs nog essentiëler, omdat er geen vaste regels zijn. Zonder cultuurhistorische voorbeelden zou smaak vervallen tot willekeur, maar blindelings navolgen zou het vermogen tot eigen oordeel ondermijnen. Het is een delicaat evenwicht: leren van traditie, terwijl men toch vrij en verantwoordelijk blijft voor zijn eigen ervaring.

Deze dynamiek doet sterk denken aan de maçonnieke benadering van symboliek. Ook daar gaat het niet om voorgeschreven betekenissen, maar om een persoonlijke, maar herkenbare interpretatie — een dialoog tussen individu en gedeelde bronnen van betekenis. Net als bij Kant is het doel niet om een "juist" antwoord te vinden, maar om te leren hoe men betekenis kan geven, binnen een kader dat zowel vrijheid als verbondenheid mogelijk maakt. Het ritueel en de symbolen fungeren als voorbeelden die niet voorschrijven, maar uitnodigen tot een eigen, doch gedeelde ervaring van waarheid en schoonheid. Zo wordt smaak, net als moreel inzicht, een oefening in autonomie binnen een gemeenschappelijke traditie.
\end{mdframed}

\section{Tweede eigenaardigheid van het smaakoordeel.}

Het smaakoordeel is geheel niet bepaalbaar door gronden van bewijs, alsof het slechts \textbf{subjectief} was.

Als iemand een gebouw, een gezicht, of een gedicht niet mooi vindt, dan laat hij, \textbf{ten eerste}, zich niet innerlijk een goedkeuring opleggen door honderd stemmen die het allemaal hoog prijzen. Hij kan zich wel zo gedragen alsof het hem ook bevalt, om niet als smaakloos te worden beschouwd, hij kan zelfs beginnen te twijfelen of hij zijn smaak voldoende heeft gevormd door kennisname van een voldoende aantal voorwerpen van een bepaalde soort, (zoals iemand die meent in de verte iets als een bos te herkennen, wat iedereen anders als een stad beschouwt, twijfelt aan het oordeel van zijn eigen ogen). Maar wat hij wel duidelijk ziet, is dit: dat de goedkeuring van anderen geen geldig bewijs levert voor het oordelen over schoonheid, dat anderen misschien voor hem kunnen zien en waarnemen, en dat wat velen op één manier hebben gezien, wat hij meent anders te hebben gezien, hem als voldoende grond van bewijs kan dienen voor een theoretisch, dus een logisch oordeel, maar dat wat anderen heeft behaagd, nooit als grond kan dienen voor een esthetisch oordeel. Het oordeel van anderen, wanneer het ongunstig is voor het onze, kan ons wel terecht aarzelingen geven over ons eigen oordeel, maar kan ons nooit overtuigen van zijn onjuistheid. Zo is er geen empirische \textbf{grond van bewijs} om het oordeel aan iemand op te leggen.

\textbf{Ten tweede}, kan een \textit{a priori} bewijs volgens bepaalde regels het oordeel over schoonheid nog minder bepalen. Als iemand mij zijn gedicht voorleest of mij meeneemt naar een toneelstuk dat uiteindelijk mijn smaak niet bevalt, dan kan hij \textbf{Batteux} of \textbf{Lessing} aanhalen, of zelfs oudere en beroemdere critici van smaak, en al hun regels die zij hebben opgesteld als bewijzen aanvoeren dat zijn gedicht mooi is; bepaalde passages, die nu juist de passages zijn die mij niet bevallen, kunnen zelfs overeenstemmen met regels van schoonheid, (zoals die daar zijn gegeven en algemeen erkend zijn): ik zal mijn oren dichthouden, naar geen redenen en argumenten luisteren, en zou liever geloven dat die regels van de critici onjuist zijn, of ten minste dat dit geen geval is voor hun toepassing, dan toelaten dat mijn oordeel bepaald wordt door middel van \textit{a priori} gronden van bewijs, aangezien het een smaakoordeel zou moeten zijn en geen oordeel van het verstand of de rede.

Het schijnt dat dit een van de hoofdredenen is, waarom aan dit vermogen van esthetisch oordelen juist de naam ‘‘smaak’’ is gegeven. Want iemand kan mij alle bestanddelen van een gerecht opsommen, en over elk opmerken dat het mij anders wel bevalt, en bovendien zelfs terecht de gezondheid van dit voedsel prijzen; toch ben ik doof voor al deze gronden, ik proef het gerecht met \textbf{mijn} tong en mijn gehemelte, en op die grond, (niet op grond van algemene principes), doe ik mijn oordeel uit.

Inderdaad, het smaakoordeel wordt altijd gedaan als een enkelvoudig oordeel over het voorwerp. Het verstand kan een universeel oordeel vellen door te vergelijken hoe bevredigend het voorwerp is met de oordelen van anderen, bijvoorbeeld: alle tulpen zijn mooi; maar in dat geval is dat geen smaakoordeel, maar een logisch oordeel, dat de verhouding van een voorwerp tot de smaak tot een predikaat maakt van dingen van een bepaalde soort in het algemeen; maar datgene waardoor ik een enkele gegeven tulp mooi vind, dat is, mijn bevrediging daarin universeel geldig acht, is het smaakoordeel alleen. Zijn eigenaardigheid bestaat echter hierin: dat het, hoewel het slechts subjectieve geldigheid heeft, niettemin een aanspraak maakt op \textbf{alle} subjecten van een soort die alleen gemaakt zou kunnen worden als het een objectief oordeel was, dat op kennissgronden berust en door middel van een bewijs afgedwongen kan worden.

\begin{mdframed}[leftmargin=0pt,rightmargin=0pt,backgroundcolor=blue!6,linecolor=white,innertopmargin=6pt,innerbottommargin=6pt]

De tweede opvallende eigenschap van het smaakoordeel is dat het laat zich niet overtuigen door bewijzen, of die nu gebaseerd zijn op de meningen van anderen of op vaste regels. Dit onderscheidt smaak fundamenteel van andere soorten oordelen, zoals theoretische of logische oordelen, die wel met argumenten kunnen worden onderbouwd.

Wanneer iemand een gebouw, landschap of gedicht niet mooi vindt, maakt het niet uit hoe veel mensen het wel waarderen. Hij zal zich niet zomaar laten overtuigen door hun lof, zelfs als hij uit beleefdheid doet alsof hij het ermee eens is. Hij kan hoogstens gaan twijfelen aan zijn eigen oordeelsvermogen, net zoals iemand die in de verte een bos meent te zien, terwijl anderen zeggen dat het een stad is, aan zijn eigen waarneming gaat twijfelen. Maar één ding staat voor hem vast: de goedkeuring van anderen is geen geldig bewijs voor schoonheid. Wat anderen zien of ervaren, kan wel helpen bij een theoretisch oordeel (bijvoorbeeld: "Dit is een klassiek gebouw"), maar nooit bij een esthetisch oordeel (bijvoorbeeld: "Dit gebouw is mooi"). Kritiek van anderen kan hem wel aan het denken zetten, maar zijn eigen smaakoordeel kan niet weerlegd worden door andermans mening. Er is simpelweg geen empirisch bewijs dat iemand kan dwingen om iets mooi te vinden.

Ook \textit{a priori} bewijzen, zoals regels uit de kunsttheorie, hebben geen vat op het smaakoordeel. Stel dat iemand een gedicht voorleest dat jou niet aanspreekt, en hij haalt beroemde critici als Batteux of Lessing aan om te bewijzen dat het wel mooi is. Zelfs als bepaalde passages voldoen aan algemeen erkende schoonheidsregels, zal dat jou niet overtuigen. Je zou eerder denken dat die regels niet kloppen, of niet van toepassing zijn op dit gedicht, dan dat je je eigen smaak aanpast. Het smaakoordeel is immers geen kwestie van verstand of rede, maar van directe, persoonlijke ervaring. Je kunt niet met argumenten worden overtuigd om iets mooi te vinden; je moet het zelf voelen.
Kant suggereert dat dit de reden is waarom we dit vermogen ‘smaak’ noemen. Net zoals je de smaak van een gerecht niet kunt bewijzen met een lijstje ingrediënten of voedingswaarden, kun je schoonheid niet aantonen met regels of theorieën. Je proeft het gerecht en oordeelt op basis van je eigen zintuiglijke ervaring, niet op grond van algemene principes. Het smaakoordeel is altijd een enkelvoudig oordeel over een specifiek voorwerp: "Deze tulp is mooi", niet "Alle tulpen zijn mooi". Dat laatste zou een logisch oordeel zijn, gebaseerd op een algemene categorie, maar een echt smaakoordeel gaat over jouw directe bevrediging bij dit specifieke voorwerp, en die claimt toch universele geldigheid.
Hier ligt de eigenaardigheid: hoewel het smaakoordeel subjectief is (het gaat om jouw ervaring), doet het alsof het objectief is. Het eist namelijk dat iedereen het ermee eens zou moeten zijn, alsof het op kennis berust en afgedwongen kan worden met bewijzen. Maar in werkelijkheid is er geen objectieve grond voor die eis, het is een vrije, maar bindende claim die het smaakoordeel vanzelf maakt.

\bigskip
§33 onthult de radicale subjectiviteit van het smaakoordeel, maar dan op een manier die juist zijn unieke kracht blootlegt. Het weigert zich te laten determineren door elke vorm van externe autoriteit, of die nu komt van de menigte, de traditie, of de rede zelf. Dit is geen zwakte, maar een wezenlijk kenmerk: smaak is geen kwestie van bewijs, maar van onmiddellijke ervaring. Hierin schuilt een diepe paradox: hoewel het smaakoordeel zich onttrekt aan objectieve gronden, maakt het toch een aanspraak die alsof het objectief is. Het is alsof de menselijke geest, wanneer ze schoonheid ervaart, niet alleen een persoonlijke voorkeur uitspreekt, maar tegelijkertijd aannemt dat die ervaring gedeeld zou moeten kunnen worden, niet omdat het bewijsbaar is, maar omdat het voelt als een waarheid die voor iedereen geldt.

De vergelijking met proeven is veelzeggend. Net zoals je niet kunt bewijzen dat een gerecht lekker is door de ingrediënten op te sommen, kun je schoonheid niet afleiden uit regels of consensus. Het is een lichaamlijke, directe kennis, niet in de zin van zintuiglijke prikkels (zoals zoet of zuur), maar als een esthetische zekerheid die voorbij redeneren gaat. Dit verklaart waarom we ons niet laten overleiden door critici of meerderheidsmeningen: smaak is geen meningsverschil, maar een soevereine ervaring. Toch is het geen louter persoonlijke gril. Het feit dat we ons wel kunnen laten beïnvloeden door twijfel ("Misschien heb ik onvoldoende voorbeelden gezien") toont aan dat smaak zich openstelt voor ontwikkeling, zonder zijn autonomie prijs te geven. Het is alsof we intuïtief weten dat onze ervaring aansluit bij een mogelijk gedeelde werkelijkheid, zelfs als we die niet kunnen bewijzen.

Deze weerstand tegen bewijzen onthult iets fundamenteels over hoe wij betekenis geven. In tegenstelling tot wetenschappelijke of morele oordelen, die berusten op principes of argumenten, is smaak een vrije, maar bindende vorm van kennis. Het is alsof Kant hier zegt: Schoonheid is geen zaak van weten, maar van herkennen, en dat herkennen is zowel persoonlijk als universeel gericht. Dit doet denken aan de maçonnieke benadering van symbolen: ook daar gaat het niet om vastomlijnde betekenissen, maar om een persoonlijke, maar herkenbare ervaring van waarheid. Een vrijmetselaar die een symbool interpreteert, doet dat niet op gezag van een leermeester of een handboek, maar uit een eigen, directe confrontatie met de betekenis ervan. Toch voelt die interpretatie alsof ze aansluit bij een groter geheel, niet omdat ze voorgeschreven is, maar omdat ze voortkomt uit een gedeelde bron van inzicht.

Wat §33 vooral blootlegt, is dat smaak een vorm van kennis is die niet kan worden afgedwongen, maar wel kan worden gescherpt. Het is geen toeval dat Kant smaak vergelijkt met proeven: net zoals een fijnproever zijn vermogen ontwikkelt door te oefenen (zonder ooit "objectief" te kunnen bewijzen wat lekker is), leert de mens schoonheid herkennen door ervaring en reflectie. De weigering om zich te laten overleiden is dus geen arrogantie, maar een voorwaarde voor echtheid. Pas wanneer we ons oordeel niet laten dicteren, kunnen we werkelijk leren zien, en dat is precies wat zowel esthetica als maçonnieke symboliek vragen: een vrije, maar gevoelige geest.

\end{mdframed}

\section{Geen objectief principe van smaak is mogelijk.}

Onder een principe van smaak zou men een fundamentele stelling verstaan, onder de voorwaarde waarvan men het begrip van een voorwerp zou kunnen subsumeren en vervolgens door middel van een redenering concluderen dat het mooi is. Maar dat is volkomen onmogelijk. Want ik moet het genot onmiddellijk in de voorstelling ervan gewaarworden, en ik kan er niet door enige bewijzen toe worden overgehaald. Zo kunnen critici, zoals Hume zegt, wel plausibeler redeneren dan koks, maar zij ondergaan toch hetzelfde lot als deze. Zij kunnen geen bepalende grond voor hun oordeel verwachten van bewijzen, maar alleen van de reflectie van het subject op zijn eigen toestand (van genot of ongenoegen), waarbij alle voorschriften en regels worden verworpen.

Wat critici niettemin kunnen en moeten redeneren, op een wijze die nuttig is voor het corrigeren en verbreden van onze smaakoordelen, is dit: niet de uiteenzetting van de bepalende grond van deze soort esthetische oordelen in een universeel bruikbare formule, wat onmogelijk is, maar het onderzoek naar de kennisvermogens en hun functies in deze oordelen en het uiteenzetten aan de hand van voorbeelden van de onderlinge subjectieve doelmatigheid, waarvan boven is aangetoond dat haar vorm in een gegeven voorstelling de schoonheid van haar voorwerp is. Zo is de kritiek van de smaak zelf alleen subjectief, met betrekking tot de voorstelling door middel waarvan een voorwerp ons gegeven wordt: dat is, het is de kunst of wetenschap om de wederzijdse verhouding van het verstand en de verbeelding tot elkaar in de gegeven voorstelling (zonder betrekking tot een voorafgaand gevoel of begrip) onder regels te brengen, en derhalve hun overeenstemming of strijd te bepalen en deze met betrekking tot haar voorwaarden vast te stellen. Het is \textbf{kunst}, als het dit alleen in voorbeelden aantoont; het is \textbf{wetenschap}, als het de mogelijkheid van een dergelijk oordelen afleidt uit de aard van dit vermogen als een kennisvermogen in het algemeen. Het is met dit laatste, als transcendentale kritiek, dat wij hier uitsluitend te maken hebben.

Zij moet het subjectieve principe van smaak ontwikkelen en rechtvaardigen als een \textit{a priori} principe van de oordeelskracht. Kritiek, als kunst, zoekt slechts de fysiologische (hier psychologische) en derhalve empirische regels toe te passen, volgens welke de smaak feitelijk te werk gaat bij het oordelen over haar voorwerpen (zonder na te denken over haar mogelijkheid), en beoordeelt de producten van de schone kunst, net zoals de \textbf{eerste} de faculteit om ze te beoordelen zelf beoordeelt.

\begin{mdframed}[leftmargin=0pt,rightmargin=0pt,backgroundcolor=blue!6,linecolor=white,innertopmargin=6pt,innerbottommargin=6pt]

Er kan geen objectief principe voor smaak bestaan. Met een principe bedoelt Kant een vaste regel of stelling waarmee je zou kunnen aantonen waarom iets mooi is, alsof je schoonheid kunt bewijzen door een voorwerp onder een algemene definitie te scharen. Dat is onmogelijk, omdat schoonheid geen kwestie is van redeneren of bewijzen, maar van directe ervaring: je moet het genot zelf voelen wanneer je iets waarneemt. Critici kunnen wel overtuigender argumenteren dan koks, maar uiteindelijk kunnen ze net als koks niet met logica of regels afdwingen dat iemand iets mooi vindt. Het oordeel over schoonheid berust niet op argumenten, maar op hoe het subject zelf reageert, op zijn persoonlijke bevrediging of afkeer.

Toch hebben critici wel een belangrijke taak. Ze kunnen niet bewijzen wat mooi is, maar ze kunnen wel helpen om onze smaakoordelen te verfijnen en verbreden. Dat doen ze niet door een algemene formule voor schoonheid te bedenken (wat niet kan), maar door te onderzoeken hoe onze kennisvermogens, met name het verstand en de verbeelding samenwerken wanneer we iets mooi vinden. Ze kunnen laten zien hoe deze vermogens in een voorstelling harmonisch samengaan (of juist niet), en dat is precies wat schoonheid uitmaakt: de subjectieve doelmatigheid van een voorstelling.

Kant maakt hierbij een onderscheid tussen twee soorten kritiek:
\begin{enumerate}
    \item \textbf{Kritiek als kunst}: Hierbij gaat het om het analyseren van voorbeelden, bijvoorbeeld door te laten zien hoe bepaalde kunstwerken de verbeelding en het verstand in evenwicht brengen.

    \item \textbf{Kritiek als wetenschap}: Dit is wat Kant zelf doet in zijn transcendentale kritiek. Hij onderzoekt niet wat mensen mooi vinden, maar hoe smaakoordelen überhaupt mogelijk zijn. Hij wil aantonen dat smaak berust op een \textit{a priori} principe van de oordeelskracht, een vermogen dat we allemaal hebben, maar dat niet op vaste regels of voorschriften berust.
\end{enumerate}

Empirische critici (die kijken naar hoe mensen feitelijk oordelen) gebruiken psychologische of fysiologische regels om kunstwerken te beoordelen. Maar Kants benadering is fundamenteel: hij wil begrijpen waarom we überhaupt in staat zijn om esthetische oordelen te vellen, zonder afhankelijk te zijn van externe regels of gevoelens.

\bigskip
§34 markeert een beslissend moment in Kants esthetica, waar hij niet alleen uitlegt wat smaak niet is, maar ook waarom dat juist haar kracht uitmaakt. Het onvermogen om schoonheid in objectieve principes te vangen is geen tekortkoming, maar een wezenlijk kenmerk van esthetische ervaring. Waar wetenschap en moraliteit berusten op redeneren en bewijzen, onttrekt smaak zich aan elke vorm van afdwingbare logica. Dit is geen zwakte, maar een eigen soort zekerheid: schoonheid is geen kwestie van weten, maar van herkennen -- en dat herkennen is zowel onmiddellijk als universeel gericht. Het is alsof Kant hier zegt: de waarheid van schoonheid ligt niet in het voorwerp zelf, noch in algemene regels, maar in de wisselwerking tussen onze geest en de wereld.

De kritische taak die Kant schetst is dubbel. Enerzijds is er de empirische benadering, die voorbeelden analyseert en probeert te beschrijven hoe smaak in de praktijk werkt. Maar dat is slechts een voorbereiding op het echte werk: de transcendentale kritiek, die onderzoekt hoe smaakoordelen überhaupt mogelijk zijn. Dit is geen louter psychologische vraag, maar een filosofisch fundament: smaak berust op de harmonie tussen verbeelding en verstand, een vermogen dat in elke mens aanwezig is, maar dat zich niet laat vangen in voorschriften. Hierin ligt de verbinding met het maçonnieke idee van symboliek: ook daar gaat het niet om vaste betekenissen, maar om het leren zien van betekenis in een vrije, maar gestructureerde ervaring. Net zoals een vrijmetselaar leert om symbolen te interpreteren zonder dat hem een ``juist'' antwoord wordt voorgeschreven, leert de mens schoonheid herkennen door de eigen geest te oefenen in het waarnemen van harmonie.

Wat §34 vooral onthult, is dat smaak een vrije, maar niet willekeurige activiteit is. Critici kunnen wel helpen om ons oordeel te verfijnen, maar de uiteindelijke bevestiging komt altijd van binnenuit -- van het subject dat zelf de overeenstemming tussen verbeelding en verstand ervaart. Dit is geen subjectivisme, maar een intersubjectieve claim: we gaan ervan uit dat anderen, als ze dezelfde vermogens bezitten, tot eenzelfde ervaring zouden kunnen komen. Zo wordt smaak een brug tussen het persoonlijke en het universele, niet door dwang, maar door een gedeelde structuur van waarnemen. Het is deze spanning -- tussen autonomie en herkenbaarheid -- die zowel Kants esthetica als de maçonnieke zoektocht naar betekenis kenmerkt.

\end{mdframed}

\section{Het principe van de smaak is het subjectieve principe van de oordeelskracht in het algemeen.}

Het smaakoordeel verschilt van het logische oordeel doordat het laatste een voorstelling subsumpeert onder begrippen van het voorwerp, maar het eerste helemaal niet subsumpeert onder een begrip, want anders zou de noodzakelijke universele instemming door bewijzen afgedwongen kunnen worden. Niettemin is het ermee verwant doordat het een universaliteit en noodzakelijkheid claimt, hoewel niet in overeenstemming met begrippen van het voorwerp, en dus slechts een subjectieve. Nu de begrippen in een oordeel de inhoud ervan uitmaken (wat betrekking heeft op de kennis van het voorwerp), maar het smaakoordeel niet bepaalbaar is door middel van begrippen, berust het slechts op de subjectieve formele voorwaarde van een oordeel in het algemeen. De subjectieve voorwaarde van alle oordelen is het vermogen om te oordelen zelf, ofwel de oordeelskracht. Deze, toegepast op een voorstelling door middel waarvan een voorwerp gegeven wordt, vereist de overeenstemming van twee voorstellingsvermogens: namelijk de verbeelding (voor de aanschouwing en de samenstelling van de veelheid der aanschouwing), en het verstand (voor het begrip als voorstelling van de eenheid van deze samenstelling). Nu er hier geen begrip van het voorwerp de grond van het oordeel is, kan het slechts bestaan in de subsumptie van de verbeelding zelf (in het geval van een voorstelling door middel waarvan een voorwerp gegeven wordt) onder de voorwaarde dat het verstand in het algemeen van aanschouwingen tot begrippen overgaat. Dat wil zeggen, aangezien de vrijheid van de verbeelding juist hierin bestaat dat zij schematiseert zonder een begrip, moet het smaakoordeel berusten op een louter gevoel van de wederzijds belevende verbeelding in haar \textbf{vrijheid} en het verstand met zijn \textbf{wetmatigheid}, dus op een gevoel dat het voorwerp toelaat te worden beoordeeld volgens de doelmatigheid van de voorstelling (door middel waarvan een voorwerp gegeven wordt) voor de bevordering van het kennisvermogen in zijn vrije spel; en de smaak, als een subjectief vermogen van oordelen, bevat een principe van subsumptie, niet van aanschouwingen onder \textbf{begrippen}, maar van het \textbf{vermogen} van aanschouwingen of voorstellingen (dat is, van de verbeelding) onder het \textbf{vermogen} van begrippen (dat is, het verstand), voor zover de eerste \textbf{in haar vrijheid} in harmonie is met de laatste \textbf{in zijn wetmatigheid}.

Nu, om deze rechtvaardigingsgrond door middel van een deductie van smaakoordelen te ontdekken, kunnen slechts de formele eigenaardigheden van deze soort oordelen, dat is, voor zover het slechts hun logische vorm is die beschouwd wordt, ons als leidraad dienen.

\begin{mdframed}[leftmargin=0pt,rightmargin=0pt,backgroundcolor=blue!6,linecolor=white,innertopmargin=6pt,innerbottommargin=6pt]

§35 onderscheidt het smaakoordeel fundamenteel van het logische oordeel. Waar een logisch oordeel een voorstelling onder een begrip van het voorwerp plaatst, doet het smaakoordeel dit niet, omdat schoonheid niet met bewijzen kan worden afgedwongen. Toch claimt het smaakoordeel wel universaliteit en noodzakelijkheid, maar dan op subjectieve grond. Omdat het smaakoordeel niet bepaald wordt door begrippen, berust het op de formele voorwaarde van oordelen in het algemeen: de wisselwerking tussen verbeelding en verstand.

De verbeelding composeert de veelheid van aanschouwingen, terwijl het verstand de eenheid van die compositie als begrip voorstelt. In het smaakoordeel is er echter geen vast begrip van het voorwerp dat als grond dient. In plaats daarvan wordt de verbeelding zelf - in haar vrije schematisering zonder begrip - onder de algemene voorwaarde gebracht dat het verstand van aanschouwingen tot begrippen overgaat. Het smaakoordeel berust daarom op het gevoel van een harmonieuze wisselwerking tussen de vrije verbeelding en het wetmatige verstand. Deze harmonie bevordert het kennisvermogen in zijn vrije spel en vormt zo het subjectieve principe van smaak: niet het onderbrengen van aanschouwingen onder begrippen, maar het onderbrengen van de verbeelding onder het verstand, voor zover beide vermogens in evenwicht zijn.

Voor de deductie van smaakoordelen moeten we ons richten op de formele eigenaardigheden van deze oordelen, met name hun logische vorm. Deze vorm is de sleutel om de rechtvaardigingsgrond van smaak te begrijpen.

\bigskip
§35 onthult de diepere structuur van het smaakoordeel als een unieke vorm van kennis die noch louter subjectief noch objectief is. Waar logische oordelen berusten op het onderbrengen van voorstellingen onder vaste begrippen, functioneert smaak juist in de vrije ruimte tussen verbeelding en verstand - een ruimte die toch universaliteit claimt zonder afdwingbare regels. Deze spanning is geen tegenspraak, maar de kern van esthetische ervaring: smaak is geen kennis in traditionele zin, maar een vrije, doch gestructureerde harmonie tussen onze kennisvermogens.

De essentie ligt in het begrip van schematisering zonder vast begrip. Dit doet denken aan de maçonnieke benadering van symbolen, waar betekenis niet voorgeschreven is, maar ontstaat in de wisselwerking tussen individuele ervaring en gedeelde structuur. De verbeelding composeert vrij, het verstand brengt ordening, en in hun samenspel ontstaat schoonheid als een ervaring die zowel persoonlijk als intersubjectief geldig is. Deze dynamiek tussen vrijheid en wetmatigheid is geen toeval, maar een noodzakelijke voorwaarde voor esthetische beoordeling.

Wat deze paragraaf vooral blootlegt, is dat smaak berust op een fundamentele menselijke capaciteit: het vermogen om in de afwezigheid van vaste regels toch betekenisvol te oordelen. Dit principe vindt zijn parallel in de maçonnieke traditie, waar symbolen niet eenduidig gedefinieerd zijn, maar binnen een rituele structuur een persoonlijke, maar herkenbare betekenis krijgen. Zoals het ritueel zowel ruimte laat voor individuele interpretatie als een kader biedt voor gedeelde ervaring, zo is smaak een oefening in het vinden van evenwicht tussen subjectieve beleving en universele aanspraak. De deductie die Kant voorbereidt, zal moeten laten zien hoe deze subjectieve harmonie toch een claim op algemene geldigheid kan maken - niet door externe dwang, maar door de gemeenschappelijke structuur van de menselijke geest.

\end{mdframed}

\section{Over het probleem voor een deductie van smaakoordelen.}

De waarneming van een voorwerp kan onmiddellijk verbonden worden met het begrip van een voorwerp in het algemeen, waaraan de eerste de empirische predikaten levert, voor een kennisoordeel, en zo kan een ervaringsoordeel tot stand komen. Dit is gegrondvest in \textit{a priori} begrippen van de synthetische eenheid van de veelheid, om het te denken als de bepaling van een voorwerp; en deze begrippen (de categorieën) vereisen een deductie, die bovendien is gegeven in de \emph{Kritiek van de Zuivere Rede}, waardoor de oplossing van het probleem "Hoe zijn synthetische \textit{a priori} kennisoordelen mogelijk?" is geboden. Dit probleem betrof dus de \textit{a priori} principes van het zuivere verstand en zijn theoretische oordelen.

Een waarneming kan echter ook onmiddellijk verbonden worden met een gevoel van genot (of ongenoegen) en een bevrediging die de voorstelling van het voorwerp begeleidt en deze in plaats van een predikaat dient, en zo kan een esthetisch oordeel ontstaan, dat geen kennisoordeel is. Een dergelijk oordeel, als het niet slechts een gevoelsoordeel is maar een formeel reflectie-oordeel, dat deze bevrediging van iedereen als noodzakelijk vereist, moet gegrondvest zijn in iets als een \textit{a priori} principe, al is het slechts een louter subjectief principe (als een objectief principe voor deze soort oordeel onmogelijk zou zijn), maar dat als zodanig principe eveneens een deductie vereist, door middel waarvan begrepen kan worden hoe een esthetisch oordeel aanspraak kan maken op noodzakelijkheid. Dit is de grond van het probleem waarmee wij ons nu bezighouden: Hoe zijn smaakoordelen mogelijk? Dit probleem betreft dus de \textit{a priori} principes van de zuivere oordeelskracht in esthetische oordelen, dat wil zeggen in die oordelen waarin zij niet (zoals in theoretische oordelen) slechts onder objectieve begrippen van het verstand hoeft te subsumeren en onder een wet staat, maar waarin zij zelf, subjectief, zowel voorwerp als wet is.

Dit probleem kan ook als volgt voorgesteld worden: Hoe is een oordeel mogelijk dat, louter uit het \textbf{eigen} gevoel van genot in een voorwerp, onafhankelijk van zijn begrip, dit genot, als verbonden met de voorstelling van hetzelfde voorwerp \textbf{in elk ander subject}, \textit{a priori}, dat wil zeggen zonder het akkoord van anderen af te hoeven wachten, beoordeelt?

Dat smaakoordelen synthetisch zijn, is gemakkelijk in te zien, omdat zij verder gaan dan het begrip en zelfs de aanschouwing van het voorwerp, en daaraan als predikaat iets toevoegen dat niet eens kennis is, namelijk het gevoel van genot (of ongenoegen). Dat dergelijke oordelen, hoewel het predikaat (van het \textbf{eigen} genot dat met de voorstelling verbonden is) empirisch is, niettemin, voor zover het gaat om de vereiste instemming \textbf{van iedereen}, \textit{a priori} oordelen zijn, of als zodanig zouden worden beschouwd, is reeds impliciet in de uitdrukkingen van hun aanspraak; en zo behoort dit probleem van de kritiek van de oordeelskracht tot het algemene probleem van de transcendentale filosofie: Hoe zijn synthetische \textit{a priori} oordelen mogelijk?

\begin{mdframed}[leftmargin=0pt,rightmargin=0pt,backgroundcolor=blue!6,linecolor=white,innertopmargin=6pt,innerbottommargin=6pt]

Hoe zijn smaakoordelen mogelijk als synthetische a priori oordelen? Terwijl kennisoordelen waarnemingen onder algemene begrippen (categorieën) brengen en hun mogelijkheid in de \textit{Kritiek van de Zuivere Rede} is aangetoond, werpt deze paragraaf de vraag op hoe esthetische oordelen - die geen kennis uitdrukken maar een gevoel van genot of ongenoegen - toch aanspraak kunnen maken op universele geldigheid.

Het probleem wordt tweevoudig geformuleerd:

\begin{enumerate}
    \item Hoe kan een oordeel dat louter gebaseerd is op persoonlijk genot in een voorwerp, zonder begrip van dat voorwerp, toch eisen dat iedereen hetzelfde genot zou moeten ervaren?
    
    \item Hoe kan zo'n oordeel, hoewel het predikaat (het genot) empirisch is, toch a priori gelden voor alle subjecten?
    
\end{enumerate}

Smaakoordelen zijn synthetisch omdat ze verder gaan dan louter waarneming of begrip: ze voegen een subjectief gevoel toe als predikaat. Hun bijzondere status ligt hierin dat ze, ondanks hun subjectieve grond, een noodzakelijke instemming van allen claimen. Dit plaatst ze binnen het algemene transcendentale vraagstuk naar de mogelijkheid van synthetische a priori oordelen, maar dan met een uniek kenmerk: in esthetische oordelen is de oordeelskracht zelf zowel het subjectieve 'voorwerp' (het gevoel) als de 'wet' die universaliteit eist, zonder zich te baseren op objectieve begrippen van het verstand.

\end{mdframed}

\section{Wat wordt er eigenlijk\textit{a priori} van een voorwerp uitgesproken in een smaakoordeel?}

Dat de voorstelling van een voorwerp onmiddellijk verbonden is met een genot kan slechts innerlijk waargenomen worden, en zou, als men niets meer dan dit zou willen aangeven, slechts een empirisch oordeel opleveren. Want ik kan een bepaald gevoel (van genot of ongenoegen) niet\textit{a priori} met enige voorstelling verbinden, behalve waar mijn grond een\textit{a priori} principe van de rede is dat de wil bepaalt; want dan is het genot (in het moreel gevoel) het gevolg daarvan, maar juist om die reden kan het geheel niet vergeleken worden met het genot in de smaak, aangezien het een bepaald begrip van een wet vereist, terwijl het smaakoordeel daartegenover onmiddellijk verbonden moet zijn met het louter oordelen, voorafgaand aan elk begrip. Derhalve zijn alle smaakoordelen ook enkelvoudige oordelen, omdat zij hun predikaat van bevrediging niet verbinden met een begrip, maar met een gegeven enkelvoudige empirische voorstelling.

Het is dus niet het genot zelf, maar \textbf{de universele geldigheid} van dit in de geest waargenomen genot, verbonden met het louter oordelen over een voorwerp, dat in een smaakoordeel als een universele regel voor de oordeelskracht, geldig voor iedereen, wordt voorgesteld. Het is een empirisch oordeel dat ik een voorwerp met genot waarneem en beoordeel. Maar het is een\textit{a priori} oordeel dat ik het mooi vind, dat wil zeggen, dat ik mag eisen dat deze bevrediging voor iedereen noodzakelijk is.

\begin{mdframed}[leftmargin=0pt,rightmargin=0pt,backgroundcolor=blue!6,linecolor=white,innertopmargin=6pt,innerbottommargin=6pt]

Het directe genot dat we bij een voorwerp ervaren is op zichzelf louter empirisch en subjectief - dit kan niet\textit{a priori} aan een voorstelling gekoppeld worden, behalve in morele oordelen waar een redelijk principe de wil bepaalt. Smaakoordelen verschillen hierin fundamenteel: ze verbinden het genot niet met een begrip of wet, maar met een concrete, enkelvoudige voorstelling.

De kern van het smaakoordeel ligt niet in het genot zelf, maar in de claim op universele geldigheid van dat genot. Waar de uitspraak "Ik vind dit aangenaam" puur empirisch is, bevat de uitspraak "Dit is mooi" een\textit{a priori} element: de eis dat deze bevrediging voor iedereen geldt. Het smaakoordeel stelt dus niet alleen vast dat ik genot ervaar, maar claimt dat dit genot noodzakelijk verbonden is met de beoordeling van het voorwerp, en dat deze verbinding voor alle subjecten zou moeten gelden. Deze universele aanspraak onderscheidt esthetische oordelen van louter persoonlijke voorkeuren.

\bigskip
§37 legt bloot wat het smaakoordeel zijn unieke karakter geeft: het is geen loutere registratie van persoonlijk genot, maar een claim op universele geldigheid die rechtstreeks voortkomt uit de structuur van ons oordeelsvermogen zelf. De paradox is treffend: hoewel het genot dat we ervaren bij een voorwerp volledig subjectief en empirisch lijkt, eist het smaakoordeel toch dat dit genot voor \emph{iedereen} zou moeten gelden. Dit is geen willekeurige aanspraak, maar een noodzakelijk kenmerk van esthetische beoordeling.

De crux ligt in het onderscheid tussen het \emph{ervaren} van genot (empirisch) en het \emph{eisen} van instemming (a priori). Waar morele oordelen hun universaliteit ontlenen aan redelijke principes die de wil binden, put het smaakoordeel zijn claim op algemene geldigheid uit de wijze waarop verbeelding en verstand in harmonie samenwerken bij het louter beoordelen van een voorstelling. Het is alsof de menselijke geest, wanneer ze in vrije reflectie een voorwerp aanschouwt, niet alleen een persoonlijke reactie heeft, maar tegelijkertijd aannemt dat deze reactie \emph{herkenbaar} zou moeten zijn voor anderen, niet omdat er een objectieve regel is, maar omdat de structuur van ons kennisvermogen zelf deze eis in zich draagt.

Dit doet denken aan de maçonnieke benadering van symbolen. Ook daar gaat het niet om persoonlijke associaties of voorgeschreven betekenissen, maar om een ervaring die zowel individueel als gedeeld is. Net zoals een vrijmetselaar een symbool interpreteert met het vertrouwen dat zijn interpretatie (hoewel subjectief) aansluit bij een gemeenschappelijke bron van betekenis, zo claimt het smaakoordeel een intersubjectiviteit die niet berust op vaste regels, maar op de gedeelde structuur van de menselijke geest. Het is deze spanning tussen persoonlijke ervaring en universele aanspraak die zowel Kants esthetica als de maçonnieke traditie kenmerkt: een vrije, maar gestructureerde zoektocht naar betekenis die zowel individueel als collectief geldig is. De schoonheid die we ervaren is dus geen eigenschap van het voorwerp, maar een openbaring van hoe onze geest werkt wanneer hij in vrije harmonie met zichzelf is.

\end{mdframed}

\section{Deductie van smaakoordelen.}

Als men toegeeft dat in een zuiver smaakoordeel de bevrediging in het voorwerp verbonden is met het louter oordelen over zijn vorm, dan is het niets anders dan de subjectieve doelmatigheid van die vorm voor de oordeelskracht die wij gewaarworden als verbonden met de voorstelling van het voorwerp in de geest. Nu de oordeelskracht met betrekking tot de formele regels van het oordelen, zonder enige stof (geen gevoelens noch begrippen), slechts gericht kan zijn op de subjectieve voorwaarden van het gebruik van de oordeelskracht in het algemeen (die noch beperkt is tot een bepaalde soort van zintuig noch tot een bepaald begrip van het verstand), en dus op dat subjectieve element dat men bij alle mensen kan veronderstellen (als vereist voor mogelijke kennis in het algemeen), moet de overeenstemming van een voorstelling met deze voorwaarden van de oordeelskracht\textit{a priori} als geldig voor iedereen kunnen worden aangenomen. Dat wil zeggen, het genot of de subjectieve doelmatigheid van de voorstelling voor de verhouding van de kennisvermogens bij het oordelen over een zintuiglijk voorwerp in het algemeen kan terecht van iedereen verwacht worden.\footnote{Om gerechtvaardigd te zijn in het claimen van universele instemming voor oordelen van de esthetische oordeelskracht die louter op subjectieve gronden berusten, is het voldoende om toe te geven:

\begin{enumerate}
    \item Bij alle mensen zijn de subjectieve voorwaarden van dit vermogen, voor zover het de verhouding van de kennisvermogens die daarin in werking gesteld worden tot een kennis in het algemeen betreft, dezelfde, wat waar moet zijn, want anders zouden mensen hun voorstellingen en zelfs kennis zelf niet kunnen mededelen.
    \item Het oordeel heeft uitsluitend deze verhouding (dus de formele voorwaarde van de oordeelskracht) in beschouwing genomen en is zuiver, dat wil zeggen, niet vermengd met begrippen van het voorwerp noch met gevoelens als bepalende gronden. Als er een fout wordt gemaakt met betrekking tot het laatste, dan betreft dat slechts de onjuiste toepassing op een bepaald geval van het gezag dat een wet ons geeft, waardoor het gezag in het algemeen niet wordt opgeheven.
\end{enumerate}

}

\begin{mdframed}[leftmargin=0pt,rightmargin=0pt,backgroundcolor=blue!6,linecolor=white,innertopmargin=6pt,innerbottommargin=6pt]

§38 biedt de deductie van smaakoordelen door aan te tonen waarom deze aanspraak kunnen maken op universele geldigheid, ondanks hun subjectieve grond. Het centrale argument is dat smaakoordelen niet berusten op persoonlijke gevoelens of specifieke begrippen, maar op de formeel-structurele voorwaarden van de oordeelskracht zelf - voorwaarden die bij alle mensen hetzelfde zijn, omdat ze noodzakelijk zijn voor elke vorm van kennis en communicatie.
Kant betoogt dat wanneer in een zuiver smaakoordeel het genot verbonden is met de loutere beoordeling van de vorm van een voorwerp, dit genot voortkomt uit de subjectieve doelmatigheid van die vorm voor de oordeelskracht. Omdat deze oordeelskracht - los van specifieke zintuiglijke indrukken of begrippen - alleen gericht is op de algemene voorwaarden waaronder kennis mogelijk is, en deze voorwaarden bij alle mensen gelijk zijn (immers, zonder gedeelde kennisstructuren zou communicatie onmogelijk zijn), kan de overeenstemming met deze voorwaarden a priori als universeel geldig worden beschouwd.

De rechtvaardiging voor de universele aanspraak van smaakoordelen vereist twee voorwaarden:
\begin{enumerate}
    \item De subjectieve voorwaarden van de oordeelskracht - met name de verhouding tussen verbeelding en verstand - zijn bij alle mensen identiek, omdat ze de basis vormen voor elke mogelijke kennis en communicatie.
    \item Het oordeel is zuiver, dat wil zeggen: het berust uitsluitend op deze formele verhouding en wordt niet beïnvloed door persoonlijke gevoelens of specifieke begrippen van het voorwerp.
\end{enumerate}

Als een oordeel niet zuiver is, wanneer het bijvoorbeeld beïnvloed wordt door persoonlijke voorkeuren of specifieke conceptuele kennis, dan gaat het niet om de geldigheid van het principe zelf, maar om een verkeerde toepassing ervan in een concreet geval. De universele aanspraak van het smaakoordeel berust dus niet op inhoudelijke overeenstemming, maar op de gedeelde structuur van de menselijke kennisvermogens.

\end{mdframed}

\section*{Opmerking}

Deze deductie is zo eenvoudig omdat het niet nodig is om enige objectieve realiteit van een begrip te rechtvaardigen; want schoonheid is geen begrip van het voorwerp, en het smaakoordeel is geen kennisoordeel. Zij stelt slechts dat wij gerechtvaardigd zijn om bij elke mens dezelfde subjectieve voorwaarden van de oordeelskracht vooronderstellen die wij bij onszelf aantreffen; en dan alleen als wij het gegeven voorwerp correct onder deze voorwaarden hebben ondergebracht. Nu dit laatste onvermijdelijke moeilijkheden met zich meebrengt die niet gelden voor de logische oordeelskracht (omdat men in de laatste onder begrippen subsumpeert, maar in de esthetische oordeelskracht onder een verhouding die louter een kwestie van gevoel is, namelijk die van de verbeelding en het verstand die wederzijds op elkaar zijn afgestemd in de voorgestelde vorm van het voorwerp, waar de subsumptie gemakkelijk bedrieglijk kan zijn), toch wordt daardoor niets afgedaan aan de rechtmatigheid van de aanspraak van de oordeelskracht op universele instemming, die slechts hierop neerkomt: de juistheid van het principe om voor iedereen geldig te oordelen op subjectieve gronden. Want voor zover het de moeilijkheid en de twijfel over de juistheid van de subsumptie onder dat principe betreft, maakt dit de rechtmatigheid van de aanspraak op deze geldigheid van een esthetisch oordeel in het algemeen, en dus het principe zelf, niet twijfelachtiger dan de eveneens (hoewel niet zo vaak en niet zo gemakkelijk) foutieve subsumptie van de logische oordeelskracht onder haar principe het laatste, dat objectief is, twijfelachtig kan maken. Maar als de vraag zou luiden: ‘‘Hoe is het mogelijk om de natuur\textit{a priori} als een totaliteit van voorwerpen van smaak aan te nemen?’’, dan heeft dit probleem betrekking op de teleologie, omdat het voortbrengen van vormen die doelmatig zijn voor ons oordeelsvermogen als een doel van de natuur beschouwd zou moeten worden dat wezenlijk tot haar begrip behoort. De juistheid van deze veronderstelling is echter nog zeer twijfelachtig, terwijl de werkelijkheid van de schoonheden der natuur voor de ervaring openligt.

\begin{mdframed}[leftmargin=0pt,rightmargin=0pt,backgroundcolor=blue!6,linecolor=white,innertopmargin=6pt,innerbottommargin=6pt]

Deze deductie is minder ingewikkeld dan andere, omdat ze niet hoeft aan te tonen dat schoonheid een eigenschap is die in de dingen zelf ligt. Schoonheid is geen objectief begrip, en een smaakoordeel is geen kennisuitspraak. Ze stelt slechts dat we ervan uit mogen gaan dat alle mensen dezelfde innerlijke voorwaarden hebben om te oordelen, dezelfde manier waarop verbeelding en verstand samenwerken, die wij bij onszelf waarnemen. Deze aanname is gerechtvaardigd, omdat zonder zo'n gedeelde basis alle communicatie en kennis onmogelijk zou zijn.

Toch is het niet altijd eenvoudig om een concreet voorwerp correct te beoordelen volgens deze voorwaarden. In tegenstelling tot logische oordelen, waar we dingen onder duidelijke begrippen brengen, gaat het hier om een gevoelsmatige afstemming tussen verbeelding en verstand in de waargenomen vorm. Daardoor kunnen we ons gemakkelijk vergissen. Maar zo'n vergissing doet niets af aan het principe zelf: het recht om universele instemming te eisen voor onze esthetische oordelen, gebaseerd op subjectieve gronden. Dat logische oordelen soms fout gaan, maakt hun objectieve principes immers ook niet ongeldig.

Een diepere vraag zou zijn: Hoe kunnen we aannemen dat de natuur vanzelf voorwerpen voortbrengt die aan onze smaak beantwoorden? Dat is een kwestie van doelmatigheid, alsof de natuur vormgeving als doel zou hebben. Maar zo'n aanname is zeer twijfelachtig. Wat we wel zeker weten, is dat schoonheid in de natuur werkelijk bestaat, zoals de ervaring ons leert. De filosofische rechtvaardiging geldt dus voor onze beoordeling van schoonheid, niet voor de natuur zelf als bron daarvan.

\bigskip
Deze opmerking onthult de diepere structuur van esthetische beoordeling als een brug tussen subjectiviteit en universaliteit. Wat hier centraal staat, is niet de objectieve realiteit van schoonheid in de dingen zelf, maar de gemeenschappelijke wijze waarop de menselijke geest werkt wanneer ze schoonheid ervaart. De deductie berust op een fundamentele aanname: dat alle mensen dezelfde innerlijke voorwaarden delen voor het oordelen - een gedeelde structuur van verbeelding en verstand die elke vorm van kennis en communicatie mogelijk maakt. Dit is geen louter psychologische observatie, maar een transcendentale voorwaarde: zonder deze gemeenschappelijke basis zou betekenisgeving zelf onmogelijk zijn.

De moeilijkheid ligt niet in het principe zelf, maar in de toepassing ervan. Waar logische oordelen werken met duidelijke begrippen, vereist esthetisch oordelen een gevoelsmatige afstemming tussen verbeelding en verstand - een proces dat inherent onzeker is, omdat het geen vaste regels kent. Toch is deze onzekerheid geen zwakte, maar juist het kenmerk van esthetische ervaring: schoonheid is geen kwestie van correcte toepassing van regels, maar van een vrije, maar gestructureerde wisselwerking tussen onze kennisvermogens. Het is alsof de menselijke geest hier een vermogen openbaart dat zowel persoonlijk als universeel is - niet omdat het objectief vastligt, maar omdat het voortkomt uit een structuur die we allemaal delen.

De verwijzing naar teleologie wijst op een belangrijke grens. Terwijl we de subjectieve geldigheid van smaakoordelen kunnen rechtvaardigen, blijft de vraag of de natuur zelf doelmatig is afgestemd op ons esthetisch vermogen onbeantwoord. Dit onderscheid is cruciaal: de deductie geldt voor onze *wijze van beoordelen*, niet voor de natuur als zodanig. Het is deze spanning tussen de noodzakelijkheid van onze beoordelingsstructuur en de contingentie van de natuurlijke werkelijkheid die de esthetische ervaring haar bijzondere karakter geeft. Net zoals in de maçonnieke traditie symbolen hun betekenis niet ontlenen aan een voorgeschreven plan, maar aan de wijze waarop de menselijke geest ze interpreteert binnen een gedeeld kader, zo is schoonheid geen eigenschap van de dingen, maar een openbaring van hoe wij als kennende en voelende wezens in de wereld staan.

\end{mdframed}

\section{Over de mededeelbaarheid van een gevoel.}

Indien het gevoel, als het reële in de waarneming, betrekking heeft op kennis, wordt het zintuiglijk gevoel genoemd; en zijn specifieke kwaliteit kan slechts volledig mededeelbaar worden geacht als men aanneemt dat iedereen een zintuig bezit dat hetzelfde is als het onze, maar dit kan in het geval van een zintuiglijk gevoel absoluut niet worden voorondersteld. Aan iemand die het reukvermogen ontbeert, kan dit soort gevoel dus niet worden meegedeeld; en zelfs als hij dit zintuig niet ontbeert, kan men nog steeds niet zeker zijn dat hij precies hetzelfde gevoel bij een bloem heeft als wij. Nog sterker: we moeten ervan uitgaan dat mensen verschillen wat betreft de \textbf{aangenaamheid} of \textbf{onaangenaamheid} van het gevoel bij één en hetzelfde voorwerp van de zintuigen; en het kan absoluut niet geëist worden dat genot in dezelfde voorwerpen door iedereen wordt toegekend. Genot van deze aard, aangezien het via de zintuigen in de geest komt en wij daarom passief ten aanzien daarvan zijn, kan het genot van \textbf{genieten} worden genoemd.

De bevrediging in een handeling omwille van haar morele kwaliteit is daartegenover geen genot van genieten, maar van zelfactiviteit en van de geschiktheid ervan voor het idee van haar roeping. Dit gevoel, dat moreel wordt genoemd, vereist echter begrippen; en toont geen vrije, maar een wetmatige doelmatigheid, en kan derhalve ook niet universeel meegedeeld worden anders dan door middel van de rede, en, indien het genot bij iedereen van dezelfde aard moet zijn, door middel van zeer bepaalde praktische begrippen van de rede.

Het genot in het verhevene in de natuur, als een genot van beschouwing dat subtiel redeneren inhoudt, doet eveneens aanspraak op universele deelname, maar veronderstelt reeds een ander gevoel, namelijk dat van zijn bovenzintuiglijke roeping, dat, hoe vaag het ook moge zijn, een morele grond heeft. Maar dat andere mensen hier rekening mee zullen houden en bevrediging zullen vinden in de beschouwing van de brute grootsheid van de natuur (die niet waarheidsgetrouw kan worden toegeschreven aan het gezicht ervan, dat eerder angstaanjagend is), is niet iets waar ik zomaar van uit mag gaan. Niettemin kan ik, in overweging van wat in aanmerking genomen moet worden bij die morele aanleg op elke passende gelegenheid, zelfs die bevrediging van iedereen verlangen, maar alleen door middel van de morele wet, die op haar beurt weer gegrond is op begrippen van de rede.

Daartegenover is het genot in het schone noch een genot van genieten, noch van een wetmatige activiteit, en zelfs niet van een beschouwing die subtiel redeneren inhoudt volgens ideeën, maar van louter reflectie. Zonder enig doel of fundamenteel principe als leidraad gaat dit genot gepaard met de gewone opvatting van een voorwerp door de verbeelding, als een vermogen van aanschouwing, in relatie tot het verstand, als een vermogen van begrippen, door middel van een procedure van de oordeelskracht, die zij ook moet uitoefenen voor de meest gewone ervaring: in het laatste geval wordt zij daartoe gedwongen omwille van een empirisch objectief begrip, terwijl in het eerste geval (bij het esthetisch oordelen) dit slechts geschiedt om de geschiktheid van de voorstelling voor de harmonische (subjectief doelmatige) bezigheid van beide kennisvermogens in hun vrijheid waar te nemen, dat wil zeggen, om de voorstellingsstaat met genot te gewaarworden. Dit genot moet noodzakelijkerwijs op dezelfde voorwaarden bij iedereen berusten, aangezien het subjectieve voorwaarden zijn van de mogelijkheid van kennis in het algemeen, en de verhouding van deze kennisvermogens die voor smaak vereist is, ook noodzakelijk is voor het gewone en gezonde verstand dat men bij iedereen mag veronderstellen. Om precies deze reden mag degene die met smaak oordeelt (zolang hij zich in dit bewustzijn niet vergist, en de stof niet voor de vorm aanziet, de bekoring niet voor schoonheid) ook de subjectieve doelmatigheid, dat wil zeggen, zijn bevrediging in het voorwerp, van elke ander verlangen, en mag hij aannemen dat zijn gevoel universeel mededeelbaar is, zelfs zonder bemiddeling van begrippen.

\begin{mdframed}[leftmargin=0pt,rightmargin=0pt,backgroundcolor=blue!6,linecolor=white,innertopmargin=6pt,innerbottommargin=6pt]

Het genot dat wij ervaren is niet allemaal van dezelfde aard, en niet alle genot kan even gemakkelijk met anderen gedeeld worden.

Het genot dat voortkomt uit onze zintuigen, zoals de geur van een bloem, is louter persoonlijk. We kunnen niet aannemen dat een ander precies hetzelfde voelt als wij, en evenmin dat wat ons aangenaam is, ook een ander zal behagen. Dit noemen we het genot van genieten: het komt tot ons via de zintuigen, en wij ondergaan het passief, zonder dat wij er invloed op hebben. Het is niet universeel, want het hangt af van hoe elk individu de wereld waarneemt.

Het genot dat wij putten uit morele handelingen is van een andere orde. Het is geen zintuiglijk genot, maar het gevoel van voldoening dat ontstaat wanneer wij handelen volgens ons morele plichtsbesef. Dit genot is wel universeel, maar alleen omdat het gebaseerd is op redelijke principes die voor iedereen gelden. Het vereist begrippen en een wetmatige ordening, en kan alleen via de rede aan anderen worden overgedragen.

Ook het genot in het verhevene, dat ontstaat bij de beschouwing van de overweldigende grootsheid der natuur, eist universaliteit, maar het veronderstelt al een moreel fundament. Wij kunnen niet zomaar aannemen dat een ander dezelfde ontroering voelt bij het aanschouwen van een stormachtige zee of een hoge berg. Toch kunnen wij dit genot van anderen verlangen, maar alleen op grond van de morele wet, die zelf weer berust op redelijke inzichten.

Het genot in het schone is echter van een geheel andere aard. Het is geen genot van de zintuigen, geen voldoening over een plichtsgetrouwe daad, en zelfs geen diepzinnige beschouwing volgens vaste ideeën. Het ontstaat uit de vrije reflectie over de vorm van een voorwerp, waarbij verbeelding en verstand in harmonie samenwerken. Deze harmonie is geen toeval, maar een voorwaarde voor alle kennis, een structuur die wij bij elk mens mogen veronderstellen. Omdat deze verhouding tussen verbeelding en verstand bij iedereen hetzelfde is, kan het genot in het schone wel universeel worden geclaimd. Wie met smaak oordeelt, en zich niet laat verleiden door louter aantrekkelijkheid, mag ervan uitgaan dat zijn bevrediging in het schone door iedereen gedeeld kan worden. Niet omdat de dingen zelf mooi zijn, maar omdat de wijze waarop wij ze aanschouwen een gemeenschappelijke grond heeft in de menselijke geest. Zo is schoonheid niet afhankelijk van persoonlijke smaak of redelijke principes, maar van de manier waarop wij allemaal, als kennende wezens, de wereld waarnemen.

/bigskip
Wat deze paragraaf ons leert, is dat schoonheid een unieke plaats inneemt in onze ervaring van de wereld. In tegenstelling tot zintuiglijk genot, dat louter subjectief is, en moreel genot, dat berust op redelijke principes, ontstaat esthetisch genot uit een vrije, maar gestructureerde wisselwerking tussen onze kennisvermogens. Het is alsof de menselijke geest hier een vermogen openbaart dat zowel persoonlijk als universeel is, niet omdat het aan vaste regels gehoorzaamt, maar omdat het voortkomt uit een structuur die wij allemaal delen.

De kern ligt in het onderscheid tussen ondergaan en beoordelen. Zintuiglijk genot ondergaan wij passief; moreel genot volgt uit onze actieve gehoorzaamheid aan de rede. Maar het genot in het schone ontstaat wanneer wij, zonder dwang of voorschrift, een voorwerp aanschouwen en daarin een harmonie ervaren tussen verbeelding en verstand. Deze harmonie is geen toeval, maar een noodzakelijke voorwaarde voor alle kennis. Daarin ligt de kracht van het smaakoordeel: het is geen kwestie van persoonlijke voorkeur of morele plicht, maar van een vrije reflectie die toch aanspraak maakt op algemene geldigheid.

Dit doet denken aan de schoonheid van een perfect uitgevoerd ritueel. Een ritueel is niet mooi omdat het strikt volgens voorschriften verloopt, of omdat het een persoonlijke voorkeur raakt. Zijn schoonheid ontstaat wanneer de deelnemers, binnen een vast kader van symbolen en handelingen, een diepere harmonie ervaren, een evenwicht tussen vrijheid en structuur. Net zoals de verbeelding en het verstand in esthetische reflectie samenwerken, zo komen in het ritueel vorm, beweging, woord en stilte samen in een ervaring die zowel persoonlijk als universeel is. De kracht ervan ligt niet in de letterlijke uitvoering, maar in de wijze waarop het de geest in evenwicht brengt: het is een vrije, maar gestructureerde ervaring die aanspraak maakt op algemene betekenis.
Het is deze spanning tussen vrijheid en structuur die zowel Kants esthetica als de rituele traditie kenmerkt. Schoonheid is hier geen kwestie van willekeurige smaak, maar van een herkenbare harmonie die ontstaat wanneer de geest zich openstelt voor de ordening van de ervaring. Zo wordt het ritueel niet alleen een persoonlijke belevenis, maar een gedeelde openbaring van hoe de menselijke geest betekenis geeft aan wat hem overstijgt.

\end{mdframed}

\section{Over smaak als een soort van \textit{sensus communis}.}

Wanneer bij het oordeelsvermogen niet zozeer de reflectie zelf opvalt, maar slechts het resultaat daarvan, wordt het vaak een gevoel genoemd, en men spreekt van een gevoel voor waarheid, voor fatsoen, voor rechtvaardigheid, enzovoort. Toch weten wij, of zouden wij althans moeten weten, dat deze begrippen niet in de zintuigen hun grond kunnen hebben, en dat een dergelijk gevoel al helemaal geen vermogen bezit om universele regels uit te drukken. Een voorstelling van waarheid, geschiktheid, schoonheid of rechtvaardigheid zou nooit in ons opkomen als wij ons niet boven de zintuigen konden verheffen tot hogere kennisvermogens. \textbf{Het gewone menselijk verstand}, dat als louter gezond (nog niet gecultiveerd) verstand beschouwd wordt en het minste is dat wij van iemand mogen verwachten die zich mens noemt, heeft daarom het ongelukkige voorrecht om met de naam gezond verstand (\textit{sensus communis}) te worden aangeduid, en wel zo dat het woord "\textbf{gemeen}" (niet alleen in onze taal, die hier een dubbelzinnigheid bevat, maar ook in vele andere talen) hetzelfde betekent als \textit{vulgair}, iets wat overal voorkomt en waarover beschikken zeker geen voordeel of eer is.

Met \textit{sensus \textbf{communis}} moet echter het idee van een \textbf{gemeenschappelijk} gevoel worden bedoeld: een vermogen om te oordelen dat in zijn reflectie (a priori) rekening houdt met de wijze waarop ieder ander denkt, om zo zijn oordeel \textbf{als het ware} aan de menselijke rede als geheel voor te leggen en daardoor de illusie te vermijden die voortkomt uit subjectieve, private voorwaarden die gemakkelijk voor objectief kunnen worden aangezien. Dit gebeurt door ons oordeel niet zozeer te toetsen aan de werkelijke oordelen van anderen, maar aan de slechts mogelijke oordelen, en ons in de plaats van ieder ander te verplaatsen door af te zien van de toevallige beperkingen die aan ons eigen oordelen kleven. Dit bereiken wij door zoveel mogelijk alles in onze voorstellingsstaat weg te laten wat stof is, dat wil zeggen: sensatie, en ons uitsluitend te richten op de formele eigenaardigheden van onze voorstelling of onze voorstellingsstaat. Misschien lijkt deze reflectie-handeling te kunstmatig om aan het vermogen toe te schrijven dat wij \textbf{gezond} verstand noemen, maar dat is alleen zo als wij het in abstracte formules uitdrukken. In werkelijkheid is niets natuurlijker dan het abstractie maken van bekoring en emotie wanneer wij een oordeel zoeken dat als universele regel moet dienen.

De volgende grondbeginselen van het gewone menselijk verstand behoren hier niet als onderdelen van de kritiek van de smaak, maar kunnen toch dienen om de fundamentele principes ervan te verduidelijken. Zij luiden: 1. Zelf denken; 2. Zich in de plaats van ieder ander denken; 3. Altijd consequent met zichzelf denken.
Het eerste is het beginsel van het \textbf{vooroordeelvrije} denken, het tweede van het \textbf{ruimdenkende} denken, het derde van het \textbf{consequente} denken. Het eerste is het beginsel van een rede die nooit \textbf{passief} is. De neiging tot het laatste, dus tot heteronomie van de rede, heet \textbf{vooroordeel}, en het grootste vooroordeel van alle is om de rede voor te stellen alsof zij niet onderworpen is aan de natuurwetten waarop het verstand haar door middel van haar eigen wezenlijke wet grondslagen verschaft: met andere woorden, \textbf{bijgeloof}. Bevrijding van bijgeloof heet \textbf{verlichting}\footnote{Men ziet gemakkelijk in dat verlichting in theorie eenvoudig is, maar in de praktijk een moeilijke zaak die slechts langzaam kan worden verwezenlijkt. Want hoewel het voor iemand die zich tot zijn wezenlijke bestemming beperkt en niet streeft naar kennis die zijn verstand te boven gaat, heel eenvoudig is om niet passief met zijn rede om te gaan, maar altijd wetgevend voor zichzelf te zijn, is het toch zeer moeilijk om het louter negatieve element (dat de ware verlichting uitmaakt) in de wijze van denken (met name die van het publiek) te handhaven of tot stand te brengen, omdat het streven naar kennis die het verstand te boven gaat nauwelijks verboden kan worden en er altijd velen zullen zijn die vol vertrouwen beloven deze kennisbehoefte te kunnen bevredigen.}. 
Hoewel deze term ook wordt gebruikt voor bevrijding van vooroordelen in het algemeen, verdient vooral het bijgeloof (\textit{in sensu eminenti}) de naam vooroordeel, omdat de blindheid waartoe het leidt, die het zelfs als plicht voorschrijft, het meest duidelijk maakt hoe noodzakelijk het is om door anderen geleid te worden, en dus de toestand van een passieve rede.
Wat het tweede beginsel van het denken betreft: wij zijn gewend iemand beperkt (\textbf{kleinburgerlijk}, in tegenstelling tot \textbf{ruimdenkend}) te noemen wiens talenten niet toereikend zijn voor enige grote taak (met name als die intensief is). Maar hier gaat het niet om het kennisvermogen zelf, maar om de \textbf{wijze van denken} die nodig is om er een doelmatig gebruik van te maken. Hoe klein de omvang en graad van iemands natuurlijke aanleg ook mogen zijn, toch toont iemand een \textbf{ruimdenkende geest} als hij zich losmaakt van de subjectieve, private voorwaarden van zijn oordeel, waaraan zovelen zich als het ware gevangen laten houden, en zijn eigen oordeel vanuit een \textbf{universeel standpunt} beziet (dat hij alleen kan bepalen door zich in het standpunt van anderen te verplaatsen).
Het derde beginsel, dat van het \textbf{consequente} denken, is het moeilijkst te bereiken en kan alleen worden verwezenlijkt door de combinatie van de eerste twee en na frequente toepassing ervan, totdat zij vanzelfsprekend zijn geworden. Men kan zeggen dat het eerste beginsel dat van het verstand is, het tweede dat van de oordeelskracht, en het derde dat van de rede.

Ik keer terug naar de draad die door deze uitweiding was losgelaten, en stel dat smaak met meer recht \textit{sensus communis} kan worden genoemd dan het gezonde verstand, en dat de esthetische oordeelskracht eerder dan de intellectuele de naam van een gemeenschappelijk gevoel verdient,\footnote{Men zou smaak kunnen aanduiden als \textit{sensus communis aestheticus}, en het gewone menselijk verstand als \textit{sensus communis logicus}.} als men tenminste het woord "gevoel" wil gebruiken voor een effect van louter reflectie op de geest, want hier bedoelt men met "gevoel" het gevoel van genot. Men zou smaak zelfs kunnen definiëren als het vermogen om te oordelen over wat ons gevoel in een gegeven voorstelling \textbf{universeel mededeelbaar} maakt, zonder bemiddeling van een begrip.

Het vermogen van mensen om hun gedachten mede te delen vereist eveneens een verhouding tussen verbeelding en verstand, om aanschouwingen met begrippen en begrippen weer met aanschouwingen te verbinden, die samen een kennis vormen. Maar in dat geval is de overeenstemming tussen deze twee geestesvermogens \textbf{wetmatig}, onder de dwang van bepaalde begrippen. Alleen daar waar de verbeelding in haar vrijheid het verstand prikkelt, en het verstand op zijn beurt, zonder begrippen, de verbeelding in een regelmatig spel zet, wordt de voorstelling niet als gedachte, maar als het innerlijke gevoel van een doelmatige gemoedstoestand medegedeeld.

Smaak is dus het vermogen om \textit{a priori} de mededeelbaarheid van de gevoelens die met een gegeven voorstelling verbonden zijn (zonder bemiddeling van een begrip) te beoordelen.
Als men zou mogen aannemen dat de louter universele mededeelbaarheid van ons gevoel op zichzelf al een belang voor ons inhoudt (wat men echter niet gerechtvaardigd is af te leiden uit de aard van een louter reflecterend oordeelsvermogen), dan zou men kunnen verklaren waarom het gevoel in het smaakoordeel van iedereen verwacht wordt alsof het een plicht is.

\begin{mdframed}[leftmargin=0pt,rightmargin=0pt,backgroundcolor=blue!6,linecolor=white,innertopmargin=6pt,innerbottommargin=6pt]

Wanneer wij spreken over een "gevoel" voor waarheid, fatsoen of rechtvaardigheid, bedoelen wij eigenlijk geen zintuiglijk vermogen, maar een vermogen om te oordelen dat rekening houdt met hoe anderen denken. Dit noemt Kant \textit{sensus communis}: niet het gewone, vulgaire verstand, maar een \textbf{gemeenschappelijk gevoel}, een vermogen om ons in de plaats van anderen te verplaatsen en ons oordeel te toetsen aan wat universeel geldig zou kunnen zijn. Dit doen wij door af te zien van persoonlijke sensaties en ons te richten op de vorm van onze voorstellingen. Niets is natuurlijker dan bekoring en emotie opzij te zetten wanneer wij een oordeel willen vellen dat voor iedereen kan gelden.

Drie grondbeginselen verhelderen dit:

\begin{compactitem}
    \item \textbf{Zelf denken} – vrij van vooroordelen;
    \item \textbf{Zich inleven in anderen} – ruimdenkend zijn;
    \item \textbf{Consequent denken} – in harmonie met onszelf.
\end{compactitem}
Deze beginselen corresponderen met het verstand, de oordeelskracht en de rede. Bijgeloof, als grootste vooroordeel, toont hoe moeilijk het is om actief en vrij te denken, verlichting is de bevrijding daarvan.

Smaak verdient de naam \textit{sensus communis} meer dan het gewone verstand, omdat het niet gaat om kennis, maar om het \textbf{gevoel van genot} dat ontstaat wanneer verbeelding en verstand in vrijheid en harmonie samenwerken. Waar kennis berust op vaste begrippen, deelt smaak een innerlijke, doelmatige gemoedstoestand mee. Smaak is dus het vermogen om a priori te beoordelen of een gevoel universeel mededeelbaar is, zonder dat wij daarbij afhankelijk zijn van begrippen.

Als wij zouden aannemen dat universele mededeelbaarheid op zichzelf al waardevol is, dan zouden wij kunnen begrijpen waarom smaak van iedereen verwacht wordt alsof het een plicht is. Het is geen kwestie van regels, maar van een vrije, maar gedeelde ervaring die voortkomt uit de structuur van de menselijke geest zelf.

\bigskip
Smaak is niet louter een kwestie is van persoonlijke voorkeur, maar van een gemeenschappelijke grond in de menselijke geest. Het is alsof Kant hier blootlegt dat schoonheid geen eigenschap is van de dingen zelf, maar een ervaring die ontstaat wanneer wij ons oordeel niet baseren op onze eigen, beperkte sensaties, maar op wat voor iedereen herkenbaar zou kunnen zijn. Dit is geen abstracte, kunstmatige handeling, maar een natuurlijk vermogen: het vermogen om ons te verheffen boven het louter subjectieve en ons in te leven in de wijze waarop anderen de wereld waarnemen. In die zin is smaak geen individuele aangelegenheid, maar een openbare aangelegenheid, een vorm van denken die ons verbindt met anderen, niet door regels op te leggen, maar door een gedeelde structuur van ervaring.

De drie grondbeginselen, zelf denken, zich inleven in anderen, consequent denken, tonen hoe smaak niet alleen een esthetische, maar ook een morele en intellectuele dimensie heeft. Het is alsof smaak ons leert om actief en vrij te oordelen, zonder ons te laten leiden door vooroordelen of bijgeloof. Net zoals verlichting ons bevrijdt van de passiviteit van de rede, zo bevrijdt smaak ons van de beperking van onze eigen, private gevoelens. Het is een vorm van vrije reflectie die toch aanspraak maakt op universaliteit, niet omdat zij berust op vaste regels, maar omdat zij voortkomt uit de wijze waarop wij allemaal, als redelijke en voelende wezens, de wereld ordenen.

Dit doet denken aan de maçonnieke praktijk, waar het niet gaat om het volgen van voorgeschreven interpretaties, maar om het vinden van betekenis binnen een gedeeld kader. Zoals een vrijmetselaar leert om symbolen te begrijpen door zich in te leven in de ervaring van anderen, zo leert smaak ons om schoonheid te herkennen door ons te verplaatsen in de wijze waarop anderen de wereld aanschouwen. Schoonheid is dus geen kwestie van individuele smaak, maar van een herkenbare harmonie die ontstaat wanneer wij ons openstellen voor de gemeenschappelijke structuur van de menselijke geest. In die zin is smaak niet alleen een esthetisch vermogen, maar een fundamentele wijze van denken die ons in staat stelt om betekenis te geven aan wat wij ervaren, niet als geïsoleerde individuen, maar als deel van een groter geheel.

\end{mdframed}

\section{Over het empirische belang bij het schone.}

Dat het smaakoordeel, waarbij iets als mooi wordt verklaard, geen belang \textbf{als bepalende grond} mag hebben, is hierboven voldoende aangetoond. Maar daaruit volgt niet dat er, nadat het als zuiver esthetisch oordeel is geveld, geen belang aan verbonden kan worden. Deze verbinding kan echter altijd alleen indirect zijn, dat wil zeggen: de smaak moet eerst met iets anders in verband worden gebracht, om aan de bevrediging die voortkomt uit louter reflectie over een voorwerp een verdere \textbf{vreugde in} zijn \textbf{bestaan} (waaruit immers alle belang bestaat) te kunnen koppelen. Want wat geldt voor kennisoordelen (over dingen in het algemeen) geldt ook voor het esthetische oordeel: er is geen geldige afleiding van mogelijkheid naar werkelijkheid.. Dit andere element kan iets empirisch zijn, namelijk een neiging die kenmerkend is voor de menselijke natuur, of iets intellectueel, zoals een eigenschap van de wil die\textit{a priori} door de rede bepaalbaar is; beide bevatten een bevrediging in het bestaan van een voorwerp en kunnen zo de grond vormen voor een belang in datgene wat al op zichzelf, zonder enig belang, bevalt.

Het schone boeit empirisch alleen in de \textbf{samenleving}; en als de drang tot samenleving als natuurlijk voor de mens wordt erkend, terwijl de geschiktheid en de neiging daartoe, dat wil zeggen: de \textbf{gezelligheid}, als noodzakelijk voor de mens worden beschouwd, als een wezen dat voor het samenleven bestemd is en dus als een eigenschap die bij de \textbf{menselijkheid} hoort, dan kan het niet anders of smaak wordt ook gezien als een vermogen om alles te beoordelen waardoor men zelfs zijn \textbf{gevoel} aan ieder ander kan meedelen, en dus als een middel om te bevorderen wat door een ieder natuurlijke neiging wordt geëist.

Alleen voor zichzelf, verlaten op een onbewoond eiland, zou een mens noch zijn hut noch zichzelf versieren, noch bloemen zoeken of planten om zich mee te tooien. Pas in de samenleving komt het bij hem op om niet alleen een mens te zijn, maar op zijn eigen wijze een gecultiveerd mens (het begin van beschaving). Zo oordelen wij over iemand die geneigd is zijn genot met anderen te delen en daarin bedreven is, en die niet tevreden is met een voorwerp als hij zijn bevrediging daarin niet in gemeenschap met anderen kan ervaren. Verder verwacht en eist ieder van een ander een streven naar universele mededeling, alsof dit voortkomt uit een oorspronkelijk verdrag dat door de menselijkheid zelf wordt voorgeschreven. Zo worden allereerst louter aantrekkelijke dingen, zoals verf voor het lichaam (rocou bij de Cariben en cinaber bij de Iroquois), bloemen, schelpen, prachtig gekleurde vogelveren, en later ook mooie vormen (zoals bij kano’s, kleding, enzovoort), die op zichzelf geen genot verschaffen, dus geen bevrediging van genieten, in de samenleving belangrijk en met groot belang verbonden. Uiteindelijk maakt de beschaving, als zij haar hoogste punt heeft bereikt, hiervan bijna het hoofdwerk van de verfijnde neiging, en sensaties krijgen waarde alleen voor zover zij universeel mededeelbaar zijn. Op dat moment is het genot dat ieder in zo’n voorwerp heeft weliswaar onbeduidend en heeft het op zichzelf geen opvallend belang, maar de gedachte aan zijn universele mededeelbaarheid vergroot zijn waarde bijna oneindig.

Dit belang, dat indirect via een neiging tot samenleving aan het schone is verbonden en dus empirisch is, heeft voor ons hier echter geen betekenis. Wij moeten die betekenis immers alleen vinden in wat\textit{a priori} met het smaakoordeel in verband kan worden gebracht, al is het maar indirect. Want zelfs als in deze laatste vorm een belang dat eraan verbonden is, aan het licht zou komen, dan zou smaak in ons oordeelsvermogen een overgang tonen van zintuiglijk genot naar moreel gevoel. En niet alleen zou men daardoor beter geleid worden in het doelmatige gebruik van smaak, maar ook zou een schakel in de keten van menselijke vermogens a priori, waarop alle wetgeving berust, als zodanig aan het licht komen. Zoveel kan zeker gezegd worden over het empirische belang bij voorwerpen van smaak en bij smaak zelf: aangezien smaak de neiging koestert, hoe verfijnd ook, laat hij zich graag vermengen met alle neigingen en passies die in de samenleving hun grootste verscheidenheid en hoogste niveau bereiken. En het belang in het schone, als het hierop gebaseerd is, zou slechts een zeer dubieuze overgang van het aangename naar het goede kunnen bieden. Maar of het goede misschien wel bevorderd zou kunnen worden door smaak, als deze in haar zuiverheid wordt genomen, hebben wij reden om te onderzoeken.

\begin{mdframed}[leftmargin=0pt,rightmargin=0pt,backgroundcolor=blue!6,linecolor=white,innertopmargin=6pt,innerbottommargin=6pt]

Wanneer wij iets als mooi beoordelen, doet dat oordeel zich voor zonder enig belang als bepalende grond. Toch betekent dit niet dat er geen belang aan verbonden kan zijn, alleen niet rechtstreeks. Het belang dat wij in schoonheid vinden, ontstaat altijd indirect: doordat smaak zich verbindt met iets buiten zichzelf. Dat "iets" kan tweevoudig zijn: ofwel een natuurlijke neiging van de mens, zoals de drang tot gezelschap, ofwel een zuiver intellectuele eigenschap, zoals de morele bestemming van onze wil. Beide brengen een bevrediging met zich mee die verder gaat dan louter esthetisch genot.

Stel u voor: een mens, alleen op een eiland, zou zijn hut niet versieren, zichzelf niet tooien met bloemen of kleuren. Pas tussen mensen ontwaakt de behoefte om niet alleen te leven, maar om te leven als een wezen dat zichzelf en zijn omgeving verfijnt. Dit is het begin van beschaving. Wat wij dan waarderen, is niet het voorwerp op zich, een schelp, een veer, een vorm, maar het vermogen om het genot dat het ons schenkt met anderen te delen. De waarde van schoonheid groeit naarmate wij haar als mededeelbaar ervaren. Een eenvoudig genot wordt oneindig waardevol wanneer wij weten dat het door iedereen zo ervaren kan worden.

Maar dit is nog slechts het empirische belang van smaak, en dat is niet waar het mij om te doen is. Wat werkelijk telt, is de vraag of smaak, in haar zuiverste vorm, meer is dan een verfijning van onze neigingen. Kan zij ons leiden van het louter aangename naar het goede? Kan zij, los van alle persoonlijke voorkeuren, een brug slaan tussen zintuiglijke ervaring en moreel inzicht? Dat is de overgang die wij moeten onderzoeken, want alleen dan toont smaak haar ware betekenis: niet als dienstmaagd van onze grillen, maar als een vermogen dat ons verheft.

/bigskip
Wat deze paragraaf ons leert, is dat schoonheid niet louter een zaak is van individuele contemplatie, maar een fundamentele menselijke behoefte aan verbinding. Het is alsof Kant hier aantoont dat esthetische ervaring niet eindigt bij het subject, maar juist haar volle betekenis krijgt in het gezelschap van anderen. Schoonheid is geen luxe, maar een uiting van onze diepste neiging om ons mens-zijn te delen, om niet alleen te voelen, maar om dat gevoel herkenbaar te maken voor een ander. Dit is geen toeval, maar een wezenlijk kenmerk van onze rede: wij zoeken niet alleen naar waarheid of goedheid, maar ook naar wat ons verbindt in een gedeelde ervaring van de wereld.

De paradox is treffend: hoewel het smaakoordeel zich aanvankelijk presenteert als zuiver en belangloos, openbaart het in de praktijk een diepe sociale functie. De bloem die wij plukken, de kleur die wij dragen, de vorm die wij bewonderen, hun waarde ligt niet in hun materiële aanwezigheid, maar in hun vermogen om een gemeenschappelijke taal te spreken. Dit doet denken aan het ritueel, waar niet de handeling zelf, maar de gedeelde betekenis ervan de kern vormt. Zo is schoonheid geen vlucht uit de wereld, maar een manier om ons in de wereld te verankeren, door ons te verbinden met anderen op een niveau dat voorbij woorden gaat.

Toch waarschuwt Kant ons hier ook: als smaak zich beperkt tot het empirische, tot het verfijnen van onze neigingen in gezelschap, blijft zij gevangen in het aangename. Haar ware kracht ligt in de mogelijkheid om ons te verheffen boven louter genot, om ons voor te bereiden op het morele. Het is alsof smaak ons leert om, voordat wij handelen volgens plicht, eerst te ervaren wat universeel betekenisvol is. In die zin is esthetische beoordeling geen doel op zich, maar een voorbereiding op een hogere vorm van vrijheid: de vrijheid om niet alleen te voelen wat ons bindt, maar om te handelen volgens wat ons verenigt.

In de loge is het niet de vorm van het ritueel op zich die telt, maar het vermogen van die vorm om een gemeenschappelijke innerlijke ervaring op te roepen. Net zoals Kant stelt dat de waarde van schoonheid ligt in haar mededeelbaarheid, zo is het maçonnieke ritueel geen mechanische handeling, maar een levende taal die alleen betekenis krijgt wanneer broeders haar samen beleven. De symbolen, de woorden, de stilte, ze zijn niet mooi omdat ze voorgeschreven zijn, maar omdat ze een ruimte openen waarin ieder individu een persoonlijke, maar herkenbare waarheid ervaart. Zo wordt het ritueel een oefening in wat Kant noemt: het universaliseerbare gevoel. Het is geen toeval dat vrijmetselaren spreken over "het licht zien", dat licht is geen individuele openbaring, maar een ervaring die alleen werkelijk schijnt wanneer zij gedeeld wordt.

Kant benadrukt dat schoonheid pas werkelijk leeft in gezelschap. Voor vrijmetselaren is dit geen abstractie, maar een dagelijkse werkelijkheid. De loge is de plaats waar de ruwe steen (het ongevormde individu) niet in isolatie wordt bewerkt, maar in het gezelschap van anderen die dezelfde zoektocht delen. De "versiering" waar Kant over spreekt, de verfijning van de geest, is precies wat er gebeurt wanneer broeders samenkomen: zij tooien zich niet met uiterlijkheden, maar met gedeelde inzichten, met de taal van symbolen die alleen betekenis hebben omdat zij door allen worden begrepen. Het is deze dynamiek die de loge onderscheidt van een eenzame meditatie: verfijning is geen solitaire daad, maar een collectieve arbeid.

Kants vraag of smaak een overgang kan zijn van het aangename naar het goede is voor vrijmetselaren geen filosofische speculatie, maar een praktische opgave. Het maçonnieke ideaal is niet om schoonheid te bewonderen omwille van zichzelf, maar om door esthetische ervaring (in ritueel, architectuur, muziek) te worden voorbereid op morele actie. Wanneer een vrijmetselaar leert om schoonheid te herkennen in de harmonie van het ritueel, oefent hij tegelijkertijd zijn vermogen om morele harmonie in de wereld te zoeken. De loge is zo een school voor smaak en voor deugd: zij leert de broeder om niet alleen te voelen wat universeel is, maar om daar ook naar te handelen. Dat is de "overgang" waar Kant naar verwijst, en precies waarom vrijmetselarij geen louter spirituele of intellectuele bezigheid is, maar een weg naar actieve verbetering, zowel van zichzelf als van de samenleving.

Wat Kant hier beschrijft, is niets minder dan de grondslag van maçonnieke broederschap: het inzicht dat ware verfijning niet bestaat in individuele verlichting, maar in het vermogen om die verlichting zo te delen dat zij anderen ook verheft. Voor vrijmetselaren is schoonheid dus geen luxe, maar een noodzakelijke voorwaarde voor de zoektocht naar waarheid en deugd. Zij herinnert hen eraan dat hun arbeid, of het nu gaat om zelfkennis, morele groei of maatschappelijke bijdrage, alleen betekenis heeft wanneer zij deze kunnen mededelen en delen. In die zin is smaak niet alleen een esthetisch vermogen, maar een morele plicht: de plicht om het licht dat men in zich draagt, zichtbaar te maken voor anderen.

\end{mdframed}

\section{Over het intellectuele belang bij het schone.}

Het was met goede bedoelingen dat zij die alle bezigheden van de mens, waartoe deze gedreven wordt door zijn innerlijke natuurlijke aanleg, graag richten op het uiteindelijke doel van de mensheid, namelijk het moreel goede, het stellen van belang in het schone in het algemeen als een teken van een goed moreel karakter hebben beschouwd. Maar zij zijn door anderen, niet zonder grond, weersproken, die wezen op de ervaring dat kenners van smaak, die niet alleen vaak maar zelfs meestal ijdel, eigenzinnig en vatbaar voor verdorven neigingen zijn, misschien nog minder dan anderen aanspraak kunnen maken op de verdienste van toewijding aan morele principes. Zo lijkt het dat het gevoel voor het schone niet alleen specifiek verschilt van het moreel gevoel (wat het inderdaad is), maar ook dat het belang dat eraan verbonden kan zijn, moeilijk met het morele belang verenigd kan worden, en zeker niet door een innerlijke verwantschap.

Ik geef graag toe dat het belang in het \textbf{schone in de kunst} (waaronder ik ook reken het kunstige gebruik van de schoonheden der natuur voor decoratie, en dus voor ijdelheid) geen bewijs levert van een gezindheid die toegewijd is aan het moreel goede, of zelfs maar daartoe geneigd. Maar ik beweer wel dat een \textbf{onmiddellijk belang} stellen in de schoonheid der natuur (niet slechts smaak hebben om haar te beoordelen) altijd een teken is van een schone ziel, en dat, als dit belang gewoonte wordt, het ten minste wijst op een gezindheid van de geest die gunstig is voor het moreel gevoel, als het graag verbonden wordt met de \textbf{aanschouwing der natuur}. Men moet echter bedenken dat ik hier strikt de schone \textbf{vormen} der natuur bedoel, en daarentegen de \textbf{bekoringen} die zij daarmee zo overvloedig verbindt, opzij zet, aangezien het belang daarin weliswaar ook onmiddellijk is, maar toch empirisch.

Iemand die alleen (en zonder enige bedoeling om zijn waarnemingen aan anderen mee te delen) de schone vorm van een wilde bloem, een vogel, een insect, enzovoort, beschouwt om zich eraan te verwonderen, ervan te houden en niet te willen dat het geheel uit de natuur verdwijnt, zelfs als het hem eerder schade dan voordeel zou berokkenen, stelt een onmiddellijk en zeker intellectueel belang in de schoonheid der natuur. Dat wil zeggen: niet alleen de vorm van haar producten, maar ook hun bestaan bevalt hem, zelfs als geen zintuiglijke bekoring hierin een rol speelt en hij geen enkel doel met dit belang verbindt.

Het is echter opmerkelijk dat, als iemand deze liefhebber van het schone heimelijk bedrogen had, door kunstbloemen te plaatsen (die zo natuurgetrouw kunnen zijn dat ze niet van echte te onderscheiden zijn) of kunstmatig gevormde vogels op de takken van bomen, en hij vervolgens de bedriegerij ontdekte, het onmiddellijke belang dat hij ervoor had, meteen zou verdwijnen. Misschien zou een ander belang, zoals ijdelheid om zijn vertrek te versieren voor de ogen van anderen, daarvoor in de plaats komen. De gedachte dat de natuur deze schoonheid heeft voortgebracht, moet de aanschouwing en reflectie begeleiden, en daarop alleen berust het onmiddellijke belang dat men erin stelt. Anders blijft er ofwel een louter smaakoordeel zonder enig belang, ofwel een belang dat slechts indirect verbonden is met de samenleving, en dit laatste biedt geen zekere aanwijzingen voor een moreel goede gezindheid.

Deze voorrang van de schoonheid der natuur boven die van de kunst, die op zichzelf al een onmiddellijk belang wekt, zelfs als de kunst de natuur in vorm zou overtreffen, stemt overeen met het verfijnde en weloverwogen denken van alle mensen die hun moreel gevoel hebben gecultiveerd. Als een mens, die genoeg smaak bezit om producten van schone kunst met de grootste juistheid en verfijning te beoordelen, graag de kamer verlaat waar zich die schoonheden bevinden die ijdelheid en hoogstens sociale vreugde voeden, en zich wendt tot de schoonheid der natuur, alsof hij daar een extase voor zijn geest hoopt te vinden in een gedachtengang die hij nooit geheel kan uitputten, dan zouden wij deze keuze met achting beschouwen en in hem een schone ziel veronderstellen, waar geen kenner of liefhebber van kunst aanspraak op kan maken op grond van het belang dat hij in zijn voorwerpen stelt. – Wat is nu het onderscheid tussen deze verschillende waarderingen van twee soorten voorwerpen, die in het louter smaakoordeel nauwelijks met elkaar zouden wedijveren om voorrang?

Wij bezitten een louter esthetisch oordeelsvermogen, om vormen te beoordelen zonder begrippen en om bevrediging te vinden in het louter oordelen over deze vormen, die wij tegelijkertijd tot een regel voor iedereen maken, zonder dat dit oordeel op een belang berust of er een voortbrengt. Daarnaast bezitten wij ook een intellectueel oordeelsvermogen, om \textit{a priori} voor louter praktische maximen (voor zover deze in zichzelf geschikt zijn voor universele wetgeving) een bevrediging vast te stellen die wij tot een wet voor iedereen maken, zonder dat ons oordeel op enig belang berust, hoewel het er wel een voortbrengt. Het genot of de onlust in het eerste oordeel noemen wij smaak, in het tweede moreel gevoel.

Maar omdat het de rede ook interesseert dat de ideeën (waarvoor zij een onmiddellijk belang in het moreel gevoel voortbrengt) ook objectieve werkelijkheid bezitten, dat wil zeggen, dat de natuur ten minste enige spoor of teken toont dat zij in zichzelf enige grond bevat voor een wetmatige overeenstemming van haar producten met onze bevrediging, die onafhankelijk is van elk belang (die wij \textit{a priori} als een voor iedereen geldende wet erkennen, zonder dit op bewijzen te kunnen gronden), moet de rede belang stellen in elke uiting in de natuur van een dergelijke overeenstemming. Bijgevolg kan de geest niet over de schoonheid der \textbf{natuur} reflecteren zonder zich daarbij ook geïnteresseerd te voelen. Door deze verwantschap is dit belang moreel, en wie zo’n belang stelt in het schone in de natuur, kan dat alleen doen voor zover hij reeds een vast belang in het moreel goede heeft gesteld. Wij hebben dus reden om ten minste een aanleg tot een goede morele gezindheid te vermoeden bij iemand die onmiddellijk belang stelt in de schoonheid der natuur.

Men zou kunnen zeggen dat deze uitleg van esthetische oordelen in termen van hun verwantschap met het moreel gevoel te gezocht lijkt om als de ware interpretatie te gelden van het raadselachtige teken waarmee de natuur in haar schone vormen figuratief tot ons spreekt. Maar ten eerste is dit onmiddellijke belang in het schone in de natuur niet algemeen, maar behoort het alleen toe aan hen wier denken ofwel reeds getraind is in het goede, ofwel bijzonder ontvankelijk is voor een dergelijke training. Ten tweede leidt, zelfs zonder duidelijke, subtiele en bewuste reflectie, de analogie tussen het zuivere smaakoordeel, dat, zonder afhankelijk te zijn van enig belang, een genot toelaat dat tegelijkertijd \textit{a priori} als passend voor de mensheid in het algemeen kan worden voorgesteld, en het morele oordeel, dat hetzelfde doet op grond van begrippen, tot een even onmiddellijk belang in het voorwerp van het eerste als in dat van het laatste. Alleen is het eerste een vrij belang, het laatste een belang dat op objectieve wetten berust. Daaraan voegen wij toe de bewondering voor de natuur, die in haar schone producten zich toont als kunst, niet toevallig, maar alsof zij doelbewust handelt, volgens een wetmatige ordening en als doelmatigheid zonder doel. Deze doelmatigheid, die wij nooit extern aantreffen, zoeken wij van nature in onszelf, en wel in datgene wat het uiteindelijke doel van ons bestaan uitmaakt: de morele roeping.

Dat de bevrediging in schone kunst in het zuivere smaakoordeel niet verbonden is met een onmiddellijk belang, zoals dat bij de schoonheid der natuur wel het geval is, is eveneens gemakkelijk te verklaren. De eerste is immers ofwel een zodanige nabootsing van de laatste dat zij bedrieglijk is, en in dat geval heeft zij het effect van natuurlijke schoonheid (waarvoor zij wordt aangezien), ofwel is het een kunst die duidelijk doelbewust gericht is op onze bevrediging. In dat geval zou de bevrediging in dit product weliswaar onmiddellijk door de smaak ontstaan, maar zou zij slechts een indirect belang opwekken in de oorzaak waarop zij berust, namelijk een kunst die alleen door haar doel kan interesseren en nooit op zichzelf.

De bekoringen in de schoonheid der natuur, die zo vaak als het ware versmolten zijn met de schone vorm, behoren ofwel tot de wijzigingen van het licht (in kleuring) ofwel tot die van het geluid (in tonen). Dit zijn immers de enige sensaties die niet alleen zintuiglijk gevoel toelaten, maar ook reflectie op de vorm van deze wijzigingen van de zintuigen, en die dus als het ware een taal bevatten die de natuur tot ons richt en die een hogere betekenis lijkt te hebben. Zo lijkt de witte kleur van de lelie de geest te stemmen tot ideeën van onschuld, en de zeven kleuren, in hun volgorde van rood tot violet, tot de ideeën van 1)~verhevenheid, 2)~durf, 3)~oprechtheid, 4)~vriendelijkheid, 5)~bescheidenheid, 6)~standvastigheid, en 7)~tederheid. Het gezang van de vogel verkondigt vreugde en tevredenheid met zijn bestaan. Ten minste, zo interpreteren wij de natuur, of zij het nu werkelijk zo bedoelt of niet.
Maar dit belang, dat wij hier stellen in schoonheid, vereist absoluut dat het de schoonheid der natuur is. Het verdwijnt geheel zodra men merkt dat men bedrogen is en dat het slechts kunst is, zozeer zelfs dat zelfs de smaak er niets moois meer in kan vinden of iets bekorends kan zien. Wat wordt er door dichters meer geprezen dan het betoverend mooie gezang van de nachtegaal, in een eenzaam bosje, op een stille zomeravond, onder het zachte maanlicht? Toch zijn er voorbeelden van gasten die, toen er geen nachtegaal te vinden was, door een grappenmaker van een herbergier werden misleid. Deze verstopte een kwajongen in het struikgewas, die met een rietje of fluitje het gezang van de vogel perfect kon nabootsen, tot volle tevredenheid van de gasten. Maar zodra men ontdekte dat het bedrog was, wilde niemand langer naar dit gezang luisteren, dat men voorheen zo betoverend vond. En hetzelfde geldt voor elk ander zangvogelgezang. Het moet de natuur zijn, of door ons als natuur beschouwd worden, wil wij een dergelijk onmiddellijk \textbf{belang} kunnen stellen in het schone. En des te meer als wij van anderen mogen verwachten dat zij dit belang ook stellen, wat inderdaad gebeurt, aangezien wij grof en laag denken wie geen \textbf{gevoel} heeft voor de schoonheid der natuur (wat wij noemen: de ontvankelijkheid voor een belang in haar aanschouwing), en die zich beperken tot het genot van louter zintuiglijke sensaties aan tafel of uit de fles.

\begin{mdframed}[leftmargin=0pt,rightmargin=0pt,backgroundcolor=blue!6,linecolor=white,innertopmargin=6pt,innerbottommargin=6pt]

Wanneer wij beweren dat belangstelling voor het schone een teken is van een goed moreel karakter, wordt ons vaak tegenspraak geboden en niet zonder reden. Kenners van kunst en schoonheid blijken immers niet zelden ijdel, eigenzinnig of vatbaar voor verdorven neigingen. Dit lijkt aan te tonen dat het gevoel voor schoonheid niet alleen wezenlijk verschilt van het morele gevoel, maar ook dat het belang dat wij eraan hechten moeilijk verenigbaar is met morele toewijding. Toch is er een cruciaal onderscheid. Het belang dat wij stellen in kunstmatige schoonheid, zoals versieringen of decoraties, levert geen bewijs van morele gezindheid. Maar wie zich onmiddellijk en oprecht interesseert voor de schoonheid der natuur, zonder enige bedoeling om anderen te imponeren of er voordeel uit te halen, toont daarmee wel degelijk iets diepers.

Iemand die stilstaat bij de schone vorm van een wilde bloem, een vogel of een insect, niet om zijn waarnemingen met anderen te delen, maar louter om zich eraan te verwonderen en ervan te houden, stelt een intellectueel belang in die schoonheid. Het is niet alleen de vorm die hem bevalt, maar het bestaan ervan. Dit belang verdwijnt direct wanneer blijkt dat de schoonheid kunstmatig is. Wat hem raakt, is de gedachte dat de natuur zelf dit heeft voortgebracht. Dit onmiddellijke belang, los van elke praktische of sociale overweging, is een teken van een schone ziel.

De schoonheid der natuur wekt een onmiddellijk belang op, op een manier die kunst niet kan evenaren. Zij spreekt tot ons als een taal die wij herkennen als een echo van onze morele roeping. Ons vermogen om schoonheid te waarderen zonder belang, louter omwille van zichzelf, vertoont een opvallende gelijkenis met ons moreel oordeel. Beide zijn universeel en vrij. De natuur nodigt ons uit om in onszelf te zoeken naar de grond van deze harmonie. Wij interpreteren haar verschijnselen als symbolen van menselijke deugden, en wie zich hierin verdiept, toont een geest die openstaat voor het morele.

Dit is geen bewijs van deugd, maar wel een belofte. Wie schoonheid in de natuur waardeert om haarzelf, heeft een ziel die gereed is voor morele verfijning. De schoonheid der natuur is geen luxe, maar een oefening in aandacht, een voorproefje van de harmonie waarnaar wij streven in ons handelen. Wie daarvoor ontvankelijk is, bewijst dat zijn hart al gericht is op wat werkelijk telt.

\bigskip
Schoonheid is geen louter esthetische ervaring, maar een morele oefening. Wat Kant hier blootlegt, is dat onze reactie op natuurlijke schoonheid geen kwestie is van smaak, maar van karakter. Het onderscheid tussen kunstmatige en natuurlijke schoonheid is niet alleen een kwestie van oorsprong, maar van betekenis. Natuurlijke schoonheid confronteert ons met iets dat ons overstijgt, zonder dat wij het kunnen controleren of bezitten. Zij vraagt niets van ons, behalve dat wij ons openstellen voor haar bestaan. In die openheid ligt de kern van morele ontvankelijkheid: het vermogen om waarde te ervaren zonder beloning, om betekenis te zien zonder voorschrift.

Wanneer wij een bloem bewonderen om wat zij is, en niet om wat zij voor ons doet, oefenen wij een vorm van belangloze aandacht die fundamenteel is voor elke morele houding. Het is alsof de natuur ons leert om te kijken zonder te grijpen, om te waarderen zonder te bezitten. Deze houding is precies wat ons in staat stelt om ook in morele kwesties niet louter ons eigen voordeel na te jagen, maar om het goede omwille van zichzelf te zoeken. De schoonheid der natuur is zo niet alleen een genot, maar een school voor de ziel. Zij leert ons om te zien wat werkelijk is, in plaats van wat wij graag zouden willen dat het is.

De minachting voor hen die alleen oog hebben voor zintuiglijk genot, is daarom geen elitair oordeel, maar een erkenning van een gemis. Wie niet ontroerd kan worden door de eenvoudige schoonheid van een bloem of het gezang van een vogel, mist niet alleen een esthetische gevoeligheid, maar ook een moreel kompas. Want wie niet kan genieten van wat geen nut heeft, zal ook moeite hebben om te handelen volgens principes die geen beloning beloven. In die zin is de liefde voor natuurlijke schoonheid geen luxe, maar een noodzakelijke voorwaarde voor ware menselijkheid. Zij herinnert ons eraan dat het goede, net als het schone, geen middel is, maar een doel op zich. En dat is misschien wel de meest diepgaande les die Kant ons hier geeft: dat ware verfijning niet bestaat in wat wij bezitten, maar in hoe wij aanschouwen.

\bigskip
Voor vrijmetselaren heeft deze paragraaf een bijzondere diepgang, omdat hij de essentie raakt van wat maçonnieke symboliek en ritueel beogen: \textbf{een school voor de ziel}. De schoonheid der natuur, die Kant hier beschrijft als een belofte van morele ontvankelijkheid, is precies wat het maçonnieke ritueel nastreeft. Net zoals de natuur ons leert om te aanschouwen zonder te bezitten, zo leert het ritueel de broeder om te zien zonder direct te begrijpen, om te ervaren zonder meteen toe te passen. Het is geen toeval dat de loge vaak wordt vergeleken met een tuin of een bos, een plaats waar de mens zich niet als meester, maar als leerling opstelt.

Wat Kant beschrijft als het ``onmiddellijke belang'' in natuurlijke schoonheid, is verwant aan de maçonnieke ervaring van symbolen. Een symbool in de loge, zoals passer en winkelhaak, het licht, of de ruwe steen, heeft geen nut in de wereldse zin. Zijn waarde ligt niet in wat het doet, maar in wat het openbaart: een waarheid die alleen zichtbaar wordt voor wie bereid is om zonder vooroordeel te aanschouwen. Net zoals de schoonheid van een bloem verdwijnt wanneer zij kunstmatig blijkt, zo verliest een symbool zijn kracht wanneer het louter als conventie wordt gezien. De kracht ervan ligt in het feit dat het, net als de natuur, een taal spreekt die niet voorschrijft, maar uitnodigt. Het vraagt om een houding van ontvankelijkheid, van ``zien zonder te begrijpen'', precies de houding die Kant hier als moreel fundamenteel beschrijft.

Bovendien is de loge zelf een microkosmos van die natuurlijke schoonheid waar Kant over spreekt. De harmonie van het ritueel, de ordening van de tempel, de stilte en het licht, het zijn allemaal uitingen van een schoonheid die niet bedoeld is om te behagen, maar om te verheffen. Wie daarvoor ontvankelijk is, oefent niet alleen zijn esthetisch vermogen, maar ook zijn morele gevoeligheid. Want zoals Kant stelt: wie leert om schoonheid om haarzelf te waarderen, leert ook om het goede omwille van zichzelf na te streven. Voor de vrijmetselaar is dit geen abstractie, maar een dagelijkse praktijk. Het ritueel is geen doel op zich, maar een middel om de geest te verfijnen, zodat hij in staat wordt om in de wereld te handelen volgens principes die geen beloning vragen.

Ten slotte ligt hierin ook de verbinding met het maçonnieke ideaal van broederschap. De schoonheid die Kant beschrijft, is geen individuele ervaring, maar een universele. Zo is ook de waarheid die in de loge wordt gezocht geen privé-bezit, maar iets wat alleen werkelijk schijnt wanneer het gedeeld wordt. Wie leert om de schoonheid der natuur, of die van het ritueel, te zien, leert tegelijkertijd om de waardigheid van elk mens te erkennen. En dat is de kern van ware vrijmetselarij: niet de zoektocht naar geheimen, maar de oefening in het herkennen van wat alle mensen verbindt.
\end{mdframed}

\section{Over de kunst in het algemeen.}

1) \textbf{Kunst} onderscheidt zich van de \textbf{natuur} zoals het doen (\textit{facere}) zich onderscheidt van het handelen of voortbrengen (\textit{agere}) in het algemeen, en het product of gevolg van de eerste onderscheidt zich als een \textbf{werk} (\textit{opus}) van het laatste als een effect (\textit{effectus}).

Met recht zou alleen voortbrenging door vrijheid, dat wil zeggen door een vermogen tot keuze dat zijn handelingen in de rede grondslagen, kunst moeten worden genoemd. Want hoewel men graag het product van de bijen (de regelmatig gebouwde raten) als een kunstwerk beschrijft, gebeurt dit alleen vanwege de analogie met het laatste; dat wil zeggen, zodra wij ons herinneren dat zij hun werk niet grondvesten op enige redelijke overweging van hun eigen, zeggen wij dat het een product van hun natuur (van instinct) is, en als kunst wordt het alleen aan hun Schepper toegeschreven.

Als iemand die door een moeras loopt, zoals soms gebeurt, een stuk bewerkt hout vindt, zegt hij niet dat het een product van de natuur is, maar van kunst; de oorzaak die het voortbracht, bedacht een doel, waaraan het hout zijn vorm te danken heeft. Ook in andere gevallen ziet men kunst in alles wat zo is samengesteld dat een voorstelling ervan in zijn oorzaak aan zijn werkelijkheid moet zijn voorafgegaan (zoals zelfs bij bijen het geval is), hoewel het misschien niet precies aan het effect heeft \textbf{gedacht}; maar als iets zonder meer een kunstwerk wordt genoemd, om het te onderscheiden van een effect van de natuur, dan wordt daarmee altijd een werk van mensen bedoeld.

2) \textbf{Kunst} als een vaardigheid van mensen onderscheidt zich ook van \textbf{wetenschap} (\textbf{het kunnen} van \textbf{het weten}), zoals een praktisch vermogen zich onderscheidt van een theoretisch, zoals techniek zich onderscheidt van theorie (zoals landmeten zich onderscheidt van meetkunde). En dus wordt datgene wat men \textbf{kan} doen zodra men weet wat gedaan moet worden, niet precies kunst genoemd. Alleen datgene wat men niet onmiddellijk vaardig kan doen, zelfs als men het volledig kent, behoort in die mate tot de kunst. 
\textbf{Camper} beschrijft heel precies hoe de beste schoen gemaakt moet worden, maar hij zou er zeker niet een hebben kunnen maken.\footnote{In mijn streek zegt de gewone man, wanneer hij geconfronteerd wordt met een probleem zoals dat van Columbus en zijn ei: ``Dat is geen kunst, dat is gewoon wetenschap.'' Dat wil zeggen: als men het weet, dan kan men het ook doen. En hij zegt hetzelfde over alle zogenaamde kunsten van de goochelaar. Maar hij zou nooit weigeren om de kunsten van de koorddanser kunst te noemen.}

3) \textbf{Kunst} onderscheidt zich ook van \textbf{ambacht}: de eerste wordt \textbf{vrije kunst} genoemd, de tweede kan ook \textbf{belonende kunst} worden genoemd. De eerste wordt beschouwd alsof zij alleen doelmatig (succesvol) kan zijn als spel, dat wil zeggen, een bezigheid die op zichzelf aangenaam is; de tweede wordt beschouwd als arbeid, dat wil zeggen, een bezigheid die op zichzelf onaangenaam (vermoeiend) is en alleen aantrekkelijk is vanwege zijn effect (bijvoorbeeld de beloning), en dus als iets dat gedwongen kan worden opgelegd. Of men in de hiërarchie van de gilden horlogemakers als kunstenaars moet tellen, maar smeden als ambachtslieden, vereist een ander standpunt dan hier wordt ingenomen, namelijk de verhouding van de talenten waarop de ene of de andere van deze bezigheden moet worden gegrondvest. Verder zal ik hier niet bespreken of onder de zogenaamde zeven vrije kunsten er misschien enkele zijn opgenomen die tot de wetenschappen moeten worden gerekend, en verschillende andere die met ambachten moeten worden vergeleken.
Maar het is niet onverstandig om te bedenken dat in alle vrije kunsten toch iets dwingends vereist is, of, zoals het wordt genoemd, een \textbf{mechanisme}, zonder welke de \textbf{geest}, die in de kunst \textbf{bevrijd} moet worden en die alleen het werk bezielt, geheel geen lichaam zou hebben en geheel zou verdampen (bijvoorbeeld in de dichtkunst, de juistheid en rijkdom van de woordkeus alsmede de prosodie en het metrum), aangezien vele moderne leraren menen dat zij een vrije kunst het best kunnen bevorderen als zij alle dwang ervan wegnemen en haar van arbeid in louter spel veranderen.

\begin{mdframed}[leftmargin=0pt,rightmargin=0pt,backgroundcolor=blue!6,linecolor=white,innertopmargin=6pt,innerbottommargin=6pt]

Kunst is niet wat de natuur vanzelf voortbrengt, maar wat de mens met vrijheid en bedoeling schept. Wanneer wij spreken van kunst, bedoelen wij geen instinctmatige handelingen, zoals die van bijen die hun raten bouwen, maar een bewuste daad die haar vorm ontleent aan een vooraf bedacht doel. Een stuk bewerkt hout dat wij in de natuur vinden, herkennen wij niet als een product van toeval, maar als het werk van een geest die een plan had. Kunst is dus altijd verbonden met menselijke vrijheid: zij vereist dat de maker weet wat hij doet, en dat hij het doet omdat hij het wil.

Maar kunst is meer dan louter kennis. Wetenschap leert ons wat gedaan moet worden, maar kunst leert ons hoe het gedaan moet worden, en dat is geen kwestie van theorie, maar van vaardigheid. Wie precies weet hoe een schoen gemaakt moet worden, kan hem nog niet noodzakelijk vervaardigen. Kunst is dus niet alleen weten, maar ook kunnen, en wel op een manier die niet louter mechanisch is.

Daarbij is er een wezenlijk onderscheid tussen vrije kunst en ambacht. Vrije kunst is een bezigheid die op zichzelf bevrediging schenkt, zoals een spel dat om het spel zelf wordt gespeeld. Ambacht daartegen is arbeid: het is vermoeiend op zichzelf en wordt alleen aantrekkelijk door het resultaat, zoals loon of nut. Toch mag men niet denken dat vrije kunst geheel zonder regels of dwang kan bestaan. Zelfs in de meest vrije kunsten, zoals de dichtkunst, is er een zekere mechaniek nodig, zoals de juiste woordkeus, het metrum, de prosodie, zonder welke de geest die het werk bezielt geen vaste vorm zou hebben. Wie alle dwang uit de kunst wegneemt, vernietigt haar wezen. Want kunst is geen willekeur, maar een vrije daad die toch aan regels gehoorzaamt, niet om te onderdrukken, maar om de geest te bevrijden.

\bigskip
Wat deze paragraaf ons leert, is dat kunst niet louter een techniek is, maar een uiting van menselijke vrijheid die tegelijkertijd gebonden is aan regels. Dit is geen tegenstrijdigheid, maar de essentie van wat kunst werkelijk maakt. De bij die haar raat bouwt, handelt uit instinct; de mens die een kunstwerk schept, handelt uit vrijheid, maar die vrijheid is geen willekeur. Zij vereist kennis, vaardigheid en een bewust doel. Hierin ligt de diepere betekenis van kunst: zij is geen louter mechanische handeling, maar een daad die de geest bevrijdt door hem te binden aan een vorm.

Het onderscheid tussen vrije kunst en ambacht onthult iets fundamenteels over de menselijke natuur. Ambacht is arbeid die beloond wordt door haar resultaat; vrije kunst is een bezigheid die beloond wordt door de daad zelf. Toch is zelfs de vrije kunst niet geheel vrij van dwang. Zonder regels, zonder de mechaniek van vorm en structuur, zou de geest die de kunst bezielt geen houvast hebben en verdampen in leegte. Dit is geen beperking, maar een voorwaarde voor ware vrijheid. Net zoals de morele wet geen ketens zijn, maar de voorwaarde voor ware morele handelingen, zo zijn de regels van de kunst geen belemmeringen, maar de grond waarop de geest zich kan verheffen.

Dit inzicht is niet alleen van belang voor de kunstenaar, maar voor ieder die streeft naar ware vrijheid. Kunst leert ons dat ware bevrijding niet bestaat in het ontbreken van regels, maar in het vermogen om binnen die regels een eigen vorm te scheppen. Zo is kunst niet alleen een bezigheid, maar een metafoor voor het menselijk streven zelf: het vermogen om, binnen de grenzen van onze wereld, een eigen betekenis te geven aan wat wij doen. In die zin is kunst niet alleen een product, maar een oefening in menselijkheid.

\bigskip
Voor vrijmetselaren heeft deze paragraaf een diepe, symbolische betekenis, omdat hij de essentie raakt van wat maçonniek werk werkelijk is: een vrije daad binnen een gebonden vorm. Het ritueel, de symbolen, de structuur van de loge, het zijn allemaal ``regels'' die niet bedoeld zijn om de geest te beperken, maar om hem te bevrijden. Net zoals Kant stelt dat ware kunst niet zonder mechaniek kan bestaan, zo is ook het maçonnieke ritueel geen leeg spel, maar een gestructureerde weg naar inzicht. De vorm is niet het doel, maar de voorwaarde: zonder de vaste ordening van het ritueel, zonder de symbolische taal die generaties lang is doorgegeven, zou de geest die erin werkt geen houvast hebben. Hij zou verdampen in willekeur, in persoonlijke interpretaties die geen verbinding meer hebben met het universele.

Wat deze paragraaf vooral duidelijk maakt, is dat vrijmetselarij geen ambacht is, maar een vrije kunst. Het is geen arbeid die beloond wordt door een extern resultaat, maar een bezigheid die op zichzelf betekenisvol is. De broeder die het ritueel uitvoert, doet dat niet voor loon of erkenning, maar omdat de daad zelf hem verheft. Toch is die daad geen willekeurige expressie. Hij is gebonden aan een traditie, aan een vorm die hem draagt en die hem in staat stelt om zijn eigen, persoonlijke zoektocht te verbinden met die van alle broeders die voor hem kwamen en na hem zullen volgen. In die zin is het maçonnieke werk een school voor de vrijheid: het leert de broeder om binnen de grenzen van de vorm zijn eigen, authentieke weg te vinden.

Bovendien onthult deze paragraaf waarom de loge zo'n belangrijke rol speelt. Net zoals kunst niet kan bestaan zonder een zekere mechaniek, zo kan de vrijmetselaar zijn innerlijke werk niet doen zonder de steun van de gemeenschap. De loge is niet alleen een plaats, maar een levend kader dat de individuele geest bevrijdt door hem te verbinden met iets groters. Zo wordt het ritueel geen mechanische handeling, maar een vrije daad die toch universeel is, precies zoals Kant beschrijft dat ware kunst zowel persoonlijk als algemeen betekenisvol moet zijn. Voor de vrijmetselaar is dit geen abstractie, maar een dagelijkse ervaring: in het ritueel leert hij om, binnen de vaste vorm, zijn eigen, unieke bijdrage te leveren aan een gedeelde zoektocht naar waarheid. Dat is de ware kunst van de vrijmetselarij: niet het volgen van regels, maar het gebruik ervan om de geest te verheffen.
    
\end{mdframed}

\section{Over de schone kunst.}

Er bestaat noch een wetenschap van het schone, maar slechts een kritiek ervan, noch een schone wetenschap, maar slechts schone kunst. Want als de eerste zou bestaan, dan zou daarin wetenschappelijk, dat wil zeggen door middel van bewijzen, bepaald worden of iets als mooi beschouwd moet worden of niet; zo zou het oordeel over schoonheid, als het tot een wetenschap behoorde, geen smaakoordeel zijn.
Wat betreft het tweede: een wetenschap die als zodanig mooi zou moeten zijn, is belachelijk. Want als men in haar, als wetenschap, naar gronden en bewijzen zou vragen, zou men worden weggestuurd met smaakvolle uitdrukkingen (\textit{bons mots}). - Wat zonder twijfel aanleiding heeft gegeven tot de gebruikelijke uitdrukking \textbf{schone wetenschappen} is niets anders dan het feit dat terecht is opgemerkt dat voor de schone kunst in haar volle volmaaktheid veel wetenschap vereist is, zoals bijvoorbeeld kennis van oude talen, uitgebreide lectuur van die auteurs die als klassiek worden beschouwd, geschiedenis, kennis van oudheden, enzovoort. En om die reden zijn deze historische wetenschappen, omdat zij de noodzakelijke voorbereiding en grond vormen voor de schone kunst, en ook gedeeltelijk omdat de kennis van de producten van de schone kunst (welsprekendheid en poëzie) zelfs daarin is opgenomen, door een verbale verwarring zelf schone wetenschappen genoemd.

Als kunst, die voldoende is voor de \textbf{kennis} van een mogelijk voorwerp, slechts de handelingen verricht die nodig zijn om het werkelijk te maken, is zij \textbf{mechanisch}; maar als zij het gevoel van genot als haar onmiddellijke doel heeft, dan wordt zij \textbf{esthetische} kunst genoemd. Deze is ofwel \textbf{aangename} ofwel \textbf{schone} kunst. Het is de eerste als haar doel is dat genot de voorstellingen als louter \textbf{sensaties} begeleidt, de laatste als haar doel is dat het deze als \textbf{soorten van kennis} begeleidt.

Aangename kunsten zijn die welke slechts op genot gericht zijn; van deze aard zijn alle bekoringen die het gezelschap aan tafel kunnen bevredigen, zoals het vertellen van vermakelijke verhalen, het aanmoedigen van het gezelschap om open en levendig te praten, het scheppen van een zekere toon van vrolijkheid door middel van grappen en gelach, waarbij, zoals men zegt, veel kan worden gekletst en niemand verantwoordelijk wordt gehouden voor wat hij zegt, omdat het alleen bedoeld is als vluchtige vermaking, niet als enige blijvende stof voor latere reflectie of discussie. (Hiertoe behoort ook de manier waarop de tafel voor het genot is gedekt, of zelfs, bij grote partijen, de tafelmuziek – een vreemd iets, dat slechts bedoeld is om de stemming van vreugde in stand te houden als een aangename herrie, en om de vrije conversatie van de ene buurman met de andere aan te moedigen zonder dat iemand de minste aandacht schenkt aan de samenstelling ervan.) Ook alle spelen die geen ander belang dienen dan het onopgemerkt laten voorbijgaan van de tijd, behoren hiertoe.

Schone kunst daartegen is een soort voorstelling die op zichzelf doelmatig is en, hoewel zonder doel, niettemin de cultivering van de geestesvermogens voor de maatschappelijke omgang bevordert.

De universele mededeelbaarheid van een genot sluit in haar begrip al in dat dit geen genot van genieten mag zijn, dat louter uit sensatie voortkomt, maar een genot van reflectie; en dus is esthetische kunst, als schone kunst, er een die het reflecterend oordeelsvermogen en niet louter sensatie als haar maatstaf heeft.

\begin{mdframed}[leftmargin=0pt,rightmargin=0pt,backgroundcolor=blue!6,linecolor=white,innertopmargin=6pt,innerbottommargin=6pt]

Schoonheid laat zich niet wetenschappelijk vastleggen. Er bestaat geen wetenschap die kan bewijzen wat mooi is, want dan zou het geen kwestie meer zijn van smaak, maar van kennis. Evenmin is er zoiets als een ``schone wetenschap''; die term is alleen ontstaan doordat men inzag dat ware kunst kennis vereist, zoals taal, geschiedenis en klassieke literatuur. Deze disciplines zijn geen kunst op zich, maar slechts de voorbereiding ertoe.

Kunst die louter gericht is op het maken van dingen, is mechanisch. Maar kunst die het genot zelf als doel heeft, noemen wij esthetische kunst. Daarbij is een wezenlijk onderscheid: aangename kunst richt zich op vluchtig plezier, zoals tafelgesprekken, grappen of muziek die alleen dient om de stemming te verheffen zonder dat men er echt naar luistert. Schone kunst daartegen is meer dan louter vermaak. Zij is op zichzelf betekenisvol, niet omdat zij een bepaald doel dient, maar omdat zij de geest verfijnt en de menselijke vermogens ontwikkelt voor het gemeenschappelijke leven. Haar waarde ligt niet in de sensatie, maar in de reflectie die zij oproept. Ware kunst is daarom geen kwestie van oppervlakkig genot, maar van een ervaring die gedeeld en overwogen kan worden, niet om de tijd te doden, maar om hem zin te geven.

\bigskip
Wat deze paragraaf ons leert, is dat ware kunst geen kwestie is van techniek of vermaking, maar van betekenis. Kant onthult hier dat schoonheid niet kan worden gereduceerd tot wetenschappelijke principes of louter zintuiglijk genot. Ware kunst is geen ambacht, geen mechanische vaardigheid, maar een daad die de geest in beweging zet. Zij is geen middel om de tijd te doden, maar een manier om de tijd te verheffen; om hem te vullen met iets dat niet alleen aangenaam is, maar ook betekenisvol.

Het onderscheid tussen aangename en schone kunst is cruciaal. Aangename kunst, zoals tafelmuziek of gebabbel, is vluchtig en oppervlakkig. Zij vraagt niets van ons behalve een passieve ontvankelijkheid. Schone kunst daartegen daagt ons uit. Zij is niet gericht op louter genot, maar op reflectie, op een ervaring die wij niet alleen ondergaan, maar ook kunnen delen en overdenken. In die zin is schone kunst geen luxe, maar een noodzaak: zij ontwikkelt in ons het vermogen om niet alleen te voelen, maar ook te begrijpen, om niet alleen te genieten, maar ook te overwegen.

Dit inzicht is niet alleen van belang voor de kunstenaar, maar voor ieder die streeft naar een verfijnd en betekenisvol leven. Ware kunst leert ons dat genot niet het hoogste doel is, maar slechts een begin. Zij toont ons dat ware bevrediging niet ligt in het vluchtige, maar in het blijvende, in wat ons niet alleen even afleidt, maar wat ons verrijkt en verbindt met anderen. In die zin is schone kunst geen vlucht uit de wereld, maar een manier om dieper in de wereld te staan: om niet alleen te leven, maar om bewust te leven. Dat is de ware betekenis van kunst: niet om ons te vermaken, maar om ons te verheffen.

\end{mdframed}

\section{Schone kunst is kunst voor zover zij tegelijkertijd als natuur lijkt.}

In een product van kunst moet men zich ervan bewust zijn dat het kunst is en geen natuur; toch moet de doelmatigheid in haar vorm zo vrij lijken van elke dwang door willekeurige regels, alsof het louter een product van de natuur was. Op dit gevoel van vrijheid in het spel van onze kennissenvermogens, dat toch tegelijkertijd doelmatig moet zijn, berust dat genot dat alleen universeel mededeelbaar is, zonder op begrippen te berusten. De natuur was schoon, als zij tegelijkertijd als kunst leek; en kunst kan alleen schoon worden genoemd als wij ons ervan bewust zijn dat het kunst is, en toch lijkt zij ons als natuur.

Want wij kunnen in het algemeen zeggen, of het nu de schoonheid van de natuur of van de kunst betreft: dat is schoon wat bevalt in het louter oordelen (noch in de sensatie, noch door een begrip). Nu heeft kunst altijd een bepaalde bedoeling om iets voort te brengen. Als deze bedoeling echter slechts een sensatie (iets louter subjectiefs) was die met genot gepaard zou moeten gaan, dan zou dit product in het oordelen alleen door het zintuiglijk gevoel behagen. Als de bedoeling gericht was op de voortbrenging van een bepaald voorwerp, dan zou, als dit door kunst werd bereikt, het voorwerp alleen door begrippen behagen. Maar in beide gevallen zou de kunst niet behagen in het louter oordelen, dat wil zeggen, zij zou niet behagen als schoon, maar als mechanische kunst.

De doelmatigheid in het product van de schone kunst, hoewel zij zeker opzettelijk is, mag toch niet opzettelijk lijken; dat wil zeggen, schone kunst moet als natuur worden beschouwd, hoewel men zich natuurlijk bewust is dat het kunst is. Een product van kunst lijkt echter op natuur, als wij vinden dat het nauwkeurig, maar niet moeizaam, overeenstemt met de regels volgens welke het product alleen kan worden wat het moet zijn, dat wil zeggen, zonder dat de academische vorm erdoorheen schemert, zonder enig teken te tonen dat de regel voor de ogen van de kunstenaar heeft gezworven en zijn geestesvermogens heeft geboeid.

\begin{mdframed}[leftmargin=0pt,rightmargin=0pt,backgroundcolor=blue!6,linecolor=white,innertopmargin=6pt,innerbottommargin=6pt]

Schone kunst bereikt haar hoogste vorm wanneer zij zo natuurlijk lijkt, dat wij bijna vergeten dat zij door mensenhanden is gemaakt. Het is essentieel dat wij weten dat het kunst is, maar de schoonheid ervan moet zo vrij en ongedwongen overkomen, alsof zij zonder moeite of dwang is ontstaan, zoals de natuur zelf. Dit gevoel van vrijheid, gecombineerd met een diepere doelmatigheid, maakt dat schoonheid universeel aanspreekt, zonder dat wij precies kunnen uitleggen waarom.

Wanneer kunst louter bedoeld is om een bepaalde sensatie op te roepen of een specifiek doel te dienen, verliest zij haar ware schoonheid. Zij wordt dan ofwel louter vermakelijk, ofwel een kwestie van techniek. Ware schone kunst daartegen onthult haar regels niet als opgedrongen, maar als vanzelfsprekend. Zij volgt haar eigen wetten, maar doet dat met zo’n lichtheid, dat het lijkt alsof er geen regels zijn; alsof de kunstenaar niet heeft nagedacht over elke stap, maar alsof het werk vanzelf zo moest zijn. Pas dan spreekt kunst niet alleen tot de zintuigen of het verstand, maar tot de geest die vrij en toch doelmatig oordeelt. Zo wordt kunst niet alleen een product, maar een ervaring die ons doet vergeten dat zij gemaakt is, en ons toch bewust maakt van haar volmaaktheid.

\bigskip
Deze paragraaf raakt aan de kern van maçonnieke symboliek en ritueel. Wat Kant hier beschrijft als de hoogste vorm van schone kunst kunst die zo natuurlijk lijkt dat zij haar eigen regels verbergt is precies wat het maçonnieke ritueel nastreeft. Het ritueel is geen mechanische handeling, maar een levende ervaring die alleen haar volle kracht ontvindt wanneer zij niet als een opgedrongen vorm wordt ervaren, maar als iets dat vanzelfsprekend en diep waar *voelt*. Net zoals schone kunst haar regels niet als dwang toont, maar als een organische ordening, zo moet ook het maçonnieke ritueel aanvoelen als een natuurlijke uiting van een hogere waarheid, niet als een stijve ceremonie.

Voor de vrijmetselaar ligt de betekenis in de spanning tussen vorm en vrijheid. Het ritueel is strikt gestructureerd, met symbolen, woorden en handelingen die precies zijn voorgeschreven. Toch is het doel niet om deze regels als beperking te ervaren, maar om ze zo te beheersen dat zij als vanzelfsprekend worden zoals een meester-muzikant die de techniek zo volledig beheerst dat de muziek *leeft*. Pas dan kan het ritueel de geest bevrijden in plaats van te boeien. Het is deze paradox die Kant hier blootlegt: ware kunst (en ware rituele beoefening) vereist een diepgaande kennis van de regels, maar de ervaring ervan moet vrij zijn, alsof het geen inspanning kost. Voor de vrijmetselaar is dit geen tegenstrijdigheid, maar de essentie van de zoektocht: de vorm is niet het doel, maar de weg naar een ervaring die zowel persoonlijk als universeel is.

Bovendien onthult deze paragraaf waarom maçonnieke symbolen zo krachtig zijn. Zij werken niet omdat zij worden "begrepen" in intellectuele zin, maar omdat zij *ervaren* worden als iets dat klopt, zonder dat men precies kan uitleggen waarom. Net zoals schone kunst ons raakt zonder dat wij haar kunnen analyseren, zo spreekt het ritueel tot het diepste in de broeder niet door uitleg, maar door beleving. De loge wordt zo een plaats waar kunst en natuur, vrijheid en ordening, samenkomen: een ruimte waar de mens leert om binnen de grenzen van de traditie zijn eigen, authentieke weg te vinden. Dat is de ware "natuur" van de vrijmetselarij: niet het volgen van regels om de regels, maar het gebruik ervan om de geest te verheffen tot iets dat groter is dan zichzelf. In die zin is het maçonnieke werk zelf een vorm van schone kunst een daad die zowel vrij als bindend is, net als de schoonheid waar Kant over spreekt.
\end{mdframed}

\section{Schone kunst is kunst van genie.}

\textbf{Genie} is het talent (natuurlijke gave) dat de regel aan de kunst geeft. Aangezien het talent, als een aangeboren productieve vermogen van de kunstenaar, zelf tot de natuur behoort, kan dit ook als volgt worden uitgedrukt: \textbf{Genie} is de aangeboren aanleg van de geest (\textit{ingenium}), \textbf{waardoor} de natuur de regel aan de kunst geeft.

Hoe het ook zij met deze definitie, en of zij louter willekeurig is of wel voldoet aan het begrip dat gewoonlijk met het woord \textbf{genie} wordt verbonden (wat in de volgende paragrafen zal worden besproken), toch kan reeds van meet af aan worden aangetoond dat, volgens de betekenis die hier aan het woord wordt gegeven, schone kunsten noodzakelijk als kunsten van \textbf{genie} moeten worden beschouwd.

Elke kunst veronderstelt immers regels die eerst de grond leggen waardoor een product dat kunstwerk moet heten, als mogelijk wordt voorgesteld. Het begrip van schone kunst staat echter niet toe dat het oordeel over de schoonheid van haar product wordt afgeleid van enige regel die een \textbf{begrip} als haar bepalende grond heeft, en dus een begrip van hoe het mogelijk is als grond heeft. Zo kan de schone kunst zelf de regel niet bedenken volgens welke zij haar product tot stand moet brengen. Maar aangezien zonder een voorafgaande regel een product nooit kunst kan worden genoemd, moet de natuur in het subject (en door middel van de aanleg van zijn vermogens) de regel aan de kunst geven, dat wil zeggen: schone kunst is alleen mogelijk als een product van genie.

Hieruit blijkt: 1)~dat genie een \textbf{talent} is om dat voort te brengen waarvoor geen bepaalde regel kan worden gegeven, geen aanleg van vaardigheid voor datgene wat volgens enige regel kan worden geleerd; bijgevolg moet \textbf{originaliteit} zijn voornaamste kenmerk zijn. 2)~Dat er, aangezien er ook oorspronkelijke onzin kan bestaan, zijn producten tegelijkertijd modellen moeten zijn, dat wil zeggen \textbf{voorbeeldig}; zij zijn dus, hoewel zelf niet het resultaat van nabootsing, toch voor anderen een maatstaf of regel voor het oordelen. 3)~Dat het genie zelf niet wetenschappelijk kan beschrijven of aangeven hoe het zijn product tot stand brengt, maar dat het de regel als \textbf{natuur} geeft; en daarom weet de maker van een product dat hij aan zijn genie te danken heeft, zelf niet hoe de ideeën ervoor bij hem opkomen, en heeft hij het ook niet in zijn vermogen om dergelijke dingen naar willekeur of volgens plan te bedenken en anderen voorschriften te geven die hen in staat zouden stellen om gelijksoortige producten voort te brengen. (Want zo is waarschijnlijk ook het woord ``genie'' afgeleid van \textit{genius}, in de zin van de bijzondere geest die een mens bij zijn geboorte wordt gegeven, die hem beschermt en leidt, en waarvan de inspiratie die oorspronkelijke ideeën voortbrengt, afstamt.) 4)~Dat de natuur door middel van genie geen regel voorschrijft aan de wetenschap, maar aan de kunst, en zelfs aan deze alleen voor zover zij schone kunst moet zijn.

\begin{mdframed}[leftmargin=0pt,rightmargin=0pt,backgroundcolor=blue!6,linecolor=white,innertopmargin=6pt,innerbottommargin=6pt]

Schone kunst ontstaat niet door het toepassen van vaste regels of technieken die men kan leren, maar door genie: een aangeboren vermogen dat de kunstenaar in staat stelt om iets te scheppen waarvoor geen vooraf bepaalde voorschriften bestaan. Genie is geen vaardigheid die men kan aanleren, maar een natuurlijke gave die de kunst haar eigen regels geeft. Het kenmerkt zich door originaliteit, want ware kunst kan niet voortkomen uit nabootsing of mechanische toepassing. Toch zijn de werken die uit genie voortkomen niet willekeurig; zij dienen als voorbeeld, als een maatstaf voor anderen, zonder dat de kunstenaar zelf precies kan uitleggen hoe hij tot zijn schepping is gekomen. Het genie handelt niet volgens een bewust plan, maar als een natuurkracht die inspiratie schenkt.

Daarom kan genie geen wetenschap voortbrengen, maar alleen kunst en dan nog wel alleen die kunst die werkelijk schoon is. Want waar wetenschap berust op begrippen en regels, daagt genie de kunstenaar uit om iets te creëren dat zowel uniek als universeel aanspreekt, zonder dat hij zelf de weg ernaartoe volledig kan verklaren. Zo is genie niet alleen een persoonlijk talent, maar een bron van inspiratie die de kunst haar diepste betekenis geeft.

\bigskip
Deze paragraaf onthult dat ware kunst niet voortkomt uit techniek of regels, maar uit een diepere bron van scheppende kracht het genie. Kant benadrukt hier dat schoonheid niet kan worden "geleerd" of "voorgeschreven", maar alleen kan ontspringen uit een natuurlijke aanleg die de kunstenaar in staat stelt om iets te scheppen dat zowel origineel als voorbeeldig is. Het genie is geen kwestie van vaardigheid, maar van inspiratie: een innerlijke stem die de kunstenaar leidt zonder dat hij zelf precies kan verklaren hoe. Dit is geen willekeur, maar een vorm van scheppende vrijheid die de kunst haar ware waarde geeft.

Wat deze paragraaf zo diepzinnig maakt, is dat hij de kunstenaar niet als een ambachtsman beschrijft, maar als een medium iemand die iets tot uiting brengt dat groter is dan zichzelf. Het genie is als een natuurkracht die door de kunstenaar heen werkt, en daarin ligt de paradox: hoewel de kunstenaar de regels niet bewust kan formuleren, schept hij toch iets dat anderen als maatstaf kunnen ervaren. Dit is geen tegenstrijdigheid, maar de essentie van ware kunst: zij is zowel persoonlijk als universeel, zowel vrij als bindend.

Voor de filosofie is dit een cruciale gedachte, want het toont aan dat schoonheid niet louter een kwestie is van smaak of techniek, maar van een scheppende geest die de mens verbindt met iets dat hem overstijgt. Het genie is niet alleen een individueel talent, maar een uiting van de natuur zelf, die door de kunstenaar heen spreekt. Zo wordt kunst niet alleen een product, maar een openbaring een manier om de wereld te zien en te begrijpen die niet in woorden of concepten is vast te leggen, maar alleen in de ervaring zelf. In die zin is genie niet alleen een gave, maar een roeping: de kunstenaar is niet de meester van zijn werk, maar de dienaar van een schoonheid die hij zelf niet volledig kan beheersen, maar alleen kan laten spreken.

\end{mdframed}

\section{Verduidelijking en bevestiging van de bovenstaande uitleg van genie.}

Iedereen is het erover eens dat genie geheel tegengesteld is aan de \textbf{geest van nabootsing}. Nu leren niets anders is dan nabootsen, geldt zelfs de grootste aanleg voor leren, de gemakkelijkheid om te leren (vermogen) als zodanig, nog niet als genie. Maar zelfs als iemand zelfstandig denkt of schrijft, en niet slechts overneemt wat anderen hebben gedacht, ja, zelfs als hij veel voor kunst en wetenschap uitvindt, is dit nog geen juiste reden om zo’n groot \textbf{verstand} (in tegenstelling tot iemand die, omdat hij nooit meer kan doen dan leren en nabootsen, een \textbf{sufkop} wordt genoemd) een \textbf{genie} te noemen, aangezien juist dit soort dingen ook geleerd \textbf{kan} worden en dus nog steeds op de natuurlijke weg van onderzoek en reflectie volgens regels ligt, en niet specifiek verschilt van datgene wat met inspanning door middel van nabootsing kan worden verworven. Alles wat Newton in zijn onsterfelijke werk over de beginselen van de natuurfilosofie heeft uiteengezet, hoe groot een geest het ook vereiste om het te ontdekken, kan nog steeds worden geleerd; maar men kan niet leren om geïnspireerde poëzie te schrijven, hoe uitputtend alle regels voor de dichtkunst en hoe uitstekend de voorbeelden ervan ook mogen zijn. De reden is dat Newton alle stappen die hij moest zetten, vanaf de eerste elementen van de meetkunde tot zijn grote en diepzinnige ontdekkingen, niet alleen voor zichzelf, maar ook voor ieder ander geheel inzichtelijk kon maken en ze zo voor het nageslacht duidelijk kon uiteenzetten; maar geen Homerus of Wieland kan aangeven hoe zijn ideeën, die fantastisch en toch rijk aan gedachten zijn, in zijn hoofd ontstaan en samenkomen, omdat hij het zelf niet weet en het dus ook niet aan een ander kan leren. In het wetenschappelijke domein verschilt de grootste ontdekker dus slechts in graad van de meest ijverige nabootser en leerling, terwijl hij in wezen verschilt van iemand die door de natuur is begaafd voor schone kunst.
Dit is niet bedoeld om deze grote mannen, aan wie het menselijk geslacht zoveel verschuldigd is, te verlagen ten opzichte van die gunstelingen der natuur met betrekking tot hun talent voor schone kunst. In hun talent om steeds grotere volmaaktheid van kennis en alle nut die daarvan afhangt na te streven, en eveneens in de opvoeding van anderen tot het verwerven van dezelfde kennis, ligt het grote voordeel van dergelijke mensen boven hen die de eer hebben genieën te worden genoemd: want voor de laatsten komt de kunst ergens tot stilstand, omdat er een grens aan is gesteld die niet overschreden kan worden, die waarschijnlijk ook al lang is bereikt en niet verder kan worden uitgebreid; en bovendien kan een dergelijke vaardigheid niet worden doorgegeven, maar wordt aan ieder onmiddellijk door de hand der natuur toebedeeld, en sterft dus met hem, totdat de natuur ooit op soortgelijke wijze een ander begiftigt, die niets meer nodig heeft dan een voorbeeld om het talent waarvan hij zich bewust is op gelijke wijze te laten werken.

Aangezien de gave der natuur de regel aan de kunst (als schone kunst) moet geven, wat voor regel is dit dan? Zij kan niet in een formule worden gevangen om als voorschrift te dienen, want dan zou het oordeel over het schone volgens begrippen bepaalbaar zijn; in plaats daarvan moet de regel uit de daad zelf, dat wil zeggen uit het product, worden afgeleid, waartegen anderen hun eigen talent kunnen toetsen, door het als model te laten dienen niet om te \textbf{kopiëren}, maar om \textbf{na te volgen}. Hoe dit mogelijk is, is moeilijk uit te leggen. De ideeën van de kunstenaar wekken soortgelijke ideeën op bij zijn leerling, als de natuur hem met een gelijksoortige verhouding van geestesvermogens heeft uitgerust. De voorbeelden van schone kunst zijn dus het enige middel om deze aan het nageslacht door te geven, wat niet kan gebeuren door louter beschrijvingen (met name niet op het gebied van de redekunsten); en zelfs in dat laatste geval zijn het alleen die in oude en dode talen, die nu alleen als geleerde talen bewaard zijn gebleven, die klassiek kunnen worden.

Hoewel mechanische en schone kunst, de eerste als louter kunst van ijver en leren, de tweede als die van genie, sterk van elkaar verschillen, is er toch geen schone kunst waarin niet iets mechanisch, dat volgens regels kan worden begrepen en gevolgd, en dus iets \textbf{academisch corrects}, de wezenlijke voorwaarde van de kunst uitmaakt. Want er moet iets in worden gedacht als een doel, anders kan men haar product aan geen enkele kunst toeschrijven; het zou louter een product van toeval zijn. Maar om een doel in het werk na te streven, zijn bepaalde regels vereist, waarvan men zich niet kan onttrekken. Nu de originaliteit van zijn talent een (maar niet het enige) essentiële element van het karakter van het genie uitmaakt, menen oppervlakkige geesten dat zij niet beter kunnen aantonen dat zij bloeiende genieën zijn dan door zich vrij te verklaren van de academische dwang van alle regels, en zij menen dat men beter pronkt op een paard met de slinger dan op een dat goed is afgericht. Genie kan slechts rijk \textbf{materiaal} leveren voor producten van kunst; de uitwerking en vorm ervan vereisen een academisch geschoold talent, om er een gebruik van te maken dat aan de kracht van het oordeel kan weerstaan. Maar wanneer iemand zelfs in zaken van de meest zorgvuldige rationele onderzoekingen spreekt en oordeelt als een genie, dan is dat volkomen belachelijk; men weet niet goed of men meer moet lachen om de kwakzalver die zulk een mist om zichzelf verspreidt dat men niet duidelijk kan oordelen, maar des te meer aan de verbeelding kan toegeven, of om het publiek, dat vertrouwend denkt dat zijn onvermogen om duidelijk te herkennen en het meesterwerk van inzicht te begrijpen, voortkomt uit het feit dat er hele massa’s nieuwe waarheden op hem worden afgeworpen, in vergelijking waarmee detail (bereikt door zorgvuldige uitleggingen en academisch correct onderzoek van grondbeginselen) louter het werk van leken lijkt.

\begin{mdframed}[leftmargin=0pt,rightmargin=0pt,backgroundcolor=blue!6,linecolor=white,innertopmargin=6pt,innerbottommargin=6pt]

Genie is geen kwestie van leren of nabootsen, maar van een aangeboren, scheppend vermogen dat de kunst haar eigen regels geeft. Een groot wetenschapper als Newton kon zijn ontdekkingen stap voor stap uitleggen, zodat anderen ze konden leren en toepassen. Maar een dichter als Homerus of Wieland kan niet verklaren hoe zijn geïnspireerde ideeën ontstaan en dus ook niet leren hoe anderen hetzelfde kunnen doen. Waar wetenschap berust op regels en methoden die men kan overdragen, is ware kunst gebonden aan een talent dat niet kan worden geleerd of voorgeschreven. Het genie schept iets origineels, iets dat als voorbeeld kan dienen, maar niet kan worden gekopieerd.

Toch is zelfs in schone kunst een zekere mechaniek, een academische grondigheid, noodzakelijk. Zonder regels en structuur zou kunst louter toeval zijn. Het genie levert het rijke materiaal, maar de vorm en uitwerking vereisen vaardigheid en kennis. Wie denkt dat genie betekent dat men zich niets hoeft aan te trekken van regels, vergist zich: zonder discipline vervalt kunst in willekeur. Echte genialiteit toont zich niet in het verwerpen van regels, maar in het vermogen om binnen die regels iets unieks en tijdloos te scheppen iets dat anderen wel kunnen bewonderen en navolgen, maar niet precies kunnen reproduceren. Zo is genie niet alleen een gave, maar ook een verantwoordelijkheid: het vereist zowel vrijheid als meesterschap.

\bigskip
Kant gaat zo diep in op genie, omdat het de sleutel vormt tot het begrijpen van ware schoonheid en scheppende vrijheid. Voor hem is genie niet alleen een persoonlijk talent, maar een brug tussen natuur en kunst, tussen vrijheid en noodzakelijkheid. Het genie toont aan dat schoonheid niet louter een kwestie is van smaak of techniek, maar van een diepere, bijna mystieke inspiratie die de kunstenaar overstijgt. Dit is cruciaal voor zijn esthetica, want het verklaart hoe kunst universeel kan aanspreken, zonder dat zij zich laat vangen in regels of begrippen. Het genie is als een natuurkracht die door de kunstenaar heen werkt, en daarin ligt de paradox: kunst is zowel vrij als bindend, zowel persoonlijk als universeel.

\bigskip
Voor de vrijmetselaar heeft dit concept een bijzondere resonantie, omdat het symboliek en ritueel op een dieper niveau raakt:

\paragraph{Inspiratie versus regels:} Net zoals het genie de kunst haar regels geeft zonder ze bewust te formuleren, zo is ook het maçonnieke ritueel geen mechanische handeling, maar een levende ervaring. De symbolen en handelingen zijn niet willekeurig, maar zij ontlenen hun kracht aan een traditie die als een "natuurkracht" door de broeder heen werkt. Het ritueel is geen stijve ceremonie, maar een vorm die de geest bevrijdt om betekenis te ervaren die niet in woorden is vast te leggen. Dit doet denken aan het maçonnieke idee dat waarheid niet wordt "geleerd", maar ervaren net zoals het genie zijn ideeën niet kan uitleggen, maar alleen kan laten spreken.

\paragraph{Originaliteit en voorbeeldigheid:} Het genie schept iets origineels dat toch als voorbeeld kan dienen. Zo is ook het maçonnieke werk: iedere broeder interpreteert de symbolen op zijn eigen manier, maar binnen een gedeeld kader. Het ritueel is geen kopie, maar een persoonlijke, doch herkenbare uiting van een universele zoektocht. Dit is precies wat Kant beschrijft: een balans tussen vrijheid en structuur, tussen individuele ervaring en gedeelde betekenis.

\paragraph{De grens van het leerbare:} Kant benadrukt dat genie niet kan worden doorgegeven via regels of voorschriften. Evenzo is de essentie van de vrijmetselarij niet te vatten in dogma’s of handleidingen. Wat in de loge wordt overgedragen, is geen kennis, maar een ervaring een manier van zien en zijn die alleen kan worden begrepen door zelf deel te nemen. Dit verklaart waarom maçonnieke kennis niet "onderwezen" wordt, maar beleid wordt: zoals het genie zijn talent alleen kan tonen door te scheppen, zo kan de vrijmetselaar zijn inzichten alleen delen door te leven volgens de principes die hij in het ritueel ervaart.

\paragraph{Discipline en vrijheid:} Kant waarschuwt dat genie zonder academische vormgeving ( zonder "mechaniek") vervalt in willekeur. Zo is ook het maçonnieke ritueel geen vrijblijvend spel: de structuur is niet bedoeld om te beperken, maar om de geest te verheffen. Zonder de "academische" precisie van het ritueel zou de symboliek haar diepgang verliezen. Voor de vrijmetselaar is dit een herinnering dat ware vrijheid niet bestaat in het verwerpen van regels, maar in het vermogen om binnen die regels een eigen, authentieke weg te vinden.

\paragraph{Tijdloosheid en klassiekers:} Kant wijst erop dat alleen voorbeelden uit de klassieke kunst (in "dode talen") echt klassiek kunnen worden omdat zij een tijdloze waarheid uitdrukken die niet afhankelijk is van mode of context. Dit doet denken aan de maçonnieke traditie, die zich baseert op eeuwenoude symbolen en verhalen die steeds opnieuw betekenis krijgen voor elke generatie. Net zoals het genie zijn werk niet kan uitleggen, maar alleen kan tonen, zo is ook de maçonnieke symboliek geen kwestie van uitleg, maar van beleving.

\end{mdframed}

\section{Over de verhouding van genie tot smaak.}

Voor het \textbf{oordelen} over schone voorwerpen als zodanig is \textbf{smaak} vereist; maar voor de schone kunst zelf, dat wil zeggen voor het \textbf{voortbrengen} van dergelijke voorwerpen, is \textbf{genie} vereist.

Als genie wordt beschouwd als een talent voor schone kunst (wat de eigenlijke betekenis van het woord impliceert), en men dit met het oog daarop analyseert in de vermogens die samen moeten komen om een dergelijk talent te vormen, dan is het nodig om eerst precies het verschil te bepalen tussen de schoonheid van de natuur, waarvan het oordelen slechts smaak vereist, en de schoonheid van de kunst, waarvan de mogelijkheid (die ook in het oordelen over een dergelijk voorwerp in aanmerking moet worden genomen) genie vereist.

Een schoonheid van de natuur is een \textbf{mooi ding}; de schoonheid van de kunst is een \textbf{mooie voorstelling} van een ding.

Om een schoonheid van de natuur als zodanig te beoordelen, hoef ik niet eerst een begrip te hebben van wat voor soort ding het voorwerp zou moeten zijn; dat wil zeggen, het is niet nodig dat ik de materiële doelmatigheid (het doel) ken, maar de loutere vorm bevalt op zichzelf in het oordelen, zonder kennis van het doel. Maar als het voorwerp als een product van kunst wordt gegeven en als zodanig als mooi moet worden verklaard, dan moet, omdat kunst altijd een doel in de oorzaak (en haar causaliteit) veronderstelt, eerst een begrip de grond zijn van wat het ding zou moeten zijn. En omdat de overeenstemming van het veelvoudige in een ding met zijn innerlijke bepaling als doel de volmaaktheid van het ding is, zal bij het oordelen over de schoonheid van kunst ook de volmaaktheid van het ding in aanmerking moeten worden genomen, wat bij het oordelen over een natuurlijke schoonheid (als zodanig) niet eens een vraag is. — Weliswaar wordt bij het oordelen, met name over levende wezens in de natuur, zoals een mens of een paard, ook vaak objectieve doelmatigheid in aanmerking genomen om hun schoonheid te beoordelen; maar in dat geval is het oordeel ook niet langer louter esthetisch, dat wil zeggen, een louter smaakoordeel. De natuur wordt dan niet langer beoordeeld zoals zij als kunst verschijnt, maar voor zover zij werkelijk kunst is (al is het dan bovenmenselijke kunst); en het teleologische oordeel dient als grond voor het esthetische en als een voorwaarde waarvan het laatste rekening moet houden. In zo’n geval, als men bijvoorbeeld zegt: ``Dat is een mooie vrouw,'' denkt men in feite niets anders dan dat de natuur in haar gestalte de doelen van het vrouwelijke lichaam mooi voorstelt, want men moet verder kijken dan de loutere vorm, naar een begrip waarmee het voorwerp op een bepaalde manier wordt gedacht via een logisch voorwaardelijk esthetisch oordeel.

Schone kunst toont haar excellentie juist door dingen die in de natuur lelijk of onaangenaam zouden zijn, mooi te beschrijven. Furieën, ziekten, verwoestingen van oorlog en dergelijke kunnen, als schadelijke dingen, zeer mooi worden beschreven, ja, zelfs in de schilderkunst worden voorgesteld; alleen één soort lelijkheid kan niet op een manier die aan de natuur beantwoordt worden voorgesteld zonder alle esthetische bevrediging te vernietigen, en dus schoonheid in de kunst, namelijk die welke \textbf{afschuw} opwekt. Want in deze vreemde sensatie, die berust op louter verbeelding, wordt het voorwerp voorgesteld alsof het het genot oplegt dat wij toch met kracht weerstaan, en daarom wordt de artistieke voorstelling van het voorwerp in ons gevoel zelf niet langer onderscheiden van de aard van het voorwerp zelf, en het wordt dan onmogelijk om de eerste als mooi te beschouwen. De beeldhouwkunst, omdat in haar producten kunst bijna verward wordt met natuur, heeft ook de voorstelling van lelijke voorwerpen uitgesloten van haar beelden, en staat bijvoorbeeld toe dat de dood (in een mooie genie) of de geest van de oorlog (in de persoon van Mars) alleen via een allegorie of kenmerken die aantrekkelijk lijken, dus slechts indirect door middel van een interpretatie van de rede, en niet voor de esthetische oordeelskracht alleen, worden voorgesteld.

Zoveel over de mooie voorstelling van een voorwerp, die in feite slechts de vorm is van de presentatie van een begrip, waardoor het laatste universeel wordt medegedeeld. — Om deze vorm aan het product van schone kunst te geven, is echter slechts smaak vereist, waaraan de kunstenaar, nadat hij deze heeft geoefend en gecorrigeerd door middel van verschillende voorbeelden van kunst of natuur, zijn werk toetst, en na vele, vaak moeizame pogingen om deze tevreden te stellen, de vorm vindt die hem bevredigt; dit is dus niet als het ware een kwestie van inspiratie of een vrije uiting van de geestesvermogens, maar een langzame en zelfs moeizame verbetering, om deze geschikt te maken voor de gedachte en toch niet schadelijk voor de vrijheid in het spel van de geestesvermogens.

Smaak is echter slechts een vermogen om te oordelen, geen productief vermogen; en wat daaraan beantwoordt, is om die reden geen werk van schone kunst, hoewel het een product kan zijn dat behoort tot een nuttige en mechanische kunst of zelfs tot de wetenschap, overeenkomstig bepaalde regels die kunnen worden geleerd en die precies moeten worden gevolgd. Maar de aangename vorm die men eraan geeft, is slechts het voertuig van mededeling en een wijze van presentatie, waarbij men tot op zekere hoogte vrij blijft, zelfs als men anders gebonden is aan een bepaald doel. Zo eist men dat tafelschikkingen, een moreel traktaat of zelfs een preek op zichzelf deze vorm van schone kunst moeten hebben, zonder \textbf{gestudeerd} te lijken; maar zij worden daarom niet werken van schone kunst genoemd. Onder de laatste worden echter een gedicht, een muziekstuk, een schilderijengalerij en dergelijke gerekend; en daar kan men in het ene zogenaamde werk van schone kunst vaak genie zonder smaak waarnemen, en in het andere smaak zonder genie.

\begin{mdframed}[leftmargin=0pt,rightmargin=0pt,backgroundcolor=blue!6,linecolor=white,innertopmargin=6pt,innerbottommargin=6pt]

Smaak is nodig om schoonheid te herkennen, maar genie is nodig om schone kunst voort te brengen. Het cruciale verschil ligt hierin: de schoonheid van de natuur bevalt om haar vorm alleen, zonder dat wij hoeven te weten wat haar doel is. Bij kunst daartegenover moet het voorwerp niet alleen mooi \textit{lijken}, maar ook een \textbf{bedoeling} hebben, een concept dat de kunstenaar vormgeeft. De natuur toont ons mooie \textit{dingen}; kunst toont ons mooie \textit{voorstellingen} van dingen, waarbij de kunstenaar een idee gestalte geeft.

Kunst kan zelfs lelijke of afschrikwekkende onderwerpen, zoals oorlog of ziekte, op een schone manier uitbeelden, mits zij geen afschuw opwekken. Afschuw vernietigt immers elke esthetische bevrediging, omdat het voorwerp dan niet meer als kunst, maar als werkelijkheid wordt ervaren. Daarom vermijdt de beeldhouwkunst bijvoorbeeld directe weergaven van lelijkheid en kiest voor symbolische of allegorische voorstellingen.

Hoewel genie de scheppende kracht levert, is het \textbf{smaak} die de vorm perfektioneert. De kunstenaar werkt zijn ideeën uit door te oefenen, te corrigeren en te verfijnen, tot de vorm zowel aan zijn bedoeling voldoet als de vrijheid van de verbeelding behoudt. Smaak alleen maakt nog geen schone kunst: een mooi opgemaakte tafel of een welsprekende preek zijn nuttig en aantrekkelijk, maar geen ware kunst. Echte schone kunst, zoals gedichten, muziek of schilderijen, vereist zowel genie als smaak. Het ene zonder het andere levert ofwel chaotische inspiratie, ofwel leeg vakmanschap op.

\bigskip
Kant benadrukt in §48 een schijnbare paradox: smaak is essentieel voor schoonheid, maar schoonheid is uiteindelijk geen kwestie van smaak. Deze spanning is geen tegenstrijdigheid, maar de kern van zijn esthetica. 

Smaak is het vermogen om schoonheid te herkennen en te beoordelen. Het is een subjectief, maar universeel oordeelsvermogen: wij ervaren schoonheid niet via regels of begrippen, maar via een direct gevoel van harmonie tussen verbeelding en verstand. Smaak is dus de toegangspoort tot schoonheid: zonder smaak kunnen wij niet eens beginnen te spreken over wat mooi is. Maar hier ligt ook de beperking: smaak is reactief. Het stelt ons in staat om te oordelen, maar niet om te scheppen. Een tafelschikking, een preek of een moreel traktaat kunnen smaakvol zijn, maar zij zijn geen schone kunst, omdat zij geen originele, scheppende daad vereisen. Smaak is als het ware de grammatica van schoonheid: noodzakelijk om te begrijpen, maar niet voldoende om te inspireren.

Kant zegt uiteindelijk dat schoonheid geen kwestie is van smaak, omdat smaak alleen de vorm beoordeelt, niet de bron van schoonheid. Ware schoonheid ontstaat niet door smaak, maar door genie, een scheppende kracht die nieuwe, onvoorziene mogelijkheden openbaart. Smaak kan een kunstwerk verfijnen, maar het kan geen kunstwerk uitvinden. Zo is smaak als de passer en winkelhaak: onmisbaar om orde te scheppen, maar niet de bron van de inspiratie zelf.

Kants radicale stelling dat schoonheid geen kwestie van smaak is, betekent dat schoonheid niet berust op persoonlijke voorkeur of conventie. Schoonheid is geen ``lekker vinden'', maar een ervaring van vrijheid: het gevoel dat iets niet alleen ons bevalt, maar moet bevallen, alsof het een universele waarheid raakt. Dit is waarom hij schoonheid verbindt met het vrije spel van de geestesvermogens. Een echt smaakoordeel claimt geldigheid voor iedereen, niet omdat het objectief is (zoals een wetenschappelijk feit), maar omdat het aansluit bij een gedeelde structuur in de menselijke geest.

Voor de vrijmetselaar resoneren deze ideeën sterk met het maçonnieke begrip van symboliek. Een symbool is niet mooi omdat het ``ons smaakt'', maar omdat het een diepere waarheid openbaart die wij herkennen zonder haar volledig te kunnen uitleggen. Net zoals smaak ons leert om schoonheid te herkennen, leert het ritueel ons om betekenis te ervaren die niet in woorden is vast te leggen. Het is geen kwestie van ``lekker vinden'', maar van ontvankelijkheid voor iets dat ons overstijgt.

Genie is het vermogen om iets te scheppen dat lijkt alsof het door de natuur zelf is voortgebracht; iets dat zowel vrij als noodzakelijk aanvoelt. Dit is precies wat Kant bedoelt wanneer hij zegt dat schoonheid geen kwestie van smaak is: genie onthult een schoonheid die niet afhangt van onze persoonlijke smaak, maar die ons juist confronteert met een ordening die wij niet zelf hebben bedacht. Het genie is als een natuurkracht die door de kunstenaar heen werkt. Zo is ook het maçonnieke ritueel geen product van willekeur, maar een vorm die de broeder ervaart als zowel persoonlijk als universeel.

Hierin ligt een diepe parallel met de vrijmetselarij:
\begin{compactitem}
\item \textbf{Smaak} is als de kennis van de rituele vormen: noodzakelijk om deel te nemen, maar niet voldoende om de ervaring te begrijpen.
\item \textbf{Genie} is als de inspiratie die het ritueel tot leven wekt: iets dat niet kan worden geleerd, maar alleen kan worden beleid. Een meester-vrijmetselaar kent de regels, maar het is zijn innerlijke houding, zijn vermogen om binnen die regels een eigen, authentieke betekenis te vinden, die het ritueel tot kunst maakt.
\end{compactitem}

Hoewel schoonheid niet berust op smaak, is smaak wel noodzakelijk om schoonheid te delen. Kant ziet smaak als een cultureel vermogen dat ons in staat stelt om ons oordeel te verfijnen en te verbreden. Zonder smaak vervalt kunst in willekeur; zonder genie vervalt zij in mechanische nabootsing. Voor de vrijmetselaar is dit een herinnering dat ware wijsheid zowel discipline als vrijheid vereist. Het ritueel is geen starre ceremonie, maar een levende traditie die alleen betekenis heeft wanneer zij zowel de vorm respecteert als de geest bevrijdt.

In de loge wordt dit zichtbaar in de spanning tussen traditie en interpretatie. De symbolen zijn vast, maar hun betekenis is altijd persoonlijk en actueel. Zo is smaak de gemeenschappelijke taal die broeders in staat stelt om met elkaar te communiceren, terwijl genie elk van hen in staat stelt om die taal op zijn eigen manier te spreken.

Kants inzicht dat schoonheid geen kwestie van smaak is, is uiteindelijk een oproep om verder dan voorkeuren en conventies te kijken. Ware schoonheid, of het nu in kunst, natuur of ritueel is,– is een openbaring van vrijheid. Voor de vrijmetselaar is dit een uitnodiging om het ritueel niet als een plicht te ervaren, maar als een vrije daad die toch gebonden is aan een hogere orde. Net zoals het genie de kunst haar regels geeft zonder ze te dicteren, zo geeft de vrijmetselarij de broeder een vorm zonder zijn vrijheid te beperken.

In die zin is smaak de voorbereiding, maar genie de bestemming: het vermogen om binnen de grenzen van de traditie iets te scheppen dat zowel persoonlijk als tijdloos is. Dat is de ware kunst van het leven, en van de vrijmetselarij.

\end{mdframed}

\section{Over de vermogens van de geest die genie constitueren.}

Men zegt van bepaalde producten, waarvan men verwacht dat zij, althans gedeeltelijk, als schone kunst aan het licht zouden moeten treden, dat zij zonder \textbf{geest} zijn, ook al vindt men niets in hen om op grond van smaak te bekritiseren. Een gedicht kan heel aardig en elegant zijn, maar zonder \textbf{geest}. Een verhaal is nauwkeurig en goed geordend, maar zonder geest. Een plechtige redevoering is grondig en tegelijkertijd bloemrijk, maar zonder geest. Veel gesprekken zijn niet zonder vermakelijkheid, maar toch zonder geest; zelfs van een vrouw kan men wel zeggen dat zij mooi, spraakzaam en bekoorlijk is, maar zonder geest. Wat wordt hier dan bedoeld met "geest"?

\textbf{Geest}, in esthetische zin, betekent het belevende beginsel in het \textbf{gemoed}. Datgene echter waardoor dit beginsel de ziel belevend maakt, het materiaal dat het hiervoor gebruikt, is datgene wat doelmatig de geestesvermogens in beweging zet, dat wil zeggen, in een spel dat zichzelf in stand houdt en zelfs de krachten daartoe versterkt.

Nu stel ik dat dit beginsel niets anders is dan het vermogen tot de voorstelling van \textbf{esthetische ideeën}; onder een esthetische idee versta ik echter die voorstelling van de verbeelding die veel denken teweegbrengt, zonder dat er een bepaald gedachte, dat wil zeggen een \textbf{begrip}, aan adequaat kan zijn, waaraan dus geen taal volledig recht kan doen of begrijpelijk kan maken. — Men ziet gemakkelijk in dat dit het tegendeel (pendant) is van een \textbf{idee van de rede}, die omgekeerd een begrip is waaraan geen \textbf{aanschouwing} (voorstelling van de verbeelding) adequaat kan zijn.

De verbeelding (als een productief kennend vermogen) is namelijk zeer krachtig in het scheppen, als het ware, van een andere natuur, uit het materiaal dat de werkelijke natuur haar verschaft. Wij vermaken ons ermée wanneer de ervaring ons te alledaags lijkt; wij transformeren deze, zeker altijd in overeenstemming met analoge wetten, maar ook in overeenstemming met beginselen die hoger in de rede liggen (en die ons even natuurlijk zijn als die volgens welke het verstand de empirische natuur begrijpt); hierbij voelen wij onze vrijheid van de wet der associatie (die geldt voor het empirisch gebruik van dat vermogen), volgens welke het materiaal ons zeker door de natuur kan worden geleverd, maar door ons in iets geheel anders kan worden omgevormd, namelijk in datgene wat boven de natuur uitstijgt.

Men kan dergelijke voorstellingen van de verbeelding \textbf{ideeën} noemen: enerzijds omdat zij ten minste streven naar iets dat buiten de grenzen van de ervaring ligt, en zo een voorstelling van begrippen van de rede (van intellectuele ideeën) trachten te benaderen, wat hun de schijn van objectieve werkelijkheid geeft; anderzijds, en wel voornamelijk, omdat geen begrip er volledig adequaat aan kan zijn, als innerlijke aanschouwingen. De dichter waagt het om zintuiglijke voorstellingen te geven van onzichtbare wezens, het rijk der zaligen, het rijk der hel, de eeuwigheid, de schepping, enzovoort, alsmede van datgene waarvan er voorbeelden in de ervaring zijn, zoals de dood, afgunst, en allerlei ondeugden, evenals liefde, roem, enzovoort, zintuiglijk voor te stellen voorbij de grenzen van de ervaring, met een volmaaktheid die alles overtreft waarvan er een voorbeeld in de natuur is, door middel van een verbeelding die het voorbeeld van de rede volgt in het streven naar een maximum; en het is werkelijk de kunst der poëzie waarin het vermogen tot esthetische ideeën zich in zijn volle omvang kan openbaren. Dit vermogen is echter, op zichzelf beschouwd, eigenlijk slechts een talent (van de verbeelding).

Als wij nu aan een begrip een voorstelling van de verbeelding toevoegen die bij zijn presentatie hoort, maar die op zichzelf zoveel denken prikkelt dat zij nooit in een bepaald begrip kan worden begrepen, en dus het begrip zelf esthetisch op onbegrensde wijze uitbreidt, dan is in dit geval de verbeelding scheppend, en zet zij het vermogen van intellectuele ideeën (de rede) in beweging, dat wil zeggen, zij geeft aanleiding tot een voorstelling die meer aan het denken geeft dan in haar kan worden begrepen en duidelijk gemaakt (hoewel zij zeker bij het begrip van het voorwerp hoort).

Die vormen die niet de voorstelling van een gegeven begrip zelf uitmaken, maar als aanvullende voorstellingen van de verbeelding slechts de bijbetekenissen die daarmee verbonden zijn en hun verwantschap met andere voorstellingen uitdrukken, worden (esthetische) \textbf{attributen} van een voorwerp genoemd, waarvan het begrip, als een idee van de rede, niet adequaat kan worden voorgesteld. Zo is de adelaar van Jupiter, met de bliksem in zijn klauwen, een attribuut van de machtige koning der hemelen, evenals de pauw van de prachtige koningin der hemelen. Zij stellen niet, zoals \textbf{logische attributen}, voor wat in onze begrippen van de verhevenheid en majesteit van de schepping ligt, maar iets anders, wat de verbeelding aanleiding geeft om zich uit te strekken over een menigte verwante voorstellingen, die meer laten denken dan men in een door woorden bepaald begrip kan uitdrukken; en zij leveren een \textbf{esthetische idee}, die in plaats van een logische voorstelling die idee van de rede dient, hoewel in feite alleen om de geest te belevendigen door het openen van het vooruitzicht op een onmetelijk veld van verwante voorstellingen. Schone kunst doet dit echter niet alleen in de schilderkunst of beeldhouwkunst (waar de namen van de attributen gewoonlijk worden gebruikt); ook de poëzie en welsprekendheid ontlenen de geest die hun werken belevend maakt uitsluitend aan de esthetische attributen van de voorwerpen, die naast de logische attributen staan, en de verbeelding een impuls geven om meer te denken, al is het op onontwikkelde wijze, dan in een begrip kan worden bevatten, en dus in een bepaalde taalkundige uitdrukking. — Omwille van de beknoptheid moet ik mij beperken tot slechts enkele voorbeelden.

Wanneer de grote koning zich in een van zijn gedichten als volgt uitdrukt: ``Laten wij het leven verlaten zonder te klagen en zonder iets te betreuren,
de wereld achterlatend vervuld van goede daden. Zo verspreidt de zon, nadat zij haar dagelijkse loop heeft voltooid, nog een zacht licht over de hemel; en de laatste stralen die zij naar de lucht zendt,
zijn haar laatste zuchten voor het welzijn der wereld,'' dan belevendigt hij zijn redelijke idee van een kosmopolitische gezindheid, zelfs aan het einde van het leven, door middel van een attribuut dat de verbeelding (bij de herinnering aan alles wat aangenaam is in een mooie zomerdag, die ten einde loopt, welke een heldere avond oproept) met die voorstelling associeert, en dat een menigte gevoelens en aanvullende voorstellingen oproept waarvoor geen uitdrukking wordt gevonden. Omgekeerd kan zelfs een intellectueel begrip dienen als attribuut van een zintuiglijke voorstelling, en zo deze verlevendigen door middel van de idee van het bovenzintuiglijke; maar alleen voor zover het esthetische, dat subjectief verbonden is aan het bewustzijn daarvan, hiertoe wordt gebruikt. Zo zegt bijvoorbeeld een bepaalde dichter in de beschrijving van een mooie morgen: ``De zon straalde uit, zoals rust uit de deugd straalt.'' Het bewustzijn van deugd, wanneer men zich, al is het maar in gedachten, in de plaats stelt van een deugdzaam persoon, verspreidt in de geest een menigte verheven en kalmerende gevoelens, en een onbegrensd vooruitzicht op een gelukkige toekomst, die geen uitdrukking die adequaat is aan een bepaald begrip volledig kan vatten.\footnote{Misschien is er niets verhevener gezegd, of een gedachte verhevener uitgedrukt, dan in het opschrift boven de tempel van \textbf{Isis} (Moeder \textbf{Natuur}): ``Ik ben alles wat is, wat was, en wat zal zijn, en mijn sluier heeft geen sterfelijke opgelicht.'' \textbf{Segner} maakte gebruik van deze gedachte door middel van een vignet, \textbf{rijk aan zin}, geplaatst aan het begin van zijn natuurleer, om bij de aanvang zijn leerling, die hij gereed was om in deze tempel te leiden, met de heilige vrees te vervullen die de geest tot plechtige aandacht moet stemmen.}

Kortom, de esthetische idee is een voorstelling van de verbeelding, verbonden met een gegeven begrip, die met zoveel deelvoorstellingen in het vrije gebruik van de verbeelding wordt gecombineerd, dat er geen uitdrukking die een bepaald begrip aanduidt voor kan worden gevonden, en die daarom toelaat dat aan een begrip veel onnoembaars wordt toegevoegd, waarvan het gevoel de kennende vermogens belevendigt en geest verbindt met de loutere letter van de taal.

De geestesvermogens die, in hun vereniging (in een bepaalde verhouding), \textbf{genie} constitueren, zijn verbeelding en verstand. Alleen in het gebruik van de verbeelding voor kennis is de verbeelding onderworpen aan de beperking van het verstand en moet zij adequaat zijn aan zijn begrip; in esthetisch opzicht echter is de verbeelding vrij om, naast die overeenstemming met het begrip, ongevraagd uitgebreid en onontwikkeld materiaal voor het verstand te verschaffen, waarnaar het verstand in zijn begrip geen acht heeft geslagen, maar dat het gebruikt, niet zozeer objectief voor kennis, als wel subjectief voor de belevendiging van de kennende vermogens, en zo ook indirect voor kennis; genie bestaat dus eigenlijk in de gelukkige verhouding, die geen wetenschap kan leren en geen ijver kan aanleren, om aan de ene kant ideeën te vinden voor een gegeven begrip en aan de andere kant de \textbf{uitdrukking} voor deze ideeën te treffen, waardoor de subjectieve stemming van de geest die daardoor wordt voortgebracht, als begeleiding van een begrip, aan anderen kan worden medegedeeld. Dit laatste talent is eigenlijk datgene wat men geest noemt, want om datgene wat onnoembaar is in de geestesgesteldheid bij een bepaalde voorstelling uit te drukken en universeel mededeelbaar te maken, of de uitdrukking nu in taal, schilderkunst of beeldende kunst bestaat — dat vereist een vermogen om het snel voorbijgaande spel van de verbeelding te vatten en tot een begrip te verenigen (dat daarom oorspronkelijk is en tegelijkertijd een nieuwe regel onthult, die niet uit enige voorafgaande beginselen of voorbeelden kan worden afgeleid), die zonder de dwang van regels kan worden medegedeeld.

\bigskip
Als wij, na deze analyses, terugblikken op de hierboven gegeven uitleg van wat \textbf{genie} wordt genoemd, dan vinden wij: \textbf{ten eerste}, dat het een talent voor kunst is, niet voor wetenschap, waarin regels die duidelijk worden gekend eerst moeten komen en de procedure daarin bepalen; \textbf{ten tweede}, dat het, als talent voor kunst, een bepaald begrip van het product als doel veronderstelt, en dus verstand, maar ook een voorstelling (zelfs als deze onbepaald is) van het materiaal, dat wil zeggen van de aanschouwing, voor de presentatie van dit begrip, en dus een verhouding van de verbeelding tot het verstand; \textbf{ten derde}, dat het zich niet zozeer toont in de uitvoering van het voorgestelde doel in de presentatie van een bepaald \textbf{begrip}, als wel in de uiteenzetting of uitdrukking van \textbf{esthetische ideeën}, die rijk materiaal voor dat doel bevatten, en dus de verbeelding, in haar vrijheid van alle leiding door regels, toch als doelmatig wordt voorgesteld voor de presentatie van het gegeven begrip; ten slotte, \textbf{ten vierde}, dat de ongevraagde en onopzettelijke subjectieve doelmatigheid in de vrije overeenstemming van de verbeelding met de wetmatigheid van het verstand een verhouding en aanleg van dit vermogen veronderstelt die niet kan worden voortgebracht door enige navolging van regels, hetzij van de wetenschap of van mechanische nabootsing, maar die alleen de natuur van het subject kan voortbrengen.

Volgens deze vooronderstellingen is genie de voorbeeldige oorspronkelijkheid van de natuurlijke aanleg van een subject voor het \textbf{vrije} gebruik van zijn kennende vermogens. Op deze wijze is het product van een genie (wat betreft datgene wat daaraan aan genie moet worden toegeschreven, niet aan mogelijke leer of scholing) een voorbeeld, niet voor nabootsing (want dan zou datgene wat genie is en de geest van het werk uitmaakt verloren gaan), maar voor navolging door een ander genie, dat daardoor wordt gewekt tot het gevoel van zijn eigen oorspronkelijkheid, om de vrijheid van dwang in zijn kunst zo uit te oefenen dat deze daardoor zelf een nieuwe regel verkrijgt, waardoor het talent zich als voorbeeldig toont. Maar aangezien het genie een gunsteling van de natuur is, wiens gelijke men slechts als een zeldzaam verschijnsel moet beschouwen, geeft zijn voorbeeld voor andere goede geesten aanleiding tot een school, dat wil zeggen, een methodische onderwijzing volgens regels, voor zover het mogelijk is geweest deze uit die producten van geest en hun individualiteit af te leiden; en voor deze is schone kunst in die mate nabootsing, waartoe de natuur de regel door een genie heeft gegeven.

Maar deze nabootsing wordt \textbf{na-apen} als de leerling alles \textbf{kopieert}, zelfs tot aan datgene toe wat het genie als een misvorming heeft moeten laten zitten, alleen omdat het niet gemakkelijk kon worden verwijderd zonder de idee te verzwakken. Deze moed is een verdienste alleen van een genie, en een zekere \textbf{stoutmoedigheid} in uitdrukking en in het algemeen enige afwijking van de gewone regel past hem wel, maar is geenszins waardig om nagebootst te worden, en blijft altijd op zichzelf een gebrek dat men moet trachten te verwijderen, maar waarvoor het genie als het ware bevoorrechtigd is, omdat datgene wat onnavolgbaar is in de kracht van zijn geest zou lijden onder angstige voorzichtigheid.
\textbf{Manierisme} is een andere soort na-apen, namelijk die van louter \textbf{individualiteit} (oorspronkelijkheid) in het algemeen, om zich zo ver mogelijk van nabootsers te onderscheiden, zonder daarbij het talent te hebben om tegelijkertijd \textbf{voorbeeldig} te zijn. — Er zijn in het algemeen, zeker, twee manieren (\textit{modus}) om gedachten in een voorstelling samen te voegen, waarvan de ene een \textbf{stijl} (\textit{modus aestheticus}) wordt genoemd en de andere een \textbf{methode} (\textit{modus logicus}), die van elkaar verschillen doordat de eerste geen andere maatstaf heeft dan het \textbf{gevoel} van eenheid in de voorstelling, terwijl de laatste bepaalde \textbf{beginselen} volgt; voor schone kunst geldt dus alleen de eerste. Maar men noemt een product van kunst \textbf{gemanierd} alleen als de voorstelling van zijn idee in dat product \textbf{gericht} is op singulariteit in plaats van adequaat te zijn aan de idee. Het opzichtige (precieuze), het gekunstelde en het geaffecteerde, alleen bedoeld om zich van het gewone volk te onderscheiden (maar zonder enige geest), zijn als het gedrag van iemand van wie men zegt dat hij graag het geluid van zijn eigen stem hoort, of die staat en beweegt alsof hij op een toneel staat, om aangegaapt te worden, wat altijd een prutser verraadt.

\begin{mdframed}[leftmargin=0pt,rightmargin=0pt,backgroundcolor=blue!6,linecolor=white,innertopmargin=6pt,innerbottommargin=6pt]

Wanneer we spreken over ware kunst, dan hebben we het niet over iets dat louter mooi is of technisch vaardig gemaakt. Nee, ware kunst ontstaat uit genie, een aangeboren vermogen dat de kunstenaar in staat stelt om iets te scheppen wat niet alleen origineel is, maar ook een diepere waarheid raakt, een waarheid die woorden of begrippen niet volledig kunnen vatten. Dit genie is geen kwestie van geleerdheid of nabootsing, maar van een scheppende kracht die de kunstenaar dwingt om voorstellingen te maken die ons denken prikkelen en onze verbeelding uitdagen. Het is alsof de kunstenaar een andere wereld schept, een wereld die boven de alledaagse werkelijkheid uitstijgt en ons een glimp geeft van iets groters, iets dat we wel kunnen ervaren, maar niet helemaal kunnen begrijpen.

Wat maakt deze scheppende kracht zo bijzonder? Het is het vermogen om esthetische ideeën te vormen. Dit zijn geen gewone gedachten of beelden, maar voorstellingen die zo rijk en diepgaand zijn, dat ze ons aan het denken zetten zonder dat we ze ooit volledig in woorden kunnen vangen. Denk aan de voorstelling van de eeuwigheid, de schepping, of deugd, ideeën die ons raken, maar die we nooit helemaal kunnen omschrijven. Een dichter kan bijvoorbeeld de zon beschrijven als iets dat "rust uitstraalt als de deugd". Dit beeld roept niet alleen een visuele voorstelling op, maar ook een gevoel van kalmte, hoop en morele diepgang. Het is alsof de dichter ons meeneemt naar een wereld die we wel kunnen voelen, maar niet helemaal kunnen verklaren.

Genie is dus geen louter talent voor techniek, maar een unieke combinatie van twee vermogens: verbeelding en verstand. De verbeelding stelt de kunstenaar in staat om nieuwe, onverwachte beelden te scheppen die boven de werkelijkheid uitstijgen. Het verstand zorgt ervoor dat deze beelden niet willekeurig zijn, maar betekenisvol en communicatief. Het is alsof de kunstenaar een brug slaat tussen wat we zien en wat we voelen, tussen de zintuiglijke wereld en de wereld van ideeën. Deze verhouding is niet te leren of na te bootsen; het is een natuurlijke gave die de kunstenaar in staat stelt om iets te creëren dat anderen wel kunnen ervaren en navolgen, maar niet precies kunnen reproduceren.

Wanneer we een werk van ware kunst aanschouwen, voelen we dat er meer in zit dan we kunnen uitleggen. Een beeld, een gedicht, of een muziekstuk kan ons raken op een manier die ons denken verrijkt en onze geest belevendigt. Dit is geen kwestie van louter techniek of vaardigheid, maar van een diepere, scheppende kracht die de kunstenaar in staat stelt om iets te maken dat zowel persoonlijk als universeel resoneren kan. Ware kunst is dus niet alleen mooi, maar ook betekenisvol, zij spreekt tot ons op een manier die ons uitdaagt om verder te denken, om dieper te voelen, en om ons te verbinden met iets dat groter is dan onszelf.

Genie is daarom geen kwestie van regels of voorschriften, maar van een vrije, scheppende daad die toch doelmatig is. Het is alsof de kunstenaar een nieuwe taal spreekt, een taal die ons in staat stelt om ideeën te ervaren die we anders niet zouden kunnen begrijpen. En juist daarom is genie zo bijzonder: het is niet alleen een talent, maar een roeping, een vermogen om de wereld op een nieuwe manier te zien en anderen uit te nodigen om hetzelfde te doen.

\bigskip

Kants diepgaande beschrijving van het genie in §49 lijkt op het eerste gezicht in te druisen tegen zijn idee van de universaliteit van schoonheid. Schoonheid, stelt hij, is iets dat iedereen kan ervaren en beoordelen via smaak, een vermogen dat niet afhankelijk is van geleerdheid of expertise, maar van een gedeelde, menselijke structuur van verbeelding en verstand. Toch is genie, volgens Kant, een zeldzame gave, slechts aan weinigen geschonken. Hoe rijmt dit met de universaliteit van schoonheid? Dit schijnbare paradox onthult juist de diepgang van Kants esthetica, en heeft zelfs een maatschappelijke en spirituele dimensie die voor de vrijmetselaar herkenbaar is.

Kant maakt een cruciaal onderscheid tussen scheppen en ervaren. Genie is het vermogen om originele, esthetische ideeën voort te brengen, ideeën die zo rijk en diepgaand zijn, dat ze niet in woorden of begrippen zijn te vatten. Dit is een scheppende daad, die niet kan worden geleerd of nagebootst. Het is als een vonk die alleen ontstaat wanneer verbeelding en verstand in een unieke, vrije verhouding samenkomen. Deze gave is zeldzaam, omdat zij niet berust op techniek of kennis, maar op een natuurlijke inspiratie die de kunstenaar zelf niet volledig kan verklaren. Smaak daarentegen is het vermogen om schoonheid te herkennen en te beoordelen. Dit is universeel, omdat het berust op een gedeelde structuur in de menselijke geest: het vrije spel van verbeelding en verstand, dat bij ieder mens op dezelfde manier functioneert. Smaak is niet afhankelijk van genie; het is een vermogen dat iedereen kan ontwikkelen door oefening en reflectie.

Het genie schept schoonheid, maar smaak ervaart en beoordeelt haar. Universaliteit ligt dus niet in het scheppen zelf, maar in het vermogen om schoonheid te herkennen. Dit is waarom Kant kan stellen dat schoonheid voor iedereen toegankelijk is, terwijl het genie dat haar voortbrengt zeldzaam blijft. Schoonheid is als een taal: iedereen kan haar begrijpen (smaak), maar slechts weinigen kunnen nieuwe, diepzinnige gedichten in die taal schrijven (genie).

Kant benadrukt dat genie een gave van de natuur is, iets dat niet kan worden doorgegeven of aangeleerd. Dit verklaart waarom genieën zo zeldzaam zijn: zij zijn als uitschieters in de natuur, zoals een zeldzame bloem die alleen onder specifieke omstandigheden ontluikt. Maar hoewel genie zelf niet kan worden geleerd, kan de smaak die nodig is om geniale werken te waarderen wel worden ontwikkeld. Dit is waarom kunst en cultuur een maatschappelijke functie hebben: zij verfijnen ons oordeelsvermogen en stellen ons in staat om de diepgang van geniale werken te ervaren.

Voor de vrijmetselaar resoneren deze ideeën met het maçonnieke begrip van traditie en individuele zoektocht. Het ritueel en de symbolen van de vrijmetselarij zijn als de ``taal'' van schoonheid: zij zijn toegankelijk voor iedereen die de smaak heeft ontwikkeld om ze te begrijpen (door studie, reflectie en ervaring). Maar de manier waarop een broeder deze symbolen beleeft en interpreteert is uniek en persoonlijk, net als het genie. Niet iedereen zal dezelfde diepgang of inspiratie uit het ritueel halen, maar iedereen kan wel leren om de schoonheid ervan te herkennen.

Kant stelt dat het werk van een genie geen regel is die anderen kunnen volgen, maar een voorbeeld dat anderen kan inspireren om hun eigen genialiteit te ontdekken. Dit is cruciaal: genie is niet exclusief in de zin dat het anderen uitsluit, maar in de zin dat het anderen uitdaagt om hun eigen, unieke vermogens te ontwikkelen. Het is als een vonk die een vuur kan aanwakkeren, maar dat vuur moet bij ieder individu zelf ontbranden.
Dit principe is diep geworteld in de vrijmetselarij. De vrijmetselarij leert geen dogma’s, maar biedt symbolen en rituelen als voorbeelden die elke broeder op zijn eigen manier kan interpreteren. Het is niet de bedoeling dat broeders elkaars ervaringen kopiëren, maar dat zij zich laten inspireren om hun eigen innerlijke licht te vinden. Zo is genie niet iets dat slechts weinigen mogen bezitten, maar iets dat iedereen kan ontwikkelen, mits men bereid is om de vrijheid en moed te tonen die daarvoor nodig zijn.

Kants idee van universaliteit is niet dat iedereen genie moet zijn, maar dat iedereen toegang heeft tot de ervaring van schoonheid. Genie is zeldzaam, maar de schoonheid die het voortbrengt is voor iedereen toegankelijk die zijn smaak heeft verfijnd. Dit is een democratisch idee: ware kunst en ware schoonheid zijn niet voorbehouden aan een elite, maar kunnen door iedereen worden ervaren die zich ervoor openstelt. Voor de vrijmetselaar betekent dit dat de waarheid die in het ritueel en de symbolen ligt opgesloten niet exclusief is voor ``genieën'' of geleerden, maar voor iedereen die bereid is om te luisteren, te reflecteren en zich te laten raken. De loge is als een school voor smaak: zij leert broeders om schoonheid en betekenis te herkennen, niet door hen regels op te leggen, maar door hen te laten ervaren hoe symbolen en rituelen hun eigen innerlijke wereld kunnen verrijken.

Kant toont aan dat genie en universaliteit geen tegenstellingen zijn, maar twee kanten van dezelfde medaille. Genie is zeldzaam omdat het een scheppende vrijheid vereist die niet kan worden geleerd. Maar de schoonheid die het voortbrengt is universeel, omdat zij aansluit bij een gedeelde menselijke structuur: het vermogen om vrij en toch doelmatig te oordelen.

Voor de vrijmetselaar is dit een krachtige metafoor voor de maçonnieke zoektocht. Genie staat voor de individuele inspiratie die elke broeder in zich draagt, het unieke vermogen om betekenis te vinden in symbolen en rituelen. Smaak staat voor de gedeelde taal van de vrijmetselarij: de structuur en traditie die het mogelijk maken om die inspiratie te ervaren en te delen. Ware vrijmetselarij is dus geen kwestie van blind volgen of mechanisch nabootsen, maar van het ontwikkelen van je eigen smaak en het ontdekken van je eigen genie, niet als een exclusief talent, maar als een roeping om binnen de vrijheid van de traditie je eigen, authentieke weg te vinden. Zo wordt schoonheid niet alleen ervaren, maar ook voortgebracht, niet door weinigen, maar door ieder die bereid is om te luisteren, te reflecteren en te scheppen.

\bigskip
Kants ideeën over kunst en genie lijken op het eerste gezicht elitair, omdat hij ware schoonheid koppelt aan het zeldzame vermogen van het genie: een aangeboren, scheppende kracht die slechts weinigen bezitten. Deze opvatting lijkt haaks te staan op het moderne, inclusieve kunstdebat, waarin ruimte is voor zowel het Concertgebouworkest als André Rieu, en waarin kunst niet alleen wordt gewaardeerd om haar diepgang, maar ook om haar toegankelijkheid, emotionele kracht en vermogen om mensen te verbinden. Toch is Kants benadering niet zozeer een afwijzing van populaire kunst, als wel een pleidooi voor authenticiteit en betekenis in alle kunst.

Centraal in Kants esthetica staat het onderscheid tussen genie en smaak. Genie is het vermogen om originele, esthetische ideeën te scheppen ideeën die zo rijk en diepgaand zijn dat ze niet in woorden of begrippen zijn te vatten. Dit is een zeldzame gave, maar smaak het vermogen om schoonheid te herkennen en te waarderen is universeel. Iedereen kan smaak ontwikkelen door reflectie, oefening en blootstelling aan kunst. Hierin ligt de democratische kern van Kants theorie: schoonheid is niet voorbehouden aan een elite, maar toegankelijk voor iedereen die bereid is om zijn geestesvermogens te verfijnen. Kunst is dus niet elitair omdat ze diepgang vereist, maar omdat ze een uitnodiging is om die diepgang te zoeken.

Het moderne debat over kunst toont aan dat kunst vele functies kan hebben: ze kan verheffen, maar ook vermaken; ze kan uitdagen, maar ook troosten; ze kan intellectueel zijn, maar ook puur emotioneel. Kants theorie herinnert ons eraan dat ware kunst of ze nu "hoog" of populair is altijd een vrije, authentieke uiting moet zijn die ons uitdaagt om betekenis te vinden. Of het nu gaat om een symfonie, een volkslied, of een ritueel: kunst heeft de kracht om ons te raken, te inspireren en te verbinden met iets dat groter is dan onszelf. In die zin is Kants visie geen elitair oordeel, maar een oproep om open te staan voor de rijke, veelzijdige ervaring die kunst ons kan bieden of ze nu komt in de vorm van een meesterwerk of een eenvoudig, maar hartverwarmend deuntje.

\end{mdframed}

\section{Over de combinatie van smaak met genie in producten van schone kunst.}

Als de vraag is of in zaken van schone kunst het belangrijker is dat genie dan wel smaak aan het licht treedt, dan is dat hetzelfde als vragen of verbeelding of oordeelskracht meer telt in deze kunst. Nu kunst juist door de eerste van deze twee het recht verdient om \textbf{geïnspireerd} te worden genoemd, maar alleen door de tweede het recht verdient om \textbf{schone} kunst te heten, is het laatste, althans als onmisbare voorwaarde (\textit{conditio sine qua non}), dus het belangrijkste waarnaar men moet kijken bij het beoordelen van kunst als schone kunst. Rijk en origineel zijn in ideeën is niet zo noodzakelijk voor de schoonheid als wel de geschiktheid van de verbeelding in haar vrijheid voor de wetmatigheid van het verstand. Want al de rijkdom van de eerste brengt, in haar wetloze vrijheid, niets anders voort dan onzin; de oordeelskracht is echter het vermogen om haar in overeenstemming te brengen met het verstand.

Smaak, zoals de oordeelskracht in het algemeen, is de discipline (of correctie) van het genie; hij knipt zijn vleugels en maakt het welgemanierd of gepolijst; maar tegelijkertijd geeft hij het genie leiding over waar en hoe ver het zich moet uitstrekken om doelmatig te blijven; en door helderheid en orde aan te brengen in de overvloed aan gedachten, maakt hij de ideeën houdbaar, in staat tot een blijvende en universele goedkeuring, en om een nageslacht te genieten bij anderen en in een steeds vorderende cultuur. Als er dus iets moet worden opgeofferd in het conflict tussen deze twee eigenschappen in één product, dan moet dat eerder aan de kant van het genie zijn; en de oordeelskracht, die in zaken van schone kunst haar uitspraken doet op basis van haar eigen principes, zal eerder schade toestaan aan de vrijheid en rijkdom van de verbeelding dan aan het verstand.

Voor schone kunst zijn dus \textbf{verbeelding}, \textbf{verstand}, \textbf{geest} en \textbf{smaak} vereist.\footnote{De eerste drie vermogens bereiken eerst hun \textbf{eenheid} door de vierde.
\textbf{Hume} geeft in zijn geschiedschrijving de Engelsen te verstaan dat zij, hoewel zij in hun werken aan geen enkel volk ter wereld toegeven wat betreft het bewijs van de eerste drie eigenschappen \textbf{afzonderlijk} beschouwd, toch in datgene wat deze verenigt, achter hun buren, de Fransen, moeten blijven.}

\begin{mdframed}[leftmargin=0pt,rightmargin=0pt,backgroundcolor=blue!6,linecolor=white,innertopmargin=6pt,innerbottommargin=6pt]

Ware schoonheid in kunst ontstaat niet door louter inspiratie of door strikte regels zij vereist een evenwicht tussen genie en smaak. Genie is de kracht die kunst haar vuur, haar originaliteit en haar rijkdom aan ideeën geeft. Maar zonder smaak, die als een strenge maar wijze leraar optreedt, zou al die vrijheid en overvloed slechts chaos voortbrengen. Smaak brengt orde, helderheid en doelmatigheid. Hij knipt de vleugels van het genie niet af, maar leert het hoe het zijn vlucht moet richten zodat de ideeën niet alleen krachtig, maar ook begrijpelijk en tijdloos worden.

Wanneer deze twee met elkaar botsen, moet het genie wijken. Want wat baat een werk dat vol vuur en originaliteit is, als het niemand raakt of overtuigt? Smaak zorgt ervoor dat kunst niet alleen indrukwekkend is voor een enkeling, maar dat zij een plaats verdient in de cultuur, dat zij blijft voortleven en anderen blijft inspireren. Schone kunst vereist dus niet alleen verbeelding en verstand, maar ook die belevende kracht die een werk tot leven wekt en bovenal smaak, die alles samenbindt en verfijnt. Zonder smaak is genie als een wild vuur; met smaak wordt het een licht dat voor iedereen schijnt.

\bigskip
§50 leert ons dat ware schoonheid niet alleen afhangt van geniale inspiratie, maar ook van het vermogen om die inspiratie zo te vormen dat zij anderen kan raken. Smaak is niet de vijand van genie, maar de partner die ervoor zorgt dat geniale kunst ook communiceerbaar, betekenisvol en tijdloos wordt. Dit is geen pleidooi voor conventie of saaiheid, maar voor een kunst die niet alleen origineel is, maar ook diepgang en verbinding biedt.

In een tijd waarin kunst vaak wordt gereduceerd tot louter originaliteit of schokwaarde, herinnert Kant ons eraan dat ware schoonheid ook orde, betekenis en communicatie vereist. Schone kunst is niet alleen een kwestie van "alles kunnen", maar ook van weten hoe je je vrijheid moet gebruiken om iets te scheppen dat niet alleen indrukwekkend is, maar ook tijdloos en universeel. Zo is smaak niet de beperking van genie, maar de voorwaarde die genie pas echt laat schitteren.

\end{mdframed}

\section{Over de indeling van de schone kunsten.}

Schoonheid (of het nu schoonheid van de natuur of van de kunst is) kan in het algemeen worden genoemd de \textbf{uitdrukking} van esthetische ideeën; alleen moet in de schone kunst deze idee worden opgeroepen door een begrip van het voorwerp, terwijl in de schone natuur louter reflectie op een gegeven aanschouwing, zonder een begrip van wat het voorwerp zou moeten zijn, voldoende is om de idee op te wekken en mede te delen, waarvan dat voorwerp als de \textbf{uitdrukking} wordt beschouwd.

Als we de schone kunsten willen indelen, kunnen we, althans als experiment, geen eenvoudiger principe kiezen dan de analogie van kunst met de wijze van uitdrukking die mensen gebruiken in hun spraak om met elkaar te communiceren, dat wil zeggen, niet alleen hun begrippen, maar ook hun gevoelens.\footnote{De lezer zal deze schets voor een mogelijke indeling van de schone kunsten niet beoordelen alsof het een doordachte theorie is. Het is slechts een van de verschillende experimenten die nog kunnen en moeten worden beproefd.} – Dit bestaat uit het \textbf{woord}, het \textbf{gebaar} en de \textbf{toon} (articulatie, gebaren en modulatie). Alleen de combinatie van deze drie soorten uitdrukking vormt de volledige communicatie van de spreker. Want gedachte, aanschouwing en gevoel worden daardoor gelijktijdig en verenigd aan de ander overgedragen.

Er zijn dus slechts drie soorten schone kunsten: de kunst van de \textbf{taal}, de \textbf{beeldende} kunst en de kunst van \textbf{het spel der sensaties} (als externe zintuiglijke indrukken). Men zou deze indeling ook als een tweedeling kunnen opstellen, zodat de schone kunst zou worden verdeeld in die van de uitdrukking van gedachten of van aanschouwingen, en de laatste weer volgens hun vorm of hun stof (van sensatie). Maar dan zou het er te abstract uitzien en niet zo geschikt zijn voor gewone begrippen.

1) De kunsten van de \textbf{taal} zijn \textbf{retorica} en \textbf{poëzie}. \textbf{Retorica} is de kunst om een zaak van het verstand te behandelen als een vrij spel van de verbeelding; \textbf{poëzie} is de kunst om een vrij spel van de verbeelding uit te voeren als een zaak van het verstand.

De \textbf{redenaar} kondigt een zaak aan en voert deze uit alsof het slechts een \textbf{spel} met ideeën is om het publiek te vermaken. De \textbf{dichter} kondigt slechts een vermakelijk \textbf{spel} met ideeën aan, en toch komt er voor het verstand evenveel uit als wanneer hij slechts de bedoeling had gehad om zijn zaak te behartigen. De combinatie en harmonie van de twee kennende vermogens, de gevoeligheid en het verstand, die weliswaar niet zonder elkaar kunnen, maar die toch niet gemakkelijk zonder dwang en wederzijdse schade kunnen worden verenigd, moeten onopzettelijk lijken en vanzelf gebeuren; anders is het geen \textbf{schone} kunst. Daarom moet alles gekunstelde en moeizame erin worden vermeden; want schone kunst moet in dubbel opzicht vrije kunst zijn: zij mag niet een kwestie van beloning zijn, een arbeid waarvan de omvang kan worden beoordeeld, afgedwongen of betaald volgens een bepaalde maatstaf; maar terwijl de geest zeker wel bezig is, moet hij zich bevredigd en gestimuleerd voelen (onafhankelijk van beloning), zonder verder te kijken naar een ander doel.

De redenaar biedt dus zeker iets wat hij niet belooft, namelijk een vermakelijk spel van de verbeelding; maar hij onthoudt ook iets van wat hij wel belooft, namelijk de doelmatige bezigheid van het verstand. De dichter daartegen belooft weinig en kondigt slechts een spel met ideeën aan, maar bereikt iets dat de moeite waard is, namelijk voedsel voor het verstand in het spel, en leven voor zijn begrippen door de verbeelding: daarom levert de eerste eigenlijk minder dan hij belooft, de laatste meer.

2) De \textbf{beeldende} kunsten, of die van de uitdrukking van ideeën in \textbf{zintuiglijke aanschouwing} (niet door voorstellingen van de loutere verbeelding, die door woorden worden opgeroepen), zijn ofwel die van \textbf{zintuiglijke waarheid} of van \textbf{zintuiglijke illusie}. De eerste worden \textbf{plastische} kunsten genoemd, de tweede \textbf{schilderkunst}. Beide maken vormen in de ruimte tot uitdrukkingen van ideeën: de eerste maakt vormen kenbaar door twee zintuigen, gezicht en tast (hoewel in het geval van de laatste weliswaar zonder acht te slaan op schoonheid), de tweede alleen voor het eerste van deze. De esthetische idee (archetype, oerbeeld) is voor beide gegrondvest in de verbeelding; de vorm echter, die haar uitdrukking vormt (ectype, afbeelding), wordt gegeven ofwel in zijn lichamelijke uitgestrektheid (zoals het voorwerp zelf bestaat) ofwel overeenkomstig de manier waarop dit laatste in het oog wordt afgebeeld (overeenkomstig zijn verschijning op een vlak); of anders, wat het eerste ook is, ofwel de betrekking tot een werkelijk doel of slechts de schijn daarvan wordt tot een voorwaarde voor reflectie gemaakt.

De \textbf{plastische} kunsten, als de eerste soort van schone beeldende kunsten, omvatten \textbf{beeldhouwkunst} en \textbf{architectuur}. De \textbf{eerste} is die welke lichamelijke begrippen van dingen voorstelt zoals zij \textbf{in de natuur zouden kunnen bestaan} (hoewel, als schone kunst, met betrekking tot esthetische doelmatigheid); de \textbf{tweede} is de kunst om, met deze bedoeling maar toch tegelijkertijd op esthetisch doelmatige wijze, begrippen van dingen voor te stellen die \textbf{alleen door kunst} mogelijk zijn, en waarvan de vorm niet de natuur, maar een vrijwillig doel als bepalende grond heeft. In de laatste is een bepaald \textbf{gebruik} van het artistieke voorwerp het hoofdzaak, waartoe de esthetische ideeën als voorwaarde worden beperkt. In de eerste is de loutere \textbf{uitdrukking} van esthetische ideeën het hoofddoel. Zo behoren beelden van mensen, goden, dieren, enzovoort, tot de eerste soort; maar tempels, prachtige gebouwen voor openbare bijeenkomsten, evenals woningen, triomfbogen, zuilen, cenotafen en dergelijke, opgericht als gedachtenissen, behoren tot de architectuur. Inderdaad kunnen alle huishoudelijke inrichtingen (het werk van de timmerman en dergelijke dingen voor gebruik) tot de laatste worden gerekend, omdat de geschiktheid van het product voor een bepaald gebruik wezenlijk is in een \textbf{werk van architectuur}, terwijl daartegenover een louter \textbf{schilderij}, dat strikt genomen alleen voor het aanschouwen is gemaakt en op zichzelf moet behagen, als een lichamelijke voorstelling slechts een nabootsing van de natuur is, hoewel met betrekking tot esthetische ideeën: waar dan de \textbf{zintuiglijke waarheid} niet zo ver mag gaan dat het ophoudt eruit te zien als kunst en een product van de keuzevermogen.

De \textbf{kunst van de schilder}, als de tweede soort van beeldende kunst, die \textbf{zintuiglijke illusie} in kunstige combinatie met ideeën voorstelt, zou ik verdelen in die van de schone \textbf{afbeelding} van de natuur en die van de schone \textbf{ordening} van haar \textbf{producten}. De eerste zou de \textbf{schilderkunst in engere zin} zijn, de tweede de \textbf{kunst van de lusttuinen}. Want de eerste geeft slechts de illusie van lichamelijke uitgestrektheid; de laatste geeft deze weliswaar in waarheid, maar geeft slechts de illusie van gebruik en nut voor doelen anders dan louter het spel van de verbeelding bij het aanschouwen van haar vormen.\footnote{Dat de kunst van de lusttuinen als een soort schilderkunst kan worden beschouwd, hoewel zij haar vormen weliswaar lichamelijk presenteert, lijkt vreemd; maar aangezien zij haar vormen inderdaad aan de natuur ontleent (de bomen, struiken, grassen en bloemen uit bossen en velden, althans om te beginnen), en in die mate geen kunst is zoals de plastische kunsten, en ook geen begrip van het voorwerp en zijn doel (zoals in de architectuur) als voorwaarde voor haar ordening heeft, maar louter het vrije spel van de verbeelding bij de beschouwing, komt zij in die mate overeen met louter esthetische schilderkunst, die geen bepaald thema heeft (die lucht, land en water door middel van licht en schaduw op een vermakelijke manier samenbrengt). – Over het algemeen moet de lezer dit slechts beschouwen als een poging om de combinatie van de schone kunsten onder één principe te beoordelen, dat in dit geval dat van de uitdrukking van esthetische ideeën (in overeenstemming met de analogie van een taal) moet zijn, en het niet beschouwen als een definitieve afleiding ervan.} 
De laatste is niets anders dan de versiering van de grond met dezelfde verscheidenheid (grassen, bloemen, struiken en bomen, zelfs water, heuvels en dalen) waarmee de natuur deze aan de aanschouwing presenteert, alleen anders gerangschikt en aangepast aan bepaalde ideeën. De schone ordening van lichamelijke dingen wordt echter ook alleen voor het oog gegeven, zoals de schilderkunst; het gevoel van tast kan immers geen aanschouwelijke voorstelling van een dergelijke vorm leveren. 
Tot de schilderkunst in ruime zin zou ik ook de versiering van kamers door middel van behang, stucwerk en allerlei soorten mooie inrichting rekenen, die slechts dienen om te worden \textbf{bekeken}; evenzo de kunst van het smaakvol kleden (ringen, pillendozen, enzovoort). Want een terras met allerlei bloemen, een kamer met allerlei versieringen (zelfs inclusief de opschik van de dames) vormen, op een prachtig feest, een soort schilderkunst, die, net als de schilderkunst in engere zin (die niet het doel heeft, bijvoorbeeld, om geschiedenis of kennis van de natuur te onderwijzen), er slechts is om te worden bekeken, om de verbeelding in vrij spel met ideeën te vermaken en de kracht van het esthetisch oordeel zonder een bepaald doel bezig te houden. Het werk in al deze versieringen kan mechanisch heel verschillend zijn en heel verschillende kunstenaars vereisen; maar het smaakoordeel over wat mooi is in deze kunst wordt op één manier bepaald: namelijk om alleen de vormen (zonder acht te slaan op een doel) te beoordelen zoals zij aan het oog worden geboden, afzonderlijk of in hun onderlinge verbinding, overeenkomstig het effect dat zij op de verbeelding hebben. – Maar hoe de beeldende kunst (door analogie) als gebaar in een taal kan worden beschouwd, wordt gerechtvaardigd door het feit dat de geest van de kunstenaar door deze vormen een lichamelijke uitdrukking geeft aan wat en hoe hij heeft gedacht, en het voorwerp zelf als het ware in mime laat spreken: een zeer gewoon spel van onze fantasie, die aan levenloze dingen, overeenkomstig hun vorm, een geest toekent die uit hen spreekt.

3) De kunst van het \textbf{schone spel der sensaties} (die van buitenaf worden opgeroepen), die toch universeel mededeelbaar moeten zijn, kan niets anders betreffen dan de verhouding van de verschillende graden van de spanning (tensie) van het zintuig waartoe de sensatie behoort, dat wil zeggen, haar toon; en in deze uitgebreide betekenis van het woord kan zij worden verdeeld in het artistieke spel der sensaties van het gehoor en het gezicht, en dus in \textbf{muziek} en de \textbf{kunst der kleuren}.
– Het is opmerkelijk dat deze twee zintuigen, naast de ontvankelijkheid voor sensaties voor zover dat nodig is om door hun middel tot begrippen van externe voorwerpen te komen, ook in staat zijn tot een speciale sensatie die daarmee verbonden is, waarover niet met zekerheid kan worden gezegd of deze haar grond heeft in het zintuig of in de reflectie; en dat deze gevoeligheid toch soms kan ontbreken, hoewel het zintuig voor zover het betrekking heeft op het gebruik voor de kennis van voorwerpen niet in het geheel defect is, maar integendeel uitzonderlijk scherp. Dat wil zeggen, men kan niet met zekerheid zeggen of een kleur of een toon (geluid) slechts aangename sensaties zijn of op zichzelf al een mooi spel van sensaties, dat als zodanig een bevrediging in de vorm bij esthetisch oordelen inhoudt. Als men bedenkt hoe snel de trillingen van het licht, of in het tweede geval, van de lucht, waarschijnlijk ver boven onze capaciteit gaan om de verhouding van de tijdsindeling direct in de waarneming te beoordelen, dan zou men moeten geloven dat alleen het \textbf{effect} van deze trillingen op de elastische delen van ons lichaam wordt gevoeld, maar dat de \textbf{tijdsindeling} door middel daarvan niet wordt opgemerkt en in het oordeel wordt betrokken, en dat bij kleuren en tonen dus alleen aangenaamheid is verbonden, niet de schoonheid van hun samenstelling. Maar als men, integendeel, \textbf{eerst} bedenkt wat men wiskundig kan zeggen over de verhouding van de trillingen in de muziek en het oordeel daarover, en de contrasten tussen kleuren beoordeelt, zoals passend is, in analogie met de laatste, en als men rekening houdt, \textbf{ten tweede} met die zeker zeldzame voorbeelden van mensen die, met het beste gezichtsvermogen ter wereld, geen kleuren kunnen onderscheiden en, met het scherpste gehoor, geen tonen kunnen onderscheiden, en ook, voor hen die dat wel kunnen, de waarneming van een gewijzigde kwaliteit (niet louter van de graad van de sensatie) in verschillende posities op de schaal van kleuren of tonen, en verder dat het aantal hiervan bepaald is voor \textbf{begrijpelijke} onderscheidingen: dan kan men zich gedwongen voelen om de sensaties van beide niet als louter zintuiglijke indrukken te beschouwen, maar als het effect van een oordeel over de vorm in het spel van vele sensaties. Het verschil tussen de ene of de andere mening in het oordeel over muziek zou echter alleen het volgende wijzigen in de definitie: dat muziek zou worden uitgelegd, zoals wij hebben gedaan, als het schone spel van sensaties (door het gehoor), of als \textbf{aangename} sensaties. Alleen bij de eerste definitie zou muziek volledig als een \textbf{schone} kunst worden voorgesteld; bij de tweede echter zou zij als een \textbf{aangename} kunst (althans gedeeltelijk) worden voorgesteld.

\begin{mdframed}[leftmargin=0pt,rightmargin=0pt,backgroundcolor=blue!6,linecolor=white,innertopmargin=6pt,innerbottommargin=6pt]
    
Wanneer we de schone kunsten willen begrijpen, moeten we allereerst inzien dat schoonheid, of ze nu in de natuur of in kunst voorkomt, altijd een uitdrukking is van esthetische ideeën. Deze ideeën zijn geen gewone gedachten die we in woorden kunnen vatten, maar voorstellingen die ons denken prikkelen en onze verbeelding verrijken, zonder dat we ze ooit helemaal kunnen omschrijven. In de kunst worden deze ideeën bewust vormgegeven door een begrip van het voorwerp, terwijl de natuur ons deze ideeën spontaan, zonder voorafgaand begrip, aanreikt. Als we de schone kunsten willen indelen, kunnen we dat het beste doen door te kijken naar hoe mensen met elkaar communiceren: via woorden, gebaren en toon. Net zoals deze drie elementen samen een volledige boodschap overbrengen, kunnen we de schone kunsten verdelen in drie hoofdgroepen, elk gericht op een andere vorm van uitdrukking.

De kunsten van de taal, zoals retorica en poëzie, gebruiken woorden om ideeën en gevoelens over te brengen. Retorica behandelt een serieuze zaak alsof het een spel van de verbeelding is, om het publiek te vermaken en te overtuigen. De redenaar belooft een zaak van het verstand, maar biedt eigenlijk een vermakelijk spel van ideeën. Poëzie doet het omgekeerde: zij presenteert zich als een vrij spel van ideeën, maar levert uiteindelijk meer op dan alleen vermaking. Een goed gedicht voedt het verstand en geeft leven aan abstracte begrippen door ze zintuiglijk en beeldend te maken. Beide kunsten moeten vrij en ongedwongen zijn, zonder dat ze als geforceerd of berekend overkomen. Ze moeten het verstand en de verbeelding in harmonie brengen, alsof dat vanzelfsprekend gebeurt.

De beeldende kunsten drukken ideeën uit via zintuiglijke aanschouwing, zonder afhankelijk te zijn van woorden. Deze kunsten kunnen worden onderverdeeld in twee soorten: die van de zintuiglijke waarheid (zoals beeldhouwkunst en architectuur) en die van de zintuiglijke illusie (zoals schilderkunst en tuinarchitectuur). Beeldhouwkunst en architectuur geven vorm aan lichamelijke begrippen van dingen zoals ze in de natuur zouden kunnen bestaan, of zoals ze door de mens bedacht zijn voor een bepaald doel. Architectuur richt zich op het praktische gebruik van het kunstwerk, terwijl beeldhouwkunst zich richt op de loutere uitdrukking van esthetische ideeën. Schilderkunst en tuinarchitectuur daartegen creëren een illusie van ruimte en vorm, die de verbeelding prikkelen zonder dat ze direct een praktisch doel dienen. Een lusttuin is als een schilderij in de openlucht: hij nodigt uit tot contemplatie en vrij spel van de verbeelding, zonder dat er een bepaald doel aan verbonden is.

Ten slotte is er de kunst van het spel der sensaties, zoals muziek en de kunst der kleuren. Deze kunsten spelen met sensaties die van buitenaf worden opgeroepen, maar die toch universeel mededeelbaar moeten zijn. Muziek en kleur kunnen niet alleen aangename sensaties oproepen, maar ook een schone samenstelling van sensaties, die ons een dieper gevoel van bevrediging geven. Hoewel we niet precies kunnen zeggen of de schoonheid van muziek of kleur ligt in de sensatie zelf of in onze reflectie daarop, is het duidelijk dat ze een speciale plek innemen in onze ervaring van schoonheid. Ze spreken tot ons op een manier die zowel zintuiglijk als geestelijk is, en ze kunnen ons een gevoel van harmonie en evenwicht geven dat boven de loutere sensatie uitstijgt.

Kortom, de schone kunsten zijn niet alleen een kwestie van techniek of vaardigheid, maar van het vermogen om esthetische ideeën uit te drukken op een manier die zowel vrij als betekenisvol is. Of het nu gaat om woorden, beelden of sensaties, ware kunst spreekt tot ons verstand en onze verbeelding, en nodigt ons uit om de wereld op een nieuwe en verrijkende manier te ervaren.

\bigskip
Kants uitgebreide onderverdeling van de schone kunsten in §51 is geen louter academische oefening, maar een fundamentele poging om te laten zien hoe kunst functioneert als een universele taal van esthetische ideeën, ideeën die ons denken en gevoel verrijken, zonder zich te laten vangen in woorden of regels. Deze indeling is essentieel voor zijn betoog om drie redenen, die ook voor vrijmetselaren diep resoneren.

Kant wil aantonen dat schoonheid niet willekeurig is, maar berust op een structuur van vermogens die samenwerken om esthetische ideeën, diepzinnige, niet-rationele voorstellingen, tot uitdrukking te brengen. Door de kunsten in te delen naar hun uitdrukkingswijze (taal, beeld, sensatie), laat hij zien dat elk type kunst een unieke manier heeft om deze ideeën communiceerbaar te maken. Dit is geen starre categorisatie, maar een kaart van hoe kunst de menselijke geest raakt:
\begin{compactitem}
\item Taal (retorica, poëzie) spreekt tot ons verstand en verbeelding via woorden.
\item Beeld (beeldhouwkunst, schilderkunst, architectuur) maakt abstracte ideeën zintuiglijk en ruimtelijk.
\item Sensatie (muziek, kleur) roept directe, maar universeel deelbare ervaringen op.
\end{compactitem}
Symbolen en rituelen doen precies hetzelfde: zij drukken niet-rationele waarheden uit (zoals morele idealen, spirituele inzichten) via vormen die zowel zintuiglijk als geestelijk werken. Een tempel is niet alleen een gebouw, maar een lichamelijke uitdrukking van een esthetische idee, net als een gedicht of een symfonie. Kants indeling helpt begrijpen waarom maçonnieke kunst (zoals architectuur, muziek, of rituele taal) zo krachtig is: zij spreekt tot alle vermogens van de menselijke geest tegelijk.

Kant benadrukt dat ware kunst vrij moet zijn, niet gebonden aan strikte regels of utilitaire doelen, maar toch gestructureerd genoeg om betekenisvol te zijn. Zijn indeling toont hoe elk type kunst deze balans bereikt:
\begin{compactitem}
\item Poëzie lijkt een vrij spel, maar voedt het verstand.
\item Architectuur dient een praktisch doel, maar drukt ook esthetische ideeën uit.
\item Muziek is een spel van sensaties, maar roept toch een gevoel van orde en harmonie op.
\end{compactitem}
Dit is precies de spanning die ook in de vrijmetselarij centraal staat: rituelen zijn gestructureerd, maar hun ware waarde ligt in de vrije, persoonlijke ervaring die ze mogelijk maken. Kants indeling laat zien dat vrijheid en vorm geen tegenstellingen zijn, maar elkaar nodig hebben, net zoals het ritueel zowel een vaste vorm als een open interpretatie vereist. Zonder structuur vervalt kunst (of ritueel) in chaos; zonder vrijheid wordt het een dode ceremonie.

Kants indeling is ook een antwoord op de vraag: \textit{Hoe kan schoonheid universeel zijn, als genie zeldzaam is}? Zijn punt is dat esthetische ideeën, de kern van schoonheid, niet afhankelijk zijn van het genie van de kunstenaar, maar van het vermogen van de toeschouwer om ze te ervaren. Door kunsten in te delen naar hun uitdrukkingswijze, laat hij zien dat elke kunstvorm een universele toegang biedt tot schoonheid:
\begin{compactitem}
\item Een gedicht spreekt tot ieders verbeelding, ook als niet iedereen dichter kan zijn.
\item Architectuur beïnvloedt ons allemaal, ook als we geen bouwmeester zijn.
\item Muziek raakt mensen direct, ook als ze geen musicus zijn.
\end{compactitem}
Symbolen en rituelen zijn niet bedoeld voor een elite, maar voor iedereen die openstaat voor hun betekenis. Net zoals Kant stelt dat schoonheid universeel is omdat ze mededeelbaar is via verschillende kunstvormen, zo is de vrijmetselarij een praktijk van gedeelde ervaring, waarin iedere broeder, ongeacht achtergrond, toegang heeft tot de diepere waarheden die in de vorm zijn verborgen.

Kants indeling laat zien hoe kunst, en dus ook maçonnieke symboliek, functioneert als een brug tussen het zintuiglijke en het geestelijke. Elke kunstvorm die hij noemt, heeft een parallel in de loge:
\begin{compactitem}
\item Poëzie en retorica: De rituele taal en allegorieën die meer zeggen dan hun letterlijke betekenis.
\item Beeldhouwkunst en architectuur: De tempel als lichamelijke uitdrukking van morele en spirituele ideeën.
\item Muziek en kleur: De sfeer en emotie van het ritueel, die direct werken op het gevoel.
\end{compactitem}
De maçonnieke praktijk is geen louter intellectuele of dogmatische aangelegenheid, maar een kunst, een manier om onzegbare waarheden uit te drukken en te ervaren. Net zoals schone kunst ons leert om met open zintuigen en een vrij verstand naar de wereld te kijken, zo leert de vrijmetselarij ons om symbolen niet als vaste betekenissen te zien, maar als uitnodigingen tot reflectie. In die zin is §51 niet alleen een theorie over kunst, maar een gids voor hoe wij betekenis kunnen vinden in datgene wat ons overstijgt.

\end{mdframed}

\section{Over de combinatie van de schone kunsten in één en hetzelfde product.}

Retorica kan worden gecombineerd met een schilderachtige voorstelling van haar onderwerpen, evenals voorwerpen in een toneelstuk; poëzie kan met muziek worden verenigd in een lied; muziek op zijn beurt kan met een schilderachtige (theatrale) voorstelling worden gecombineerd in een opera; het spel der sensaties in een muziekstuk kan worden verenigd met het spel der vormen in dans, enzovoort. Verder kan de voorstelling van het verhevene, voor zover het tot de schone kunst behoort, met schoonheid worden verenigd in een tragisch gedicht in verzen, een leerzaam gedicht, een oratorium; en in deze combinaties is de schone kunst des te kunstzinniger, hoewel men in sommige van deze gevallen kan twijfelen of zij ook mooier is (aangezien zoveel verschillende soorten bevrediging elkaar kruisen).
Toch bestaat in alle schone kunst het wezenlijke uit de vorm, die doelmatig is voor waarneming en oordeel, waarbij het genot tegelijkertijd cultuur is en de geest tot ideeën stemt, en hem daardoor ontvankelijk maakt voor verschillende soorten genot en vermaak – niet in de stof van de sensatie (de bekoring of de emotie), waar het louter gericht is op genot, dat niets in het idee achterlaat, en de geest verdooft, het voorwerp op den duur weerzinwekkend maakt, en het gemoed, omdat het zich bewust is dat zijn gesteldheid in het oordeel van de rede contraproductief is, ontevreden met zichzelf en humeurig.

Als de schone kunsten niet, hetzij nauw hetzij van verre, met morele ideeën worden gecombineerd, die als enige een op zichzelf staande bevrediging met zich meebrengen, dan is dat hun uiteindelijke lot. Zij dienen dan slechts voor afleiding, die men des te meer nodig heeft naarmate men ze gebruikt om de ontevredenheid van de geest met zichzelf te verdrijven, waardoor men zichzelf steeds nuttelozer en ontevredener met zichzelf maakt. Over het algemeen zijn de schoonheden van de natuur het meest verenigbaar met het eerste doel, als men er vroeg aan gewend is geraakt om ze waar te nemen, te beoordelen en te bewonderen.

\begin{mdframed}[leftmargin=0pt,rightmargin=0pt,backgroundcolor=blue!6,linecolor=white,innertopmargin=6pt,innerbottommargin=6pt]

Wanneer verschillende schone kunsten in één werk worden gecombineerd, zoals retorica met schilderkunst in een toneelstuk, poëzie met muziek in een lied, of muziek met dans in een opera, dan wordt de kunst weliswaar ingewikkelder en kunstzinniger, maar dat betekent niet automatisch dat ze ook mooier wordt. Ware schoonheid ligt niet in de hoeveelheid sensaties of de complexiteit van de vorm, maar in de doelmatige eenheid die onze geest prikkelt en veredelt. Als kunst louter gericht is op genot, op het opwekken van emoties of bekoring zonder diepgang, dan verdooft ze uiteindelijk de geest. Ze wordt saai, voorspelbaar, en laat ons leeg achter, omdat ze niets achterlaat om over na te denken.

De echte kracht van kunst ligt in haar vermogen om ons te verheffen, om ons niet alleen te vermaken, maar ook te laten reflecteren, om onze geest open te stellen voor ideeën die verder gaan dan het zintuiglijke. Als kunst zich niet verbindt met morele of spirituele betekenis, met iets dat ons uitdaagt om na te denken over wat echt waardevol is, dan vervalt ze tot louter afleiding. Ze wordt een vlucht, een manier om onze ontevredenheid met onszelf even te vergeten, zonder ons echt te veranderen. Maar wanneer kunst ons leert om met aandacht en bewondering naar de wereld te kijken, zoals de schoonheid van de natuur dat doet, dan wordt ze een bron van ware cultuur. Ze scherpt onze geest, maakt ons ontvankelijk voor diepere ervaringen, en helpt ons om niet alleen te voelen, maar ook te begrijpen. Dat is het verschil tussen kunst die ons slechts even afleidt en kunst die ons echt raakt.

\bigskip
In §52 komt Kants esthetische idealisme opnieuw naar voren, maar het is niet per se elitair in de zin van een afwijzing van populaire of toegankelijke kunst. Wat wel naar voren komt, is zijn strenge eis aan ware kunst: zij moet meer zijn dan louter vermaak of sensatie. Dit kan als elitair overkomen, omdat het suggereert dat niet alle kunst even waardevol is; alleen die kunst die de geest veredelt en tot reflectie aanzet, verdient volgens Kant de naam ``schone kunst''. Maar dit is geen afwijzing van toegankelijkheid, maar een pleidooi voor diepgang.

Kant waarschuwt dat kunst die zich beperkt tot sensatie (zoals louter emotionele bekoring of oppervlakkig vermaak) op den duur verdooft in plaats van verrijkt. Dit is geen afwijzing van genot, maar een waarschuwing dat kunst die alleen op genot is gericht, uiteindelijk leeg aanvoelt. Het is alsof je je voedt met suiker: het smakt, maar het vult niet. Dit lijkt elitair, omdat het suggereert dat ``lage'' kunst (zoals puur vermaak) minderwaardig is. Maar Kants punt is niet dat kunst moeilijk of intellectueel moet zijn, maar dat zij betekenis moet hebben: iets moet achterlaten om over na te denken.

Voor moderne lezers, en voor vrijmetselaren, die waarde hechten aan toegankelijkheid en verbinding, is dit een uitnodiging om na te denken over wat kunst werkelijk doet. Moet kunst ons alleen afleiden, of moet zij ons ook veranderen? Kants kritiek is niet gericht op populaire kunst, maar op kunst die niets achterlaat, die ons niet uitdaagt om dieper te kijken.

Kant stelt dat kunst haar ware waarde alleen bereikt wanneer zij morele ideeën raakt, ideeën die ons doen reflecteren over wat echt belangrijk is in het leven. Dit lijkt elitair, omdat het suggereert dat kunst een ``hogere'' bestemming moet hebben. Maar zijn punt is niet dat kunst moeilijk of serieus moet zijn, maar dat zij menselijk moet zijn. Kunst die ons alleen maar even afleidt, zonder ons te confronteren met vragen over waarheid, schoonheid of deugd, mist volgens hem haar ware potentieel.

Voor vrijmetselaren resoneren deze ideeën sterk. Het maçonnieke ritueel is niet bedoeld als loutere ceremonie, maar als een uitnodiging tot reflectie, een manier om betekenis te vinden in symbolen die verder gaan dan hun oppervlak. Kant lijkt hier te zeggen dat ware kunst hetzelfde moet doen: zij moet ons niet alleen vermaken, maar ook verheffen. Dit is geen elitisme, maar een oproep tot authenticiteit. Kunst die ons alleen maar even afleidt, is als een maaltijd zonder voedingswaarde: lekker op het moment, maar zonder blijvend effect.

Kants kritiek op oppervlakkige kunst is geen afwijzing van toegankelijkheid, maar van leegte. Hij erkent dat schoonheid in de natuur, die voor iedereen toegankelijk is, de hoogste vorm van kunst benadert, omdat zij ons leert om met aandacht en bewondering naar de wereld te kijken. Zijn punt is niet dat kunst alleen voor een elite is, maar dat ware kunst ons moet leren om dieper te kijken, of we nu naar een symfonie luisteren, een gedicht lezen, of een tuin bewonderen.

Voor vrijmetselaren is dit een herkenbare gedachte. Het maçonnieke ideaal is niet om kunst of symbolen te reserveren voor een kleine groep "ingewijden", maar om iedereen uit te nodigen om met open ogen en hart naar de wereld te kijken. Kants waarschuwing is dus geen elitisme, maar een opleiding tot aandacht: ware schoonheid vraagt om betrokkenheid, niet om passieve consumptie.

Kants §52 is geen elitair manifest, maar een pleidooi voor kunst die ertoe doet. Zijn kritiek is niet gericht op toegankelijke kunst, maar op kunst die niets achterlaat, die ons niet uitdaagt om na te denken, om te groeien, of om ons te verbinden met iets groters dan onszelf. Ware kunst, volgens Kant, is niet iets wat ons even afleidt, maar iets wat ons verandert, wat ons leert om met meer aandacht, meer diepgang, en meer bewondering in het leven te staan.

Voor vrijmetselaren is dit een uitnodiging om ritueel en symbolen niet als loutere versiering te zien, maar als wegen naar inzicht. Ware schoonheid ligt niet in de complexiteit of de exclusiviteit, maar in het vermogen om ons te raken en te veranderen. In die zin is Kants boodschap niet elitair, maar menselijk: kunst is niet bedoeld om ons te vluchten, maar om ons thuis te brengen bij onszelf.

\end{mdframed}

\section{Vergelijking van de esthetische waarde van de schone kunsten onderling.}

De \textbf{kunst van de poëzie} (die haar oorsprong bijna geheel aan het genie te danken heeft en het minst door voorschrift of voorbeeld wordt geleid) eist de hoogste plaats op van alle kunsten. Zij verbreedt de geest door de verbeelding vrij te laten en, binnen de grenzen van een gegeven begrip en te midden van de onbegrensde veelheid van vormen die ermee overeen kunnen komen, die ene vorm te presenteren die haar voorstelling verbindt met een volheid van gedachten waaraan geen taalkundige uitdrukking volledig recht kan doen, en verheft zich zo esthetisch tot het niveau van ideeën. Zij versterkt de geest door hem te laten voelen dat hij in staat is om de natuur, als verschijning, vrij, zelfstandig en onafhankelijk van bepaling door de natuur, te beschouwen en te beoordelen volgens gezichtspunten die de natuur zelf in de ervaring noch voor de zintuigen noch voor het verstand presenteert, en zo de natuur te gebruiken als het ware als het schema van het bovenzintuiglijke. Zij speelt met de illusie die zij naar believen oproept, zonder daarbij bedrieglijk te zijn; want zij verklaart zelf haar bezigheid tot louter spel, dat niettemin doelmatig door het verstand voor zijn eigen zaken kan worden gebruikt. - Retorica, voor zover daarmee de kunst van het overtuigen wordt bedoeld, dat wil zeggen, het misleiden door middel van een mooie illusie (als een \textit{ars oratoria}), en niet louter de vaardigheid in spreken (welsprekendheid en stijl), is een dialectiek, die slechts zoveel leent van de kunst van de poëzie als nodig is om de geesten te winnen voor het voordeel van de spreker voordat zij kunnen oordelen, en om hen hun vrijheid te ontnemen; daarom kan zij noch voor de rechtbank noch voor de preekstoel worden aanbevolen. Want wanneer het gaat om burgerlijke wetten die de rechten van individuele personen betreffen, of om de blijvende onderrichting en bepaling van de geesten tot correcte kennis en gewetensvolle naleving van hun plicht, dan is het beneden de waardigheid van een dergelijke belangrijke zaak om zelfs maar een spoor van overvloed aan geestigheid en verbeelding te tonen, laat staan de kunst van het overtuigen en het misleiden voor het voordeel van een ander. Want zelfs als deze kunst soms kan worden toegepast op doelen die op zichzelf legitiem en lofwaardig zijn, blijft het bezwaarlijk dat de maximen en gezindheden subjectief worden gecorrumpeerd, zelfs als de daad objectief rechtvaardig is: want het is niet genoeg om het juiste te doen, het moet ook uitsluitend om die reden worden gedaan dat het juist is.
Bovendien heeft het louter duidelijke begrip van deze soorten menselijke aangelegenheden, gecombineerd met een levendige presentatie in voorbeelden, en zonder aanstoot te geven aan de regels van welklank in de spraak of van gepastheid in de uitdrukking, voor ideeën van de rede (die samen welsprekendheid uitmaken), op zichzelf al voldoende invloed op de menselijke geesten, zonder dat het ook nodig is om de machinerie van overtuiging in te zetten, die, omdat zij ook kan worden gebruikt om ondeugd en dwaling te verdoezelen of te verbergen, nooit geheel de diepgewortelde verdenking van listige kunstgrepen kan wegnemen. In de poëzie gaat alles eerlijk en oprecht toe. Zij verklaart dat zij slechts een vermakelijk spel met de verbeelding zal voeren, en wel wat de vorm betreft, in overeenstemming met de wetten van het verstand, en eist niet dat het verstand door zintuiglijke voorstelling wordt misleid en in verwarring wordt gebracht.\footnote{Ik moet bekennen dat een mooi gedicht mij altijd zuiver genot heeft gegeven, terwijl het lezen van de beste redevoering van een Romeinse volkstribuun of een hedendaagse parlements- of kerkspreker altijd gemengd was met het onaangename gevoel van afkeer van een bedrieglijke kunst, die weet hoe zij mensen, als machines, tot een oordeel in belangrijke zaken kan bewegen, dat alle gewicht voor hen moet verliezen bij kalme reflectie. Welsprekendheid en spraakvaardigheid (samen: retorica) behoren tot de schone kunst; maar de kunst van de redenaar (\textit{ars oratoria}), als de kunst om de zwakheid van mensen voor eigen doeleinden te gebruiken (hoe goedbedoeld of zelfs werkelijk goed deze ook mogen zijn), is geen enkele \textbf{achting} waard. Bovendien bereikte zij zowel in Athene als in Rome haar hoogtepunt slechts in een tijd waarin de staat op zijn ondergang afstevende en een waarlijk vaderlandslievend denken was uitgedoofd. Hij die, naast een helder inzicht in de feiten, de taal in al haar rijkdom en zuiverheid beheerst, en die, met een vruchtbare verbeelding die in staat is zijn ideeën voor te stellen, een levendige sympathie voelt voor het ware goede, is de \textit{vir bonus dicendi peritus}, de spreker zonder kunstgrepen maar vol kracht, zoals \textbf{Cicero} hem voor ogen had, hoewel hij zelf niet altijd trouw bleef aan dit ideaal.}

Na de poëzie zou ik, \textbf{als het gaat om bekoring en beweging van de geest}, datgene plaatsen wat het dichtst bij haar staat onder de kunsten van de taal en ook zeer natuurlijk met haar kan worden verenigd, namelijk de \textbf{toonkunst}. Want hoewel zij zeker spreekt door louter sensaties zonder begrippen, en dus, in tegenstelling tot poëzie, niets achterlaat voor reflectie, beweegt zij de geest op meer veelzijdige en, hoewel slechts tijdelijke, diepere wijzen; maar zij is weliswaar meer genot dan cultuur (het spel van gedachten dat erdoor wordt opgeroepen is slechts het effect van een als het ware mechanische associatie); en zij heeft, beoordeeld door de rede, minder waarde dan enige andere van de schone kunsten. Daarom eist zij, zoals elk ander genot, frequente verandering, en kan zij frequente herhaling niet verdragen zonder weerzin op te wekken. Haar bekoring, die zo universeel mededeelbaar is, lijkt hierop te berusten: dat elke uitdrukking van taal, in haar context, een toon heeft die past bij haar betekenis; dat deze toon meer of minder een affect van de spreker aanduidt en omgekeerd ook een dergelijk affect bij de hoorder oproept, die vervolgens op zijn beurt de idee die in de taal door middel van zo'n toon wordt uitgedrukt, bij de hoorder oproept; en dat, net zoals modulatie als het ware een taal van sensaties is die universeel begrijpelijk is voor elk menselijk wezen, de toonkunst deze taal voor zichzelf in al haar kracht in praktijk brengt, namelijk als een taal van de affecten, en zo, volgens de wet van associatie, universeel de esthetische ideeën mededeelt die van nature daaraan verbonden zijn; echter, aangezien die esthetische ideeën geen begrippen noch bepaalde gedachten zijn, dient de vorm van de samenstelling van deze sensaties (harmonie en melodie) slechts, in plaats van de vorm van een taal, om door middel van een evenredige ordening daarvan (die, aangezien in het geval van tonen zij berust op de verhouding van het aantal trillingen van de lucht in dezelfde tijd, voor zover de tonen tegelijkertijd of achtereenvolgens worden gecombineerd, wiskundig onder bepaalde regels kan worden gebracht) de esthetische ideeën van een samenhangend geheel van een onuitsprekelijke volheid van gedachten, overeenkomstig een bepaald thema, dat de overheersende affect in het stuk vormt, uit te drukken. Op deze wiskundige vorm, hoewel niet voorgesteld door bepaalde begrippen, berust uitsluitend de bevrediging die de loutere reflectie op zo'n veelheid van sensaties die elkaar begeleiden of opvolgen, verbindt met dit spel ervan als voorwaarde voor haar schoonheid die voor iedereen geldt; en het is slechts volgens deze vorm dat smaak voor zichzelf het recht kan opeisen om vooraf het oordeel van iedereen te bepalen.

De wiskunde heeft echter zeker geen enkel aandeel in de bekoring en de beweging van de geest die muziek voortbrengt; integendeel, zij is slechts de onmisbare voorwaarde (\textit{conditio sine qua non}) van die verhouding van de indrukken, in hun combinatie evenals in hun afwisseling, waardoor het mogelijk wordt om ze samen te vatten en te voorkomen dat ze elkaar vernietigen, zodat ze in plaats daarvan overeenstemmen in een voortdurende beweging en belevendiging van de geest door middel van harmonieuze affecten en hierdoor in een aangenaam zelfgenot.

Als men daarentegen de waarde van de schone kunsten schat naar de cultuur die zij de geest bieden en als maatstaf neemt de vergroting van de vermogens die in de oordeelskracht moeten samenwerken ten behoeve van de kennis, dan neemt muziek in die mate de laagste plaats in onder de schone kunsten (net zoals zij misschien de hoogste plaats inneemt onder die kunsten die naar hun aangenaamheid worden geschat), omdat zij slechts speelt met sensaties. De beeldende kunsten gaan haar daarom in dit opzicht ver uit de hoogte; want terwijl zij de verbeelding in een vrij spel zetten dat niettemin ook geschikt is voor het verstand, brengen zij tevens een zaak tot stand door een product voort te brengen dat dient als een duurzaam en zelfaanbevelend voertuig voor de vereniging van het verstand met de zintuiglijkheid en zo als het ware voor de bevordering van de beschaafdheid van de hogere kennende vermogens. De twee soorten kunsten volgen geheel verschillende paden: de eerste van sensaties naar onbepaalde ideeën, de laatste van bepaalde ideeën naar sensaties. De laatste hebben een \textbf{blijvende} indruk, de eerste slechts een \textbf{voorbijgaande}. De verbeelding kan de eerste terugroepen en zich aangenaam ermée vermaken; maar de laatste worden of geheel uitgedoofd of, als zij onwillekeurig door de verbeelding worden teruggeroepen, zijn zij eerder lastig dan aangenaam voor ons.
Bovendien is er een zekere gebrek aan beschaafdheid in muziek, doordat zij, voornamelijk vanwege het karakter van haar instrumenten, haar invloed verder uitstrekt (tot in de omgeving) dan gewenst is, en zo zich als het ware opdringt, waardoor zij de vrijheid van anderen, buiten de muzikale kring, aantast, wat de kunsten die tot het oog spreken niet doen, aangezien men slechts de ogen hoeft af te wenden als men hun indruk niet wil toelaten. Het is bijna hetzelfde als bij het genot van een wijdverspreide geur. Iemand die zijn geparfumeerde zakdoek uit zijn zak haalt, dwingt iedereen in de omgeving er ongewild aan mee te doen, en dwingt hen, als zij willen ademen, er tegelijkertijd van te genieten; daarom is het ook uit de mode geraakt.\footnote{Degenen die het zingen van geestelijke liederen als onderdeel van de huiskamer-godsdienst hebben aanbevolen, hebben er niet bij stilgestaan dat zij met deze luidruchtige (en juist daardoor meestal schijnheilige) vorm van eredienst een groot ongemak aan het publiek opleggen. Zij dwingen de omgeving immers óf mee te zingen, óf hun eigen gedachtengang op te geven.} - Onder de beeldende kunsten zou ik de palm toekennen aan de \textbf{schilderkunst}, gedeeltelijk omdat zij, als de kunst van het tekenen, de basis vormt van alle andere beeldende kunsten, gedeeltelijk omdat zij veel verder kan doordringen in het rijk der ideeën en ook het veld van de aanschouwing in overeenstemming met deze veel verder kan uitbreiden dan voor de rest mogelijk is.

\begin{mdframed}[leftmargin=0pt,rightmargin=0pt,backgroundcolor=blue!6,linecolor=white,innertopmargin=6pt,innerbottommargin=6pt]
    
De kunst van de poëzie verdient de hoogste plaats onder alle schone kunsten, omdat zij het meest afhankelijk is van genie en het minst door regels of voorbeelden wordt beperkt. Zij verbreedt onze geest door de verbeelding vrij te laten en binnen de grenzen van een gegeven begrip die ene vorm te vinden die een volheid van gedachten uitdrukt, gedachten die geen taal ooit volledig kan vatten. Poëzie verheft ons boven de natuur, doordat zij ons leert om de wereld niet alleen te zien, maar ook te bedenken op manieren die de natuur zelf niet aanreikt. Zij speelt met illusies, maar bedriegt niet; zij verklaart zelf dat haar werk louter spel is, en toch dient dit spel als een brug naar diepere inzichten.

Retorica daartegen, als zij wordt misbruikt om te overtuigen door middel van mooie, maar misleidende woorden, is geen ware kunst. Zij steelt de vrijheid van het oordeel door de luisteraar te manipuleren voordat hij zelf kan nadenken. Dit is onwaardig voor belangrijke zaken zoals rechtspraak of prediking, waar het gaat om waarheid en plichtsbesef. Een echte spreker hoeft geen kunstgrepen; wie helder denkt, zuiver spreekt en oprecht voelt, overtuigt zonder bedrog.

Na de poëzie komt de muziek, die ons het meest direct raakt door haar sensaties, zonder dat zij ons concrete begrippen biedt om over na te denken. Muziek beweegt de geest diep en veelzijdig, maar haar effect is vluchtig en meer gericht op genot dan op cultuur. Haar schoonheid berust op wiskundige verhoudingen in klanken, die een gevoel van harmonie en evenwicht oproepen, een taal van affecten die iedereen kan begrijpen, maar die geen duidelijke gedachten achterlaat. Daarom heeft muziek, ondanks haar kracht, minder waarde voor de geest dan andere kunsten: zij prikkelt, maar verrijkt niet.

De beeldende kunsten, zoals schilderkunst en architectuur, gaan verder dan muziek, omdat zij niet alleen de verbeelding in een vrij spel zetten, maar ook een tastbaar resultaat voortbrengen dat het verstand voedt. Een schilderij of gebouw blijft bestaan als een blijvende uitdrukking van ideeën, terwijl muziek verdwijnt zodra zij ophoudt te klinken. Schilderkunst verdient onder de beeldende kunsten de hoogste plaats, omdat zij de basis vormt voor alle andere en het vermogen heeft om diep in het rijk der ideeën door te dringen, ver voorbij wat woorden of tonen alleen kunnen bereiken.

Kortom: poëzie staat bovenaan omdat zij de geest het meest verrijkt, muziek raakt ons het meest direct maar laat weinig blijvends achter, en de beeldende kunsten bieden een duurzame verbinding tussen zintuig en verstand. Ware kunst is niet alleen mooi, maar verheft ons, zij leert ons om met vrijheid en diepgang naar de wereld te kijken.

\bigskip
Kants ordening van de kunsten in §53 is geen willekeurige rangschikking, maar een diepgaande poging om te laten zien hoe kunst functioneert als een weg naar menselijke verrijking. Deze hiërarchie is essentieel voor zijn betoog over schoonheid om drie fundamentele redenen.

\paragraph{Schoonheid als weg naar vrijheid en cultuur} Kant wil aantonen dat schoonheid niet louter een kwestie is van aangename sensaties, maar van geestelijke groei. Door de kunsten te ordenen naar hun vermogen om de geest te verheffen, laat hij zien dat ware kunst ons niet alleen vermaakt, maar ons ook leert om vrij en zelfstandig te denken. Poëzie staat bovenaan omdat zij het meest geniaal is: zij daagt ons uit om verder te kijken dan de oppervlakte, om de wereld te zien door een lens die de natuur zelf niet biedt. Dit is geen elitisme, maar een pleidooi voor diepgang. Kants ordening herinnert ons eraan dat kunst niet alleen mooi moet zijn, maar ook betekenisvol, dat zij ons moet helpen om ons verstand en onze verbeelding te verfijnen.

\paragraph{De balans tussen verbeelding en verstand} Kants ordening toont hoe elke kunstvorm een unieke balans vindt tussen verbeelding (vrijheid, creativiteit) en verstand (structuur, betekenis). Poëzie combineert deze het best: zij is vrij, maar niet willekeurig; zij is diep, maar niet abstract. Muziek daartegen raakt ons diep, maar zonder duidelijke gedachten achter te laten, zij is meer genot dan cultuur. De beeldende kunsten bieden een tastbare, blijvende uitdrukking van ideeën, maar missen soms de directe emotionele kracht van muziek of de diepgang van poëzie.
Deze ordening is geen waardeoordeel, maar een kaart van hoe kunst werkt. Zij laat zien dat ware schoonheid niet alleen afhangt van wat kunst doet (vermaken, raken, verheffen), maar ook van hoe zij dat doet: door een evenwicht te vinden tussen vrijheid en structuur, tussen gevoel en gedachte. 

\paragraph{Kunst als middel tot morele en intellectuele verrijking} Kants hiërarchie is uiteindelijk een pleidooi voor kunst als cultuur, niet als afleiding. Hij waarschuwt dat kunst die louter gericht is op genot (zoals muziek zonder diepgang of retorica zonder oprechtheid) ons verdooft in plaats van verrijkt. Ware kunst daagt ons uit om na te denken, om ons verstand en onze zintuigen in harmonie te brengen. Poëzie doet dit het best, omdat zij ons leert om de wereld te zien door de ogen van de verbeelding en het verstand tegelijk.
Voor vrijmetselaren is dit een cruciale les. Het maçonnieke ideaal is niet om symbolen of rituelen te gebruiken als loutere traditie, maar als middelen tot inzicht. Kants ordening van de kunsten herinnert ons eraan dat ware schoonheid, of het nu in kunst, natuur of ritueel is, ons altijd uitnodigt om dieper te kijken, om vrijer te denken, en om ons te verbinden met iets dat groter is dan onszelf. Zijn hiërarchie is dus geen starre ranglijst, maar een gids voor hoe wij kunst kunnen gebruiken om ons leven te verrijken.

Kant brengt deze ordening aan omdat hij wil laten zien dat schoonheid niet willekeurig is, maar berust op een structuur van vermogens die samenwerken om ons te verheffen. Zijn hiërarchie is geen afwijzing van bepaalde kunstvormen, maar een uitnodiging om kunst te zien als een weg naar wijsheid. 

\end{mdframed}

\section*{Opmerking}

Tussen wat \textbf{louter in het oordelen bevalt} en wat \textbf{gratie geeft} (bevalt in de sensatie) bestaat, zoals we vaak hebben aangetoond, een wezenlijk verschil. Dat laatste is iets wat men, in tegenstelling tot het eerste, niet van iedereen kan verlangen. Gratie (zelfs als de oorzaak ervan in ideeën kan liggen) lijkt altijd te bestaan in een gevoel van bevordering van het totale leven van de mens, en dus ook van lichamelijk welzijn, dat wil zeggen: van gezondheid. Epicurus, die alle gratie uiteindelijk als lichamelijke sensatie beschouwde, had in die zin misschien niet ongelijk, en heeft zich alleen vergist toen hij ook intellectuele en zelfs praktische bevrediging als gratie beschouwde. Als men dit onderscheid voor ogen houdt, kan men verklaren hoe een gratie zelfs degene die haar ervaart kan misnoegen (zoals de vreugde van een behoeftige maar rechtvaardige persoon over de erfenis van zijn liefdevolle, maar gierige vader), of hoe een diepe pijn toch degene die haar lijdt kan behagen (zoals de droefheid van een weduwe bij de dood van haar lofwaardige echtgenoot). Ook kan een gratie bovendien behagen (zoals in de wetenschappen die we beoefenen) of een pijn (bijvoorbeeld haat, afgunst of wraakzucht) ons bovendien misnoegen. De bevrediging of ontevredenheid berust hier op de rede en is hetzelfde als \textbf{goedkeuring} of \textbf{afkeuring}; gratie en pijn kunnen echter alleen berusten op het gevoel of de vooruitzicht (wat de grond ervan ook moge zijn) van een mogelijke toestand van \textbf{welzijn} of \textbf{onwelzijn}.

Elk veranderlijk, vrij spel van sensaties (dat niet op enige bedoeling berust) gratieert, omdat het het gevoel van gezondheid bevordert, ongeacht of we genoegen scheppen in het rationele oordeel over het voorwerp of zelfs in deze gratie; en deze gratie kan uitgroeien tot een affect, ook al hebben we geen belangstelling voor het voorwerp zelf, of in elk geval niet in de mate die evenredig is aan de graad van dat affect. We kunnen dit verdelen in het \textbf{spel van toeval}, het \textbf{spel van toon} en het \textbf{spel van gedachten}. Het \textbf{eerste} vereist een \textbf{belangstelling}, of het nu ijdelheid of eigenbelang is, die echter ver afstaat van de belangstelling voor de manier waarop we proberen deze te bevredigen; het \textbf{tweede} vereist slechts de verandering van \textbf{sensaties}, die elk hun relatie tot affect hebben, maar niet de graad van een affect, en roept esthetische ideeën op; het \textbf{derde} ontstaat louter uit de verandering in de voorstellingen, in het oordeelsvermogen, waardoor weliswaar geen enkele gedachte die enig belang met zich meebrengt wordt gegenereerd, maar de geest toch wordt geanimeerd.

Hoe gratievol deze spelen moeten zijn, zonder dat er behoefte is om ze te gronden in een geïnteresseerde bedoeling, blijkt uit al onze avondse sociale bijeenkomsten, want zonder spelen vindt bijna niemand deze vermakelijk. Maar hier zijn de affecten van hoop, angst, vreugde, woede en minachting in het spel, die elk moment van rol wisselen en zo levendig zijn dat daardoor het hele bedrijf van het lichamelijke leven, als een innerlijke beweging, lijkt te worden bevorderd, wat blijkt uit de opgewektheid van de geest die daardoor ontstaat, ook al is er niets gewonnen of geleerd. Omdat kansspelen echter geen mooi spel zijn, zullen we deze hier buiten beschouwing laten. Muziek en stof tot lachen zijn daarentegen twee soorten spel met esthetische ideeën of zelfs voorstellingen van het verstand, waarbij uiteindelijk niets wordt gedacht, en die slechts door hun verandering kunnen gratieren, en dat toch op een levendige manier doen; waardoor ze duidelijk maken dat de animatie in beide gevallen louter lichamelijk is, hoewel ze wordt opgeroepen door ideeën van de geest, en dat het gevoel van gezondheid dat voortkomt uit een beweging van de ingewanden die overeenkomt met dat spel, de hele gratie uitmaakt in een levendig gezelschap, dat als zo verfijnd en geestig wordt geprezen.
Het is niet het oordeel over de harmonieën in tonen of geestige opmerkingen, die met hun schoonheid slechts als noodzakelijk voertuig dienen, maar de bevordering van het levensbedrijf in het lichaam, het affect dat de ingewanden en het middenrif in beweging brengt, met andere woorden: het gevoel van gezondheid (dat anders niet gevoeld kan worden zonder een dergelijke prikkel), dat de gratie uitmaakt waarin men ontdekt dat men zelfs via de ziel bij het lichaam kan komen en deze kan gebruiken als de arts voor de eerste.

In muziek gaat dit spel van de lichamelijke sensatie naar esthetische ideeën (van de voorwerpen voor affecten), en vervolgens weer terug, maar met verenigde kracht, naar het lichaam. In de grap (die, net als muziek, eerder als aangename dan als schone kunst beschouwd moet worden) begint het spel met gedachten die, als geheel, voor zover ze zintuiglijk moeten worden uitgedrukt, ook het lichaam in beslag nemen; en omdat het verstand in deze voorstelling niet vindt wat het verwachtte, ontspant het zich plotseling, en men voelt het effect van deze ontspanning in het lichaam door de trilling van de organen, die de herstelling van hun evenwicht bevordert en een gunstige invloed heeft op de gezondheid.

In alles wat een levendige, bulderende lach moet opwekken, moet iets onzinnigs zitten (waarbij het verstand op zichzelf geen bevrediging kan vinden). \textbf{Lachen is een affect dat voortkomt uit de plotselinge transformatie van een verhoogde verwachting in niets}. Deze transformatie, die zeker niets aangenaams is voor het verstand, is toch indirect aangenaam en, voor een moment, zeer levendig. De oorzaak moet dus bestaan in de invloed van de voorstelling op het lichaam en de wederkerige werking daarvan op de geest; zeker niet voor zover de voorstelling objectief een voorwerp van gratie is (want hoe kan een teleurgestelde verwachting gratievol zijn?), maar louter doordat zij, als louter spel van voorstellingen, een evenwicht in de levenskrachten van het lichaam teweegbrengt.

Als iemand dit verhaal vertelt: Een Indiër, aan de tafel van een Engelsman in Surat, ziet een fles bier worden geopend en al het bier, veranderend in schuim, eruit stromen, en toont zijn grote verbazing met vele uitroepen. Op de vraag van de Engelsman: ``Wat is hier zo verbazingwekkend aan?'' antwoordt hij: ``Ik ben niet verbaasd dat het eruit komt, maar hoe jullie het allemaal erin hebben gekregen,'' dan lachen we, en het geeft ons een hartelijk genoegen: niet omdat we ons slimmer voelen dan deze onwetende persoon, of vanwege iets anders aangenaams dat het verstand ons hierin laat opmerken, maar omdat onze verwachting werd verhoogd en plotseling in niets verdween. Of als de erfgenaam van een rijke verwant een waardige begrafenis voor hem wil regelen, maar klaagt dat hij het niet helemaal voor elkaar krijgt, omdat (zo zegt hij) ``hoe meer geld ik mijn rouwenden geef om er treurig uit te zien, hoe vrolijker ze eruitzien,'' dan lachen we hardop, en de reden is dat een verwachting plotseling in niets verandert.
Merk op dat het niet in het positieve tegendeel van een verwacht voorwerp mag veranderen – want dat is altijd iets, en kan vaak verdrietig zijn – maar in niets. Want als iemand ons een verhaal vertelt en een grote verwachting oproept, en we bij de afloop onmiddellijk de onwaarheid ervan inzien, dan is dat onaangenaam, zoals bijvoorbeeld het verhaal van mensen wier haar in één nacht grijs zou zijn geworden door groot verdriet. Als daarentegen, als reactie op zo'n verhaal, een andere grappenmaker een zeer uitgewerkt verhaal vertelt over het verdriet van een koopman die, terugkerend van India naar Europa met al zijn fortuin in koopwaar, gedwongen werd alles overboord te gooien in een vreselijke storm, en zo van streek was dat zijn \textbf{pruik} in dezelfde nacht grijs werd, dan lachen we en het geeft ons gratie, omdat we een tijdje heen en weer worden geslingerd als een bal tussen onze eigen misvatting over een voorwerp dat ons verder onverschillig laat, of liever ons eigen idee dat we achterna hebben gejaagd, terwijl we alleen maar probeerden het vast te grijpen en stevig vast te houden. Het is niet het wegsturen van een leugenaar of een domoor dat hier de gratie opwekt, want zelfs op zichzelf zou het laatste verhaal, met een schijnbare ernst verteld, een gezelschap tot schaterlachen kunnen bewegen, en het eerste zou niet eens gewoonlijk de moeite van de aandacht waard zijn.

Het is opmerkelijk dat in al dergelijke gevallen de grap altijd iets moet bevatten dat voor een moment kan misleiden: vandaar dat, wanneer de illusie plotseling in niets verdwijnt, de geest terugkijkt om het nog eens te proberen, en zo heen en weer wordt geslingerd door snel opvolgende toename en afname van spanning, en in trilling wordt gebracht; wat, omdat datgene wat als het ware de snaar heeft doen trillen plotseling terugveert (niet door een geleidelijke verslapping), een beweging van de geest en een daaraan beantwoordende innerlijke lichamelijke beweging moet teweegbrengen, die onwillekeurig voortduurt en vermoeidheid, maar tegelijkertijd ook opgewektheid veroorzaakt (de effecten van een beweging die gunstig is voor de gezondheid).

Want als we aannemen dat al onze gedachten tegelijkertijd harmonieus verbonden zijn met een bepaalde beweging in de organen van het lichaam, dan zal men een redelijk inzicht hebben in hoe aan die plotselinge verschuiving van de geest, eerst naar het ene en dan naar het andere standpunt voor het beschouwen van zijn voorwerp, een wederkerige spanning en ontspanning van de elastische delen van onze ingewanden kan corresponderen, die zich voortplant naar ons middenrif (zoals datgene wat kietelige mensen voelen), zodat de longen de lucht met snel opvolgende pauzes uitstoten, en zo een beweging teweegbrengen die bevorderlijk is voor de gezondheid, die alleen, en niet wat er in de geest gebeurt, de echte oorzaak is van een gratie in een gedachte die in wezen niets voorstelt. –
Voltaire zei dat de hemel ons twee dingen heeft gegeven als tegenwicht tegen de vele lasten van het leven: \textbf{hoop} en \textbf{slaap}. Hij had er ook het \textbf{lachen} bij kunnen voegen, als de middelen om het bij redelijke mensen op te wekken maar zo gemakkelijk voorhanden waren, en de geestigheid of de oorspronkelijke verbeeldingskracht die daarvoor nodig is niet zo zeldzaam was als het talent frequent is om werken te maken die het \textbf{hoofd breken}, zoals die van mystieke dromers, of de \textbf{nek}, zoals die van een genie, of het \textbf{hart}, zoals die van sentimentale romanciers (of wat dat betreft moralisten van dezelfde soort).

Men kan dus, dunkt mij, aan Epicurus toegeven dat alle gratie, zelfs als zij door begrippen wordt veroorzaakt die esthetische ideeën oproepen, \textbf{dierlijk} is, dat wil zeggen: lichamelijke sensatie, zonder daarmee het minste krenken van het \textbf{geestelijke} gevoel van respect voor morele ideeën, dat geen gratie is, maar zelfrespect (van de menselijkheid in ons) dat ons boven de behoefte aan gratie verheft, zonder zelfs schade toe te brengen aan het minder edele gevoel van \textbf{smaak}.

Iets met een vleugje van beide vinden we in de \textbf{naïviteit}, die de weerstand is van de oorspronkelijke, natuurlijke oprechtheid van de mens tegen de kunst van de schijn die een tweede natuur is geworden. Men lacht om de eenvoud die nog niet weet hoe zij moet veinzen, en toch verheugt men zich ook over de eenvoud van de natuur die hier die kunst te slim af is. Men verwacht de gebruikelijke kunstmatige uitdrukkingswijze die zorgvuldig is gericht op een mooie illusie, en zie! het is onbedorven, onschuldige natuur, waar men helemaal niet op voorbereid was en die degene die haar even liet doorschemeren niet eens bedoelde bloot te stellen. Dat de mooie, maar valse illusie, die gewoonlijk zoveel betekent in ons oordeel, hier plotseling in niets verandert, dat als het ware de grappenmaker in onszelf wordt ontmaskerd, brengt de opeenvolgende beweging van de geest in twee tegengestelde richtingen teweeg, die tegelijkertijd het lichaam een gezonde schudbeurt geeft. Maar dat er in de menselijke natuur iets oneindig beters is dan elke aangenomen gewoonte, namelijk de zuiverheid van gedachte (of ten minste de aanleg daartoe), niet geheel is uitgedoofd, voegt ernst en hoge achting toe aan dit spel van het oordeelsvermogen. Omdat het echter een verschijning is die zich slechts voor korte tijd manifesteert, en het gordijn van de kunst van de schijn al snel weer wordt gesloten, bevat het ook een element van spijt, dat een gevoel van tederheid is, dat heel goed als spel kan worden gecombineerd met zo'n goedhartige lach, en dat inderdaad meestal ermee wordt gecombineerd, en tegelijkertijd meestal degene die de stof daartoe levert compenseert voor de gênante situatie om niet scherp te zijn in de omgang met mensen. - Een kunst om \textbf{naïef} te zijn is dus een tegenspraak; maar het is zeker mogelijk om naïviteit in een fictieve persoon voor te stellen, en dit is een mooie, hoewel ook zeldzame kunst. Naïviteit mag niet worden verward met openhartige eenvoud, die de natuur niet kunstmatig verbergt, simpelweg omdat zij niet begrijpt wat de kunst van het sociale leven is.

Naast wat vrolijk is, nauw verwant aan de gratie van het lachen, en deel uitmaakt van de oorspronkelijkheid van de geest, maar daarom niet van het talent voor schone kunst, kan ook de \textbf{grillige} wijze worden gerekend. \textbf{Grilligheid} in de goede zin betekent het talent om zich naar believen in een bepaalde geestesgesteldheid te kunnen verplaatsen, waarin alles heel anders wordt beoordeeld dan gebruikelijk (zelfs geheel omgekeerd), en toch in overeenstemming met bepaalde principes van de rede in zo'n geestesgesteldheid. Iemand die onwillekeurig aan dergelijke veranderingen onderhevig is, is \textbf{onderhevig aan grilligheid}, maar iemand die ze vrijwillig en doelbewust kan aannemen (voor een levendige presentatie door middel van een lachwekkend contrast), zo iemand en zijn optreden worden \textbf{grillig} genoemd. Deze wijze behoort echter meer tot de aangename dan tot de schone kunst, omdat het voorwerp van de laatste altijd enige waardigheid in zichzelf moet tonen, en daarom een zekere ernst in de presentatie vereist, net als smaak in zijn oordeel.

\begin{mdframed}[leftmargin=0pt,rightmargin=0pt,backgroundcolor=blue!6,linecolor=white,innertopmargin=6pt,innerbottommargin=6pt]

Er is een diep verschil tussen wat ons bevalt omdat het onze geest aanspreekt, zoals schoonheid, die ons doet nadenken en ons verrijkt, en wat ons simpelweg gratie geeft, omdat het ons zintuigen streelt of ons een gevoel van welbehagen schenkt. Dat laatste is geen kwestie van smaak of inzicht, maar van louter genot: een gevoel van levendigheid, van gezondheid, alsof ons hele wezen even wordt opgefrist. Het is alsof ons lichaam zelf ons vertelt dat het goed gaat, los van wat ons verstand ervan vindt.

Soms genieten we van dingen die ons eigenlijk niet bevallen, zoals een erfenis van een gierigaard, die ons weliswaar rijk maakt, maar waar we ons schuldig over voelen. Of we vinden troost in pijn, zoals de droefheid van een weduwe die haar man mist, maar die toch een zekere voldoening vindt in haar verdriet, omdat het getuigt van haar liefde. Ons genot of onze pijn hangt niet alleen af van wat we voelen, maar ook van hoe we het beoordelen: ons verstand kan ons laten genieten van iets wat ons lichaam afwijst, of ons laten lijden onder iets wat ons eigenlijk deugd doet.

Er zijn drie soorten spel die ons gratie geven zonder dat we er bewust naar hoeven te streven: het spel van toeval (zoals gokken), het spel van klanken (muziek) en het spel van gedachten (humor). Deze spelen vermaken ons niet omdat ze ons iets leren, maar omdat ze ons lichaam en onze geest in beweging zetten. Bij een gezellig samenzijn, een feest, een avond met vrienden, zijn het juist deze spelen die ons vermaken. We lachen, we voelen hoop, angst, vreugde of ergernis, en die wisselende gevoelens geven ons een levendig gevoel, alsof ons hele wezen even wordt wakker geschud. Muziek en humor zijn hierbij het krachtigste: ze roepen geen diepe gedachten op, maar hun ritme, hun klanken, hun onverwachte wendingen geven ons een direct, lichamelijk genot.

Muziek werkt als een wisselwerking tussen lichaam en geest: de tonen raken ons, roepen beelden en gevoelens op, en die versterken op hun beurt weer de lichamelijke sensatie. We voelen de muziek niet alleen met ons hoofd, maar met ons hele lichaam, alsof onze ziel even meedeint op de maat. Bij humor is het andersom: een gedachte, een verwachting, wordt plotseling ontkracht, en die ontlading voelen we als een lichamelijke schok. We lachen niet om de grap zelf, maar om het feit dat onze geest even in de war wordt gebracht en dan weer tot rust komt. Dat spel van spanning en ontspanning, alsof we even uit balans worden gebracht en dan weer terugvallen, brengt een trilling teweeg in ons lichaam, een gevoel van levendigheid, alsof we even helemaal wakker zijn.

Neem bijvoorbeeld de verbazing van een Indiaan die voor het eerst een fles bier ziet openen: hij snapt niet hoe al dat schuim erin past, en zijn verwondering is zo groot dat we erom moeten lachen. Of stel je voor: een erfgenaam klaagt dat zijn rouwenden er juist vrolijker uitzien naarmate hij ze meer betaalt om treurig te kijken. In beide gevallen lachen we niet om de inhoud, maar om de plotselinge leegte die ontstaat wanneer onze verwachting in duigen valt. Onze geest wordt even op het verkeerde been gezet, en dat geeft ons een gevoel van opluchting, alsof we even bevrijd zijn van de ernst van het leven.

Een bijzonder soort humor is naïviteit, de onschuldige eenvoud van iemand die de gebruikelijke spelregels van de samenleving niet kent. We lachen om de naïeve persoon omdat hij de wereld nog ziet zoals die is, zonder schijn of bedrog, en dat confronteert ons met onze eigen kunstmatigheid. Maar tegelijkertijd raakt zijn oprechtheid ons, omdat ze ons herinnert aan een zuiverheid die we zelf misschien zijn kwijtgeraakt. Naïviteit is geen kunstmatige deugd, maar een natuurlijke kwaliteit die ons zowel doet lachen als ontroeren. Het is een moment waarop de schijnwereld even wordt doorbroken, en dat geeft ons een gevoel van verlichting en tederheid.

Ten slotte is er nog de grillige geest, iemand die speels en bewust met perspectieven schuift, die dingen anders ziet dan anderen en ons daardoor verrast. Dit is geen ware kunst, maar wel een vorm van geestigheid die ons vermakt door onverwachte contrasten. Het is geen diepgaande schoonheid, maar wel een spel dat ons even uit onze dagelijkse sleur haalt.

Al deze vormen van gratie, muziek, humor, naïviteit, grilligheid, laten zien dat ons genot niet alleen een kwestie is van de geest, maar van ons hele wezen: lichaam en ziel. Ze herinneren ons eraan dat we niet alleen denkende wezens zijn, maar ook voelende, levende wezens, die soms gewoon nodig hebben om te lachen, om mee te deinen op de maat, of om even verrast te worden. Dat is geen oppervlakkigheid, maar een teken van levendigheid, van het vermogen om het leven niet alleen serieus te nemen, maar ook om ervan te genieten.

\bigskip
Deze opmerking vormt een verrassend en diepzinnig scharnier in het betoog over schoonheid en esthetica. Ze staat als een tussenspel tussen de analyse van schone kunsten en de overgang naar morele en teleologische vraagstukken, en dat is geen toeval. Hier onthult zich namelijk de menselijke kern van Kants esthetica: schoonheid en genot zijn niet louter kwesties van smaak of verstand, maar raken ons in ons hele wezen, als lichamelijke en geestelijke wezens. Deze paragraaf doet drie dingen die cruciaal zijn voor het begrijpen van Kants filosofie als geheel, en die ook voor vrijmetselaren diep resoneren.

Kant laat hier zien dat onze ervaring van de wereld niet alleen een zaak is van redelijke beoordeling (zoals bij schoonheid), maar ook van lichamelijke sensatie (zoals bij gratie, humor of muziek). Dit is een radicale erkenning dat de mens niet alleen een denkend, maar ook een voelend wezen is. Schoonheid raakt ons via ons oordeel, maar genot, of het nu komt van muziek, een grap, of een gezellig samenzijn, raakt ons via ons lichaam. Dat is geen mindere ervaring, maar een andere laag van ons bestaan. Voor vrijmetselaren is dit herkenbaar: rituelen en symbolen werken niet alleen op het verstand, maar ook op het gevoel en de zintuigen. Een tempel is niet alleen een gebouw, maar een ruimte die ons lichamelijk en geestelijk aanspreekt; muziek in de loge is niet alleen mooi, maar roert ons diep vanbinnen.
Deze opmerking herinnert ons eraan dat ware wijsheid niet alleen een kwestie is van denken, maar ook van ervaren, van het vermogen om ons hele wezen, lichaam en ziel, open te stellen voor de wereld. Dat is precies wat het maçonnieke ideaal van integrale ontwikkeling nastreeft: niet alleen het verstand, maar het hele menselijke wezen verrijken.

De nadruk op spel, of het nu gaat om muziek, humor, of het vrije spel van gedachten, is geen toeval. Kant laat hier zien dat genot vaak voortkomt uit beweging, verandering en onverwachte wendingen: het is een dynamisch proces, geen statische toestand. Dit is een sleutelidee: ware levendigheid ontstaat niet uit vastomlijnde regels, maar uit vrijheid en spontaniteit. Muziek en humor werken omdat ze ons verrassen, omdat ze ons even uit ons normale denkpatroon halen en ons dwingen om anders te kijken. Dat is precies wat ook het maçonnieke ritueel doet: het doorbreekt de allang, het biedt een tijdelijke ruimte waarin we de wereld anders kunnen ervaren.
Deze opmerking is dus niet alleen een uitstapje over genot, maar een diepe reflectie op vrijheid zelf. Ware cultuur, of het nu in kunst, filosofie of spirituele praktijk is, vereist het vermogen om spelenderwijs met ideeën en ervaringen om te gaan. Voor vrijmetselaren is dit een herkenbaar principe: het ritueel is geen starre ceremonie, maar een levende ervaring, die elke keer opnieuw kan verrassen en inspireren.

Deze opmerking staat niet toevallig tussen de bespreking van schone kunsten en de overgang naar morele en teleologische vraagstukken. Kant bereidt hier de lezer voor op een belangrijke verschuiving: schoonheid is niet alleen een kwestie van smaak, maar ook een brug naar morele reflectie. Het genot dat hij hier beschrijft, of het nu komt van muziek, humor of naïviteit, is lichamelijk, maar het opent ook de deur naar iets diepers: het vermogen om onbevangen en vrij te ervaren, zonder direct een doel of belang te zoeken. Dat is precies de houding die nodig is voor morele reflectie.
Voor vrijmetselaren is dit een cruciale gedachte. Het maçonnieke pad is niet alleen gericht op persoonlijke verrijking, maar ook op morele groei en het zoeken naar zin in een groter geheel. Kants opmerking laat zien hoe esthetische ervaringen, zoals humor, muziek of het aanschouwen van schoonheid, ons kunnen voorbereiden op die morele reflectie. Ze leren ons om open en vrij te staan in de wereld, zonder ons direct te laten leiden door vooroordelen of eigenbelang. Dat is geen kleine zaak: het is de grondhouding die nodig is voor ware wijsheid.

Kant besteedt speciale aandacht aan naïviteit en grilligheid, twee kwaliteiten die op het eerste gezicht oppervlakkig lijken, maar die in werkelijkheid diepgaande lessen bevatten. Naïviteit is de onschuldige eenvoud van iemand die nog niet heeft geleerd om te veinzen of te manipuleren. We lachen erom, maar tegelijkertijd raakt het ons, omdat het ons herinnert aan een zuiverheid die we zelf misschien zijn kwijtgeraakt. Grilligheid daartegen is het vermogen om speels en vrij met perspectieven om te gaan, om dingen anders te zien dan anderen. Beide kwaliteiten laten zien dat ware wijsheid niet altijd serieus hoeft te zijn, dat ze soms juist schuilt in het vermogen om lichtvoetig en open te staan in de wereld.
Voor vrijmetselaren is dit een belangrijke les. Het maçonnieke pad is niet alleen een kwestie van ernst en diepgang, maar ook van speelsheid en openheid. Het ritueel is geen starre ceremonie, maar een levende ervaring, die ruimte laat voor verrassing, humor en onverwachte inzichten. Kants opmerking herinnert ons eraan dat ware wijsheid niet alleen komt van studeerkamerfilosofie, maar ook van het vermogen om met verwondering en vrijheid in het leven te staan.



\end{mdframed}